\documentclass[11pt]{article} 
\usepackage{etoolbox} \newtoggle{perso}
%\toggletrue{perso} % Comment for student version
\usepackage[mathletters]{ucs}
\usepackage[utf8x]{inputenc}

\voffset=-2cm
%\voffset=-3cm
\hoffset=-1.5cm
\textwidth=16cm
\textheight=23cm

\usepackage{amsthm}
\usepackage{amssymb,amsmath}
\usepackage{todonotes}
\usepackage{enumerate}
\usepackage{abraces}
\newcommand\overpar[1]{\aoverbrace[L1R]{#1}}
\usepackage{hyperref}
% Generated by scripts/lean_pages/leancheck.py from the Lean files: do not edit by hand.
% Usage: % Generated by scripts/lean_pages/leancheck.py from the Lean files: do not edit by hand.
% Usage: % Generated by scripts/lean_pages/leancheck.py from the Lean files: do not edit by hand.
% Usage: \input{leananchors} in the preamble, after hyperref; then \leancheck{<label>} next to
% a formalized statement prints a "Lean" mark linking to its Lean formalization.
% Website: \newcommand\leanmapurl{<address>} before \input{leananchors} overrides the default.
% Mark: likewise \newcommand\leancheckmark{...}.
\providecommand\leanmapurl{https://mirebeau.github.io/lean-proof-maps/M1Analyse/}
\providecommand\leancheckmark{\textup{\textsf{\scriptsize Lean}}\ensuremath{^\checkmark}}
\makeatletter
\DeclareRobustCommand\leancheck[1]{%  robust: usable in titles
  \@ifundefined{leancheck@#1}%
    {\PackageWarning{leananchors}{No Lean declaration for the label `#1'}}%
    {\href{\leanmapurl\#label=#1}{\leancheckmark}}}
% label -> number of Lean declarations (their names)
\@namedef{leancheck@def:dense}{2}% def_dense, def_separable
\@namedef{leancheck@def:interieur}{2}% def_interieur, def_adherence
\@namedef{leancheck@def:topo}{3}% def_topo_axioms, def_topo_of_axioms, def_topo_induite
\@namedef{leancheck@eq:finter_open}{1}% eq_finter_open
\@namedef{leancheck@eq:int_adh_compl}{1}% eq_int_adh_compl
\@namedef{leancheck@eq:int_inter_adh_cup}{1}% eq_int_inter_adh_cup
\@namedef{leancheck@rem:fermes}{1}% rem_fermes
\@namedef{leancheck@rem:unionVide}{1}% rem_unionVide
\@namedef{leancheck@subsubsec:ouvertsFermes}{1}% M1Analyse.TopologieGenerale.OuvertsFermes
\makeatother
 in the preamble, after hyperref; then \leancheck{<label>} next to
% a formalized statement prints a "Lean" mark linking to its Lean formalization.
% Website: \newcommand\leanmapurl{<address>} before % Generated by scripts/lean_pages/leancheck.py from the Lean files: do not edit by hand.
% Usage: \input{leananchors} in the preamble, after hyperref; then \leancheck{<label>} next to
% a formalized statement prints a "Lean" mark linking to its Lean formalization.
% Website: \newcommand\leanmapurl{<address>} before \input{leananchors} overrides the default.
% Mark: likewise \newcommand\leancheckmark{...}.
\providecommand\leanmapurl{https://mirebeau.github.io/lean-proof-maps/M1Analyse/}
\providecommand\leancheckmark{\textup{\textsf{\scriptsize Lean}}\ensuremath{^\checkmark}}
\makeatletter
\DeclareRobustCommand\leancheck[1]{%  robust: usable in titles
  \@ifundefined{leancheck@#1}%
    {\PackageWarning{leananchors}{No Lean declaration for the label `#1'}}%
    {\href{\leanmapurl\#label=#1}{\leancheckmark}}}
% label -> number of Lean declarations (their names)
\@namedef{leancheck@def:dense}{2}% def_dense, def_separable
\@namedef{leancheck@def:interieur}{2}% def_interieur, def_adherence
\@namedef{leancheck@def:topo}{3}% def_topo_axioms, def_topo_of_axioms, def_topo_induite
\@namedef{leancheck@eq:finter_open}{1}% eq_finter_open
\@namedef{leancheck@eq:int_adh_compl}{1}% eq_int_adh_compl
\@namedef{leancheck@eq:int_inter_adh_cup}{1}% eq_int_inter_adh_cup
\@namedef{leancheck@rem:fermes}{1}% rem_fermes
\@namedef{leancheck@rem:unionVide}{1}% rem_unionVide
\@namedef{leancheck@subsubsec:ouvertsFermes}{1}% M1Analyse.TopologieGenerale.OuvertsFermes
\makeatother
 overrides the default.
% Mark: likewise \newcommand\leancheckmark{...}.
\providecommand\leanmapurl{https://mirebeau.github.io/lean-proof-maps/M1Analyse/}
\providecommand\leancheckmark{\textup{\textsf{\scriptsize Lean}}\ensuremath{^\checkmark}}
\makeatletter
\DeclareRobustCommand\leancheck[1]{%  robust: usable in titles
  \@ifundefined{leancheck@#1}%
    {\PackageWarning{leananchors}{No Lean declaration for the label `#1'}}%
    {\href{\leanmapurl\#label=#1}{\leancheckmark}}}
% label -> number of Lean declarations (their names)
\@namedef{leancheck@def:dense}{2}% def_dense, def_separable
\@namedef{leancheck@def:interieur}{2}% def_interieur, def_adherence
\@namedef{leancheck@def:topo}{3}% def_topo_axioms, def_topo_of_axioms, def_topo_induite
\@namedef{leancheck@eq:finter_open}{1}% eq_finter_open
\@namedef{leancheck@eq:int_adh_compl}{1}% eq_int_adh_compl
\@namedef{leancheck@eq:int_inter_adh_cup}{1}% eq_int_inter_adh_cup
\@namedef{leancheck@rem:fermes}{1}% rem_fermes
\@namedef{leancheck@rem:unionVide}{1}% rem_unionVide
\@namedef{leancheck@subsubsec:ouvertsFermes}{1}% M1Analyse.TopologieGenerale.OuvertsFermes
\makeatother
 in the preamble, after hyperref; then \leancheck{<label>} next to
% a formalized statement prints a "Lean" mark linking to its Lean formalization.
% Website: \newcommand\leanmapurl{<address>} before % Generated by scripts/lean_pages/leancheck.py from the Lean files: do not edit by hand.
% Usage: % Generated by scripts/lean_pages/leancheck.py from the Lean files: do not edit by hand.
% Usage: \input{leananchors} in the preamble, after hyperref; then \leancheck{<label>} next to
% a formalized statement prints a "Lean" mark linking to its Lean formalization.
% Website: \newcommand\leanmapurl{<address>} before \input{leananchors} overrides the default.
% Mark: likewise \newcommand\leancheckmark{...}.
\providecommand\leanmapurl{https://mirebeau.github.io/lean-proof-maps/M1Analyse/}
\providecommand\leancheckmark{\textup{\textsf{\scriptsize Lean}}\ensuremath{^\checkmark}}
\makeatletter
\DeclareRobustCommand\leancheck[1]{%  robust: usable in titles
  \@ifundefined{leancheck@#1}%
    {\PackageWarning{leananchors}{No Lean declaration for the label `#1'}}%
    {\href{\leanmapurl\#label=#1}{\leancheckmark}}}
% label -> number of Lean declarations (their names)
\@namedef{leancheck@def:dense}{2}% def_dense, def_separable
\@namedef{leancheck@def:interieur}{2}% def_interieur, def_adherence
\@namedef{leancheck@def:topo}{3}% def_topo_axioms, def_topo_of_axioms, def_topo_induite
\@namedef{leancheck@eq:finter_open}{1}% eq_finter_open
\@namedef{leancheck@eq:int_adh_compl}{1}% eq_int_adh_compl
\@namedef{leancheck@eq:int_inter_adh_cup}{1}% eq_int_inter_adh_cup
\@namedef{leancheck@rem:fermes}{1}% rem_fermes
\@namedef{leancheck@rem:unionVide}{1}% rem_unionVide
\@namedef{leancheck@subsubsec:ouvertsFermes}{1}% M1Analyse.TopologieGenerale.OuvertsFermes
\makeatother
 in the preamble, after hyperref; then \leancheck{<label>} next to
% a formalized statement prints a "Lean" mark linking to its Lean formalization.
% Website: \newcommand\leanmapurl{<address>} before % Generated by scripts/lean_pages/leancheck.py from the Lean files: do not edit by hand.
% Usage: \input{leananchors} in the preamble, after hyperref; then \leancheck{<label>} next to
% a formalized statement prints a "Lean" mark linking to its Lean formalization.
% Website: \newcommand\leanmapurl{<address>} before \input{leananchors} overrides the default.
% Mark: likewise \newcommand\leancheckmark{...}.
\providecommand\leanmapurl{https://mirebeau.github.io/lean-proof-maps/M1Analyse/}
\providecommand\leancheckmark{\textup{\textsf{\scriptsize Lean}}\ensuremath{^\checkmark}}
\makeatletter
\DeclareRobustCommand\leancheck[1]{%  robust: usable in titles
  \@ifundefined{leancheck@#1}%
    {\PackageWarning{leananchors}{No Lean declaration for the label `#1'}}%
    {\href{\leanmapurl\#label=#1}{\leancheckmark}}}
% label -> number of Lean declarations (their names)
\@namedef{leancheck@def:dense}{2}% def_dense, def_separable
\@namedef{leancheck@def:interieur}{2}% def_interieur, def_adherence
\@namedef{leancheck@def:topo}{3}% def_topo_axioms, def_topo_of_axioms, def_topo_induite
\@namedef{leancheck@eq:finter_open}{1}% eq_finter_open
\@namedef{leancheck@eq:int_adh_compl}{1}% eq_int_adh_compl
\@namedef{leancheck@eq:int_inter_adh_cup}{1}% eq_int_inter_adh_cup
\@namedef{leancheck@rem:fermes}{1}% rem_fermes
\@namedef{leancheck@rem:unionVide}{1}% rem_unionVide
\@namedef{leancheck@subsubsec:ouvertsFermes}{1}% M1Analyse.TopologieGenerale.OuvertsFermes
\makeatother
 overrides the default.
% Mark: likewise \newcommand\leancheckmark{...}.
\providecommand\leanmapurl{https://mirebeau.github.io/lean-proof-maps/M1Analyse/}
\providecommand\leancheckmark{\textup{\textsf{\scriptsize Lean}}\ensuremath{^\checkmark}}
\makeatletter
\DeclareRobustCommand\leancheck[1]{%  robust: usable in titles
  \@ifundefined{leancheck@#1}%
    {\PackageWarning{leananchors}{No Lean declaration for the label `#1'}}%
    {\href{\leanmapurl\#label=#1}{\leancheckmark}}}
% label -> number of Lean declarations (their names)
\@namedef{leancheck@def:dense}{2}% def_dense, def_separable
\@namedef{leancheck@def:interieur}{2}% def_interieur, def_adherence
\@namedef{leancheck@def:topo}{3}% def_topo_axioms, def_topo_of_axioms, def_topo_induite
\@namedef{leancheck@eq:finter_open}{1}% eq_finter_open
\@namedef{leancheck@eq:int_adh_compl}{1}% eq_int_adh_compl
\@namedef{leancheck@eq:int_inter_adh_cup}{1}% eq_int_inter_adh_cup
\@namedef{leancheck@rem:fermes}{1}% rem_fermes
\@namedef{leancheck@rem:unionVide}{1}% rem_unionVide
\@namedef{leancheck@subsubsec:ouvertsFermes}{1}% M1Analyse.TopologieGenerale.OuvertsFermes
\makeatother
 overrides the default.
% Mark: likewise \newcommand\leancheckmark{...}.
\providecommand\leanmapurl{https://mirebeau.github.io/lean-proof-maps/M1Analyse/}
\providecommand\leancheckmark{\textup{\textsf{\scriptsize Lean}}\ensuremath{^\checkmark}}
\makeatletter
\DeclareRobustCommand\leancheck[1]{%  robust: usable in titles
  \@ifundefined{leancheck@#1}%
    {\PackageWarning{leananchors}{No Lean declaration for the label `#1'}}%
    {\href{\leanmapurl\#label=#1}{\leancheckmark}}}
% label -> number of Lean declarations (their names)
\@namedef{leancheck@def:dense}{2}% def_dense, def_separable
\@namedef{leancheck@def:interieur}{2}% def_interieur, def_adherence
\@namedef{leancheck@def:topo}{3}% def_topo_axioms, def_topo_of_axioms, def_topo_induite
\@namedef{leancheck@eq:finter_open}{1}% eq_finter_open
\@namedef{leancheck@eq:int_adh_compl}{1}% eq_int_adh_compl
\@namedef{leancheck@eq:int_inter_adh_cup}{1}% eq_int_inter_adh_cup
\@namedef{leancheck@rem:fermes}{1}% rem_fermes
\@namedef{leancheck@rem:unionVide}{1}% rem_unionVide
\@namedef{leancheck@subsubsec:ouvertsFermes}{1}% M1Analyse.TopologieGenerale.OuvertsFermes
\makeatother
 % \leancheck{<label>} : lien vers la formalisation Lean (généré par scripts/lean_pages/leancheck.py)
\usepackage{pifont} \def\facultatif{\ding{40}} % avion

\def\cA{\mathcal{A}}
\def\cB{\mathcal{B}}
\def\cD{\mathcal{D}}
\def\cE{\mathcal{E}}
\def\cF{\mathcal{F}}
\def\cG{\mathcal{G}}
\def\cK{\mathcal{K}}
\def\cO{\mathcal{O}}
\def\cP{\mathcal{P}}
\def\cS{\mathcal{S}}
\def\cU{\mathcal{U}}
\def\cV{\mathcal{V}}
\def\cW{\mathcal{W}}

\def\bC{\mathbb{C}}
\def\bF{\mathbb{F}}
\def\bK{\mathbb{K}}
\def\bN{\mathbb{N}}
\def\bP{\mathbb{P}}
\def\bQ{\mathbb{Q}}
\def\bR{\mathbb{R}}
\def\bS{\mathbb{S}}
\def\bT{\mathbb{T}}
\def\bU{\mathbb{U}}
\def\bV{\mathbb{V}}
\def\bX{\mathbb{X}}
\def\bY{\mathbb{Y}}
\def\bZ{\mathbb{Z}}

\def\sm{\setminus}
\def\vp{\varphi}
\def\ve{\varepsilon}
\def\<{\langle} \def\>{\rangle}
\usepackage{bbm} \def\bbone{\mathbbm{1}}
\def\bi{\boldsymbol{i}}

\iftoggle{perso}{
\usepackage[notref]{showkeys}
\newcommand\todoi[1]{\todo[inline]{#1}}
\newcommand\notei[1]{\todo[inline,color=yellow]{#1}}
\newcommand\hiddenproof[1]{\begin{proof}[Preuve (brouillon seul)] #1 \end{proof}}

\newcommand\inventaire[1]{}
%\newcommand\inventaire[1]{\paragraph{Inventaire (pour le cours)\\} #1 \hfill $\square$}

\newcommand\formalizationInstruction[1]{\todo[inline,color=yellow]{#1}}
\newcommand\claudereport[1]{\todo[inline,color=green]{#1}}
}{
\newcommand\todoi[1]{}
\newcommand\notei[1]{}
\newcommand\hiddenproof[1]{}
\newcommand\inventaire[1]{}
\newcommand\formalizationInstruction[1]{}
\newcommand\claudereport[1]{}
}

\usepackage[french]{cleveref} %capitalize,
\newtheorem{theorem}{Th\'eor\`eme}[section]
\newtheorem{proposition}[theorem]{Proposition}
\newtheorem{remark}[theorem]{Remarque}
\newtheorem{lemma}[theorem]{Lemme}
\newtheorem{corollary}[theorem]{Corollaire}
\newtheorem{definition}[theorem]{D\'efinition}
\newtheorem{notation}[theorem]{Notation}
\newtheorem{assumption}[theorem]{Hypoth\`ese}
\newtheorem{exo}[theorem]{Exercice}
\newtheorem{exemple}[theorem]{Exemple}
\newtheorem*{retenir}{\`A retenir absolument}
\renewcommand*{\proofname}{Preuve}
%\crefname{exemple}{exemple}{exemples} % Incompatible avec TexPad, donc un exemple à la fois...

\usepackage{nameref} 
\newcommand{\nref}[1]{\cref{#1}\space(\nameref{#1})}
\newcommand{\Nref}[1]{\Cref{#1}\space(\nameref{#1})}

% The clever way to do it, below, kills pdfLaTeX. https://tex.stackexchange.com/a/109864
%\makeatletter
%\newcommand{\nref}[1]{\cref{#1}\mynameref{#1}{\csname r@#1\endcsname}}
%\newcommand{\Nref}[1]{\Cref{#1}\mynameref{#1}{\csname r@#1\endcsname}}
%\def\mynameref#1#2{%
%  \begingroup
%    \edef\@mytxt{#2}%
%    \edef\@mytst{\expandafter\@thirdoffive\@mytxt}%
%    \ifx\@mytst\empty\else
%    \space(\nameref{#1})\fi
%  \endgroup
%}
%\makeatother

\def\diff{\mathrm{d}}

\DeclareMathOperator\Adh{Adh}
\DeclareMathOperator\supp{supp}
\DeclareMathOperator\diver{div}
\DeclareMathOperator\Lip{Lip}
\DeclareMathOperator\Ext{Ext}
\DeclareMathOperator\Ker{Ker}
\DeclareMathOperator\Image{Im}
\DeclareMathOperator\Diag{Diag}
\DeclareMathOperator\Dom{Dom}
\DeclareMathOperator\Aff{Aff}
\DeclareMathOperator\Relint{ri}
\DeclareMathOperator\Proj{Proj}
\DeclareMathOperator\Id{Id}
\DeclareMathOperator\Tr{Tr}
\DeclareMathOperator\Hull{Hull}
\DeclareMathOperator\Vect{Vect}
\DeclareMathOperator\sign{sign}
\DeclareMathOperator*\argmin{argmin}

\def\loc{\mathrm{loc}}
\DeclareMathOperator\ordre{ordre}
\DeclareMathOperator\Leb{Leb}
\def\kC{\mathfrak{C}}
\def\weakto{\rightharpoonup}

\newcommand{\vertiii}[1]{{\left\vert\kern-0.25ex\left\vert\kern-0.25ex\left\vert #1 
    \right\vert\kern-0.25ex\right\vert\kern-0.25ex\right\vert}}
\usepackage{stmaryrd}
\newcommand\braii[1]{{\left\llbracket #1 \right\rrbracket}}


\def\eqdef{\overset {\tiny \mathrm{def}} \Leftrightarrow}
\def\csubset{\subset_c}


%latex_to_unicode=False
% The script latex_to_unicode.py will not tamper with this part of the document.
%latex_to_unicode=True

\title{Analyse fonctionnelle}

\author{Jean-Marie Mirebeau\footnote{
Universit\'e Paris-Saclay, ENS Paris-Saclay, CNRS, Centre Borelli, F-91190 Gif-sur-Yvette, France.
}}

\begin{document}
\maketitle
\begin{abstract}
Ce cours est une introduction à l'analyse et à l'analyse fonctionnelle, et présente un certain nombre de résultats fondateurs : théorèmes de point fixe (Picard, Brouwer) et autres conséquences classiques de la complétude et de la compacité, dualité et topologie faible, espaces de Sobolev.
Le cours se conçoit comme une base solide pour ceux qui se destinent à étudier l'analyse ou l'analyse numérique, et pourra également intéresser les candidats à l'agrégation. 
\end{abstract}


\tableofcontents

\paragraph{Références}

Ces \href{https://www.dropbox.com/scl/fi/d3vfdgnaje0iwxyfyejqw/M1_Analyse_Shared.pdf?rlkey=qfngap7sa3uh1xo5ukxt8ddpi&dl=0}{notes de cours (lien vers la dernière version)} sont fortement inspirées du 
\href{http://talazard.perso.math.cnrs.fr/Livre.pdf}{livre de Thomas Alazard (lien)} 
\cite{alazard2017analyse} et 
\href{https://www.dropbox.com/scl/fi/t7tj1pm7pxcc5t67gyizj/Desvillettes_2015_AF.pdf?rlkey=8k1gvfgcmiczqqjxxvwblnugn&dl=1}{du polycopié de Laurent Desvillettes (lien)} 
\cite{desvilettes2015analyse}, qui ont précédemment enseigné ce cours
%\footnote{Lien de téléchargement des polycopiés précédents : 
%\url{http://talazard.perso.math.cnrs.fr/Livre.pdf}\newline
%\href{https://www.dropbox.com/scl/fi/t7tj1pm7pxcc5t67gyizj/Desvillettes_2015_AF.pdf?rlkey=8k1gvfgcmiczqqjxxvwblnugn&dl=1}{Laurent Desvillettes}
%\url{http://laurent.desvillettes.free.fr/analyse_fonctionnelle79.pdf} % Site mort en 2024
%https://dl.dropbox.com/s/2xi2jzwc4htfjw5/Desvillettes_2015_AF.pdf
%}. 

Le contenu du cours étant assez classique, de nombreuses autres références de qualité et en français sont disponibles, par exemple le livre de Daniel Li \cite{li2013cours}, ou le 
\href{https://www.ceremade.dauphine.fr/~carlier/poly2010.pdf}{polycopié plus ancien du même cours par Guillaume Carlier (lien)} \cite{carlier2009analyse}.

Je tiens également à mettre en avant les 
\href{https://www.dropbox.com/scl/fi/3gir3uhkt842idl9sfyc3/main_Shared.pdf?rlkey=xruzgtc6kzfedhjx0uecku0yq&dl=0}{notes de cours indépendantes prises par un groupe d'étudiants mené par Félix Yvonnet en 2023 (lien)}, particulièrement agréables à lire et accompagnées de figures. 


\paragraph{Points importants, ou annexes.}
Chaque sous-chapitre contient une liste de résultats à retenir absolument, fondamentaux pour les applications en analyse fonctionnelle. À l'inverse le symbole \facultatif\ signale des points facultatifs, d'illustration et d'ouverture vers d'autres domaines mathématiques.

\newpage
\section{Topologie générale}
\label{sec:topo}
Cette section est principalement constituée de rappels de topologie générale, qui forment une base indispensable pour l'analyse. Les espaces fonctionnels rencontrés en analyse sont en général des espaces vectoriels normés, ou munis de familles de semi-normes, voir les \cref{subsec:evn,subsec:evtlc}.


\begin{remark}
	Dans cette section, les espaces topologiques, métriques, ... sont présentés comme une paire de symboles tels que $(X,\bU)$ ou $(X,d)$, où $X$ désigne un ensemble de points, et $\bU$ ou $d$ une structure associée (topologie, distance, ...). Par abus de language, et en l'absence d'ambiguïté, on se permettra également par la suite de donner seulement le symbole correspondant à l'ensemble de points (par exemple ``soit $X$ un espace topologique''). 
\end{remark}


\subsection{Espaces topologiques}

\begin{retenir}
La définition d'une topologie, de la topologie engendrée, et d'une base de voisinages. 
La caractérisation séquentielle de l'adhérence \emph{si} l'espace topologique a une base dénombrable de voisinages. 
L'axiome de séparation de Hausdorff, la définition et la caractérisation de la continuité.
(\Cref{def:topo,def:topoEngendree,def:base_vois,prop:adhSuite,def:separe,def:continuite,prop:caract_cont}.)
\end{retenir}

\inventaire{
\Nref{subsubsec:ouvertsFermes} \\
\Nref{def:topo}\\
\Nref{rem:unionVide}\\
\Nref{rem:fermes}\\
\Nref{def:interieur}\\
\Nref{def:dense}\\ 

\Nref{subsubsec:compTopo} \\
\Nref{def:finesseTopo} \\
\Nref{def:topoEngendree} \\
\Nref{prop:baseOuverts} \\
\Nref{corol:baseOuverts} (Identification de la topologie engendrée par une base d'ouverts) \\
\Nref{ex:topoOrdre} \\
\Nref{ex:topoProduit} \\

\Nref{subsubsec:Vois}\\
\Nref{def:filtre_vois} \\
\Nref{prop:adh_vois} \\
\Nref{def:base_vois} \\
\Nref{rem:topoVois} \\
\Nref{def:denVois} \\
\Nref{def:cvSuite} \\
\Nref{prop:adhSuite} \\
\Nref{prop:adhSousSuite} \\
\Nref{rem:axiomFiltre} \\

\Nref{subsubsec:separation} \\
\Nref{def:separe} \\
\Nref{rem:uniqueLim} \\
\Nref{rem:T1} \\
\Nref{ex:T1nonT2} \\

\Nref{subsubsec:continuite} \\
\Nref{def:continuite} \\
\Nref{prop:caract_cont} \\
\Nref{rem:topoFunc} \\
\Nref{def:semiContinuite} \\
\Nref{prop:surGrapheSCI} (Sci ssi sur-graphe fermé)\\

\Nref{subsec:metric} \\
\Nref{def:distance} \\
\Nref{def:boule} \\
\Nref{rem:distInfinie} \\
\Nref{prop:baseVoisMetric} (Base de voisinages dans un espace métrique) \\
\Nref{def:moduleCont} \\
\Nref{lem:regModuleCont} \\
\Nref{lem:contDistance} \\
\Nref{def:unifContinu} \\
\Nref{prop:distCvSimple} \\
\Nref{def:dist_unif} \\
\Nref{rem:equicont} \\

\Nref{subsec:evn} \\
\Nref{def:norme} \\
\Nref{def:normeOperateur} \\
\Nref{def:normeHolder} \\
\Cref{def:normeCnb} ($C_b^n$)\\
\Nref{def:normeLp} \\

\Nref{subsec:evtlc} \\
\Nref{def:evtlc} \\
\Nref{lem:evtlc_separe} \\
\Nref{lem:evtlcLinCont} \\
\Nref{rem:sommeMaxEvtlc} \\
\Nref{rem:evtlcGradue} \\
\Nref{prop:distanceEvtlcs} \\
\Nref{def:lineairBorne} \\
\Nref{def:cinf_bornee} \\
\Nref{def:classeSchwartz} \\
\Nref{prop:exhaustComp} \\
\Cref{def:lploc} ($L^p_{\loc}$)\\
\Nref{def:cont_nonbornee} \\
\Nref{def:distribution} \\
\Nref{prop:distribCont} \\
\Nref{def:topoFaible} \\
\Nref{prop:bouleMetrisable} 
} %inventaire


\subsubsection{Ouverts et fermés}
\label{subsubsec:ouvertsFermes}

\notei{Conventions : 
Élément $x$, ensemble $X$, ensemble d'ensembles $\bX$
Espace vectoriel $E$.
Exceptions : $\cP(X)$ parties, pour ne pas confondre avec $\bP(X)$ les probas .\newline
Indices $i$, $n$, de $1$ à $I$ ou $N$}

On note $\cP(X)$ l'ensemble des parties d'un ensemble $X$. 

\begin{definition}[Espace topologique]
\label{def:topo}
\leancheck{def:topo}%
	Un espace topologique est une paire $(X, \bU)$ où $X$ est un ensemble, et où $\bU\subset \cP(X)$ appelé \emph{ensemble des ouverts} satisfait les axiomes suivants : 
	\begin{enumerate}[(i)]
		\item $\emptyset$ et $X$ sont ouverts.
		\begin{equation*}
			\emptyset, X \in \bU
		\end{equation*}
		\item Une union quelconque d'ouverts est un ouvert. 
		\begin{equation*}
		\label{eq:union_open}
			\forall \cU \subset \bU,\ \bigcup_{U \in \cU} U \in \bU.
		\end{equation*}
		\item L'intersection de deux ouverts est un ouvert.
		\begin{equation*}
		\label{eq:inter_open}
			\forall U,V \in \bU,\ U \cap V \in \bU.
		\end{equation*}
	\end{enumerate}
Sur toute partie $A \subset X$ donnée, on dispose de la topologie induite $\bU_A := \{U \cap A \mid U \in \bU\}$. 
\end{definition}


On déduit immédiatement du point (iii) de la \cref{def:topo} que l'intersection d'un nombre fini positif d'ouverts est un ouvert. 
\begin{equation}
	\label{eq:finter_open}
	\forall I \geq 1,\ \forall U_1,\cdots,U_I \in \bU,\ \bigcap_{1 \leq i \leq I} U_i \in \bU.
\leancheck{eq:finter_open}%
\end{equation}
Pour le point (ii), plutôt qu'un ensemble $\cU$ d'ouverts, on pourrait de manière équivalente considérer l'union $\cup_{i \in I} U_i$ d'une famille d'ouverts $(U_i)_{i \in I}$ indexée par un ensemble $I$ quelconque. 

\begin{remark}[Union et intersection vides]
\label{rem:unionVide}
\leancheck{rem:unionVide}%
L'union vide, correspondant à $\cU = \emptyset$ dans \eqref{eq:union_open}, est définie comme l'ensemble vide. Ainsi, l'axiome $\emptyset \in \bU$ est redondant. 

L'intersection vide, cas hypothétique $I=0$ dans \eqref{eq:finter_open}, n'est pas définie en théorie des ensembles en général.
Cependant, si l'on considère l'intersection d'une famille de parties d'un ensemble $X$ \emph{fixé a priori}, alors il est naturel de définir par convention l'intersection vide comme l'ensemble $X$ dans sa totalité.
%(On peut  à moins que le type d'ensemble soit fixé a priori. Si c'est le cas, alors on peut étendre \eqref{eq:finter_open} au cas $I=0$ obtenir l'ensemble $X$ dans sa totalité, ce qui rend l'axiome $X \in \bU$ redondant. Dans la suite, cependant, `ìntersection quelconque sous entend'' d'une famille non-vide d'ensembles. 
\end{remark}

\begin{remark}[Fermés]
\label{rem:fermes}
\leancheck{rem:fermes}%
On appelle fermé, de l'espace topologique $(X,\bU)$, le complémentaire d'un ouvert. Il s'ensuit (i) que $\emptyset$ et $X$ sont des fermés, et que l'ensemble des fermés $\bF := \{X \sm U\mid U \in \bU\}$ est (ii) stable par intersection quelconque, et (iii) stable par union finie. 
\end{remark}

On peut définir une topologie par la donnée de ses fermés, pourvu qu'ils satisfassent les axiomes de la \cref{rem:fermes}, voir l'\cref{ex:T1nonT2} (Topologie de Zariski).

La stabilité des ouverts par union, et des fermés par intersection, permet de parler du plus grand ouvert, et du plus petit fermé satisfaisant une propriété (qui doit être stable par union ou intersection respectivement). On définit ainsi l'intérieur et l'adhérence d'une partie.

\begin{definition}[Intérieur, adhérence]
\label{def:interieur}
\leancheck{def:interieur}%
Soit $A$ une partie d'un espace topologique $(X,\bU)$. Son intérieur, noté $\mathring A$, est le plus grand ouvert qu'il contient. Son adhérence, notée $\overline A$, est le plus petit fermé le contenant. Ces ensembles peuvent être caractérisés par 
\begin{align}
\label{eqdef:int_adh}
	\mathring A &:= \cup \{U\in \bU \mid U \subset A\}, &
	\overline A &:= \cap \{F \in \bF \mid A \subset F\}.
\end{align}	
\end{definition}

On note, par passage au complémentaire, que
\begin{align}
\label{eq:int_adh_compl}
	X \sm \mathring A &= \overline {X \sm A}, &
	\mathring {\overpar{X \sm A}} &= X \sm \overline A. 
\leancheck{eq:int_adh_compl}%
\end{align}
Également, pour tous $A,B \subset X$,
\begin{align}
\label{eq:int_inter_adh_cup}
	\mathring{\overpar{A \cap B}} &= \mathring A \cap \mathring B,&
	\overline {A \cup B} = \overline A \cup \overline B.
\leancheck{eq:int_inter_adh_cup}%
\end{align}

\hiddenproof{
On note que l'intersection \eqref{eqdef:int_adh} n'est pas vide car $A\subset X \in \bF$.
	Étant donné $U \in \bU$ ouvert et $F := X \sm U$ son complémentaire, on a $A \supset U$ ssi $F\supset A$, d'où (\ref{eq:int_adh_compl}, gauche). En appliquant ce résultat à $Y := X \sm A$ on obtient $X \sm \mathring{\overpar{X\sm A}} = \overline {X \sm (X \sm A)} = \overline A$, d'où (\ref{eq:int_adh_compl}, droite) par passage au complémentaire.
}

\begin{definition}[Partie dense, séparabilité]
\label{def:dense}
\leancheck{def:dense}%
	Une partie $A \subset X$ est dite dense si $\overline A = X$.
	Un espace topologique est dit séparable si il admet une partie dénombrable dense.
	(\emph{Attention}: ne pas confondre séparable et séparé, voir la \cref{def:separe})
\end{definition}

\subsubsection{Comparaison de topologies}
\label{subsubsec:compTopo}

%\notei{Manque une intro. Topologie grossière : permet d'avoir de la compacité. Topologie fine : permet de déduire des propriétés fortes de la convergence d'une suite.}

Il est courant en analyse de disposer de plusieurs topologies sur un même ensemble, comme les  topologies forte et faible sur un evn, voir la \cref{subsec:evtlc}. Heuristiquement, l'intérêt d'une topologie fine est que beaucoup de fonctions sont continues,
% (par exemple la norme de evn), 
tandis que l'intérêt d'une topologie grossière est que beaucoup d'ensembles sont compacts, voir le \nref{th:BanachAlaoglu}.

\begin{definition}[Finesse d'une topologie]
\label{def:finesseTopo}
\leancheck{def:finesseTopo}%
Une topologie $\bU_1$ sur un ensemble $X$ est dite \emph{plus fine} qu'une autre $\bU_2$ si elle possède plus d'ouverts, i.e. 
$
	\bU_1 \supset \bU_2. 
$
On définit en particulier :
%C'est une relation d'ordre sur les topologies, inversée par rapport à l'inclusion, dont l'élément minimal et maximal est respectivement:
\begin{itemize}
	\item La topologie fine, définie par $\bU := \cP(X)$.
	\item La topologie grossière, définie par $\bU := \{\emptyset, X\}$. 
\end{itemize}	
\end{definition}

On vérifie facilement que l'intersection d'une famille quelconque de topologies reste une topologie. On peut ainsi considérer la plus petite topologie satisfaisant une propriété, pourvu que celle-ci soit compatible avec l'intersection de topologies. En particulier : 

\begin{definition}[Topologie engendrée]
\label{def:topoEngendree}
\leancheck{def:topoEngendree}%
	Soit $X$ un ensemble. La topologie engendrée par $\bU_0 \subset \cP(X)$ est la plus petite topologie $\bU$ sur $X$ contenant $\bU_0$. 
\end{definition}

Les éléments de la topologie engendrée sont $X$ lui-même, ainsi que les unions quelconques d'intersection finies d'elements de $\bU_0$. Ces dernières s'écrivent : 
$$
	U = \bigcup_{i \in I} \bigcap_{1 \leq j \leq J(i)} U_{ij},
$$
où $I$ est un ensemble quelconque d'indices, $J : I \to \bN^*$, et $U_{ij} \in \bU_0$ pour tous $i \in I$, $1 \leq j \leq J(i)$. 
La notion de base d'ouverts permet de caractériser plus simplement la topologie engendrée. 

\begin{proposition}[Base d'ouverts]
\label{prop:baseOuverts}
\leancheck{prop:baseOuverts}%
Une base d'ouverts d'un ensemble $X$, est une partie $\bU_0 \subset \cP(X)$ satisfaisant aux hypothèses suivantes : 
	\begin{enumerate}[(i)]
		\item Couverture de $X$:
		\begin{equation*}
			\bigcup_{U \in \bU_0} U = X.
		\end{equation*}
		\item Stabilité par intersection en tout point:
		\begin{equation*}			
		\forall U,V \in \bU_0,\ \forall x \in U \cap V,\ \exists W \in \bU_0,\ x\in W \text{ et } W \subset U \cap V.
		\end{equation*}
	\end{enumerate}
	Alors la topologie $\bU$ engendrée par $\bU_0$ consiste en les unions quelconques d'éléments de $\bU_0$: 
	\begin{equation}
	\label{eq:topo_engendree}
		\bU = \Big\{ \bigcup_{U \in \cU} U \mid \cU \subset \bU_0\Big\}.
	\end{equation}
\end{proposition}

\begin{proof}
	Notons $\bigcup \cU := \bigcup_{U \in \cU} U$ pour tout $\cU \subset \bU_0$. On a $\bigcup \emptyset = \emptyset \in \bU$, et par (i) $\bigcup\bU_0 = X \in \bU$.  
	De plus, étant donnée une famille $(\cU_i)_{i \in I}$ de parties de $\bU_0$, où $I$ est un ensemble quelconque d'indices, on a $\cup_{i\in I} \bigcup \cU_i = \bigcup \cU_*$ avec $\cU_* := \cup_{i \in  I} \cU_i \subset \bU_0$, d'où la stabilité de $\bU$ par union quelconque. 
	
	Fixons $U,V \in \bU_0$, et pour tout $x \in U \cap V$ choisissons $W_x\in \bU_0$ satisfaisant à la condition (ii). Alors 
	$$
		U \cap V = \bigcup_{x \in U \cap V} W_x,
	$$
	ce qui montre que $U \cap V \in \bU$. 	
	Finalement pour tous $\cU, \cV \subset \bU_0$ on a 
	\begin{equation*}
		\big(\bigcup \cU) \cap \big(\bigcup \cV) = \bigcup_{U\in\cU, V \in \cV} U \cap V 
	\end{equation*}
	ce qui prouve la stabilité de $\bU$ par intersection, et conclut la preuve.
\end{proof}

\begin{corollary}
\label{corol:baseOuverts}
\leancheck{corol:baseOuverts}%
Soit $(X,\bU)$ un espace topologique, et soit $\bU_0 \subset \bU$ un sous-ensemble d'ouverts tel que pour tous $U \in \bU$ et $x \in U$, il existe $U_x \in \bU_0$ tel que $x \in U_x \subset U$. 
Alors $\bU_0$ est une base d'ouverts pour la topologie $\bU$. 
%	Soit $(X,\bU)$ un espace topologique, et soit $\bU_0 \subset \bU$ une partie satisfaisant aux axiomes (i) et (ii) de la \cref{prop:baseOuverts}, ainsi que : 
%	(iii) pour tout $U \in \bU$ et tout $x \in U$, il existe $V\in \bU_0$ tel que $x\in V$ et $V \subset U$. Alors $\bU$ est la topologie engendrée par $\bU_0$.
\end{corollary}

\begin{proof}
	Choisissons, comme suggéré par l'énoncé, $U_x \in \bU_0$ tel que $x \in U_x \subset U$ pour tous $x \in U \in \bU$. 
	On a $X \in \bU$ par définition, et donc $X = \cup_{x \in X} U_x \subset \cup_{U \in \bU_0} U$ comme souhaité pour le point (i). 
	Le point (ii) s'obtient en notant que $U \cap V \in \bU$, pour tous $U,V \in \bU_0$, et en choisissant $W = [U\cap V]_x$. 
	Finalement, pour tout $U \in \bU$ on a $U = \cup_{x \in U} U_x$, donc tout élément de $\bU$ s'écrit comme union d'éléments de $\bU_0$, et réciproquement $\bU$ est bien une topologie contenant $\bU_0$. 
%	Par définition $X \in \bU$, et donc par hypothèse tout point $x \in X$ est contenu dans un certain $U_0 \in bU_0$, ce qui établit le point (i) de la \cref{prop:baseOuverts}. Si $U,V \in \bU_0$, alors $U \cap V \in \bU$ et donc l'hypothèse assure que (ii) est également satisfait. 
%	Finalement, tout ouvert $U = 
%	La topologie engendrée par $\bU_0$ est contenue 
\end{proof}

\begin{exemple}[topologie de l'ordre \facultatif]
\label{ex:topoOrdre}
\leancheck{ex:topoOrdre}%
	Soit $(X,\leq)$ un ensemble totalement ordonné ayant au moins deux points. On appelle topologie de l'ordre la topologie engendrée par la base d'ouverts suivante : pour tous $a,b \in X$
	\begin{align}
		&\{x\in X\mid x<b\},&
		&\{x\in X\mid a<x<b\},&
		&\{x\in X\mid a<x\}, 
		%& X
	\end{align}
	Ces ensembles sont aussi notés $]-\infty,b[$, $]a,b[$, $]a,\infty[$.
\end{exemple}

\begin{proof}
	Avoir deux points $a<b \in X$ permet de satisfaire l'axiome de couverture avec $]-\infty,b[$ et $]a,\infty[$. 
	La stabilité par intersection découle de ce que l'intersection de deux intervalles ouverts est un intervalle ouvert (éventuellement vide) : $]a,b[ \cap ]c,d[ = ]\max\{a,c\},\min\{b,d\}[$, pour tous $a,b,c,d$, éventuellement infinis et avec les conventions naturelles.
\end{proof}


\begin{exemple}[topologie produit]
\label{ex:topoProduit}
\leancheck{ex:topoProduit}%
	Soit $(X_i, \bU_i)_{i \in I}$ une famille d'espaces topologiques. Alors le produit $\Pi_{i \in I} X_i$ est muni de la topologie dont une base d'ouverts sont les produits 
	$\Pi_{i \in I} U_i$ d'ouverts $U_i \in \bU_i$, $i \in I$, avec \emph{$U_i = X_i$ sauf pour un nombre fini d'indices}. 
\end{exemple}

\begin{proof}
	L'axiome de couverture découle du choix $U_i = X_i$, $i \in I$, et la stabilité de ce que 
	$(\Pi_{i \in I} U_i) \cap (\Pi_{i \in I} V_i) = \Pi_{i \in I} (U_i\cap V_i)$.
\end{proof}

Un cas particulier est celui où $X_i = X$ pour tout $i \in I$, et où le produit s'identifie à l'ensemble $X^I$ des fonctions de $I$ dans $X$. Dans ce cas, la topologie produit correspond à la convergence simple des fonctions. 

\subsubsection{Voisinages}
\label{subsubsec:Vois}

\begin{definition}[Filtre des voisinages]
\label{def:filtre_vois}
\leancheck{def:filtre_vois}%
	Soit $(X,\bU)$ un espace topologique. Un voisinage de $x \in X$ est une partie $V\subset X$ contenant un ouvert contenant $x$, en d'autres termes telle que $x \in \mathring V$. On note $\cV_x$ l'ensemble des voisinages de $x$:
	\begin{equation}
	\label{eqdef:Vx}
		\cV_x := \{V \subset X \mid x \in \mathring V\}.
\leancheck{eqdef:Vx}%
	\end{equation}
\end{definition}
On note qu'une intersection finie de voisinages d'un point donné reste un voisinage, par \eqref{eq:int_inter_adh_cup}.

\begin{proposition}[Caractérisation de l'adhérence]
\label{prop:adh_vois}
\leancheck{prop:adh_vois}%
\begin{align*}
	\overline A &= \{x \in X \mid \forall V \in \cV_x,\ A\cap V \neq \emptyset\},&
	\mathring A &= \{x \in X \mid \exists V \in \cV_x,\ V \subset A\}.
\end{align*}
\end{proposition}

\begin{proof}
La caractérisation de $\mathring A$ est évidente. Celle de $\overline A = X \sm (\mathring{\overpar{X \sm A}})$ s'en déduit.
\end{proof}



\begin{definition}[Base de voisinages]
\label{def:base_vois}
\leancheck{def:base_vois}%
	Une base de voisinages de $x\in X$ est une partie $\cW_x \subset \cV_x$ contenant des voisinages arbitrairement petits, dans le sens où : 
	\begin{equation}
		\forall V \in \cV_x,\, \exists W \in \cW_x,\, W \subset V.
	\end{equation}
\end{definition}

\begin{remark}[Définition d'une topologie à partir d'une base de voisinages]
\label{rem:topoVois}
\leancheck{rem:topoVois}%
Soit $X$ un ensemble, et pour tout $x \in X$ soit $\cW_x$ un ensemble non-vide de parties de $X$, chacune contenant $x$.
Si $\forall x \in X$, $\forall U,V \in \cW_x$, $\exists W \in \cW_x$, $W \subset U \cap V$, alors 
$$\bU := \{U \subset X \mid \forall x \in U,\, \exists V \in \cW_x,\, V \subset U\}$$ 
est une topologie sur $X$.
Si de plus $\forall x\in X$, $\forall U \in \cW_x$, $\forall y\in U$, $\exists V \in \cW_y$, $V \subset U$, alors $(\cW_x)_{x\in X}$ constitue une base de voisinages de $\bU$, et $\cup_{x \in X} \cW_x$ est une base d'ouverts de $\bU$.
\end{remark}

\begin{proof}
	\emph{Premier point :} les axiomes d'une topologie (i) $\emptyset, X \in \bU$ et (ii) de stabilité par réunion, sont facilement vérifiés, tandis que (iii) la stabilité par intersection de deux ouverts découle de l'axiome imposé sur $\cW_x$, pour tout $x \in X$. \emph{Second point :} la nouvelle condition de compatibilité imposée sur les $(\cW_x)_{x \in X}$ assure que chaque $U \in \cW_x$, pour tout $x \in X$, est un ouvert de la topologie $\bU$. De plus $\cup_{x \in X} \cW_x$ est une base d'ouverts de $\bU$ au vu du \ref{corol:baseOuverts}, ainsi qu'une base de voisinages par construction.
\end{proof}


%\hiddenproof{
%\emph{Premier point :} vérifions les axiomes d'une topologie.
%Par construction (i) $\emptyset \in \bU$ et $X \in \bU$, puisque $\cW_x$ est non-vide pour tout $x \in X$. La stabilité par réunion quelconque (ii) est claire, et celle par intersection de deux ensembles (iii) découle du premier axiome satisfait par $\cW_x$ : si $x \in U \subset \cU \in \bU$ et $x \in V \subset \cV \in \bU$ avec $U,V \in \cW_x$, alors $x \in W \subset U \cap V \subset \cU \cap \cV$ où $W$ est choisi comme dans l'énoncé. Ainsi $\bU$ est une topologie sur $X$.  
%\emph{Second point :} le second axiome assure que chaque $W \in \cW_x$, où $x \in X$, appartient à $\bU$, i.e.\ $W$ est un ouvert, comme annoncé. Réciproquement, la définition de la topologie $\bU$ assure que $(\cW_x)_{x \in X}$ est une base de voisinages. Noter que $\bU$ est la plus petite topologie ayant cette propriété.
%}

\begin{definition}[Topologie à base dénombrable]
\label{def:denVois}
\leancheck{def:denVois}%
	Une topologie sur un ensemble $X$ est dite : 
	\begin{itemize}
		\item A \emph{base de voisinages dénombrable (first countable)} si tout point $x\in X$ admet une base de voisinages dénombrable.
		\item A \emph{base dénombrable} (second countable) si elle admet une base dénombrable d'ouverts. 
	\end{itemize}
\end{definition}

Les espaces métriques sont à base dénombrable de voisinages, et les espaces métriques séparables sont à base dénombrable, voir la \cref{prop:baseVoisMetric}.
Ces notions autorisent la caractérisation séquentielle de propriétés topologiques. 


\todoi{Éventuellement mettre le contre exemple avec la suite $\sin(nx)$ sur $[0,1]$ pour la convergence simple \cite{desvilettes2015analyse} ... }
%(Voir  \cite{desvilettes2015analyse} pour une contre exemple sans cette hypothèse.)

\begin{definition}[Convergence d'une suite]
\label{def:cvSuite}
\leancheck{def:cvSuite}%
Une suite $(x_n)_{n \in \bN}$ a valeurs dans un espace topologique $(X,\bU)$ converge vers un point $x \in X$ si et seulement si tout voisinage de ce point contient tous les termes de la suite à partir d'un certain rang:
\begin{equation}
	\lim_{n \to \infty} x_n = x \quad \eqdef \quad \forall V \in \cV_x,\ \exists N \in \bN,\ \forall n \geq N,\ x_n \in V.
\end{equation}
	
\end{definition}

\begin{proposition}[Caractérisation séquentielle de l'adhérence]
\label{prop:adhSuite}
\leancheck{prop:adhSuite}%
	Soit $A$ une partie d'un espace topologique à \emph{base dénombrable de voisinages}, alors 
	$$
		\overline A = \{ \lim_{n \to \infty} x_n \mid (x_n)_n \in A^\bN \text{ est une suite convergente} \}.
	$$
\end{proposition}

\begin{proof}
	Preuve de l'inclusion $\subset$ :
	Soit $x \in \overline A$, et soit $(V_n)_{n \in \bN}$ une base de voisinages de $x$, que l'on peut supposer décroissante pour l'inclusion. D'après la \cref{prop:adh_vois} il existe $x_n \in A \cap V_n$. La suite $(x_n)_{n \in \bN}$ converge alors vers $x$. Preuve de l'inclusion $\supset$ : découle de la caractérisation de l'adhérence de la \cref{prop:adh_vois}. (Ce sens n'utilise pas la base dénombrable).
\end{proof}

\begin{proposition}[Convergence d'une sous suite vers toute valeur d'adhérence d'une suite]
\label{prop:adhSousSuite}
\leancheck{prop:adhSousSuite}%
	L'ensemble des valeurs d'adhérence d'une suite $(x_n)_{n \in \bN}$ est défini par 
%	Soit $(x_n)_{n \in \bN}$ une suite à valeurs dans un espace topologique à \emph{base dénombrable de voisinages}. L'ensemble de ses valeurs d'adhérence est 
	\begin{equation}
	\label{eqdef:adherence_suite}		
		\Adh((x_n)_{n \in \bN}) := \bigcap_{N \geq 0} \overline {\{x_n \mid n \geq N\}}.
\leancheck{eqdef:adherence_suite}%
	\end{equation}
	Si l'espace topologique est à base dénombrable de voisinages, alors l'adhérence d'une suite est constituée des limites de suites extraites
	\begin{equation}
		\Adh((x_n)_{n \in \bN}) = \{x \in X \mid \exists \vp : \bN\to \bN, \text{ strictement croissante}, \lim_{n \to \infty} x_{\vp(n)} = x\}.
	\end{equation}	
%	pour tout $x\in A := \Adh((x_n)_{n \in \bN})$ il existe $\vp : \bN\to \bN$ strictement croissante telle que $x_{\vp(n)} \to x$ lorsque $n \to \infty$. 
\end{proposition}

\begin{proof}
	Soit $x \in \Adh((x_n)_{n \in \bN})$, et soit $(V_n)_{n \in \bN}$ une base de voisinages de $x$ décroissante pour l'inclusion. On choisit par récurrence $x_{\vp(n)} \in \{x_m \mid m> \vp(n-1)\} \cap V_n$, avec la convention $\vp(-1):=-1$; cette extractrice satisfait aux conditions. 
	
	Réciproquement, s'il existe une extractrice $\vp$ telle que $x_{\vp(n)} \to x$ lorsque $n \to \infty$, alors tout voisinage $V \in \cV_x$ contient tous les termes $(x_{\vp(n)})_{n \geq N(V)}$ à partir d'un certain rang $N(V)$. Ainsi $V$ intersecte $\{x_n\mid n \geq N\}$, quel que soit l'entier $N\in \bN$, donc $x\in \overline{\{x_n\mid n \geq N\}}$ ce qui conclut. (Ce sens n'utilise pas la base dénombrable.)
\end{proof}

\begin{remark}[Séparabilité d'une partie]
\label{rem:separable_partie}
\leancheck{rem:separable_partie}%
	Si un espace métrique est à base dénombrable d'ouverts, alors ses parties le sont aussi (prendre l'intersection des ouverts de la base avec la partie concernée), et il est séparable (prendre un point arbitraire dans chaque ouvert de la base). En particulier, toute partie d'un espace métrique séparable est séparable, voir la \cref{prop:baseVoisMetric}.
\end{remark}


Finalement, expliquons le terme ``filtre'' en entête de la \cref{def:filtre_vois}.

\begin{remark}[Axiomes de Filtre \facultatif]
\label{rem:axiomFiltre}
\leancheck{rem:axiomFiltre}%
Un filtre d'un ensemble ordonné $(\cP, \preceq)$ est une partie $\cV \subset \cP$, non-vide, et satisfaisant aux axiomes: 
\begin{itemize}
	\item (Upper closure) Si $V \in \cV$ et $W \in \cP$ satisfont $V \preceq W$, alors $W \in \cV$.
	\item (Downward directed) Si $V_0,V_1 \in \cV$, alors il existe $W \in \cV$ tel que $W\preceq V_0$ et $W \preceq V_1$.
\end{itemize}
Le filtre $\cV = \cP$ est dit trivial. 
L'ensemble des voisinages $\cV_x$ d'un point $x\in X$ d'un espace topologique $(X, \bU)$ est un filtre de $(\cP(X),\subset)$. 
%Réciproquement, on se donner pour tout $x \in X$ un filtre $\cV_x$ de parties contenant $x$, et définir la topologie $\bU := \{U \subset X\mid \forall x \in U,\ U \in \cV_x\}$ (un ouvert est un voisinage de chacun de ses points). Attention, pour la consistance il faut vérifier que $\cV_x$ admet une base formée d'ouverts c
%Attention, pour que $\cV_x$ soit bien la base de voisinages de l'espace topologique $(X,\bU)$, il faut vérifier que $\cV_x\cap \bU$ est une base de voisinages de $X$, ce n'est pas automatique. 
\end{remark}



\subsubsection{Séparation}
\label{subsubsec:separation}

\begin{definition}[Axiome de séparation de Hausdorff, T2, espace séparé]
\label{def:separe}
\leancheck{def:separe}%
Un espace topologique $(X, \bU)$ est dit séparé si deux points distincts admettent toujours des voisinages disjoints : 
\begin{equation}
	\forall x\in X,\ \forall y \in X \sm \{x\},\ \exists U \in \cV_x,\ \exists V \in \cV_y,\ U \cap V = \emptyset. 
\end{equation}
\end{definition}

\begin{remark}[Unicité de la limite]
\label{rem:uniqueLim}
\leancheck{rem:uniqueLim}%
Dans un espace séparé, toute suite admet au plus une limite.	
\end{remark}

La plupart des espaces rencontrés en analyse et analyse fonctionnelle, et donc dans ce cours, sont T2. À l'inverse, la topologie de Zariski, fondamentale en algèbre, satisfait seulement la propriété T1 plus faible.

\begin{remark}[Axiome T1 de Fréchet \facultatif] % Dit d'accessibilité
\label{rem:T1}
\leancheck{rem:T1}%
Un espace topologique $(X, \bU)$ satisfait l'axiome T1 si étant donné deux points distincts, chacun admet un voisinage ne contenant pas l'autre. 
Cet axiome ne \emph{garantit pas l'unicité} de la limite. 
\end{remark}

\begin{exemple}[Espaces T1 mais pas T2 \facultatif]
\label{ex:T1nonT2}
\leancheck{ex:T1nonT2}
\begin{itemize}
	\item (Topologie cofinie) Les entiers naturels, muni de la topologie dont les fermés sont les parties finies de $\bN$, ainsi que $\bN$ lui-même. 
	\item (Topologie de Zariski) L'ensemble $\bC^d$, où $d \geq 1$, muni de la topologie dont les fermés sont les ensembles algébriques (i.e.\ les zéros communs d'une famille de polynômes). Formellement, un fermé de Zariski $F \subset \bC^d$ est un ensemble de la forme
	\begin{equation}
		F = \{x \in \bC^d \mid P_1(x) = \cdots = P_I(x) = 0\},
	\end{equation}
	où $I\geq 0$ est un entier et où $P_1,\cdots,P_I \in \bC[X_1,\cdots,X_d]$. 
\end{itemize}
\end{exemple}

\begin{proof}[Preuve que les fermés de Zariski satisfont aux axiomes de la \cref{rem:fermes}]
L'ensemble vide s'obtient en choisissant $I=1$ et $P_1=1$, et l'ensemble total en posant $I=0$. 
Si le fermé de Zariski $F$ est associé aux polynômes $(P_i)_{1 \leq i \leq I}$, et $G$ aux polynômes $(Q_j)_{1 \leq j \leq J}$, alors $F \cup G$ est associé à la famille double $(P_i Q_j)^{1 \leq i \leq I}_{1 \leq j \leq J}$. Si $F_\alpha$ est associé aux polynômes $(P^\alpha_i)_{1 \leq i \leq I(\alpha)}$, pour tout paramètre $\alpha \in A$, alors $\cap_{\alpha\in A} F_\alpha$ est associé à une famille génératrice finie quelconque de l'idéal engendré par $\{P^\alpha_i \mid \alpha \in A,\, 1 \leq i \leq I(\alpha)\}$ dans l'anneau $\bC[X_1,\cdots,X_d]$, qui est Noetherien d'après le théorème de la base de Hilbert.
\end{proof}


\subsubsection{Continuité}
\label{subsubsec:continuite}

\begin{definition}[Continuité en un point]
\label{def:continuite}
\leancheck{def:continuite}%
	Soient $(X,\bU)$, $(Y,\bV)$, des espaces topologiques. Une application $f:X \to Y$ est continue en $x \in X$ si l'image réciproque de tout voisinage de $f(x)$ est un voisinage de $x$. 
	\begin{equation}
	\label{eq:cont_at_x}
		\forall V \in \cV_{f(x)},\ f^{-1}(V) \in \cV_x. 
	\end{equation}
	On dit que $f$ est continue (sur $X$), si elle est continue en tout point de $X$. 
\end{definition}

On peut reformuler \eqref{eq:cont_at_x} en termes d'image directe:
\begin{equation}
\label{eq:cont_at_x_direct}
	\forall V \in \cV_{f(x)},\ \exists U \in \cV_x,\ f(U) \subset V,
\leancheck{eq:cont_at_x_direct}%
\end{equation}
c'est-à-dire que tout voisinage de $f(x)$ contient l'image par $f$ d'un voisinage de $x$.
Si $f$ est continue en $x$, et $g$ est continue en $f(x)$, alors $g \circ f$ est continue en $x$. 

\notei{Supprimé la définition de l'équicontinuité ici, pas vraiment utile. (On la reprend dans les espaces métriques, avec sa variante uniforme.)
}

%Cette reformulation permet de définir l'équicontinuité d'une familles de fonctions:
%\begin{definition}[Équicontinuité]
%Une famille de fonctions $(f_i)_{i\in I}$ est dite équicontinue en $x \in X$ ssi
%\begin{equation}
%\label{eqdef:equicont}
%	\forall V \in \cV_{f(x)},\ \exists U \in \cV_x,\ \forall i \in I,\ f(U) \subset V.
%\end{equation}
%La famille est dite équicontinue ssi elle est équicontinue en tout point de $X$. 
%\end{definition}


%\todoi{La composée de fonctions continues est continue, et l'image par une fonction continue d'une suite convergente est une suite convergente. A localiser en $x\in X$.}

\begin{proposition}[Caractérisations de la continuité]
\label{prop:caract_cont}
\leancheck{prop:caract_cont}%
Soient $(X,\bU)$, $(Y,\bV)$, des espaces topologiques, et $f : X \to Y$. Les propriétés suivantes sont équivalentes:
\begin{enumerate}[(i)]
	\item $f$ est continue (en tout point)
	\item L'image réciproque de tout ouvert est un ouvert. 
	\item L'image réciproque de tout fermé est un fermé. 
	\item $\forall A \subset X$, $f(\overline A) \subset \overline {f(A)}$.  
\end{enumerate}
\end{proposition}

\begin{proof}
	(i) $\Leftrightarrow$ (ii) est classique. 
	(ii) $\Leftrightarrow$ (iii) s'obtient par passage au complémentaire. 
	(iii) $\Rightarrow$ (iv) : $f^{-1}(\overline {f(A)})$ est par (iii) un fermé, et cet ensemble contient $A$ donc aussi son adhérence $\overline A$. (iv) $\Rightarrow$ (iii) : rappelons que $f(f^{-1}(B)) \subset B$ pour toute partie $B\subset Y$. Si $B$ est fermé alors $\overline {f(f^{-1}(B))} \subset B$, puis en posant $A := f^{-1}(B)$ on déduit $\overline {f^{-1}(B)} = \overline A \subset f^{-1}(\overline {f(A)}) = f^{-1}(\overline {f(f^{-1}(B))}) \subset f^{-1}(B)$. % ce qui conclut. % et $f^{-1}(f(f^{-1}(B)))=f^{-1}(B)$
\end{proof}



\begin{remark}[Topologie rendant continue une famille de fonctions]
\label{rem:topoFunc}
\leancheck{rem:topoFunc}
Soit $X$ un ensemble, et pour tout $i \in I$ appartenant à un ensemble quelconque d'indices soit $f_i : X \to Y_i$ une fonction vers un espace topologique $(Y_i,\bU_i)$.
On peut considérer la topologie $\bU$ sur $X$ engendrée par les ensembles $\{f_i^{-1}(U_i) \mid i \in I,\ U_i \in \bU_i\}$.

C'est la topologie la moins fine rendant continue ces fonctions. Une suite $(x_n)_{n \in \bN}$ converge vers $x_*\in X$ pour $\bU$ si et seulement si $f_i(x_n) \to f_i(x_*)$ lorsque $n \to \infty$, pour tout $i \in I$.
\end{remark}

Nous terminons cette sous-section avec la notion de semi-continuité, qui intervient dans les problèmes d'optimisation, voir  l'\nref{sec:LegendreFenchel}.

\begin{definition}[Semi-continuité]
\label{def:semiContinuite}
\leancheck{def:semiContinuite}%
	Soit $(X,\bU)$ un espace topologique, et $f : X \to \overline \bR$. 
	\begin{itemize}
		\item $f$ est semi-continue inférieurement (s.c.i.) si $\{f \leq M\}$ est fermé pour tout $M \in \overline \bR$.
		\item $f$ est semi-continue supérieurement (s.c.s) si $\{f \geq M\}$ est fermé pour tout $M \in \overline \bR$. 
	\end{itemize}	
\end{definition}

Noter que $f$ est s.c.s.\ ssi $-f$ est s.c.i.
Soit $(f_i)_{i \in I}$, $f_i : X \to \overline \bR$, une famille de fonctions s.c.i.\ alors $f(x) := \sup_{i \in I} f_i(x)$ est s.c.i. De même, l'infimum d'une famille de fonctions s.c.s. est s.c.s.
Si $f,g : X \to ]-\infty,\infty]$ sont s.c.i, alors $f+g$ est s.c.i.\ par ouverture de l'ensemble suivant
\begin{equation}
\label{eq:somme_sci}
	\{f+g>M\} = \bigcup_{\alpha \in \bR} \big(\{f > M-\alpha\} \cap \{g > \alpha\}\big).
\end{equation}
De même, une somme de fonctions s.c.s.\ est s.c.s.


\begin{proposition}
\label{prop:surGrapheSCI}
\leancheck{prop:surGrapheSCI}
	Soit $(X,\bU)$ un espace topologique, et $f : X \to \overline \bR$. Alors $f$ est s.c.i.\ ssi son sur-graphe $\cG := \{(x,t) \in X \times \bR \mid f(x) \leq t\}$ est fermé. 
\end{proposition}

\begin{proof}
Si $\cG$ est fermé, alors $\{f \leq \lambda\} = \{x\in X \mid (x,\lambda) \in \cG\}$ aussi, pour tout $\lambda \in \bR$. Comme de plus $\{f \leq \infty\} = X$ et $\{f \leq -\infty\} = \cap_{\lambda \in \bR} \{f\leq \lambda\}$ sont fermés, on obtient que $f$ est s.c.i. La réciproque découle de l'égalité suivante, facilement vérifiée 
\begin{equation*}
	\{(x,t) \in X \times \bR \mid f(x) \leq t\} 
	= \bigcap_{\alpha \in \bR} \big(\{f \leq \alpha\} \times ]-\infty,\alpha]\big) \cup \big(X \times [\alpha, \infty[\big).
\end{equation*}
(Si $f(x)\leq t$, distinguer les cas $\alpha \leq t$ et $\alpha \geq t$. À l'inverse, si $f(x)>t$, considérer $\alpha \in ]t,f(x)[$.)
Si $f$ est s.c.i., alors $\cG$ est une intersection de fermés, donc fermé. 
\end{proof}

%\hiddenproof{
%Considérons $(x,t)$ tels que $f(x) \leq t$, et $\alpha\in \bR$. 
%Si $\alpha \leq t$ alors $(x,t) \in X \times [\alpha,\infty[$, et si $\alpha >t$ alors $(x,t) \in \{f \leq \alpha\} \times ]-\infty,\alpha]$).
%À l'inverse, supposons que $f(x) > t$, considérons $\alpha \in ]t,f(x)[$, de sorte que $x \notin \{f\leq \alpha\}$ et $t \notin [\alpha,\infty[$.
%}


\subsection{Espaces métriques}
\label{subsec:metric}

Les espaces métriques sont en général plus simples à manipuler que les espaces topologiques, grâce entre autres à l'analyse séquentielle. Ils permettent également d'introduire des notions quantitatives et renforcées de la notion de continuité : uniforme continuité, exposant de Hölder, constante de Lipschitz.

\begin{retenir}
Les axiomes de distance, la caractérisation de la continuité via un module de continuité, et le fait que l'on puisse choisir celui-ci croissant et lisse. La définition de l'uniforme continuité et ses cas particuliers de Lipschitz et Hölder.
(\Cref{def:distance,lem:contDistance,lem:regModuleCont,def:unifContinu})
\end{retenir}


\begin{definition}[Espace métrique]
\label{def:distance}
\leancheck{def:distance}%
	Un espace métrique est une paire $(X,d)$ où la fonction $d : X \times X \to \bR$ appelée \emph{distance} satisfait aux axiomes suivants:
	\begin{itemize}
		\item (Séparation) Pour tous $x,y \in X$, on a $d(x,y) \geq 0$, avec égalité ssi $x=y$. 
		\item (Symétrie) Pour tous $x,y \in X$, on a $d(x,y) = d(y,x)$. 
		\item (Inégalité triangulaire) Pour tous $x,y,z \in X$, on a $d(x,z) \leq d(x,y) + d(y,z)$. 
	\end{itemize}
\end{definition}

Si $x\in X$ et $A\subset X$ sont un élément et une partie non-vide d'un espace métrique $(X,d)$, on se permet de noter $d(x,A) := \inf\{d(x,y) \mid y \in A\}$.

\begin{definition}[Boules, topologie métrique]
\label{def:boule}
\leancheck{def:boule}%
	Soit $(X,d)$ un espace métrique. On définit pour tout $x \in X$ et $r > 0$ les boules ouverte et fermée:
	\begin{align}
		B(x,r) &:= \{y \in X\mid d(x,y) < r\}, &
		B'(x,r) &:= \{y\in X\mid d(x,y) \leq r\}.
	\end{align}
	On munit $X$ de la topologie dite métrique, dont les boules $\{B(x,r) \mid x \in X, r>0\}$ constituent une base d'ouverts.
\end{definition}

\hiddenproof{
(Vérifications des axiomes de la \nref{prop:baseOuverts})
\emph{Axiome de couverture (i).} Les boules recouvrent $X$ car chacune contient son centre. 
\emph{Axiome de stabilité (ii).} Soit $z \in B(x,r) \cap B(y,s)$, où $x,y \in X$, $r,s>0$. Notons $t := \min\{ r-d(x,z), s-d(y,z)\}$. Alors $B(z,t) \subset  B(x,r) \cap B(y,s)$. 
}

\begin{remark}[$\infty$-distance]
\label{rem:distInfinie}
\leancheck{rem:distInfinie}%
Une $\infty$-distance $d_\infty$ satisfait les mêmes axiomes qu'une distance, mais peut prendre la valeur $+\infty$. Cette généralisation est parfois naturelle ou pratique. Noter que $d := \min\{1,d_\infty\}$ est alors une distance définissant la même topologie que $d_\infty$. 
\end{remark}


Une partie d'un espace métrique est dite bornée si elle est incluse dans une boule. 
Sur toute partie $A \subset X$ d'un espace métrique on dispose de la distance induite $d_{|A \times A}$, qui correspond bien à la topologie induite. 
La proposition suivante assure que l'on peut utiliser les caractérisations séquentielles de l'adhérence et de la continuité dans un espace métrique, via les \cref{prop:adhSuite,prop:adhSousSuite}.

\begin{proposition}
\label{prop:baseVoisMetric}
\leancheck{prop:baseVoisMetric}%
	Soit $(X,d)$ un espace métrique. 
	Alors tout point $x \in X$ admet la base dénombrable de voisinages 
	\begin{equation}
	\label{eqdef:base_vois_metric}
		\cW_x := \{B(x,1/n) \mid n \in \bN^*\}.		
	\end{equation}
	Si $X$ est séparable, et si $(x_n)_{n \in \bN}$ est une suite dense, alors on a la base dénombrable d'ouverts 
	$$
		\{B(x_n,r) \mid n\in \bN, r \in \bQ_+^*\}.
	$$
	\notei{Alternativement, on pouvait se contenter de $r=2^{-m}$, $m \in \bN$.}
\end{proposition}

\begin{proof}
	\emph{Premier point :} soit $U\subset X$ un ouvert contenant un point $x\in X$. Par construction, il existe $y\in X$ et $r>0$ tels que $x \in B(y,r)$ et $B(y,r) \subset U$. En choisissant $n \in \bN^*$ tel que $1/n < r-d(x,y)$ on a donc $B(x,1/n) \subset B(y,r) \subset U$ comme souhaité.
	\emph{Second point :}
	soit $U\subset X$ un ouvert contenant un point $x\in X$. Par le premier point, il existe $n \in \bN^*$ tel que $B(x,1/n) \subset U$. Par densité, il existe $k\in \bN$ tel que $d(x_k,x) \leq 1/(3n)$. Ainsi $x \in B(x_k,2/(3n)) \subset B(x,1/n)\subset U$ comme souhaité.	
%	Nous passons sur le premier point que est trivial, et vérifions le second via la condition du \cref{corol:baseOuverts}. Soit $U$ un ouvert, de la topologie métrique, et soit $x \in U$. Par construction, il existe une boule ouverte telle que $x \in B(y,r) \subset U$, où $y \in X$ et $r>0$. 
%	Posons $\delta := r-d(x,y)$ et choisissons $n \in \bN$ tel que $d(x,x_n) < \delta/3$, et $r\in \bQ_+^*$ tel que $\delta/3\leq r \leq 2 \delta/3$. Alors $x\in B(x_n,r)\subset B(y,r) \subset U$ comme souhaité, ce qui conclut. 
%	Axiome de recouvrement (i) : $X = \cup_{r \in \bN^*} B(x_0,r)$. Axiome de stablilité (ii) : considérons $x\in B(x_m,r) \cap B(x_n,s)$ où $m,n \in \bN$ et $r,s \in \bQ_+^*$. Alors on peut choisir un rationnel $0<\ve\leq \min\{r-\|x-x_m\|, s-\|x-x_n\|\}$, de sorte que $B(x,\ve) \subset B(x_m,r) \cap B(x_n,s)$. Par densité de la suite, il existe $k \in \bN$ tel que $x_k \in B(x,\ve/2)$, et donc $x \in B(x_k, \ve/2) \subset B(x,\ve) \subset B(x_m,r) \cap B(x_n,s)$ de sorte que $B(x_k, \ve/2)$ satisfait la condition recherchée. Par un argument similaire, et le critère (iii) du \cref{corol:baseOuverts}, la topologie engendrée par ces ouverts est bien celle de $d$.
\end{proof}

La notion de module de continuité permet, comme annoncé, de reformuler plus simplement la continuité d'une fonction, et d'en introduire des variante renforcées, voir les \cref{lem:contDistance,def:unifContinu}.

\begin{definition}[Module de continuité]
\label{def:moduleCont}
\leancheck{def:moduleCont}%
Un module de continuité est une application $\omega : \bR_+ \to [0,\infty]$, telle que $\omega(s) \to 0$ lorsque $s \to 0$.
On dit qu'il est sous-additif si $\omega(s+t) \leq \omega(s) + \omega(t)$ pour tous $s,t \geq 0$. 
\end{definition}
Nous avons fait dans la \cref{def:moduleCont} des hypothèses minimales concernant la régularité de $\omega$. Sans nuire à la généralité, on pourrait toutefois les renforcer, comme le montre le  \cref{lem:regModuleCont} ci-dessous. 

\begin{remark}
\label{rem:dist_of_mod_cont}
\leancheck{rem:dist_of_mod_cont}%
	Un module de continuité sous-additif est toujours à valeurs finies. S'il est de plus croissant et non-identiquement nul, et si $d$ est une distance sur un espace métrique, alors $\omega \circ d$ est une distance définissant la même topologie.
\end{remark}

\hiddenproof{
Finitude : fixons $s>0$, et notons que $\omega(s/n) \to 0$ lorsque $n \to \infty$, donc il existe $n\in \bN^*$ tel que $\omega(s/n) < \infty$. Alors $\omega(s) \leq n\omega(s/n) < \infty$ par sous-additivité.

Résultat intermédiaire : si $\omega$ est non-identiquement nul, sous-additif et croissant, alors il est positif partout sur $]0,\infty[$. En effet, si $\omega(s_0)>0$ (hypothèse non-identiquement nul), alors $\omega(s_0/n) \geq \omega(s_0)/n>0$ par sous-additivité pour tout $n\in \bN^*$, et on conclut par croissance de $\omega$. 

Propriétés de distance: l'inégalité triangulaire découle de la sous-additivité et de la croissance de $\omega$, la symétrie est claire, et la séparation découle du résultat intermédiaire.
}

\begin{lemma}[Régularisation d'un module de continuité]
\label{lem:regModuleCont}
	Soit $\omega_0 : \bR_+ \to \bR_+$ un module de continuité localement borné. Alors il existe un module de continuité $\omega_*$ qui le majore, est croissant, et est $C^\infty$ sur $\bR_+^*$. 
	
	Si $\omega_0$ est à croissance sous linéaire à l'infini : $\omega_0(t) \leq \alpha + \beta t$, $\forall t \in \bR_+$, pour des constantes $\alpha,\beta \in \bR_+$, alors on peut construire $\omega_*$ qui est de plus sous-additif : $\omega_*(s+t) \leq \omega_*(s) + \omega_*(t)$ pour tous $s,t\geq 0$.
\end{lemma}

\begin{proof}
	L'application $\omega_1 : \bR_+ \to \bR_+$ définie par $\omega_1(t) := \sup\{\omega_0(s) \mid 0 \leq s \leq t\}$ tend vers $0$ en $0$, majore $\omega_0$, et est croissante donc mesurable. 
	
	Si $\omega_0$ est à croissance sous linéaire à l'infini, alors c'est aussi le cas de $\omega_1$. On pose dans ce cas 
	\begin{equation}
		\omega_2(s) := \sup_{0 < \lambda \leq 1} \lambda \omega_1(s/\lambda) 
		= \sup_{\mu \geq s} \tfrac s \mu \omega_1(\mu). 
	\end{equation}
	L'égalité s'obtient par changement de variable $\mu = s/\lambda$. 
	La première expression montre que $\omega_2$ est croissant et majore $\omega_1$, et on déduit de la seconde qu'il est sous-additif (ce qui implique également sa continuité). 
%	Par hypothèse, on a $\omega_1(\mu) \leq o(1) + \beta\mu$, où le premier terme tend vers $0$ lorsque $\mu \to 0$, donc $W_1(s) := \sup_{\mu \geq s}\omega_1(\mu)/\mu = o(1/s)+\beta$, donc $\omega_2(s) = \beta s + o(1)$, de sorte que $\omega_2(s) \to 0$ lorsque $s\to 0$ comme souhaité.	
	Supposons $\lambda_n \in ]0,1]$ et $s_n \to 0$ lorsque $n\to \infty$, et montrons que $\lambda_n \omega_1(s_n/\lambda_n) \to 0$, ce qui assure que $\lim_{s \to 0} \omega_2(s) = 0$; quitte à extraire, on peut distinguer (i) le cas où $\lambda_n \leq \sqrt{s_n}$, et donc $\lambda_n \omega_1(s_n/\lambda_n) \leq \lambda_n (\alpha + \beta s_n/\lambda_n) = \alpha \lambda_n + \beta s_n \to 0$, et (ii) le cas où $\lambda_n \geq \sqrt{s_n}$, et donc $\lambda_n \omega_1(s_n/\lambda_n) \leq \omega_1(\sqrt {s_n}) \to 0$ lorsque $n \to \infty$. 
	On note dans la suite $\omega_\times = \omega_2$ dans le cas où $\omega_0$ est à croissance sous-linéaire, ou $\omega_\times = \omega_1$ dans le cas contraire.
	
%\todoi{La preuve de $\omega_2 \to 0$ en $0$ pourrait être rendue plus limpide.}
	
%Considérons deux suites $0<s_n \leq \mu_n$ telles que $s_n \to 0$. Si $\mu_n$ est bornée inférieurement, alors $\omega_1(\mu_n)/\mu_n \leq \alpha/\mu_n+ \beta$ est borné supérieurement, et donc $s_n \omega_1(\mu_n)/\mu_n \to 0$. Si $\mu_n \to 0$, alors $(s_n/\mu_n) \omega_1(\mu_n) \leq \omega_1(\mu_n) \to 0$. 
	
	Considérons maintenant une fonction bosse $\vp \in C^\infty(\bR_+, \bR_+)$ telle que $\supp(\vp) \subset [1,2]$ et $\int_0^\infty \vp = 1$, voir l'\cref{sec:bosse} pour une construction. 
	Posons pour tout $t>0$
	\begin{equation}
	\label{eq:reg_mod_cont}
		\omega_*(t) := \int_0^\infty \omega_\times(t s) \vp(s) \diff s = \frac 1 t \int_0^\infty \omega_\times(s) \vp(s/t) \diff s.
	\end{equation}
	La première expression intégrale montre que $\omega_\times (s) \leq \omega_*(s) \leq \omega_\times(2s)$, donc que $\omega_*$ est un module de continuité majorant $\omega_0$, et que $\omega_*$ est croissant, et sous-additif si $\omega_\times$ l'est. La seconde expression, qui s'obtient par un changement de variable, montre  par dérivation sous le signe intégral que $\omega_*$ est $C^\infty$ sur $\bR_+^*$, comme annoncé, ce qui conclut. 
\end{proof}

\begin{lemma}[Continuité entre espaces métriques]
\label{lem:contDistance}
\leancheck{lem:contDistance}%
Soit $f:X \to Y$ une application entre espaces métriques $(X,d_X)$ et $(Y,d_Y)$, et soit $x \in X$. 
Les conditions suivantes sont équivalentes
\begin{enumerate}[(i)]
	\item $f$ est continue en $x$.
	\item $\forall \ve>0,\ \exists \delta>0,\ \forall x' \in X,\ d_X(x,x') < \delta \Rightarrow d_Y(f(x),f(x')) < \ve$.
	\item $\forall \ve>0,\ \exists \delta>0,\ \forall x' \in X,\ d_X(x,x') \leq \delta \Rightarrow d_Y(f(x),f(x')) \leq \ve$.
	\item Il existe un module de continuité $\omega_x$ tel que pour tout $x'\in X$,
\begin{equation}
\label{eq:mod_cont_x}
	d_Y(f(x),f(x')) \leq \omega_x(d_X(x,x')).
\end{equation}
	\item Pour toute suite $(x_n)_{n \in \bN} \in X^\bN$ telle que $x_n \to x$ lorsque $n \to \infty$, on a $f(x_n) \to f(x)$.
\end{enumerate}
\end{lemma}

\begin{proof}
	L'équivalence (i) $\Leftrightarrow$ (ii) est classique et découle de \eqref{eq:cont_at_x_direct} et de la base de voisinages \eqref{eqdef:base_vois_metric}.
	L'équivalence (ii) $\Leftrightarrow$ (iii) est triviale. 
	L'implication (iii) $\Rightarrow$ (iv) s'obtient en posant $\omega_x(\delta) := \sup \{d_Y(f(x),f(x')) \mid x' \in B'(x,\delta)\}$, et réciproquement (iv) $\Rightarrow$ (iii) en utilisant que $\omega_x$ tend vers $0$ en $0$. 
	L'implication (iv) $\Rightarrow$ (v) est claire. Montrons que non (iii) $\Rightarrow$ non (v) : en ce cas, par négation des quantificateurs, il existe $\ve>0$ et pour tout $n\geq 1$ un point $x_n \in X$ tel que $d(x,x_n)\leq 1/n$ et $d(f(x),f(x_n)) \geq \ve$, de sorte que $x_n \to x$ et $f(x_n) \not \to f(x)$ lorsque $n\to \infty$.
\end{proof}

Dans les espaces métriques (à la différence des   espaces topologiques), il est possible de considérer des modules de continuités uniformes, ou ayant une expression particulière. 

\begin{definition}[Uniforme continuité]
\label{def:unifContinu}
\leancheck{def:unifContinu}%
	Une application $f : X \to Y$ entre espaces métriques $(X,d_X)$ et $(Y,d_Y)$ est dite:
	\begin{itemize}
		\item \emph{Uniformément continue} s'il existe un module de continuité $\omega$ tel que 
		\begin{equation}
		\label{eq:mod_cont_unif}			
			\forall x,x' \in X,\ d_Y(f(x),f(x')) \leq \omega(d_X(x,x')).
		\end{equation}
		De manière équivalente, $\forall \ve>0$, $\exists\delta>0$, $\forall x,x' \in X$, $(d_X(x,x')\leq \delta \Rightarrow d_Y(f(x),f(x')) \leq \ve)$.
		\item \emph{$\alpha$-Hölderienne}, où $\alpha \in ]0,1[$, si il existe une constante $K\in \bR$ telle que 
		$$
			\forall x,x' \in X,\ d_Y(f(x),f(x')) \leq K d_X(x,x')^\alpha.
		$$
		\item \emph{($K$-)Lipschitzienne}, si il existe une constante $K \in \bR$ telle que 
		$$
			\forall x,x' \in X,\ d_Y(f(x),f(x')) \leq K d_X(x,x').
		$$
		Une application $K$-Lipschitzienne pour un certain $K<1$ est dite contractante. 
	\end{itemize}
	Ces ensembles de fonctions sont respectivement notés $\mathrm{UC}(X,Y)$, $C^\alpha(X,Y)$ et $\Lip(X,Y)$.
\end{definition}


\paragraph{Convergence uniforme et équicontinuité.}

Distinguons d'abord deux distances sur un produit d'ensembles, qui sont a-priori différentes dès que le nombre de termes est infini. 

\begin{proposition}[Distance associée à la convergence simple]
\label{prop:distCvSimple}
\leancheck{prop:distCvSimple}%
%\label{prop:distProduit}
	Soit $(X_n,d_n)_{n \in \bN}$ une famille dénombrable d'espaces métriques. La topologie produit sur $\Pi_{n \in \bN} X_n$ est métrisable, pour la distance
	\begin{equation*}
		d( (x_n)_{n \in \bN}, (y_n)_{n \in \bN} ) := \max_{n \in \bN} \min\{2^{-n},\ d_n(x_n,y_n)\}.
	\end{equation*}
\end{proposition}

\begin{definition}[Distance associée à la convergence uniforme]
\label{def:dist_unif}
\leancheck{def:dist_unif}%
	Soit $(X_i,d_i)_{i \in I}$ une famille d'espaces métriques. La topologie de la convergence uniforme sur le produit $\prod_{i \in I} X_i$ est donnée par la distance $\min\{1,d_\infty\}$ (voir la \cref{rem:distInfinie}) où:
	\begin{equation*}
		d_\infty( (x_i)_{i \in I}, (y_i)_{i \in I} ) := \sup_{i \in I} d_i(x_i,y_i).
	\end{equation*}
\end{definition}

On dispose de la notion d'équi-continuité, pour une famille de fonctions, qui se reformule en termes de convergence uniforme.

\begin{remark}[Équicontinuité]
\label{rem:equicont}
\leancheck{rem:equicont}%
Une famille $(f_i)_{i \in I}$ de fonctions entre espaces métriques $(X,d_X)$ et $(Y,d_Y)$ est dite:
\begin{enumerate}[(i)]
	\item équicontinue en $x \in X$ s'il existe un module de continuité $\omega_x$ tel que : 
	\begin{equation*}
	\forall x' \in X,\ \forall i \in I,\ d_Y(f_i(x),f_i(x')) \leq \omega_x(d_X(x,x')). 		
	\end{equation*}
	\item uniformément équicontinue s'il existe un module de continuité $\omega$ tel que :
	\begin{equation*}
	\forall x,x' \in X,\ \forall i \in I,\ d_Y(f_i(x),f_i(x')) \leq \omega(d_X(x,x')). 		
	\end{equation*}
\end{enumerate}
De manière équivalente $F : X \to Y^I : x \mapsto (f_i(x))_{i \in I}$ est continue en $x$ (resp.\ uniformément continue) pour la topologie de la convergence uniforme sur $Y^I = \Pi_{i \in I} Y$, voir la \cref{def:dist_unif}. 
\end{remark}


\subsection{Espaces vectoriels normés}
\label{subsec:evn}

Les espaces vectoriels normés (evn) possèdent une distance compatible avec leur structure linéaire, et constituent un cadre privilégié pour l'analyse linéaire; cependant leur généralisation les evtlcs, présentés dans la sous-section suivante, est nécessaire pour décrire certains espaces fonctionnels. 
Dans la suite, on note $\bK$ le corps $\bR$ ou $\bC$. 


\begin{retenir}
Les axiomes de norme et de semi-norme, la norme d'opérateur, les normes usuelles sur les espaces $C_b^\alpha(X)$, où $0 \leq \alpha < 1$, et $L^p(X,\mu)$, $1 \leq p \leq \infty$. 
(\Cref{def:norme,def:normeOperateur,def:normeHolder,def:normeLp}.)
\end{retenir}


\begin{definition}[Norme]
\label{def:norme}
\leancheck{def:norme}%
	Une norme sur un $\bK$-espace vectoriel $E$ est une application $\|\cdot\| : E \to \bR_+$ satisfaisant aux axiomes suivants : 
	\begin{itemize}
		\item (Séparation) $\forall x \in E\sm \{0\}, \|x\|>0$.
		\item (Homogénéité) $\forall x\in E, \forall \lambda \in \bK, \|\lambda x\| = |\lambda| \|x\|$.
		\item (Inégalité triangulaire) $\forall x,y \in E, \|x+y\|\leq \|x\| + \|y\|$. 
	\end{itemize}
	Une semi-norme est une application $|\cdot| : E \to \bR_+$ satisfaisant aux axiomes d'homogénéité et d'inégalité triangulaire. 
\end{definition}

Un $\bK$-espace vectoriel normé (evn) est muni de la distance $d(x,y) := \|x-y\|$ et de la topologie associée. 

\begin{definition}[Evn des applications linéaires continues]
\label{def:normeOperateur}
\leancheck{def:normeOperateur}%
	Soient $E,F$ des $\bK$-evn. On note $L(E,F)$ l'ensemble des applications $\bK$-linéaires \emph{continues} $u : E \to F$. On rappelle que $L(E,F)$ est un $\bK$-evn, muni de la norme 
	\begin{equation}
		\vertiii{u} := \sup\{ \|u (x)\|_F \mid \|x\|_E\leq 1\},
	\end{equation}
	qui est aussi la constante de Lipschitz de $u$. 
	En particulier, le dual $E^* := L(E,\bK)$ d'un $\bK$-evn est aussi un $\bK$-evn. 
\end{definition}

\subsubsection{Exemples d'espaces fonctionnels}

\begin{definition}[Fonctions Lipschitziennes, Hölderiennes]
\label{def:normeHolder}
\leancheck{def:normeHolder}%
	L'ensemble $C^0_b(X,\bK)$ des fonctions continues bornées sur un espace topologique $X$ est un $\bK$-evn muni de la norme 
		\begin{equation}
			\|f\|_\infty := \sup_{x\in X} |f(x)|. 
		\end{equation}
L'ensemble $C^\alpha_b(X,\bK)$ des fonctions $\alpha$-Hölderiennes bornées sur un espace métrique $(X,d)$, où $\alpha \in ]0,1[$, est un evn muni de la norme 
		\begin{align}
			\|f\|_{C^\alpha} &:= \|f\|_\infty + |f|_{C^\alpha}, &
			\text{où}\quad
			|f|_{C^\alpha} &:= \sup_{x\neq y} \frac{|f(x)-f(y)|}{d(x,y)^\alpha}
		\end{align}
De même les fonctions Lipschitziennes bornées, avec la norme $\|f\|_\infty + \Lip(f)$, où 
\begin{align*}
	\Lip(f) := \sup_{x \neq y} \frac{|f(x)-f(y)|}{d(x,y)}.
\end{align*}
\end{definition}


\begin{definition}[Espace $C^n_b(\Omega)$ des fonctions ayant $n$ dérivées continues bornées.]
\label{def:normeCnb}
\leancheck{def:normeCnb}%
	Soit $\Omega \subset \bR^d$ un ouvert, et soit $n \in \bN$.
	On note $C^n_b(\Omega)$ l'ensemble des fonctions $n$ fois continument différentiables, dont les dérivées jusqu'à l'ordre $n$ sont bornées, muni de la norme
	\begin{equation}
		\|f\|_{C^n_b} := \sum_{|\alpha|_1\leq n} \| \partial_\alpha f\|_\infty,
	\end{equation}
	où $|\alpha|_1 := \alpha_1+ \cdots+\alpha_d$ pour tout $\alpha = (\alpha_1,\cdots,\alpha_d) \in \bN^d$. 
	On note $C^n_b(\overline \Omega)$ le sous ensemble de $C^n_b(\Omega)$ formé des fonctions dont les dérivées jusqu'à l'ordre $n$ s'étendent continument à $\overline \Omega$. 
\end{definition}


La notion de dérivée au bord d'un domaine pose des questions non-triviales. Le théorème d'extension de Seeley affirme que si $\Omega$ a un bord lisse, alors les éléments de $C^n(\overline \Omega)$ admettent des extensions $C^n$ à l'espace entier $\bR^d$. 
Des conditions de compatibilité plus restrictives sont nécessaires dans le cas d'un ouvert général, traité par le théorème d'extension de Whitney. 
\todoi{ref ?}

Étant donné un espace mesuré $(X,\mu)$, on dit qu'une propriété des points de $X$ a lieu $\mu$ presque partout ($\mu$-p.p.) si l'ensemble des points qui ne la satisfont pas est de mesure nulle. 

Désignons par $L^*(X,\mu)$ (notation non standard) l'ensemble des fonctions mesurables de $X$ dans $\bK$, quotienté par la relation d'équivalence $f\sim g \eqdef f=g$ $\mu$-p.p. 
En pratique, on manipule les éléments de $L^*(X,\mu)$ comme des fonctions usuelles, en choisissant implicitement un représentant et en vérifiant que les résultats obtenus ne dépendent pas de ce choix. % en prenant seulement garde de ne pas considérer éventuellement des propriétés qui dépendent d'ensembles de mesure nulle. 

\begin{definition}[Espaces de Lebesgue $L^p(X,\mu)$]
\label{def:normeLp}
\leancheck{def:normeLp}%
	Soit $(X,\mu)$ un espace mesuré, et soit $f \in L^*(X,\mu)$. On note pour $p\in [1,\infty[$
	\begin{align}
		\|f\|_p &:= \Big(\int_X |f|^p \diff \mu\Big)^\frac 1 p, &
		\|f\|_\infty &:= \inf\{\lambda \geq 0 \mid |f| \leq \lambda, \ \mu\text{-p.p.}\}.
	\end{align}
	On définit $L^p(X,\mu) := \{f \in L^*(X,\mu) \mid \|f\|_p < \infty\}$, qui est un evn, pour tout $p \in [1,\infty]$. Si $X$ est fini ou dénombrable, on note parfois cet espace $l^p(X,\mu)$.
\end{definition}

\begin{proof}
L'homogénéité et la séparation sont claires.
Le seul point difficile est l'inégalité triangulaire dans le cas $p\in [1,\infty[$, appelée inégalité Minkowski.
Étant donnés $f,g \in L^p(X,\mu)$, montrons que $\|f+g\|_p \leq \|f\|_p + \|g\|_p$. Par homogénéité, et en éliminant les cas triviaux, on peut supposer que $\|f\|_p + \|g\|_p = 1$ et $0 < \|g\|_p < 1$, et poser $\lambda := \|g\|_p$. 

Posons $F := f/(1-\lambda)$ et $G := g/\lambda$. Alors par convexité de $s\in \bR_+ \mapsto s^p$ on obtient
\begin{equation}
	|(1-\lambda) F(x)+\lambda G(x)|^p \leq (1-\lambda) |F(x)|^p+\lambda |G(x)|^p,
\end{equation}
pour tout $x \in X$, d'où l'on conclut par intégration
\begin{equation*}
	\|f+g\|_p^p = \int_X |(1-\lambda) F(x)+\lambda G(x)|^p \diff \mu \leq (1-\lambda) \int_X |F|^p \diff \mu + \lambda \int_X |G|^p \diff \mu = 1.
\qedhere
\end{equation*}
\end{proof}

\subsection{Espaces vectoriels topologiques localement convexes} 
\label{subsec:evtlc}

Les evtlcs permettent de traiter des espaces fonctionnels plus généraux que les evn, et notamment les espaces de fonction test en théorie des distributions, tout en gardant une structure maniable : définie par des semi-normes, et souvent une distance associée. % - bénéficiant selon les cas de résultats de complétude et de compacité. 


\begin{retenir}
La définition d'un evtlc, le critère de séparation, de continuité des formes linéaires, la distance associée dans le cas d'une famille dénombrable de semi-normes. La topologie des espaces de fonction test $\cS(\bR^n)$ et $\cD_K(\Omega)$, la définition d'une distribution. La définition de la topologie faible, *faible, et la métrisabilité de cette dernière restreinte à un borné. 
(\Cref{def:evtlc,lem:evtlc_separe,eq:cont_lin_evtlc_evn,prop:distanceEvtlcs,def:cinf_bornee,def:classeSchwartz,def:distribution,def:topoFaible,prop:bouleMetrisable}.)
\end{retenir}



\begin{definition}[evtlc]
\label{def:evtlc}
\leancheck{def:evtlc}%
%	Une semi-norme sur un espace vectoriel $E$ est une application $|\cdot| : E \to \bR_+$ satisfaisant aux axiomes d'homogénéité et d'inégalité triangulaire d'une norme (mais pas forcément l'axiome de séparation).	
	Un \emph{espace vectoriel topologique localement convexe} est un espace vectoriel $E$ muni d'une famille de semi-normes $(|\cdot|_i)_{i \in I}$, où $I$ est un ensemble quelconque d'indices. On lui associe la topologie la plus grossière qui rend continues les applications $y \in E \mapsto |y-x|_i$, pour tout $x \in E$ et tout $i \in I$; une base de voisinages d'un point $x \in E$ est formée par les ensembles 
\begin{equation}	
\label{eq:evtlc_vois}
	B_{I_0}(x,\ve) := \{y \in E \mid \forall i \in I_0,\ |x-y|_i < \ve\},
\leancheck{eq:evtlc_vois}%
\end{equation}
pour tout $\ve>0$ et toute partie finie $I_0 \subset I$.
\end{definition}

Une suite $x_n$ d'un evtlc $(E,(|\cdot|_i)_{i \in I})$ converge vers un point $x$ ssi $|x-x_n|_i \to 0$ lorsque $n \to \infty$ pour tout $i \in I$.
Tout evn est de manière triviale un evtlc, dont la famille de semi-normes est réduite à un élément. 
Noter que la topologie sur un evtlc n'est pas automatiquement séparée; cela doit être justifié, typiquement par le lemme suivant. On abrège \emph{evtlc séparé} en evtlcs.

\begin{lemma}[Séparation dans les evtlc]
\label{lem:evtlc_separe}
\leancheck{lem:evtlc_separe}%
	Un evtlc $E$ muni d'une famille de semi-normes $(|\cdot|_i)_{i \in I}$, est séparé ssi pour tout $x \in E$, $(\forall i \in I, |x|_i = 0 ) \Rightarrow x=0$.
\end{lemma}

\begin{proof}
	Si la condition n'est pas satisfaite, alors il existe un point $x\in E\sm\{0\}$ tel que $|x|_i=0$ pour tout $i \in I$, donc appartenant à tous les voisinages de $0$ au vu de leur définition \eqref{eq:evtlc_vois}, donc l'espace n'est pas séparé.

	Réciproquement, supposons la condition satisfaite, et soient $x,y \in E$, $x \neq y$, dont on veut construire des voisinages disjoints. Comme $y-x \neq 0$, il existe $i \in I$ tel que $\ve := |x-y|_i >0$. Dès lors, $U := \{z \in E\mid |x-z|_i < \ve/2\}$ et $V := \{z \in E\mid |y-z|_i < \ve/2\}$ sont des voisinages disjoints de $x$ et $y$. 
\end{proof}


\begin{lemma}[Continuité linéaire entre evtlc]
\label{lem:evtlcLinCont}
\leancheck{lem:evtlcLinCont}%
Une application linéaire $u : E \to F$ entre evtlc $(E,(|\cdot|_i)_{i \in I})$ et $(F,(|\cdot|_j)_{j \in J})$ est continue si et seulement si pour tout $j \in J$ il existe $I_j \subset I$ fini et une constante $C_j$ tels que 
\begin{equation}
\label{eq:cont_lin_evtlc}
	|u(x)|_j \leq C_j \max_{i \in I_j} |x|_i.
\leancheck{eq:cont_lin_evtlc}%
\end{equation}	
\end{lemma}

\begin{proof}
	Supposons $u$ continue, et considérons $j \in J$. Alors le voisinage $\{y \in F \mid |y|_j < 1\}$ de l'origine de $F$ contient l'image d'un voisinage de l'origine de $E$, que l'on peut choisir dans la base \eqref{eq:evtlc_vois}, donc de la forme $\{x \in E \mid \forall i \in I_j, |x|_i < \ve\}$ avec $I_j \subset I$ fini et $\ve > 0$. 
	Par homogénéité on en déduit \eqref{eq:cont_lin_evtlc} avec $C_j := \ve^{-1}$. 
	
	Réciproquement, supposons la condition satisfaite. Par linéarité, il suffit de montrer que $u$ est continue en $0$. Considérons un voisinage de l'origine de $F$, que l'on peut choisir dans la base \eqref{eq:evtlc_vois}, donc de la forme $\{y \in F \mid \forall j \in J_0,\ |y|_j < \ve\}$ avec $J_0 \subset J$ fini et $\ve>0$. Alors il contient l'image de $\{x \in E \mid \forall i \in \cup_{j \in J_0} I_j,\ |x|_i < \ve/\max_{j \in J_0} C_j\}$, qui est bien un voisinage de l'origine. 
\end{proof}

On en déduit qu'une application linéaire d'un evtlc $(E,(|\cdot|_i)_{i \in I})$ vers un evn $(F,\|\cdot\|_F)$, ou vers le corps de base $\bK$, est continue ssi il existe $I_0\subset I$ fini et une constante $C$ tels que 
\begin{equation}
\label{eq:cont_lin_evtlc_evn}
	\|u(x)\|_F \leq C \max_{i \in I_0} |x|_i. 
\leancheck{eq:cont_lin_evtlc_evn}%
\end{equation}

\begin{remark}[Somme ou maximum de semi-normes]
\label{rem:sommeMaxEvtlc}
\leancheck{rem:sommeMaxEvtlc}%
	Certains auteurs remplacent le maximum par une somme dans \eqref{eq:cont_lin_evtlc}, \eqref{eq:cont_lin_evtlc_evn} ainsi que \eqref{eqdef:evtlc_dist} et \eqref{eqdef:met_faible_boule}. Les notions obtenues sont équivalentes, quitte à modifier certaines constantes. L'utilisation d'un maximum, préférée ici, donne souvent des preuves un peu plus simples. 
\end{remark}


\begin{remark}[Suite graduée de semi-normes]
\label{rem:evtlcGradue}
\leancheck{rem:evtlcGradue}%
Soit $(E,(|\cdot|_n)_{n \in \bN})$ un evtlc muni d'une famille dénombrable de semi-normes. On dit qu'elle est graduée si pour tout $x \in E$ la suite $(|x|_n)_{n \in \bN}$ est croissante. 
On peut faire cette hypothèse sans nuire à la généralité, quitte à considérer la suite graduée de semi-normes suivante $(|\cdot|'_N)_{N \in \bN}$ qui définit la même topologie : pour tout $N \in \bN$,
\begin{equation*}
	|x|_N' := \max_{0 \leq n \leq N} |x|_n % \{ |x|_n \mid n \leq N\}.
\end{equation*}
Noter que $|\cdot|_n = |\cdot|_n'$ si la famille initiale était déjà graduée. 
Une base de voisinages de $x\in E$ est donnée par les ensembles suivants : pour tout $N \in \bN$ et tout $\ve>0$
\begin{equation}
\label{eqdef:boule_graduee}
	B_N(x,\ve) :=  \{y \in E\mid \forall n \leq N, |y-x|_n < \ve\} = \{y \in E\mid |y-x|'_N < \ve\}.
\leancheck{eqdef:boule_graduee}%
\end{equation}

\notei{Famille filtrante, similaire dans le cas non dénombrable. Abstrait et pas vraiment nécessaire.}
\end{remark}

Les evtlcs munis d'une famille dénombrable de semi-normes sont métrisables, ce qui permet le cas échéant d'utiliser la notion de complétude, voir la \cref{def:complet}.
\begin{proposition}[Distance sur un evtlcs]
\label{prop:distanceEvtlcs}
\leancheck{prop:distanceEvtlcs}%
Soit $(E,(|\cdot|_n)_{n \in \bN})$ un evtlcs muni d'une famille dénombrable de semi-normes. Alors la topologie de $E$ est métrisable, pour la distance définie par : $\forall x,y \in E$ 
\begin{equation}
\label{eqdef:evtlc_dist}
	d(x,y) := \max_{n \in \bN} \min\{2^{-n}, |x-y|_n\}.
\leancheck{eqdef:evtlc_dist}%
\end{equation}
En notant $\log_2 \eta := \ln \eta/ \ln 2$ le logarithme en base $2$ d'un réel $\eta>0$, on a de plus
\begin{equation*}
\label{eq:boule_evtlc}
	B_d(x,\eta) = \{y \in E \mid \forall n \leq -\log_2 \eta,\ |y-x|_n < \eta\}.
\leancheck{eq:boule_evtlc}%
\end{equation*}
\end{proposition}

\begin{proof}
L'expression de la boule est claire, en remarquant que $\min\{2^{-n},|x-y|_n\}<\eta$ est toujours satisfait si $n > -\log_2 \eta$, et est équivalent à $|x-y|_n < \eta$ sinon. 
On en déduit les inclusions suivantes, pour tout $N \in \bN$ et tous $\epsilon,\eta>0$, montrant que $d$ définit une base de voisinages équivalente à \eqref{eqdef:boule_graduee}, donc la même topologie
\belowdisplayskip=-12pt
\begin{align*}
	B_d(x, \min \{\ve,2^{-N-1}\}) &\subset B_N(x,\ve), &
	B_{\lfloor -\log_2 \eta \rfloor}(x,\eta) &\subset B_d(x,\eta). 
\end{align*}
\end{proof}

\notei{Suppression de l'uniforme continuité des applications linéaires, assez triviale et inutile.}
%\begin{remark}
%Soient $(E,(|\cdot|_n)_{n \in \bN})$ et $(F,(|\cdot|_m)_{m \in \bN})$ des evtlc munis de familles dénombrables et graduées de semi-normes. Alors toute application linéaire continue $u \in L(E,F)$ est uniformément continue pour les distances associées.
%\todoi{Vraiment utile ?}
%\end{remark}
%
%\begin{proof}
%	Notons $\omega_0$ le module de continuité de $u$ en $0\in E$, qui existe car $u$ est continue. 
%	On vérifie alors que c'est un module de continuité uniforme, en utilisant l'invariance par translation : $d_F(u(x),u(y)) = d_F(0,u(y-x))  \leq \omega_0 (d_E(0,y-x)) = d_E(x,y)$. 
%\end{proof}

\begin{definition}[Partie linéairement bornée]
\label{def:lineairBorne}
\leancheck{def:lineairBorne}%
	Une partie $A$ d'un evtlcs $(E,(|\cdot|_i)_{i \in I})$ est dite linéairement bornée ssi chaque semi-norme est bornée sur $A$ : $\forall i \in I$, $\sup\{|x|_i\mid x\in A\}$ est borné. 	
	Si la famille de semi-normes est dénombrable, et si $d$ désigne la distance \eqref{eqdef:evtlc_dist}, alors on a la caractérisation équivalente : $\forall \ve>0$, $\exists \delta>0$, $\delta A \subset B_d(0,\ve)$. 
\end{definition}

\subsubsection{Exemples d'espaces fonctionnels}

Nous présentons plusieurs exemples d'espaces fonctionnels dont la topologie n'est pas décrite par une seule norme, mais par une famille de semi-normes. C'est souvent nécessaire lorsque l'on doit contrôler une infinité de dérivées, ou une fonction non-bornée. 

Soit $\alpha = (\alpha_1,\cdots,\alpha_d) \in \bN^d$, et $k = |\alpha|_1 := \alpha_1+\cdots+\alpha_d $. Si une fonction $f$ est $k$ fois dérivable au voisinage de $x \in \bR^d$, alors on pose
\begin{equation*}
	\partial_\alpha f(x) := \frac{\partial^{|\alpha|_1} f}{\partial x_1^{\alpha_1}\cdots \partial x_d^{\alpha_d}}.
\end{equation*}

La notation $K \csubset \Omega$ signifie que $K$ est un compact inclus dans $\Omega$. On étudiera la compacité en détail \cref{sec:compact}, et on utilise seulement ici le fait que les fermés bornés de $\bR^d$ sont compacts, et qu'une fonction réelle continue sur un compact atteint ses bornes.

\begin{definition}[Fonctions $C^\infty$ bornées sur un domaine]	
\label{def:cinf_bornee}
\leancheck{def:cinf_bornee}%
On note $C^\infty_b(\Omega)$, où $\Omega\subset \bR^d$ est un ouvert, l'ensemble des fonctions $f$ indéfiniment dérivables sur $\Omega$ et bornées ainsi que toutes leurs dérivées. Cet espace est muni de la famille graduée de (semi-)normes suivante : pour tout $k \in \bN$
\begin{equation*}
	|f|_k := \max_{|\alpha|_1 \leq k} \|\partial_\alpha f\|_\infty.
\end{equation*}
Pour tout compact $K \csubset \Omega$ on note $\cD_K(\Omega) := \{f \in C^\infty_b(\Omega) \mid \supp(f) \subset K\}$ le sous espace fermé des fonctions à support inclus dans $K$. 
\end{definition}

La classe de Schwartz est un espace de fonctions indéfiniment dérivables à décroissance rapide, intervenant dans l'étude de la transformée de Fourier. 
(Note orthographique : cette classe de fonctions est nommée d'après Laurent Schwartz, alors que l'inégalité de Cauchy-Schwarz est associée à Augustin-Louis Cauchy et Hermann Schwarz.)

\begin{definition}[Classe de Schwartz]
\label{def:classeSchwartz}
\leancheck{def:classeSchwartz}%
	Pour $f \in C^\infty(\bR^d)$, $\alpha \in \bN^d$, $n \in \bN$, on pose 
	\begin{equation}
		|f|_{n,\alpha} := \sup_{x \in \bR^d} (1+|x|^2)^\frac n 2 |\partial_\alpha f(x)|
	\end{equation}
	On note $\cS(\bR^d)$ le sous-ensemble des fonctions sur lequel toutes ces semi-normes sont finies, qui est donc un evtlcs métrisable. On obtient une suite graduée de (semi-)normes en posant (par exemple) $|f|_n := \max_{|\alpha|_1 \leq n} |f|_{n,\alpha}$ pour tout $n \in \bN$.
\end{definition}


On utilise dans les \cref{def:lploc,def:cont_nonbornee} une suite exhaustive de compacts d'un domaine $\Omega \subset \bR^d$, pour définir une suite graduée de semi-normes sur un espace de fonctions sur $\Omega$. L'existence d'une telle suite exhaustive est assurée par la proposition suivante, et on vérifie facilement que la topologie obtenue ne dépend pas de la suite choisie.
\notei{On peut les supposer connexes si $\Omega$ est connexe}

\begin{proposition}[Suite exhaustive de compacts]
\label{prop:exhaustComp}
\leancheck{prop:exhaustComp}%
Étant donné un ouvert $\Omega\subset \bR^d$, il existe une suite $(K_n)_{n \in \bN}$ de parties compactes telles que 
\begin{align}
\label{eqdef:exhaustive_compacts}
\cup_{n \in \bN} K_n &= \Omega,&
\forall n \in \bN,\, K_n &\csubset \mathring K_{n+1}. 
\leancheck{eqdef:exhaustive_compacts}%
\end{align}
Si $K \csubset \Omega$, alors $K \csubset \mathring K_n$ pour tout $n$ assez grand.
\end{proposition}

\begin{proof}
	\emph{Existence}: prendre par exemple $K_n := \{x \in B'(0,n)\mid d(x,\bR^d \sm \Omega) \geq \frac 1{n+1}\}$. 
	\emph{Inclusion $K \csubset \mathring K_n$}: Remarquer que $K$ est couvert par l'union croissante d'ouverts $\cup_{n\in \bN} \mathring K_{n+1} = \Omega$, puis appliquer l'axiome de Borel-Lebesgue, voir la \cref{def:compact}.
\end{proof}

Par convention $\|f\|_{L^p(K)}$ désigne $\|f_{|K}\|_{L^p(K)}$ si $f$ est définie sur un domaine contenant $K \subset \bR^d$, $f$ et $K$ étant supposés mesurables. 

\begin{definition}[Evtlcs $L^p_\loc$]
\label{def:lploc}
\leancheck{def:lploc}%
Soit un ouvert $\Omega\subset \bR^d$ et un exposant $p\in [1,\infty]$. 
On note $L^p_\loc(\Omega)$ l'evtlcs des fonctions mesurables, modulo l'égalité presque partout, pour lesquelles la famille graduée de semi-normes suivantes prend des valeurs finies :
\begin{equation*}
	|f|_n := \|f\|_{L^p(K_n)},
\end{equation*}
où $(K_n)_{n \in \bN}$ est une suite exhaustive de compacts de $\Omega$. 
\end{definition}

\begin{definition}[Evtlcs $C^k(\Omega)$ et $C^\infty(\Omega)$, formés de fonctions non-bornées a-priori]	
\label{def:cont_nonbornee}
\leancheck{def:cont_nonbornee}%
	Soit $\Omega \subset \bR^d$ un ouvert, muni d'une suite exhaustive de compacts $(K_n)_{n \in \bN}$. On pose pour tout $\alpha \in \bN^d$
	\begin{equation*}
		|f|_{n,\alpha} := \|\partial_\alpha f\|_{L^\infty(K_n)} %\sup_{x \in K_n} \max_{|\alpha|_1 \leq k} | \partial_\alpha f(x)|.
	\end{equation*}
	Pour tout $k \in \bN \cup \{\infty\}$, on a l'evtlcs $(C^k(\Omega), (|\cdot|_{n,\alpha})_{n \in \bN}^{|\alpha|_1\leq k})$. On obtient une famille graduée de semi-normes en posant, par exemple, $|f|_n := \max_{|\alpha|_1 \leq \min\{n,k\}} |f|_{n,\alpha}$.
\end{definition}

Les espaces fonctionnels précédents de cette section sont des evtlcs métrisables, car leur topologie est définie par une famille dénombrable de semi-normes. 
Ce n'est en revanche pas le cas de la topologie de $\cD(\Omega)$, décrite ci-dessous.
Une description plus abstraite, la topologie dite limite inductive, est présentée dans \cite{carlier2009analyse}.


\begin{definition}[Fonctions test $\cD(\Omega)$ et distributions]
\label{def:distribution}
\leancheck{def:distribution}%
	On note $\cD(\Omega) = C^\infty_c(\Omega)$, où $\Omega \subset \bR^d$, l'ensemble des fonctions indéfiniment dérivables à support compact dans $\Omega$. 
	On dit qu'une forme linéaire $\vp$ sur $\cD(\Omega)$ est une distribution si sa restriction à $\cD_K(\Omega)$ est continue pour tout $K \csubset \Omega$, i.e. 
	\begin{equation}
	\label{eq:distrib_cont_loc}
		\forall K \csubset \Omega,\, \exists C_K,\, \exists n_K \in \bN,\, 
		\forall f \in \cD_K(\Omega),\, |\vp(f)| \leq C_K \max_{|\alpha|_1 \leq n_K} \|\partial_\alpha f\|_{L^\infty(K)}.
	\end{equation}
\end{definition}
\notei{Mentionner qu'on pourrait se restreindre à un voisinage compact de tout point, en utilisant des partitions de l'unité.}

L'entier $n_K$ minimal dans \eqref{eq:distrib_cont_loc} est appelé ordre de la distribution $\vp$ sur $K$. 
L'ordre de la distribution est $n(\vp) := \sup \{n_K \mid K \csubset \Omega\} \in \bN \cup \{\infty\}$.

Les distributions d'ordre zero sont aussi appelées \emph{mesures signées}. Le théorème de représentation de Riesz–Markov–Kakutani établit que celles qui sont positives, i.e.\ $\<u,\vp\> \geq 0$ dès que $\vp \geq 0$ sur $\Omega$, correspondent aux mesures de Radon\footnote{Mesure intérieurement régulière et localement finie sur la tribu des Boréliens d'un espace topologique séparé.} sur $\Omega$. Noter qu'il existe des distributions qui ne sont pas d'ordre fini. 

On rappelle que pour $\rho,f\in C^\infty(\bR^d)$, $\alpha \in \bN^d$, on a la formule du binome pour la dérivation
\begin{align}
\label{eq:deriv_binom}
	\partial_\alpha(\rho f) &= \sum_{\beta \preceq \alpha} 
	\binom \alpha \beta \partial_\beta \rho \, \partial_{\alpha-\beta} f,&
	\text{où }
	\binom \alpha \beta := \prod_{1 \leq i \leq d} \frac{\alpha_i!}{\beta_i! (\alpha_i-\beta_i)!},
\end{align}
et où $\beta \preceq \alpha \eqdef \forall 1 \leq i \leq d, \beta_i \leq \alpha_i$.
La proposition suivante montre que l'ensemble des distributions est le dual topologique de l'espace des fonctions test, pour une topologie d'evtlc adéquate, raison pour laquelle on le notera $\cD(\Omega)^*$. Elle est précédée d'un lemme facile dont la démonstration est omise. 



\notei{Supprimé le produit d'une distribution et d'une fonction, assez trivial.}
%\todoi{La def suivante est elle appropriée ici, ou dans le début du cours ?}
%\begin{exemple}[Produit d'une distribution et d'une fonction]
%	Soit $T \in \cD'(\Omega)$ et $\phi \in C^\infty(\Omega)$. On définit le produit $\phi T \in \cD'(\Omega)$ par $\<\phi T,\psi\> := \<T,\phi \psi\>$ pour tout $\psi \in \cD'(\Omega)$. 
%
%On vérifie facilement que $\phi T \in \cD'(\Omega)$, que $(\phi,T) \in C^\infty \times \cD' \mapsto \phi T \in \cD'$ est continue, que $\ordre(\phi T) \leq \ordre(T)$, et que $\psi(\phi T) = (\psi\phi)  T$. De plus $\phi U_f = U_{\phi f}$, où $U_f$ est la distribution canoniquement associée à une fonction $f \in L^1_\loc$, de sorte que la multiplication de distributions prolonge celle des fonctions.
%\end{exemple}

\begin{lemma}
\label{lem:omegaNmuN}
\leancheck{lem:omegaNmuN}%
	Soit $(\omega_n)_{n \in \bN} \in \bR_+^\bN$ une suite à valeurs positives, et $(\mu_n)_{n \in \bN}\in (\bR_+^*)^\bN$ strictement positive et telle que $\sum_{n\in \bN} \mu_n=1$.
	Alors $\sum_{n \in \bN} \omega_n \leq \sup_{n \in \bN} \omega_n/\mu_n$. % \frac {\omega_n} {\mu_n}$. 
\end{lemma}
%\begin{proof}
%	$\sum_{n \in \bN} (\omega_n/\mu_n) \mu_n
%\end{proof}

\begin{proposition}[Continuité des distributions \facultatif]
\label{prop:distribCont}
Soit $\Omega\subset \bR^d$ ouvert. On munit $\cD(\Omega)$ de la topologie d'evtlcs définie par la famille de semi-normes : pour toutes fonctions $\omega,\eta \in C^0(\Omega, \bR_+)$
\begin{equation*}
	|f|_{\omega,\eta} := \max_{x \in \Omega} \Big(\omega(x) \max_{|\alpha|_1 \leq \eta(x)} |\partial_\alpha f(x)|\Big).
\end{equation*}
Pour une forme linéaire $\vp$ sur $\cD(\Omega)$, sont équivalentes (i) $\vp$ est continue sur $\cD(\Omega)$, soit encore 
	\begin{equation}
	\label{eq:distrib_cont}
		\exists \omega, \eta \in C^0(\Omega,\bR_+),\ \forall f \in \cD(\Omega),\ |\vp(f)| \leq |f|_{\omega,\eta},
	\end{equation}
et (ii) $\vp$ est une distribution, soit encore \eqref{eq:distrib_cont_loc}. 
\end{proposition}

\begin{proof}
%L'équivalence de (i) avec (iii) est claire au vu du critère \eqref{eq:cont_lin_evtlc_evn}.
\emph{Preuve de (i) $\Rightarrow$ \eqref{eq:distrib_cont}.}
 Si $\vp$ est continue, alors il existe  une constante $C$, un entier $N$ et des fonctions $(\omega_n,\eta_n)_{n=1}^N$ tels que $|\vp(f)|\leq C \max_{n=1}^N |f|_{\omega_n,\eta_n}$. On conclut en posant $\omega := C\max_{n=1}^N \omega_n$ et $\eta := \max_{n=1}^N \eta_n$.
	La réciproque \eqref{eq:distrib_cont} $\Rightarrow$ (i) est évidente. 
%Preuve de (i) $\Leftrightarrow$ \eqref{eq:distrib_cont_loc}. Il suffit de remarquer que $(|\cdot|_n)_{n \in \bN}$ est une suite graduée de semi-normes. 

\emph{Preuve de \eqref{eq:distrib_cont} $\Rightarrow$ \eqref{eq:distrib_cont_loc}.} Poser $\omega_K := \sup_K \omega$ et $\eta_K:=\sup_K \eta$. Passons à la preuve de la réciproque \eqref{eq:distrib_cont_loc} $\Rightarrow$ \eqref{eq:distrib_cont}, qui demande un peu plus de travail.  

%	Soit $(\Omega_n)_{n \in \bN}$ une suite d'ouverts tels que $\overline {\Omega_n} \csubset \Omega_{n+1}$ et $\cup_{n \in \bN} \Omega_n = \Omega$. 
Soit $K_n$ une suite exhaustive de compacts de $\Omega$.	
Notons $\omega_n := C_{K_n}$ et $\eta_n := n_{K_n}$ pour la lisibilité. 	Soit $\chi_n$, $C^\infty$, positive, telle que $\chi_n=1$ sur $K_{n-1}$ et $\supp(\chi_n) \subset \mathring K_n$, puis $\rho_n := \chi_n - \chi_{n-1}$, avec les conventions $K_{-1} = K_{-2} = \emptyset$ et $\chi_{-1} = \chi_{-2} =0$.
	On obtient
	\begin{align*}
%	\label{eq:distrib0_lin}
		|\vp(f)| &= \Big| \sum_{n \in \bN} \vp(\rho_n f) \Big|\\
%	\label{eq:distrib0_cont_loc}
		&\leq \sum_{n\in \bN} \omega_n \max_{|\alpha|_1 \leq \eta_n} \|\partial_\alpha (\rho_n f)\|_\infty\\
%	\label{eq:distrib0_binom}
		&\leq \sum_{n \in \bN} \tilde \omega_n \max_{|\alpha|_1 \leq \eta_n} \max_{x \in \supp(\rho_n) }|\partial_\alpha f(x)|\\
%	\label{eq:distrib0_sup}
		&\leq \sup_{n \in \bN} \hat \omega_n \max_{|\alpha|_1 \leq \eta_n} %\|f^{(k)}\|_\infty 
		\max_{x \in \supp(\rho_n) }|\partial_\alpha f(x)|\\
		& \leq |f|_{\omega,\eta}.
	\end{align*}
	On a utilisé successivement à chaque ligne (a) la linéarité de $\vp$, sachant que la somme a un nombre fini de termes non-nuls, (b) l'hypothèse \eqref{eq:distrib_cont_loc}, (c) la formule du binôme pour la dérivation \eqref{eq:deriv_binom}, avec une constante $\tilde \omega_n$ appropriée, (d) les nouvelles variables $\hat \omega_n = \tilde \omega_n /\mu_n$, où la suite $(\mu_n)_{n \in \bN}$ satisfait les conditions du \cref{lem:omegaNmuN}, par exemple $\mu_n := c_\alpha (1+n)^{-\alpha}$ avec $\alpha>1$ et une constante $c_\alpha>0$ appropriée, (e) des fonctions $\omega,\eta \in C^0(\Omega,\bR_+)$ telles que $\omega \geq \hat \omega_n$ et $\eta \geq \eta_n$ sur $\supp(\rho_n)\subset \mathring K_n \sm \mathring K_{n-2}$; par exemple $\omega := \max_{n\in \bN} \hat \omega_n (1-\chi_{n-2})$ et de même pour $\eta$. 
\end{proof}


\subsubsection{La topologie faible}

La topologie faible, sur un evn ou un evtlcs, est comme son nom l'indique une topologie plus grossière - possédant moins d'ouverts - que la topologie usuelle, en dimension infinie. L'un de ses intérêts principaux, on le verra, est le \cref{th:BanachAlaoglu} de Banach-Alaoglu, montrant la compacité *faible de la boule unité d'un espace dual en dimension infinie, ce qui est exclu dans le cas de la topologie usuelle d'un evn d'après le \cref{th:bouleCompRiesz} de Riesz.

\begin{definition}[Topologie faible et *faible]
\label{def:topoFaible}
\leancheck{def:topoFaible}%
Soit $E$ un evtlcs (ou un evn). 
\begin{itemize}
	\item La topologie faible sur $E$ est définie par la famille de semi-normes $x\in E\mapsto |l(x)|$, pour tout $l \in E^*$. 
	\item La topologie *faible sur $E^*$ est définie par la famille de semi-normes $l \in E^* \mapsto |l(x)|$, pour tout $x\in E$. 
\end{itemize}	
\end{definition}

La topologie faible (resp.\ *faible) est aussi la topologie la plus grossière qui rend continues toutes les formes linéaires $l\in E^*$ (resp.\ de la forme $l\in E^*\mapsto l(x)$ où $x \in E$).
La séparation de la topologie *faible est évidente par construction, tandis que la séparation de la topologie faible découle du Théorème de Hahn-Banach, voir la \cref{sec:dual}.

Les topologies faible et *faible ne sont en général pas métrisables, voir \cite[exercice]{alazard2017analyse}.
\todoi{ref exacte}
En revanche, sous certaines conditions, leur restriction à la boule unité, ou à n'importe quelle partie bornée, est métrisable. 
\begin{proposition}[Métrisabilité de la topologie *faible restreinte à un borné]
\label{prop:bouleMetrisable}
\leancheck{prop:bouleMetrisable}%
	Soit $E$ un evn séparable, soit $(x_n)_{n \in \bN}$ une suite dense dans la boule unité de $E$, et soit $B^* := B'_{E^*}(0,1)$ la boule unité de l'espace dual pour la norme d'opérateur. 
	Alors la restriction à $B^*$ de la topologie *faible est métrisable, pour la distance
	\begin{equation}
	\label{eqdef:met_faible_boule}
		d(u,v) := \max_{n \in \bN} \min \{2^{-n}, |u(x_n)-v(x_n)|\}.
\leancheck{eqdef:met_faible_boule}%
	\end{equation}
\end{proposition}
	
\begin{proof}
	Une base de voisinages de $u \in B^*$ pour la topologie *faible est donnée par les ensembles de la forme
	\begin{equation}
	\label{eq:vois_faible_boule_evtlc}
		V := \{v\in B^* \mid |u(y_k)-v(y_k)| < \ve \text{ pour tout } 1 \leq k \leq K\}.
\leancheck{eq:vois_faible_boule_evtlc}%
	\end{equation}
	où $y_1,\cdots,y_K \in E$, et où $\ve>0$. Quitte à réduire $\ve$, on peut supposer que $y_1,\cdots,y_K\in B := B'_E(0,1)$.
	Soient $n_1,\cdots,n_K\in \bN$ tels que $\|x_{n_k}-y_k\| \leq \ve/3$ pour tout $1 \leq k \leq K$. Soit $N := \max\{n_1,\cdots,n_K\}$ et soit $\eta := \min \{\ve/3, 2^{-N-1}\}$. Posons le voisinage
	\begin{equation}
		W := B_d(u,\eta) \cap B^* \subset \{v \in B^* \mid \forall 1 \leq k \leq K, |u(x_{n_k})-v(x_{n_k})| < \ve/3\}. 
	\end{equation}
	On conclut que $W \subset V$ en notant que pour tout $v \in W$ et tout $k \leq K$, 
	\begin{equation}
		|v(y_k)-u(y_k)| \leq |v(y_k)-v(x_{n_k})| + |v(x_{n_k})-u(x_{n_k})| + |u(x_{n_k}) - u(y_k)| < 3 \ve /3 = \ve.
	\end{equation}
	Réciproquement $B_d(u,\eta) \cap B^*$ est exactement de la forme \eqref{eq:vois_faible_boule_evtlc}, au vu de \eqref{eq:boule_evtlc}, en choisissant $\ve=\eta$, $K := \lfloor -\log_2 \eta\rfloor$, et $y_k := x_k$ pour $k\leq K$. 
\end{proof}

%\todoi{On peut supposer que $E$ est un evtlc, mais il faut adapter la notion de bornitude. Par exemple supposer que $\forall x \in E$, $\forall l \in B$, $|l(x)| \leq C \sum_{i \in I_0} |x|_i$ où $I_0$ est fini. Cela revient en gros à se ramener à un evn, avec la norme $x\mapsto \sum_{i \in I_0} |x|_i$. 
%}

%\todoi{Un opérateur linéaire continu, est continu pour les topologies faibles}

\begin{remark}
Soit $A \in L(E,F)$ un opérateur linéaire continu entre evn, et soit $A^*\in L(F^*,E^*)$ son adjoint. Alors $A \in L( (E,\text{faible}), (F,\text{faible}))$ et $A^*\in L( (F^*,\text{*faible}), (E^*, \text{*faible}))$.
\end{remark}

\begin{proof}
	Rappelons que $A^*(l) := l\circ A$ est un opérateur linéaire continu, puisque classiquement $\|A^*(l)\|_{E^*} = \|l\circ A\|_{E^*} \leq \|l\|_{F^*} \|A\|_{L(E,F)}$.
	Il s'agit d'appliquer le critère du \cref{lem:evtlcLinCont}. Notons $|l(x)| = |l|_x = |x|_l$ pour la lisibilité. Premier point: pour tous $l \in F^*$, $x \in E$, on a $|A(x)|_l = |l(A(x))| = |x|_{l\circ A}$.  Second point: pour tous $l \in F^*$, $x\in E$, on a $|A^*(l)|_x = |l(A(x))| = |l|_{A(x)}$.
\end{proof}



\newpage
\section{Complétude}
\label{sec:complet}

La complétude est un outil fondamental en analyse, assurant l'existence d'une limite unique, ce qui permet de prolonger de manière unique une fonction uniformément continue, de montrer l'existence d'un point fixe unique via le théorème de Picard, ou encore d'établir des résultats d'analyse linéaire abstraits et puissants via le théorème de Baire. 

\inventaire{
\Nref{subsec:Cauchy}
\Nref{def:suiteCauchy} \\
\Nref{def:complet}\\
\Nref{lem:serie_Banach}\\
\Nref{prop:fermeComplet} (Les fermés d'un complet sont complets)\\
\Nref{prop:LEFcomplet} ($L(E,F)$ est complet si $F$ complet)\\

\Nref{subsec:fctsComplets}\\
\Nref{ex:cont_bornee_Banach} ($C^0_b$ est un Banach)\\
\Nref{ex:Holder_Banach}\\
\Nref{ex:Ck_bornee_Banach} ($C^k_b(\Omega)$ et $C^k_b(\overline\Omega)$ sont des Banach)\\
\Nref{ex:LpBanach} ($L^p$ est un Banach)\\
\Nref{ex:testDistribFrechet} (Les fonctions test des distributions forment un Fréchet)\\
\Nref{ex:DKOmegaComplet} ($\cD_K(\Omega)$ est un Fréchet)\\
\Nref{ex:classeSchwartzFrechet} (La classe de Schwartz est un Fréchet)\\

\Nref{subsec:prolongements}\\
\Nref{prop:prolongement_complet}\\
\Nref{prop:espaceComplete}\\
\Nref{ex:extUnifCont} \\

\Nref{subsec:Picard} \\
\Nref{th:fixe_Picard} \\
\Nref{lem:Gronwall} \\
\Nref{th:Cauchy_Lipschitz} \\
\Nref{corol:CauchyLipAuton} \\
\Nref{ex:EDretard} \\
\Nref{ex:eqFunctIntegr} \\

\Nref{subsec:Baire} \\
\Nref{lem:fermeEmboite} \\
\Nref{th:BaireOuvert} \\
\Nref{corol:BaireFerme} \\
\Nref{rem:terminologieBaire} \\
\Nref{ex:GrasQ} (Ensemble gras de mesure nulle, contenant $\bQ$)\\
\Nref{ex:grasNonDen} (Un ensemble gras, dans un espace sans points isolé, est indénombrable)\\

\Nref{subsec:BaireLin} \\
\Nref{th:BanachSteinhauss} \\
\Nref{corol:limSimpleLin} (Limite simple d'applications linéaires) \\
\Nref{corol:BanachSteinhaussK} \\
\Nref{ex:FourierDiverge}	 (Une fonction continue dont la série de Fourier diverge en zéro) \\
\Nref{th:BanachSteinhaussFrechet} \\
\Nref{th:appliOuverte} \\
\Nref{corol:isomBanach} \\
\Nref{corol:equivNormes} \\
\Nref{th:grapheFerme} 
}


\subsection{Critère de Cauchy}
\label{subsec:Cauchy}

\begin{retenir}
	La définition d'une suite de Cauchy, la complétude d'une partie fermée d'un complet, et de l'ensemble des applications linéaires dans un complet.
	
	(\Cref{def:suiteCauchy,prop:fermeComplet,prop:LEFcomplet})
\end{retenir}


\begin{definition}[Suite de Cauchy]
\label{def:suiteCauchy}
	Soit $(X,d)$ un espace métrique. Une suite $(x_n)_{n\in \bN}$ est de dite de \emph{Cauchy} si elle satisfait au critère suivant:
	\begin{equation*}
		\forall \ve>0,\ \exists N \in \bN,\ \forall p,q \geq N,\ d(x_p,x_q) \leq \ve. 
	\end{equation*}
\end{definition}

Une suite de Cauchy est bornée, et admet au plus une limite. Une suite de Cauchy qui admet une sous-suite convergente est convergente. Réciproquement, toute suite convergente est de Cauchy. 
Noter que la propriété de Cauchy, et donc la notion de complétude, n'ont de sens que dans un espace métrique. On peut reformuler ainsi le critère de Cauchy : il existe une suite $(\ve(n))_{n \in \bN}$ tendant vers $0$, et telle que pour tous $p,q\in \bN$
\begin{equation}
\label{eqdef:suiteCauchyMin}
	d(x_p,x_q) \leq \ve(\min(p,q))
\end{equation}
\begin{definition}[Complétude]
\label{def:complet}
Un espace métrique $(X,d)$ est dit \emph{complet} si toutes ses suites de Cauchy y sont convergentes. 
Un espace de Banach est un evn complet. 
Un espace de Fréchet est evtlcs muni d'une famille dénombrable de semi-normes, et complet pour la distance associée.
\end{definition}


La complétude d'un evn peut également s'exprimer en termes de séries, plutôt que de suites.
Rappelons qu'une série à valeurs dans un evn, écrite formellement ``$\sum_{n \in \bN} x_n$'' est dite \emph{absolument convergente} ssi $\sum_{n\geq 0} \|x_n\| < \infty$, et est dite \emph{convergente} ssi la suite des sommes partielles $(\sum_{n \leq N} x_n)_{N \in \bN}$ admet une limite lorsque $n\to \infty$.

\begin{lemma}[Séries dans un Banach]
\label{lem:serie_Banach}
Un evn $E$ est complet ssi toute série absolument convergente y est convergente.   	
\end{lemma}

\begin{proof}
	\emph{Supposons $E$ complet.} Soit $(x_n)_{n \in \bN}$ les termes d'une série absolument convergente, et soit $y_N := \sum_{n<N} x_n$ pour tout $N>0$. 
	Alors pour $p,q \geq N$, disons $p\leq q$, on a $\|y_p-y_q\| \leq \sum_{p \leq n < q} \|x_n\| \leq r_N := \sum_{n \geq N} \|x_n\|$. Par sommabilité, $r_N \to 0$ lorsque $N \to \infty$. Ainsi $y_n$ est de Cauchy, donc converge, comme annoncé. 
	
	\emph{Supposons maintenant que les séries absolument convergentes sont convergentes.}
	Soit $(x_n)_{n \in \bN}$ une suite de Cauchy. 
	Pour tout $k \in \bN$, soit $N_k$ croissant et tel que $\| x_p-x_q \| \leq 2^{-k}$ pour tous $p,q \geq N_k$. 
	Posons $y_k := x_{N_{k+1}}-x_{N_k}$, de sorte $\| y_k\| \leq 2^{-k}$, et donc que $\sum_{k \in \bN} \|y_k\|$ converge. 
	Par hypothèse, les sommes partielles $\sum_{k < K} y_k = x_{N_K}-x_{N_0}$ convergent lorsque $K \to \infty$. Donc la sous-suite $(x_{N_K})_{K\geq 0}$ converge, et donc aussi $(x_n)_{n\in \bN}$ comme annoncé car elle est de Cauchy et admet une sous-suite convergente. 
\end{proof}


Plusieurs constructions mathématiques héritent de la propriété de complétude.


\begin{proposition}
\label{prop:fermeComplet}
	Une partie $A$ d'un espace métrique complet $(X,d)$, est complète ssi elle est fermée.	
\end{proposition}

\begin{proof}
	On utilise la caractérisation ($A$ fermé ssi $\overline A=A$), ainsi que la caractérisation séquentielle de l'adhérence \cref{prop:adhSuite}.

	Si $A$ n'est pas fermé, alors en considérant une suite convergeant vers un point de $\overline A \sm A$, on obtient une suite de Cauchy non-convergente dans $A$, qui n'est donc pas complet.
	
	Si $A$ fermé, alors toute suite de Cauchy à valeurs dans $A$ est convergente dans $X$, car celui-ci est complet, et sa limite appartient à $A$ par fermeture, qui est donc complet.
\end{proof}



\begin{proposition}
\label{prop:LEFcomplet}
Soient $E$ un evn et $F$ un Banach. Alors $L(E,F)$ est un Banach, et en particulier $E^*$ est un Banach. 	
\end{proposition}

\begin{proof}
	Soit $(u_n)_{n \in \bN}$ de Cauchy dans $L(E,F)$. On vérifie que $(u_n(x))_{n \in \bN}$ est également de Cauchy, pour tout $x \in E$, par continuité de l'évaluation puisque 
%$\|u_n(x)-u_m(x)\|_F \leq \|u_n-u_m\|_{L(E,F)} \|x\|_E$, 
$\|u_n(x)-u_m(x)\| \leq \vertiii{u_n-u_m} \|x\|$, 
	ce qui permet de définir une application limite $u_\infty(x) := \lim_{n\to \infty} u_n(x)$. On vérifie ensuite que $u_\infty$ est linéaire, par continuité de l'addition et de la multiplication scalaire. Par le critère de Cauchy on a $\vertiii{u_p-u_q}\leq \ve_N$ pour tous $p,q\geq N$, où la suite $\ve_N \to 0$ lorsque $N \to \infty$. Ainsi $\|u_\infty(x)-u_N(x)\| \leq \ve_N$ par passage à la limite pour tout $x \in B := B'_E(0,1)$, d'où le caractère continu de $u_\infty$ et la convergence $\vertiii{u_\infty-u_N} \leq \ve_N \to 0$ lorsque $N \to \infty$.
\end{proof}

Attention : il n'y a pas de semi-norme naturelle sur le dual d'un Fréchet, qui est seulement muni de la topologie *faible a priori. Le dual d'un Fréchet n'est donc pas un Fréchet en général, à la différence des Banach.
%A fortiori, on n'a pas de semi-norme naturelle sur $L(E,F)$ où $F$ est un Fréchet ou un evn. 

\subsection{Espaces fonctionnels complets}
\label{subsec:fctsComplets}

Les plupart espaces fonctionnels classiques, qui sont métrisables, sont complets.

\begin{retenir}
Les espaces de fonctions bornées qui sont continues, Holder, Lipschitz, ou continument dérivables jusqu'à un certain ordre, sont des Banach, de même que les espaces de Lebesgue $L^p$. Les espaces de fonction test $\cD_K(\Omega)$ et $\cS(\bR^d)$ sont des Frechet. 

(Exemples \ref{ex:cont_bornee_Banach}, \ref{ex:Holder_Banach}, \ref{ex:Ck_bornee_Banach}, \ref{ex:LpBanach}, \ref{ex:testDistribFrechet}, \ref{ex:classeSchwartzFrechet})
\end{retenir}

\paragraph{Espaces de Banach.}

\begin{exemple}
\label{ex:cont_bornee_Banach}
	Soit $(X,\bU)$ un espace topologique, $E$ un Banach. Alors $C^0_b(X,E)$ est un Banach.
\end{exemple}

\begin{proof}
	Soit $(u_n)_{n \in \bN}$ de Cauchy dans $C^0_b(X,E)$, qui est muni de la norme uniforme $\|\cdot\|_\infty$. Alors pour tout $x \in X$, $(u_n(x))_{n \in \bN}$ est de Cauchy, donc convergente, et on note $u_*(x)$ sa limite. 
	
	La limite $u_* : X\to E$ est continue (resp.\ bornée) en tant que limite uniforme de fonctions continues (resp.\ bornées). Finalement, $\|u_n - u_*\|_\infty \leq \max_{m \geq 0} \|u_n-u_m\|_\infty \to 0$ par le critère de Cauchy, ce qui conclut.  
\end{proof}

\begin{exemple}[Complétude des espaces de Lipschitz et Holder]
\label{ex:Holder_Banach}
	Soit $(X,d)$ un espace métrique, $E$ un Banach, et $\omega : \bR_+ \to \bR_+$ un module de continuité strictement positif hors de $0$. 
	On pose 
	\begin{align}
		|f|_\omega &:= \sup_{x\neq y} \frac{|f(x)-f(y)|}{\omega(d(x,y))},&
		C_b^\omega(X,E) &:= \{f \in C_b(X,E) \mid |f|_\omega < \infty\}. 
	\end{align}
	Alors $C^\omega_b(X,E)$ est un Banach pour la norme $\|\cdot\|_\omega := \|\cdot\|_\infty + |\cdot|_\omega$. 
	L'espace de Hölder $C^\alpha_b(X,E)$ correspond à $\omega(s):=s^\alpha$, où $0 < \alpha < 1$, et $\Lip_b(X,E)$ à $\omega(s):=s$.
\end{exemple} 

\begin{proof}
	On note que $|\cdot|_\omega$ positivement homogène et sous-additive, donc $C_b^\omega$ est bien un espace vectoriel. De plus $\|\cdot\|_\omega$ est une norme sur $C^\omega_b$ car $\|\cdot\|_\infty$ satisfait l'axiome de séparation. 
	
	Soit $(f_n)_{n \in \bN}$ une suite de Cauchy. Par complétude de $(C_b(X,E),\|\cdot\|_\infty)$, $f_n$ converge en norme $\|\cdot\|_\infty$ vers une limite $f_\infty \in C_b(X,E)$.
	Comme $f_n$ est de Cauchy, $|f_n|_\omega$ est bornée par une constante notée $C$, donc pour tous $x \neq y$ on a $\|f_\infty(x)-f_\infty(y)\| = \lim_{n\to \infty} \|f_n(x)-f_n(y)\| \leq C \omega(d(x,y))$, ce qui montre que $|f_\infty|_\omega \leq C$. 
	De même $\| (f_\infty(x)-f_N(x)) - (f_\infty(y)-f_N(y))\| = \lim_{n \to \infty} \| (f_n(x)-f_N(x)) - (f_n(y)-f_N(y))\| \leq \ve_N \omega(d(x,y))$, où $\ve_N\to 0$ est donnée par le critère de Cauchy, d'où $|f_\infty-f_N|_\omega \leq \ve_N$, ce qui conclut.
\end{proof}


\begin{exemple}
\label{ex:Ck_bornee_Banach}
	Soit $\Omega \subset \bR^d$ un ouvert. 
	Alors $C^k_b(\Omega)$ et $C^k_b(\overline \Omega)$ sont des Banach, $\forall k \in \bN$. 
\end{exemple}

\begin{proof}
	Notons $A := \{\alpha\in \bN^d \mid |\alpha|_1\leq k\}$. 
	Soit $(f_n)_{n \in \bN}$ de Cauchy dans $C^k_b(\Omega)$, de sorte que $\partial_\alpha f_n \in C^0_b(\Omega)$ est de Cauchy pour tout $\alpha \in A$.
	On note $f^\alpha := \lim_{n\to \infty} \partial_\alpha f_n \in C^0_b(\Omega)$, et il s'agit de montrer que $f^\alpha = \partial_\alpha f_0$. Par récurrence, il suffit de montrer que $f^{\alpha+e_i} = \frac{\partial}{\partial x_i} f^\alpha$ pour tout $\alpha\in A$ avec $|\alpha|_1 < k$ et tout $1 \leq i \leq d$, où $(e_i)_{1 \leq i \leq d}$ désigne la base canonique. 
	
	Pour établir cette dernière identité, on observe par passage à la limite de la convergence uniforme sous l'intégrale que 
	$$
		f^\alpha(x+\rho e_i) - f^\alpha(x) = \int_0^\rho f^{\alpha+e_i}(x+t e_i) \diff t,
	$$
	pour tout $x \in \Omega$ et tout $\rho>0$ tel que $[x,x+\rho e_i] \subset \Omega$.
	
	Le cas de $C^k_b(\overline \Omega)$ s'en déduit, en notant qu'une limite uniforme de fonctions continues jusqu'au bord reste continue jusqu'au bord. 
\end{proof}



\begin{exemple}
\label{ex:LpBanach}
	Soit $(X,\mu)$ un espace mesuré. Alors $L^p(X,\mu)$ est un Banach, $\forall p \in [1,\infty]$. 
\end{exemple}

\begin{proof}[Preuve \facultatif]
	La preuve est fondée sur des arguments de théorie de la mesure.
	On traite seulement le cas $p\in [1,\infty[$. 
	Compte tenu du \cref{lem:serie_Banach}, on considère $(f_n)_{n \in \bN}$ le terme général d'une série absolument convergente dans $L^p(X,\mu)$. 
	Soit $S_N(x) := \sum_{n<N} |S_n(x)|$, et $S_\infty(x) := \sum_{n\in \bN} |S_n(x)| \in [0,\infty]$, pour tout $x \in X$.
	Alors, par le Théorème de convergence monotone (Beppo Levi) et l'inégalité de Minkowski, voir la \cref{def:normeLp}, on a  
	\begin{equation}		
		\int_X S_\infty^p \diff \mu 
		= \lim_{N \to \infty} \int_X S_N(x)^p \diff \mu(x)
		\leq \lim_{N \to \infty} \Big(\sum_{n < N} \|f_n\|_p \Big)^p
		= \Big(\sum_{n \in \bN} \|f_n\|_p \Big)^p < \infty
	\end{equation}
	On en déduit que $S_\infty(x) < \infty$ pour $\mu$-p.p.\ $x$. 

	On peut alors définir $g_N(x) := \sum_{n <N} f_n(x)$, et $g_\infty(x) := \sum_{n \in \bN} f_n(x)$ par complétude de $\bK$, la somme étant absolument convergente pour $\mu$ presque tout $x$.
	Alors $|g_\infty(x)-g_N(x)| \leq\sum_{n\geq N} |f_n(x)|$, puis par intégration $\|g_\infty-g_N\|_p \leq  r_N := \sum_{n\geq N} \|f_n\|_p$, et $r_N \to 0$ lorsque $N \to \infty$ car c'est le reste d'une série convergente. 
	
\end{proof}


\paragraph{Espaces de Fréchet.}

\begin{remark}
	Considérons un Fréchet, c.a.d un evtlcs $\big(E,(|\cdot|_k)_{k \in \bN}\big)$ muni d'une famille dénombrable de semi-normes, et complet pour la distance $d(x,y) := \max_{k \in \bN} \min\{2^{-k}, |x-y|_k\}$. Une suite $(x_n)_{n \in \bN}$ est de Cauchy (pour la distance $d$) ssi 
	\begin{equation*}
		\forall k \in \bN,\ \forall\ve>0,\ \exists N \in \bN,\ \forall p,q\geq N,\ |x_p-x_q|_k \leq \ve.
	\end{equation*}
	Elle converge vers une limite $x$ ssi $\lim_{n \to \infty} |x-x_n|_k = 0$, pour tout $k \in \bN$. 
\end{remark}


\begin{exemple}
\label{ex:testDistribFrechet}
	Soit $\Omega \subset \bR^d$ un ouvert, et $k \in \bN \cup \{\infty\}$. Alors $C^k(\Omega)$ est un Fréchet muni des semi-normes \eqref{def:cont_nonbornee}. 
\end{exemple}

\begin{proof}
	Soit $(f_n)_{n \in \bN}$ une suite de Cauchy de $C^k(\Omega)$, et soit $(K_r)_{r \in \bN}$ une suite exhaustive de compacts de $\Omega$, et $\Omega_r := \mathring K_r$. % comme dans l'\cref{def:cont_nonbornee}.
	Soit $k' = k$ si $k< \infty$, et soit $k' \in \bN$ arbitraire si $k=\infty$, et soit $r \in \bN$. Alors $(f_{n|\Omega_r})_{n \in \bN}$ converge uniformément ainsi que ses dérivées vers un élément de $C^{k'}_b(\overline{\Omega_r})$, par les arguments de l'\cref{ex:cont_bornee_Banach}.
	Ceci permet de définir une fonction limite $f \in C^k(\Omega)$, puisque $\cup_{r \in \bN} \Omega_r = \Omega$ et que $k'$ est arbitraire.
	On vérifie ensuite que $|f-f_n|_{r,\alpha} \to 0$ lorsque $n \to \infty$, pour tous $r \in \bN$ et $\alpha \in \bN^d$ avec $|\alpha|_1 \leq k$, ce qui conclut. 
%	\cref{def:cont_nonbornee}
\end{proof}

Les preuves des exemples suivants sont similaires et laissées en exercice. 
%Nous donnons deux autres exemples, dont les preuves se déduisent de l'e

\begin{exemple}
\label{ex:DKOmegaComplet}
	Pour tout ouvert $\Omega\subset \bR^d$, et $K \csubset \Omega$,  $C^\infty_b(\Omega)$ et $\cD_K(\Omega)$ sont des Fréchets. 
\end{exemple}

\begin{exemple}
\label{ex:classeSchwartzFrechet}
	La classe de Schwartz $\cS(\bR^d)$ est un Fréchet. 
\end{exemple}

%\begin{proof}
%	Similaire au cas précédent.  
%\end{proof}

\begin{exemple}
\label{ex:LplocFrechet}
	$L^p_\loc(\Omega)$ est un Fréchet, pour tout $p\in [1,\infty]$, $\Omega \subset \bR^d$ ouvert. 
\end{exemple}
	
\notei{Exemple peut être mal choisi. On pourrait traiter $C^\infty_b$ et $L^p_\loc$, pour distinguer le problème du support de celui de l'ordre des dérivées.}


\subsection{Prolongements}
\label{subsec:prolongements}

La notion de complétude fournit des extensions naturelles aux fonctions uniformément continues définies sur des sous-espaces denses, et en particulier aux fonctions linéaires continues. 
Noter que l'image d'une suite de Cauchy par une fonction uniformément continue, est toujours une suite de Cauchy. 

\begin{retenir}
Le principe de prolongement d'une application uniformément continue, et de complétion d'un espace.
(\Cref{prop:prolongement_complet,prop:espaceComplete}.)
\end{retenir}


\begin{proposition}[Prolongement des fonctions uniformément continues]
\label{prop:prolongement_complet}
	Soit $(X,d)$ un espace métrique et $(Y,d_Y)$ un complet. Soit $A \subset X$ une partie dense, et $f:A \to Y$ uniformément continue, pour un module de continuité $\omega$ que l'on peut supposer continu. Alors $f$ admet une unique extension continue $F : X \to Y$, qui de plus admet le même module d'uniforme continuité. 
\end{proposition}

\begin{proof}
	Soit $x\in X$, et soit $(x_n)_{n \geq 0} \in A^\bN$ une suite convergeant vers $x$. Alors $(x_n)_{n \geq 0}$ est de Cauchy, donc $(f(x_n))_{n \geq 0}$ l'est aussi car $f$ est uniformément continue, donc $f(x_n)$ converge lorsque $n \to \infty$ vers une limite notée $y\in Y$. 
	
	Si $(x'_n)_{n \geq 0} \in A^\bN$ est une autre suite arbitraire convergeant vers $x$, alors en notant que $d(f(x_n),f(x'_n)) \leq \omega(d(x_n,x'_n)) \leq \omega(d(x_n,x)+d(x'_n,x))$ on obtient que $f(x'_n) \to y$. On peut donc définir sans ambiguité l'extension par $F(x) = y$. 
	
	Finalement si $x_n \to x$ et $x'_n \to x'$ lorsque $n \to \infty$, avec $x_n,x'_n \in A$ et $x,x'\in X$, alors $d(F(x),F(x')) = \lim_{n \to \infty} d(f(x_n),f(x'_n)) \leq \limsup_{n \to \infty} \omega(d(x_n,x'_n)) = \omega(d(x,x'))$. En effet, on a supposé $\omega$ continu, ce qui ne restreint pas la généralité compte tenu du \cref{lem:regModuleCont}.
\end{proof}


La notion de complétude permet également d'étendre l'espace lui même. 

\begin{proposition}[Complété d'un espace métrique]
\label{prop:espaceComplete}
Soit $(A,d)$ un espace métrique. Alors il existe un espace métrique complet $X$ dont une partie dense est en bijection isométrique avec $A$. De plus $X$ est unique à isométrie près. 
\end{proposition}

\begin{proof}
%	Etant données deux suites de Cauchy, on pose $(x_n)_n \sim (y_n)_n \eqdef d(x_n,y_n) = 0$
	\emph{Existence :} On définit l'espace quotient de suites de Cauchy % , sur l'ensemble des suites de Cauchy, la relation
	\begin{align*}
		X &:= \{(x_n)_{n\in \bN} \in A^\bN \mid (x_n)_n \text{ est de Cauchy}\} / \sim,
	\end{align*}
	muni de la relation d'équivalence $(x_n)_n \sim (y_n)_n  \eqdef \lim_{n \to \infty} d(x_n,y_n) = 0$, et de la distance $d( (x_n)_n, (y_n)_n) := \lim_{n \to \infty} d(x_n,y_n)$. 
	On vérifie facilement qu'il s'agit d'un espace métrique, qui est complet.
	On a l'injection isométrique $A \to X : x \to (x)_{n\in \bN}$ associant à un point la classe d'équivalence de la suite constante.
	
	\emph{Unicité :} Etendre via la \cref{prop:prolongement_complet} l'isométrie définie sur la partie dense $A$. 
\end{proof}

\hiddenproof{
Notons que la bonne définition de la distance découle du critère de Cauchy : 
\begin{equation*}
	|d(x_n,x'_n)-d(x_p,x'_p)| \leq d(x_n,x_p) + d(x'_n,x'_p) \leq \ve(n\wedge p) + \ve'(n\wedge p),
\end{equation*}
de sorte que la suite $(d(x_n,x'_n))_{n \in \bN}$ est de Cauchy dans $\bR$, que l'on sait complet, donc convergente. On a noté $p\wedge q := \min(p,q)$, et introduit les modules de Cauchy $\ve$ et $\ve'$ associés aux deux suites via \eqref{eqdef:suiteCauchyMin}.

Montrons maintenant la complétude de $X$. Considérons une suite de Cauchy $(x^p)_{p \in \bN}$ de $X$, dont des représentants sont les suites $(x^p_n)_{n \in \bN}$, elles mêmes de Cauchy dans $A$. On déduit de ces hypothèses que 
\begin{equation*}
	d(x^p_r,x^q_s) \leq \ve_p(r) + \ve_q(s) + \ve(p \wedge q),
\end{equation*}
pour tous $p,q,r,s \in \bN$,
où $\ve(n) \to 0$ et $\ve_k(n) \to 0$ lorsque $n \to \infty$, quel que soit $k \in \bN$. 
Etant donné $n \in \bN^*$, choisissons $N(n)$ tel que $\ve(n) \leq 1/n$, et $M(n)$ tel que $\ve_{N(n)}(M(n)) \leq 1/n$, puis posons $x^\infty_n := x^{N(n)}_{M(n)}$. Cette suite est de Cauchy car 
\begin{equation*}
	d(x^\infty_p,x^\infty_q) = d(x^{N(p)}_{M(p)}, q^{N(q)}_{M(q)}) \leq \ve_{N(p)}(M(p)) + \ve_{N(q)}(M(q)) + \ve(N(p) \wedge N(q)) \leq \frac 1 p + \frac 1 q + \frac 1 {p\wedge q}.
\end{equation*}
Par ailleurs
\begin{equation*}
	d(x^k_r, x^\infty_r) = d(x^k_r, x^{N(r)}_{M(r)}) \leq \ve_k(r) + \ve_{N(r)}(M(r)) + \ve(k\wedge N(r)) = \ve(k) + o(1),
\end{equation*}
lorsque $r \to \infty$. Ainsi $d_X(x^k ,x^\infty)\leq \ve(k)$ et donc $x^k \to x^\infty$ dans $X$. 
}


On peut également étendre des fonctions uniformément continues définies sur des parties non-denses, comme l'illustre le résultat suivant, qui s'applique en particulier aux fonctions Holderiennes et Lipschitziennes. L'hypothèse de sous-additivité du module de continuité n'est pas exagérément restrictive, voir le \cref{lem:regModuleCont}.

\begin{exemple}[Extension de fonctions uniformément continues]
\label{ex:extUnifCont}
	Soit $(X,d)$ un espace métrique borné, soit $A \subset X$ non-vide, et soit $f : A \to \bR$ uniformément continue, dont le module de continuité $\omega$ est sous-additif. 
	Alors la fonction $F : X \to \bR$ définie par 
	\begin{equation}
		F(x) := \inf_{y \in A} f(y) + \omega(d(x,y))
	\end{equation}
	prolonge $f$ à $X$ et admet le module de continuité $\omega$. 
\end{exemple}

\begin{proof}
	On note que $F$ est à valeurs finies car $\omega$ est sous additif et $X$ est borné. 
	On a $F(x) \leq f(x)$ pour tout $x\in A$ en choisissant $y=x$, mais par ailleurs $f(x)-f(y) \leq \omega(d(x,y))$ pour tout $y\in A$ donc $F(x) \geq f(x)$ sur $A$, d'où $F_{|A} = f$. 
	
	Comme $\omega(d(x,z)) \leq \omega(d(x,y)) + \omega(d(y,z))$ par sous-additivité, 
	on obtient $F(z) = \inf_{y \in A} f(y) + \omega(d(y,z)) \leq F(x) + \omega(d(x,z))$, donc $|F(x)-F(z)| \leq \omega(d(x,z))$ par symétrie, ce qui conclut. 
\end{proof}


\subsection{Point fixe de Picard}
\label{subsec:Picard}

Le thèorème de point fixe de Picard permet affirme qu'une application contractante d'un espace complet dans lui même admet un point fixe. Il assure également son unicité de ce point fixe, et donne un moyen de l'approcher (par itération) et d'estimer l'erreur commise. C'est un résultat extrêmement robuste, à comparer au \nref{th:Brouwer} qui assure seulement l'existence d'un point fixe. Une application immédiate et fondamentale est le \nref{th:Cauchy_Lipschitz} d'existence des solutions d'EDOs. 


\begin{retenir}
L'énoncé et la preuve du théorème de point fixe de Picard et du lemme de Grönwall. L'énoncé du théorème de Cauchy-Lipschitz. 
(\Cref{th:fixe_Picard,lem:Gronwall,th:Cauchy_Lipschitz}.)	
\end{retenir}

\begin{theorem}[Point fixe de Picard]
\label{th:fixe_Picard}
	Soit $(X,d)$ un espace métrique complet, et soit $f : X \to X$ une fonction $K$-Lipschitzienne avec $K<1$ (on dit aussi contractante). 
	
	Alors $f$ admet un unique point fixe $x_* \in X$, et pour tout $x_0 \in X$ et tout $n \in \bN$ on a 
	\begin{align}
		d(x_n,x_*) &\leq K^n \frac{d(x_0,x_1)}{1-K},&
		\text{où} \quad x_n := f^{\circ n}(x_0).
	\end{align}
\end{theorem}

\begin{proof}
Par récurrence immédiate on obtient $d(x_n,x_{n+1}) \leq K^n d(x_0,x_1)$, pour tout $n \geq 0$. Ainsi, pour tout $N \geq 0$ et tous $p,q\geq N$, on a  $d(x_p,x_q) \leq \sum_{n\geq N} K^n d(x_0,x_1) = K^n d(x_0,x_1)/(1-K)$. On en déduit que $(x_n)_{n \geq 0}$ satisfait le critère de Cauchy, et est donc convergente. Puis $d(x_N,x_*) = \lim_{n \to \infty} d(x_N,x_n) \leq \sum_{n\geq N} K^n d(x_0,x_1) = K^N d(x_0,x_1)/(1-K)$. Noter que $x_*$ est bien un point fixe par continuité de $f$, et qu'il est unique par contractivité.
\end{proof}

Le \nref{th:fixe_Picard} permet d'établir celui de Cauchy-Lipschitz, c.à.d.\ l'existence, l'unicité, et la stabilité d'une solution $C^1$ de l'équation différentielle ordinaire (EDO)
\begin{equation}
\label{eq:EDO}
	u'(t)=f(t,u(t)),
\end{equation}
avec donnée initiale $u(0)$ fixée, et $f$ assez régulière. On en rappelle la preuve ci-dessous.

\begin{lemma}[Grönwall]
\label{lem:Gronwall}
	Soit $f \in C^0([0,T], \bR_+)$, et $A,B \geq 0$ tels que pour tout $t \in [0,T]$
	\begin{equation}
		f(t) \leq A \int_0^t f(s) \diff s + B.	
	\end{equation}
	Alors $f(t) \leq B \exp(A t)$ pour tout $t \in [0,T]$. 
\end{lemma}

\begin{proof}
	Posons $F(t) := A \int_0^t f(s) \diff s + B$, $t \in [0,T]$, de sorte que $f \leq F$ sur $[0,T]$. On note que $F\in C^1$ et que $F'(t) = A f(t) \leq A F(t)$, donc $t\in [0,T] \mapsto F(t)\exp(-A t)$ est décroissante. Ainsi $f(t) \leq F(t) \leq F(0) \exp(At) = B \exp(At)$, pour tout $t \in [0,T]$, comme annoncé. 
\end{proof}

La notation ``$K \csubset \Omega$'' signifie que $K$ est une partie compacte d'un espace topologique $\Omega$. 

\begin{theorem}[Cauchy-Lispchitz, Picard-Lindelöf]
\label{th:Cauchy_Lipschitz}
Soit $\Omega\subset \bR^d$ ouvert, et soit $f : \bR_+ \times \Omega \to \bR^d$, continue et localement Lipschitzienne en sa seconde variable, i.e.\
	\begin{equation}
	\label{eq:loc_lip2}
		\forall T \geq 0,\ \forall K \subset_c \Omega,\ \exists C,\ \forall t\in [0,T],\ \forall x,y \in K,\ 
		\|f(t,x) - f(t,y)\| \leq C\|x-y\|.
	\end{equation}
	Alors:
	\begin{itemize}
		\item (Existence) Pour tout $x_0 \in \Omega$, il existe $t_0>0$ et une solution $u \in C^1( [0,t_0], \Omega)$ de l'EDO \eqref{eq:EDO}, avec condition initiale $u(t_0) = x_0$. 
		\item (Stabilité) Si $u,v \in C^1( [0,T], K)$ sont des solutions de \eqref{eq:EDO}, où $T>0$ et $K\csubset \Omega$ est un compact, alors $\|u(t)-v(t)\| \leq e^{C t} \|u(0)-v(0)\|$  où la constante $C$ dépend seulement de $T$ et $K$. 
	\end{itemize}
\end{theorem}

\begin{proof}
	\emph{Existence}.
	Soit $r>0$ tel que $B'(x_0,r) \subset \Omega$. Soit $t_1 > 0$ et soit $C_{\mathrm{lip}}$ une constante de Lipschitz pour $f$ par rapport à la seconde variable sur $[0,t_1] \times B'(x_0,r)$, et $C_\infty$ une borne pour $f$ sur cet ensemble. 
	Choisissons $t_0 \in ]0,t_1]$ tel que $t_0 C_\infty \leq r$ et $t_0 C_{\mathrm{lip}} <1$. 
	Considérons l'espace métrique complet $X := C^0([0,t_0],B'(x_0,r))$, muni de la norme uniforme, et posons pour tout $u \in X$ 
	\begin{equation*}
		F (u) (t) := x_0 + \int_0^t f(s,u(s)) \diff s.
	\end{equation*}
	Les estimations suivantes montrent que $F : X \to X$ est bien définie, et est contractante : 
	\begin{align*}
		\|F u(t) -x_0 \| &\leq \int_0^t \|f(s,u(s))\| \diff s \leq t_0 C_\infty \leq r,\\
		\|F u(t) - F v(t)\| &\leq \int_0^t \|f(s,u(s)) - f(s,v(s))\|\diff s \leq t_0 C_{\mathrm{lip}} \|u-v\|_\infty.
	\end{align*}
	Par le \cref{th:fixe_Picard}, $F$ admet un point fixe, qui est la solution recherchée. 
	
	\emph{Stabilité}.
	Pour tout $t \in [0,T]$ on a 
	\begin{equation*}
		\| u(t)-v(t)\| - \|u(0)-v(0)\| \leq \int_0^t  \|u'(s)-v'(s) \| \diff s  = \int_0^t \|f(t,u(t))-f(t,v(t))\|\diff t \leq C \int_0^t \|u-v\| \diff t, 
	\end{equation*}
	où $C = C(T,K)$ est une constante de Lipschitz de $f$ par rapport à la deuxième variable sur $[0,T] \times K$. La stabilité s'en déduit via le Lemme de Grönwall appliqué à $t\in [0,T] \mapsto \| u(t)-v(t)\|$.
\end{proof}

On déduit du \cref{th:Cauchy_Lipschitz} qu'il existe une solution dite maximale\footnote{L'ensemble des solutions est muni de l'ordre (a priori partiel) associé au prolongement. Le point (stabilité) du \cref{th:Cauchy_Lipschitz} montre que cet ordre est en fait total, ce qui permet de construire un plus grand élément.}  $u \in C^1([0,T_{\max}[,\Omega)$ de \eqref{eq:EDO}, où $T_{\max} \in ]0,\infty]$, dont toutes les autres solutions sont des restrictions.
%que toutes les solutions de \eqref{eq:EDO} sont des restrictions d'une solution définie sur un interval maximal. 
Noter qu'un autre résultat, le \nref{th:CauchyPeano} qui est fondé sur le point fixe de Brouwer, permet d'établir l'existence - mais pas l'unicité, ni a fortiori la stabilité - d'une solution de \eqref{eq:EDO} sans l'hypothèse \eqref{eq:loc_lip2}. L'énoncé précédent peut être complété dans le cas des systèmes autonomes.



\begin{corollary}[Alternative dans les systèmes autonomes]
\label{corol:CauchyLipAuton}
	Soit $\Omega \subset \bR^d$ ouvert, et soit $f : \Omega \to \bR^d$ localement Lipschitzienne, i.e.\ $\forall K\csubset \Omega, \exists C, \forall x,y \in K, \|f(x)-f(y)\| \leq C \|x-y\|$.
	Soit $x_0 \in \Omega$, et soit $u : [0,T_{\max}[ \to \Omega$ la solution maximale de l'EDO $u'(t) = f(u(t))$ (c'est à dire ne pouvant pas être prolongée), où $T_{\max} \in \bR_+ \cup \{\infty\}$.
	Alors on a l'alternative:
	\begin{itemize}
		\item (Existence globale) Le temps d'existence est $T_{\max} = \infty$. 
		\item (Sortie de tout compact) Pour tout compact $K \csubset \Omega$ il existe $0 \leq T_K < T_{\max}$ tel que $u(t) \notin K$ pour tout $t \in [T_K,T_{\max}[$. 
	\end{itemize} 
\end{corollary}

\begin{proof}
%\todoi{N'utilise pas particulièrement la complétude, plutôt la compacité, en particulier pour l'existence de $r$ (distance atteinte, compacité locale). Déplacer.}
Soit $K \csubset \Omega$ un compact non-vide, et soit $r>0$ tel que $\tilde K := \{x \in \bR^d \mid d(x,K) \leq r\} \csubset \Omega$.
En notant $C_{\infty}$  une borne uniforme pour $f$ sur $\tilde K$, et $C_{\mathrm{lip}}$ une constante de Lispchitz pour $f$ sur $\tilde K$, on obtient par le raisonnement précédent un temps d'existence $t_K = \min(r/C_\infty,1/(2C_{\mathrm{lip}}))>0$ pour l'EDO considérée, valable pour toute condition initiale appartenant à $K$. 
En particulier, si le temps d'existence $T_{\max}$ est borné, alors $u(t) \notin K$ pour tout $t \geq T_{\max}-t_K$, comme annoncé.
\end{proof}

Le schéma de preuve du \nref{th:Cauchy_Lipschitz} est robuste et permet de traiter de nombreuses variantes, comme des EDO avec retard.
\begin{exemple}[Equation différentielle avec retard]
\label{ex:EDretard}
	Il existe une unique $f\in C^1([0,1],\bR)$ telle que $f(0)=1$ et $f'(t)=f(t-t^2)$ pour tout $t \in [0,1]$. 
\end{exemple}

\begin{proof}
On pose pour tout $f \in E := C^0([0,1/2],\bR)$ et tout $t\leq 1/2$. 
\begin{equation}
\label{eq:EDO_retard_fixe}
A(f)(t) := 1+\int_0^t f(s-s^2) \diff s
\end{equation}
On vérifie facilement que $A$ est $1/2$-Lipschitz sur le Banach $E$, muni de la norme $\|\cdot\|_\infty$. Ainsi elle admet un unique point fixe par le \nref{th:fixe_Picard}, que dans la suite on note $f\in E$. Observons que $f\in C^1$ en tant que primitive d'une fonction continue.
Finalement, on étend $f$ sur $]1/2,1]$ via la formule explicite (\ref{eq:EDO_retard_fixe}, droite), qui ne dépend que de $f_{|[0,1/2]}$.
\end{proof}


De même, on montre l'existence et l'unicité de solutions à certaines équations fonctionnelles. 
\begin{exemple}[Equation fonctionnelle intégrale]
\label{ex:eqFunctIntegr}
	Soient $k \in C^0([0,1]\times [0,1],]-1,1[)$ et $\vp \in C^0([0,1])$. Alors il existe une unique solution $f\in C^0([0,1])$ à l'équation fonctionnelle
	\begin{equation*}
		 \forall x \in [0,1],\ f (x) = \int_0^1 \frac{k(x,y) f(y)}{1+f(y)^2} \diff y + \vp(x)
	\end{equation*}
%	\begin{align*}
%		f &= A(f),&
%		\text{ où } \forall x \in [0,1], A(f)(x) &:= \int_0^1 \frac{k(x,y) f(y)}{1+f(y)^2} \diff y + \vp(x).
%	\end{align*}
\end{exemple}

\begin{proof}
	Notons $A(f)(x)$ le membre de droite, et $E := C^0([0,1])$ qui est complet pour la norme $\|\cdot\|_\infty$. 
	Soit également $K := \|k\|_{L^\infty([0,1]^2)}$, notons que $K<1$ par compacité, et montrons que $A : E \to E$ est $K$-Lipschitzienne donc contractante, ce qui permet de conclure par le \nref{th:fixe_Picard}. 
	La fonction $\psi(s):= s/(1+s^2)$ satisfait $\psi'(s) = (1-s^2)/(1+s^2)^2$, donc est $1$-Lipschitzienne. On en déduit comme annoncé que 
	\begin{equation*}
		| A(f)(x)-A(g)(x)| \leq \int_0^1 |k(x,y)| |\psi(f(y))-\psi(g(y))| \diff y 
		\leq \|k\|_\infty \|f-g\|_\infty.
		\qedhere
	\end{equation*} 	
\end{proof}


\subsection{Théorème de Baire}
\label{subsec:Baire}

Le théorème de Baire fournit une notion d'ensemble \emph{maigre}, stable par union dénombrable, dans tout espace complet. On peut faire l'analogie avec les parties de mesure nulle d'un espace mesuré, mais ces notions sont en fait assez différentes comme l'illustre l'\cref{ex:GrasQ}, et le théorème de Baire a l'avantage de s'appliquer en dimension infinie, voir la \cref{subsec:BaireLin}.

\begin{retenir}
Le lemme des fermés emboités, le théorème de Baire, et son corollaire s'appliquant aux fermés d'intérieur vide. 
(\Cref{lem:fermeEmboite,th:BaireOuvert,corol:BaireFerme}.)	
\end{retenir}


\begin{lemma}[Fermés emboités]
\label{lem:fermeEmboite}
Soit $(X,d)$ un espace métrique complet, et soit $(F_n)_{n \in \bN}$ une suite de fermés non-vides, décroissants pour l'inclusion, et dont le diamètre $\delta_n := \sup \{d(x,y) \mid x,y \in F_n\}$ tend vers $0$. Alors $\bigcap_{n \in \bN} F_n$ est un singleton. 
\end{lemma}

\begin{proof}
	Choisissons $x_n \in F_n$ arbitraire pour tout $n \geq 0$. Alors pour tout $N \geq 0$, $p,q\geq N$, on a $x_p,x_q \in F_N$ et donc $d(x_p,x_q) \leq \delta_N$. Comme $\delta_N\to 0$ lorsque $N \to \infty$, le critère de Cauchy est satisfait.
\end{proof}

L'hypothèse que le diamètre tend vers $0$ est nécessaire dans le \cref{lem:fermeEmboite}, comme l'illustre le contre exemple $F_n := [n,\infty[ \subset \bR$. 
Un résultat analogue pour les fermés d'un espace compact permet en revanche de s'en passer, voir la \cref{rem:fermeBorelLeb}. La propriété établie au \cref{lem:fermeEmboite} est en fait équivalente à la complétude de l'espace.

\begin{theorem}[Baire]
\label{th:BaireOuvert}
Soit $(X,d)$ un espace métrique complet, et soit $(U_n)_{n \in \bN}$ une suite d'ouverts denses. 
Alors leur intersection $\bigcap_{n \in \bN} U_n$ est dense. 	
\end{theorem}

\begin{proof}
	Soit $x_0\in U_0$ et soit $\ve_0>0$ tel que $B'(x_0,\ve_0) \subset U_0$. On choisit par récurrence, pour tout $n\in \bN$: (i) $x_{n+1} \in B(x_n,\ve_n) \cap U_{n+1}$, et (ii) $\ve_{n+1} \in ]0,2^{-n}]$ tel que $B'(x_{n+1},\ve_{n+1}) \subset B(x_n,\ve_n) \cap U_{n+1}$. 

	On note que les fermés $B'(x_n,\ve_n)$ sont emboités et de diamètre au plus $2\ve_n\to 0$, donc leur intersection contient un point $x_*$ par le \cref{lem:fermeEmboite}.
	De plus $x_* \in \bigcap_{n \in \bN} U_n$, et $\|x_0-x_*\| \leq \ve_0$. Comme $U_0$ est dense, et  $x_0\in U_0$ et $\ve_0>0$ sont arbitraires, ceci conclut. 
\end{proof}

Par passage au complémentaire, on obtient:

\begin{corollary}[Baire, complémentaire]
\label{corol:BaireFerme}
	Soit $(X,d)$ un espace métrique complet, et soit $(F_n)_{n \in \bN}$ une suite de fermés d'intérieur vide. Alors $\bigcup_{n \in \bN} F_n$ est d'intérieur vide. 
\end{corollary}

\begin{remark}[Terminologie de Baire]
\label{rem:terminologieBaire}
Une intersection dénombrable d'ouverts est appelée un $G_\delta$, et une réunion dénombrable de fermés un $F_\sigma$. Un ensemble est dit maigre s'il est contenu dans un $F_\sigma$ d'intérieur vide, et son complémentaire est dit co-maigre (ou gras) : il contient un $G_\delta$ dense. 	
\end{remark}
 
Attention : le fait d'être maigre ou gras, pour une partie de $\bR$, n'a pas d'implication sur la mesure de celle-ci, comme le montre l'exemple suivant. 
\begin{exemple}
\label{ex:GrasQ}
	Soit $(q_k)_{k >0}$ une énumération de $\bQ$, et pour tout $n >0$ soit 
	\begin{equation}
	U_n := \bigcup_{k >0} \Big]q_k-\frac 1 {nk^2}, q_k + \frac 1 {n k^2}\Big[.	
	\end{equation}
	Alors $U_n$ est un ouvert dense de $\bR$, de mesure $|U_n| \leq \pi^2/(6n)$. En particulier, $U_\infty :=  \cap_{n \in \bN} U_n$ est un ensemble gras de mesure nulle, contenant $\bQ$, et son complémentaire est un ensemble maigre de mesure pleine. 
\end{exemple}

\begin{exemple}	
\label{ex:grasNonDen}
	Soit $(X,d)$ un espace métrique complet et sans points isolés, et soit $A \subset X$ un ensemble gras. Alors $A$ est indénombrable. 
\end{exemple}

\begin{proof}
	Comme $X$ est sans points isolés, les complémentaires de singleton $X\sm\{x\}$ sont des ouverts denses. Si $A$ était dénombrable, alors $A \cap \bigcap_{x \in A} X \sm \{x\}$ serait à la fois vide et gras, ce qui contredit \nref{th:BaireOuvert}. 
\end{proof}


\subsection{Applications de Baire aux opérateurs linéaires continus} %Conséquences du th.\ 
\label{subsec:BaireLin}

Le théorème de Baire a plusieurs conséquences importantes pour les applications linéaires entre espaces de Banach.

\begin{retenir}
	Le théorème de Banach-Steinhauss, de l'application ouverte, et du graphe fermé.
	(\Cref{th:BanachSteinhauss,th:appliOuverte,th:grapheFerme}.)
\end{retenir}


\begin{theorem}[Banach-Steinhauss]
\label{th:BanachSteinhauss}
Soit $E$ un Banach, $F$ un evn, et $\cA \subset L(E,F)$ une famille d'applications linéaires continues, que l'on suppose simplement bornée : 
\begin{equation}
	\forall x \in E,\ \sup_{L\in \cA} \|Lx\| < \infty.
\end{equation}
Alors $\cA$ est uniformément bornée :
%\begin{equation}
$	\sup_{L \in \cA} \vertiii{L} < \infty.$
%\end{equation}
\end{theorem}

\begin{proof}
	Posons, pour tout $k \in \bN$, 
	\begin{equation}
		E_k := \{x \in E \mid \forall L \in \cA, \|Lx\| \leq k\}.
	\end{equation}
	Alors $E_k$ est un fermé, et $\bigcup_{k \in \bN} E_k = E$, donc par le \nref{corol:BaireFerme}, l'un au moins des $E_k$ est d'intérieur non-vide. Par homogénéité, $E_1$ est d'intérieur non-vide, donc contient une boule ouverte $B(x,r)$, où $x \in E$ et $r>0$. Par symétrie, $B(-x,r) \subset E_1$, et par convexité $B(0,r) \subset E_1$. Par définition de la norme d'opérateur, on a donc $\vertiii{L} \leq r^{-1}$ pour tout $L \in \cA$, ce qui conclut.
\end{proof}

\begin{corollary}
\label{corol:limSimpleLin}
	Soient $E$ un Banach, $F$ un evn, et soit $L_n \in L(E,F)$, $n\in \bN$, une suite d'applications linéaires continues telles que $Lx := \lim_{n \to \infty} L_n x$ existe pour tout $x \in E$. Alors $L\in L(E,F)$ est linéaire continue.
\end{corollary}

\begin{proof}
	La linéarité de la limite $L$ est claire. Comme $(L_n x)_{n \in \bN}$ est convergente, pour tout $x \in E$, elle est bornée, et donc $(\vertiii{T_n})_{n \in \bN}$ est bornée par le \nref{th:BanachSteinhauss}, par une constante $C$. Par passage à la limite on obtient $\|Tx\| \leq C \|x\|$, pour tout $x \in E$, donc $\vertiii{T} \leq C$ comme annoncé.
\end{proof}


Le corollaire suivant applique Banach-Steinhauss aux formes linéaires. Il admet un analogue dual  le \cref{corol:faibleSepare}, dont la preuve requiert le \nref{th:HahnBanach} démontré ultérieurement.  

\begin{corollary}[Banach-Steinhauss, pour les formes linéaires]
\label{corol:BanachSteinhaussK}
	Soient $E$ un Banach, et $\cA \subset E^*$ un ensemble de formes linéaires continues. Alors $\cA$
	 %est linéairement borné pour la topologie *faible 
	est *faiblement borné
	(i.e.\ $\forall x \in E, \{l(x)\mid x \in A\}$ est borné) ssi $\cA$ est borné. 
\end{corollary}

\begin{proof}
Choisir $F := \bK$ dans le \nref{th:BanachSteinhauss}.
%	La \cref{def:lineairBorne} de bornitude dans un evtlc, et celle \cref{def:topoFaible} des semi-normes définissant la topologie faible, nous placent précisément dans le cadre du \cref{th:BanachSteinhauss}.
\end{proof}

\begin{exemple}
\label{ex:FourierDiverge}
	Il existe une fonction $f\in C^0(\bT, \bC)$, où $\bT := \bR/(2 \pi \bZ)$, dont la série de Fourier diverge en zéro. 
\end{exemple}

\begin{proof}
	La somme partielle de la série de Fourier de $f$ en $0$ est donnée par 
	\begin{equation*}
		L_n(f) 
		:= \sum_{|n| \leq N} \int_\bT f(t) \exp(\bi n t) \diff t
		= \int_\bT f(t) D_n(t) \diff t, 
	\end{equation*}
	où $\bi \in \bC$ est le nombre imaginaire unité, et 
	où on a noté
	\begin{equation*}
		D_n(t) := \sum_{|n| \leq N} \exp(\bi n t) = \frac{\sin((2N+1) t/2)}{\sin (t/2)}.
	\end{equation*}
	Rappelons que l'espace des fonctions continues périodiques $C^0(\bT,\bC)$ est un Banach, muni de la norme $\|\cdot\|_\infty$. Ainsi
	\begin{equation*}
		\vertiii{L_n} = \|D_n\|_{L^1(\bT)} = \int_{-\pi}^{\pi} \frac{|\sin((2n+1)t/2)|}{|\sin(t/2)|} \diff t
		\geq \int_{-\pi}^{\pi} |\sin((2n+1)t/2)| \frac{\diff t} {|t|}
		= \int_{-(2n+1)\pi/2}^{(2n+1)\pi/2} |\sin(s)| \frac{\diff s} {|s|},
	\end{equation*}
	et la dernière intégrale diverge clairement lorsque $t \to \infty$ lorsque $n \to \infty$. 
	Le résultat se découle de la contraposée du \cref{corol:BanachSteinhaussK}. 
\end{proof}


Le \nref{th:BanachSteinhauss} s'étend aux les espaces de Fréchet, voir \cite[Theorem 1.33]{alazard2017analyse} pour une discussion plus approfondie.

\begin{theorem}[Banach-Steinhauss dans des Fréchets \facultatif]
\label{th:BanachSteinhaussFrechet}
	Soient $(E,(|\cdot|_i)_{i \in I})$ un Fréchet, et $(F,(|\cdot|)_{j \in J})$ un evtlcs, avec $I$ et $J$ dénombrables. Soit $\cA\subset L(E,F)$ une famille d'applications linéaires continues, supposée simplement linéairement bornée : $\forall x\in E$, $\forall j \in J$, $\sup_{u \in \cA} |u(x)|_j < \infty$. Alors $\cA$ est équicontinue : il existe un module de continuité $\omega$ tel que $\forall x,y \in E$, $\forall u \in \cA$, $d_F(u(x),u(y)) \leq \omega(d_E(x,y))$, où les distances sont définies par \eqref{eqdef:evtlc_dist}.
\end{theorem}

\begin{proof}
	La preuve mime celle du cas d'un Banach : étant fixé $j \in J$, on note $E_k := \{x\in E \mid \forall L \in \cA,\ |Lx|_j \leq k\}$ qui est un fermé. Le théorème de Baire implique que l'un de ces fermés est d'intérieur non-vide, disons $E_{k_j}$. Par symmétrie et convexité, il contient un voisinage de l'origine, que l'on peut choisir de la forme $V := \{x \in E \mid \forall i \in I_j,\ |x|_i \leq \ve_j\}$ où $\ve_j>0$ et $I_j \subset I$ est fini. On en déduit que $\forall L \in \cA$, $\forall x \in E$, $\forall j \in J$, $|Lx|_j \leq (k_j/\ve_j) \max_{i \in I_j} |x|_i$. 
	
	Compte tenu de l'expression \eqref{eqdef:evtlc_dist} des distances, on obtient un module de continuité $\omega$ tel que $\forall x \in E$, $\forall u \in \cA$, $d_F(u(x),0) \leq \omega(d_E(x,0))$. Par linéarité de $u$, il satisfait aux conditions de l'énoncé. 
\end{proof}

\hiddenproof{
Détaillons la construction du module de continuité.
Disons que $|u(x)|_j \leq C_j \max_{i \leq N_j} |x|_i$, et quitte à augmenter $C_j$ supposons que $C_j 2^{-N_j} \geq 2^{-j}$.  
On choisit le module de continuité $\omega(s) := \max_{j \in \bN} \min\{2^{-j}, C_j s\}$, qui tend bien vers $0$ lorsque $s \to 0$.
Considérons $x \in E$, $j\in \bN$, et posons $w := \omega(d_E(x,0))$. Si $w \geq 2^{-j}$, alors $\min\{2^{-j}, |u(x)|_j\} \leq 2^{-j} \leq w$ comme souhaité. Si $w < 2^{-j}$, alors $C_j \min \{2^{-i},|x|_i\} < w$ pour tout $i\in \bN$ par définition de $d_E$, donc $C_j |x|_i < w$ pour tout $i \leq N_j$, donc $\min\{2^{-j}, |u(x)|_j\} \leq |u(x)|_j \leq w$ comme souhaité également. Ainsi $d_F(u(x),0) := \max_{j\in \bN} \min \{2^{-j}, |u(x)|_j\} \leq w$ ce qui conclut. 
}


Le théorème de l'application ouverte est une seconde application du \nref{th:BaireOuvert} aux applications linéaires, qui utilise la complétude des espaces de départ \emph{et} d'arrivée. 

\begin{theorem}[Application ouverte]
\label{th:appliOuverte}
Soient $E$ et $F$ des Banach et soit $L \in L(E,F)$ linéaire continue et surjective. 
Alors $\exists \delta>0$, $B_F(0,\delta) \subset L(B_E(0,1))$, et donc l'image par $L$ de tout ouvert est un ouvert. %; de manière équivalente, $\exists C$, $\forall y\in E$, $\exists x \in E$, $Lx = y$ et $\|x\| \leq C \|y\|$.

%\todoi{Ou encore : $T(B_E(0,1))\supset B_F(0,\delta)$ avec $\delta>0$.}
\end{theorem}

\begin{proof}
	Par surjectivité, on a $\cup_{n \in \bN} \overline {L(B_E(0,n))} = F$. Par le \nref{corol:BaireFerme} l'un de ces fermés est d'intérieur non-vide, et par symétrie et convexité il contient un voisinage de l'origine. 
	Il existe donc $n \in \bN$ et $r>0$ tels que $B_F(0,r) \subset \overline{L(B_E(0,n))}$; notons $C_0 := n/r$.
	Par homogénéité, il existe pour tout $y\in F$ et tout $\ve>0$, un point $x \in E$ tel que $\|x\| \leq C_0 \|y\|$ et $\|y-Lx\| \leq \ve$; la seconde partie de la preuve permet de remplacer cette dernière approximation par une égalité.
	
	Etant donné $y_0 \in F\sm \{0\}$, pour tout $n \geq 0$ par récurrence, (i) on choisit $x_n\in E$ tel que $\|y_n-Lx_n\| \leq \|y_n\|/2$ et $\|x_n\| \leq C_0 \|y_n\|$, puis (ii) on pose $y_{n+1} := y_n-L x_n$. Par une récurrence immédiate, $\|y_n\| \leq 2^{-n} \|y_0\|$, donc $\|x_n\| \leq C_0 2^{-n} \|y_0\|$. Par complétude de $E$, la série $x_* := \sum_{n \in \bN} x_n$ converge, et par sommation des normes $\|x_*\| \leq 2 C_0 \|y_0\|$.
	Finalement, on note que $\|y_0 - L\sum_{n<N} x_n\| = \|y_N\| \to 0$ lorsque $N \to \infty$, donc $\|y_0-Lx_*\|=0$ par continuité de $L$.
	
	Nous avons obtenu : $\forall y\in E$, $\exists x \in E$, $Lx = y$ et $\|x\| \leq C \|y\|$, où $C := 2 C_0$. 
	On en déduit que $L(B(0,1)) \supset B(0,1/C)$ comme annoncé. Puis $L(B(x,r)) \supset B(Lx,r/C)$ par linéarité, et donc l'image d'un voisinage de $x\in E$ arbitraire est un voisinage de $Lx$, ce qui conclut. 
\end{proof}


\begin{corollary}[Isomorphisme de Banach]
\label{corol:isomBanach}
	Soient $E,F$ des Banachs et soit $L \in L(E,F)$ linéaire continue et bijective. Alors $L^{-1} \in L(F,E)$ est continue.
\end{corollary}

\begin{proof}
Par le \cref{th:appliOuverte}, l'image directe par $L$ d'un ouvert est un ouvert, donc l'image réciproque par $L^{-1}$ d'un ouvert est un ouvert. On conclut par la \nref{prop:caract_cont}.
\end{proof}

\begin{corollary}[Equivalence des normes]
\label{corol:equivNormes}
	Soit $E$ un espace vectoriel, qui est complet pour deux normes $\|\cdot\|$ et $\|\cdot\|'$ données. 
	S'il existe $C$ tel que $\forall x \in E$, $\|x\|' \leq C \|x\|$, alors il existe $c>0$ tel que $\forall x \in E$, $\|x\|' \geq c \|x\|$.
\end{corollary}

\begin{proof}
	Appliquer le \cref{corol:isomBanach} à $\Id : (E,\|\cdot\|) \to (E,\|\cdot\|')$.
\end{proof}


\begin{theorem}[Graphe fermé]
\label{th:grapheFerme}
	Soient $(E,\|\cdot\|_E)$ et $(F,\|\cdot\|_F)$ des Banach, et soit $L : E \to F$ linéaire. Alors $L$ est continue ssi son graphe $G := \{ (x,Lx) \mid x \in E\}$ est une partie fermée de $E \times F$. 
\end{theorem}

\begin{proof}
Rappelons que $E \times F$ est un Banach muni de la norme $\|(x,y)\|_{E \times F} := \|x\|_E+\|y\|_F$.
Si $L$ est continue, alors $G = \{ (x,y) \in E \times F	\mid y - Lx = 0\}$ est fermé en tant qu'image réciproque de $\{0\}$ par l'application linéaire continue $(x,y)\in E \times F \mapsto y-Lx \in F$. 


Réciproquement, supposons $G$ fermé; c'est donc un Banach. La projection $\Pi : (x,y) \in G \mapsto x \in E$ est continue et bijective, donc par le \nref{corol:isomBanach} sa réciproque est  continue. Ainsi il existe une constante $C$ tel que $\forall x \in E$, $C \|x\|_E \geq \|\Pi^{-1} x\|_{E \times F} = \|x\|_E + \|Lx \|_F$, ce qui implique $\|Lx\| \leq (C-1)\|x\|$ et conclut.
\end{proof}

\begin{corollary}
	Soit $E$ un Banach, et soit $E = V\oplus W$ une décomposition en sev supplémentaires. La projection $P : E\to V$ parallèlement à $W$ est continue ssi $V$ et $W$ sont fermés.
\end{corollary}

\begin{proof}
	Si $P$ est continue, alors $W := \Ker P$ et $V := \Ker (\Id-P)$ sont fermés. Réciproquement, supposons $V,W$ fermés.
	Soit $(x_n,v_n)$ une suite du graphe de $P$, donc $x_n \in E$ et $v_n := P x_n$, convergeant vers $(x,v)\in E$. Alors $x_n \in V$ donc $v\in V$ par fermeture de cet espace, et par ailleurs $x_n-v_n \in W$ donc $x-v\in W$ par fermeture également. On en déduit que $v=P x$, et donc que $P$ est continue d'après le théorème du graphe fermé.
\end{proof}


\newpage
\section{Compacité}
\label{sec:compact}
La compacité permet fondamentalement d'obtenir des résultats d'existence, via le théorème des bornes atteintes, ou le théorème de point fixe de Brouwer; l'unicité des objets construits devra, le cas échéant, être établie par d'autres méthodes. La compacité est également au coeur de certains résultats d'approximation, par des fonctions simples dans le cadre du théorème de Stone-Weierstrass, ou par des opérateurs de rang fini dans le cadre linéaire. 


\inventaire{
\Nref{subsec:axiomBorelLeb} \\
\Nref{def:compact}\\
\Nref{rem:fermeBorelLeb}\\
\Nref{lem:compImplFerme} (Une partie compacte est fermée)\\
\Nref{corol:compSsiFerme} (Une partie d'un compact est compacte ssi elle est fermée)\\
\Nref{prop:compactImage}\\
\Nref{corol:homeoCompact} (Une bijection continue entre compacts est un homéomorphisme)\\
\Nref{def:locCompact}\\
\Nref{def:compactAlex}\\

\Nref{subsec:compactMetric}\\
\Nref{def:precompact}\\
\Nref{th:BolzanoWeierstrass}\\
\Nref{th:Heine}\\
\Nref{prop:noAdhCv} (Si une seule valeur d'adhérence possible, alors convergence)\\


\Nref{subsec:compactRd}\\
\Nref{prop:compSsiFermeBorne} (Un fermé borné de $\bR^d$ est compact)\\
\Nref{corol:bornesAtteintes}\\
\Nref{corol:normesEquivRd}\\
\Nref{lem:normaleRiesz}\\
\Nref{th:bouleCompRiesz}\\

\Nref{subsec:compactProd}\\
\Nref{th:compProdTycho}\\
\Nref{ex:formuleLogique}\\
\Nref{th:BanachAlaoglu}\\
\Nref{th:AscoliLip}\\
\Nref{corol:AscoliEquicont}\\
\Nref{corol:AscoliCk}\\
\Nref{corol:AscoliCkNonBorne}\\
\Nref{prop:hullCompact}\\
\Nref{th:Caratheodory}\\

\Nref{subsec:Brouwer}\\
\Nref{th:Brouwer}\\
\Nref{th:chgVarIntegr}\\
\Nref{lem:chgVarLax}	\\
\Nref{rem:chgVarDiffeo} (Formule du changement de variable sans difféomorphisme)\\
\Nref{lem:noretractSmooth} (Lemme de non-rétractation, version régulière)\\
\Nref{corol:noretract}\\

\Nref{subsec:Schauder}\\
\Nref{th:Schauder}\\
\Nref{rem:SchauderEvt} (Schauder dans les evt)\\
\Nref{th:CauchyPeano}\\
\Nref{rem:preuveCauchyPeano}\\

\Nref{subsec:StoneWeierstrass}\\
\Nref{th:StoneWeierstrass}\\
\Nref{lem:StoneUnite} ($\overline A$ est une algèbre unitaire complète)\\
\Nref{lem:StoneSqrt} ($\overline A$ contient $\sqrt f, |f|, 1/f$)\\
\Nref{corol:StoneMinMax} ($\overline A$ contient $\max(f,g)$ et $\min(f,g)$).\\
\Nref{lem:StoneDirac} (Approximation d'un Dirac dans $A$)\\
\Nref{corol:StoneComplexe} \\
\Nref{rem:polDens}\\
\Nref{prop:StoneCtrEx}\\

\Nref{subsec:linCompact}\\
\Nref{def:linCompact} (Opérateur compact)\\
\Nref{prop:opCompactComposition}\\
\Nref{prop:opCompactLimite}\\
\Nref{prop:opCompactComplCont}\\
\Nref{def:linSpectre} (Spectre d'un opérateur)\\
\Nref{th:Fredholm}\\
\Nref{lem:convK_compact}\\
\Nref{corol:eqFuncNoyau} (Équation fonctionnelle à noyau)\\
\Nref{corol:FredholmPrecise}\\ % -> th ?
%\Nref{def:propriApprox}\\
%\Nref{prop:limRangFini}
}

\subsection{Axiome de Borel-Lebesgue} % caractérisation topologique
\label{subsec:axiomBorelLeb}


\begin{retenir}
	L'axiome de Borel-Lebesgue et sa reformulation pour les fermés, l'image directe d'un compact par une application continue est un compact, la définition d'un espace localement compact séparé.
	(\Cref{def:compact,rem:fermeBorelLeb,prop:compactImage,def:locCompact})
\end{retenir}


\begin{definition}[Compact]
\label{def:compact}
\leancheck{def:compact}%
	Un espace topologique $(X, \bU)$ est dit compact s'il est \emph{séparé} et satisfait l'\emph{axiome de Borel-Lebesgue} : de toute couverture de $X$ par des ouverts on peut extraire une sous-couverture finie, i.e.\ pour tout $\cU \subset \bU$
	\begin{equation*}
	%\forall \cU \subset \bU, \Big(
	\cup\, \cU=X \Rightarrow 
	\exists\, \cU_0\subset \cU,\ \cU_0 \text{ est fini et } \cup\cU_0 = X.
	%\Big).
	\end{equation*}
\end{definition}

Une partie d'un espace topologique est dite compacte si c'est un compact pour la topologie induite. Elle est dite \emph{relativement compacte} si son adhérence est compacte (dans un espace métrique complet, c'est équivalent à la \emph{précompacité}, voir les \cref{def:precompact} et \cref{th:BolzanoWeierstrass}).

Par passage au complémentaire de l'axiome de Borel-Lebesgue, on obtient une propriété des fermés. 
\begin{remark}[Borel-Lebesgue pour les fermés]
\label{rem:fermeBorelLeb}
\leancheck{rem:fermeBorelLeb}%
Si $(X,\bU)$ est compact, et si $(F_i)_{i \in I}$ est une famille de fermés dont l'intersection est vide, alors elle admet une sous famille finie dont l'intersection est vide. 	

En particulier, une intersection décroissante de fermés non-vides d'un compact est non-vide (comparer avec \nref{lem:fermeEmboite} dans un complet). Notamment, l'ensemble des valeurs d'adhérence \eqref{eqdef:adherence_suite} d'une suite à valeurs dans un compact est toujours non-vide.
\end{remark}

Le \cref{corol:compSsiFerme} suivant illustre l'importance de l'hypothèse de séparation dans la \cref{def:compact}, voir la \nref{def:separe}. 

\begin{lemma}
\label{lem:compImplFerme}
\leancheck{lem:compImplFerme}%
	Soit $(X,\bU)$ espace topologique séparé. Si $F\subset X$ est compact, alors $F$ est fermé.
\end{lemma}

\begin{proof}
Montrons que le complémentaire de $F$ est ouvert. Soit $x \in X \sm F$, et pour tout $y \in F$ soient $V_y\ni x$ et $W_y \ni y$ des ouverts disjoints de $X$, par l'axiome de séparation.
Comme $F \subset \cup_{y\in F} W_y$, il existe par compacité $y_1,\cdots,y_N \in F$ tels que $F \subset \cup_{1 \leq n \leq N} W_{y_n}$. On note alors que $\cap_{1 \leq n \leq N} V_{y_n}$ est un ouvert disjoint de $F$ et contenant $x$.
\end{proof}

\begin{corollary}	
\label{corol:compSsiFerme}
\leancheck{corol:compSsiFerme}%
	Soit $(X, \bU)$ espace compact. Une partie $F \subset X$ est compacte ssi elle est fermée. 
\end{corollary}

\begin{proof}
	Supposons $F$ fermé, l'implication réciproque ayant été traitée dans le \cref{lem:compImplFerme}. Soit $\cU$ une couverture de $F$ par des ouverts. Alors $\cU \cup \{X\sm F\}$ est une couverture de $X$ par des ouverts, dont on peut extraire une sous-couverture finie, d'où l'on tire une couverture finie de $F$ en omettant le dernier ensemble $X \sm F$. 
\end{proof}


Alors que les ouverts et fermés sont stables par image \emph{réciproque} continue, voir la \nref{prop:caract_cont}, les compacts sont eux stables par image \emph{directe} continue.
Il en découlera une observation élémentaire mais fondamentale en optimisation : le \nref{corol:bornesAtteintes}. La notation $K \csubset X$ signifie que $K$ est une partie compacte d'un espace topologique $X$. 

\begin{proposition}[Image continue directe d'un compact]
\label{prop:compactImage}
\leancheck{prop:compactImage}%
	Soient $(X, \bU)$ et $(Y,\bV)$ des espaces séparés, et soit $f \in C^0(X,Y)$. 
	Alors pour tout $K \csubset X$, l'image directe $f(K)$ est compacte. 
\end{proposition}

\begin{proof}
	La séparation de $f(K)$ s'hérite de celle de $Y$. Par ailleurs si $\cV\subset \bV$ est une couverture de $f(K)$ par des ouverts, alors $\cU := \{f^{-1}(V) \mid V \in \cV\}$ est par continuité de $f$ une couverture de $K$ par des ouverts. 
	On peut donc en extraire une sous-couverture finie $(f^{-1}(V_i))_{i=1}^I$, d'où l'on tire une couverture finie $(V_i)_{i=1}^I$ de $f(K)$. 
\end{proof}

\begin{corollary}
\label{corol:homeoCompact}
\leancheck{corol:homeoCompact}%
	Soient $(X, \bU)$ et $(Y,\bV)$ des espaces compacts, et soit $f : X \to Y$ une bijection continue. Alors $f^{-1}$ est continue. 
\end{corollary}

\begin{proof}
Étant donné $F\subset X$ fermé, on note que $F$ est compact par le \cref{corol:compSsiFerme}, donc $f(F)$ est compact par la \cref{prop:compactImage}, donc $f(F)$ est fermé par le sens réciproque du \cref{corol:compSsiFerme}. 
Le résultat s'ensuit par la \nref{prop:caract_cont}. 
%Au vu de la caractérisation des
\end{proof}


Certains espaces (séparés) sont \emph{localement compacts}, l'exemple fondamental étant celui des ouverts de $\bR^d$, compte tenu de la compacité des boules fermées de $\bR^d$, voir la  \cref{prop:compSsiFermeBorne}.

L'hypothèse de séparation est importante dans \cref{def:locCompact}, tout comme dans \nref{def:compact}, et en son absence les propriétés (i-ii) ne sont pas équivalentes. (Il n'y a d'ailleurs pas de consensus sur l'éventuelle notion d'espace localement compact non séparé.)


\begin{definition}[Espace localement compact séparé]
\label{def:locCompact}
\leancheck{def:locCompact}%
Un espace topologique \emph{séparé} est dit localement compact s'il satisfait l'une des conditions suivantes, qui sont équivalentes:
\begin{enumerate}[(i)]
	\item Tout point admet un voisinage compact.
	\item Tout point admet une base de voisinages compacts. 
\end{enumerate}
\end{definition}

\begin{proof}[Preuve de l'équivalence]
	L'implication (ii) $\Rightarrow$ (i) étant claire, nous supposons donc (i) et montrons (ii).
	Soit $x \in X$, $K$ un voisinage compact de $x$, et $V$ un voisinage quelconque de $x$. On peut supposer $V\subset K$, quitte à considérer son intersection avec $K$, et on peut supposer $V$ ouvert, quitte à considérer son intérieur. 
	Pour tout $y \in K \sm \{x\}$, soient $V_y$ et $W_y$ des voisinages ouverts disjoints de $x$ et $y$; noter que $V_y \cap W_y = \emptyset$. 
	Alors $\{V\} \cup \{W_y \mid y \in K \sm \{x\} \}$ est une couverture de $K$, dont on peut extraire une couverture finie $K \subset V \cup_{i=1}^I W_{y_i}$. 
	Ainsi $K \sm \cup_{i=1}^I W_{y_i} \supset V \cap \cap_{i=1}^I V_{y_i}$ est un voisinage compact de $x$ contenu dans $V$, ce qui conclut. 
\end{proof}

%L'hypothèse de séparation est importante dans \cref{def:locCompact}


On vérifie facilement qu'un espace compact privé d'un point est localement compact. Récipro\-quement, 
un espace localement compact peut être compactifié en lui rajoutant seulement un point, ``à l'infini''. 

%\notei{Loc compact ssi compact privé d'un point}

\begin{definition}[Compactifié d'Alexandrov]
\label{def:compactAlex}
\leancheck{def:compactAlex}%
	Soit $(X, \bU)$ un espace topologique localement compact séparé. 
	Son compactifié d'Alexandrov est l'espace $\tilde X := X \sqcup \{\omega\}$, où $\omega$ est un élément additionnel arbitraire, muni de la topologie $\tilde \bU := \bU \cup \{\tilde X \sm K \mid K \csubset X\}$. 
\end{definition}

\hiddenproof{
Un compact $X$ privé d'un point $x$ est localement compact. D'abord, $X\sm \{x\}$ est séparé. Ensuite, étant donné $y \in X \sm \{x\}$, choisissons deux voisinages ouverts disjoints $U$ et $V$ de $x$ et $y$, par séparation dans $X$; ainsi $X \sm U \supset V$ est un voisinage compact de $y$ ne contenant pas $x$, comme souhaité pour la condition (i). 
}

\hiddenproof{
Vérifions les axiomes de topologie. (i) $\emptyset\in \bU\subset \tilde \bU$ et $\tilde X = \tilde X \sm \emptyset \in \tilde \bU$. (ii) La topologie $\bU$ est stable par union, et $\cup_{i \in I} (\tilde X\sm K_i) = \tilde X \sm \cap_{i \in I} K_i$ est de la forme souhaitée car une intersection de compacts est un compact. Finalement $U \cup (\tilde X \sm K) = \tilde X \sm (K\sm U)$ est de la forme souhaitée car $K \sm U$ est un fermé de $K$ donc un compact.
(iii) La topologie $\bU$ est stable par intersection de deux éléments, et $(\tilde X \sm K) \cap (\tilde X \sm K') = \tilde X \sm (K \cup K')$ est de la forme souhaitée puisque $K \cup K'$ est compact. Enfin $U \cap (\tilde X \sm K) = U \cap (X \sm K)$ est l'intersection de deux ouverts de $X$ (on a vu que tout compact est fermé, donc son complémentaire dans $X$ est ouvert). 

Vérifions que $\tilde X$ est compact, via l'axiome de Borel-Lebesgue. Considérons une couverture de $\tilde X$ par des ouverts. Notons que l'un d'entre eux contient $\omega$, donc prend la forme $X \sm K$. On peut supposer, quitte à retirer $\omega$ s'il est présent, que les autres membres de la couverture sont des ouverts $(U_i)_{i \in I}$ de $X$. Alors $K \subset \cup_{i \in I} U_i$, donc on peut extraire une sous-couverture finie $U_{i_1},\cdots,U_{i_N}$ de $K$. On conclut en notant que $U_{i_1},\cdots,U_{i_N}, \tilde X \sm K$ est une couverture finie de $\tilde X$.  
}


%\todoi{
%La remarque sur propriété locale vs sur tout compact ne s'applique pas bien à la Lipschitzianité, car on a besoin de chemins de longueur correctement bornée passant par l'intersection. 
%Du coup, je ne sais pas si elle est intéressante. 
%}

%\begin{remark}[Propriété valant localement, ou sur tout compact]
%	Une propriété $P$ sur un ensemble $X$ est une fonction booléenne $P : \cP(X) \to \{0,1\}$. 
%	On supposera $P$ stable par union : $P(A) = P(B) = 1 \Rightarrow P(A \cup B)$, et décroissante pour l'inclusion : $P(A)$ et $B\subset A$ $\Rightarrow$ $P(B)$. 
%	Si $X$ es localement compact, alors les conditions suivantes sont équivalentes : 
%	\begin{itemize}
%		\item $P$ est satisfaite localement, i.e. tout point admet un voisinage satisfaisant $P$.
%		\item $P$ est satisfaite sur tout compact. 
%	\end{itemize}
%\end{remark}
%
%\begin{proof}
%	Trivial. Mais pas clair que ça s'applique aux exemples d'intérêt (localement Lipschitz)
%\end{proof}

\subsection{Compacts métriques}
\label{subsec:compactMetric}
La présence d'une métrique apporte de nouvelles caractérisations de la compacité, en particulier la caractérisation séquentielle de Bolzano-Weierstrass qui est particulièrement maniable. % et que nous illustrons par le théorème de Heine d'uniforme continuité des applications continues sur un compact. 

\begin{retenir}
La triple caractérisation des compacts métriques, le théorème de Heine.
(\Cref{th:BolzanoWeierstrass,th:Heine}.)
\end{retenir}

%On dispose dans un espace métrique de caractérisations alternatives de la compacité.

\begin{definition}[Précompacité]
\label{def:precompact}
\leancheck{def:precompact}%
	Un espace métrique $(X,d)$ est dit pré-compact ssi pour tout $\ve>0$ il admet une couverture finie par des boules ouvertes de rayon $\ve$. 
\end{definition}


\begin{remark}
\label{rem:precompact_separable}
\leancheck{rem:precompact_separable}%
	Tout espace précompact est séparable. %, i.e.\ il admet une partie dénombrable et dense.
	
	Un espace est précompact ssi pour tout $\ve>0$ il peut être couvert par un nombre fini de parties, chacune de diamètre au plus $\ve$. En particulier, lorsque l'on étudie la précompacité d'une partie $X \subset \bX$ d'un espace métrique englobant, il n'est pas nécessaire que les boules formant la couverture soient centrées en des points de $X$.
\end{remark}
\begin{proof}
	\emph{Preuve du premier point.} Considérer pour tout $n\geq 0$ une couverture par un nombre fini de boules de rayon $2^{-n}$, puis collecter les centres de ces boules, ce qui forme une partie dénombrable et dense.
	
	\emph{Preuve du second point.} Une boule de rayon $\ve$ constitue par l'inégalité triangulaire une partie de diamètre au plus $2 \ve$. Réciproquement, une partie de diamètre $\ve$ est contenue dans une boule fermée de rayon $\ve$, centrée en un point quelconque de cette partie, donc dans une boule ouverte de rayon $2 \ve$ centrée au même point. Le résultat s'ensuit.
\end{proof}
 
\begin{theorem}[Compacité via Bolzano-Weierstrass]
\label{th:BolzanoWeierstrass}
\leancheck{th:BolzanoWeierstrass}%
Dans un espace métrique $(X,d)$, les propriétés suivantes sont équivalentes : 
\begin{enumerate}[(i)]
	\item $X$ est compact (i.e.\ vérifie l'axiome de Borel-Lebesgue, voir la \cref{def:compact}).
	\item Toute suite à valeurs dans $X$ admet une valeur d'adhérence (et donc une sous-suite convergente); c'est l'axiome de Bolzano-Weierstrass.
	\item $X$ est pré-compact et complet.
\end{enumerate}
\end{theorem}

\begin{proof}
	Notons que $X$ est séparé en tant qu'espace métrique. 
	
	\emph{Preuve de (i) $\Rightarrow$ (ii).} L'existence de valeurs d'adhérence est observée dans la \cref{rem:fermeBorelLeb} (valable dans un espace compact général), et la convergence d'une sous-suite vers celles-ci dans la \cref{prop:adhSousSuite} (utilise la métrique). 

	\emph{Preuve de (ii) $\Rightarrow$ (iii).} \emph{Complétude :} Si une suite est de Cauchy, alors elle admet une valeur d'adhérence par (ii), donc une sous suite convergente, donc est convergente par les propriétés générales des suites de Cauchy. 
	\emph{Précompacité :}
	Soit $\ve>0$, et pour tout $N \geq 0$ tant que c'est possible choisissons un point  
	$x_N \in X \sm \cup_{n<N} B(x_n,\ve)$ de manière récurrente. Si la suite $(x_n)_{n\in \bN}$ est définie en totalité, alors elle admet une sous-suite convergente, donc deux termes à distance strictement inférieure à $\ve$, ce qui est une contradiction. On en déduit que la suite ne peut pas être construite au delà d'un certain rang $N$, et donc que $X = \cup_{n<N} B(x_n,\ve)$. 
	
	\emph{Preuve de (iii) $\Rightarrow$ (ii).}
%	Soit $A := \{y_r \mid r \in \bN\}$ l'ensemble des valeurs d'une suite
	Soit une suite $(y_r)_{r \in \bN}$ à valeurs dans $X$, et soit $A := \{y_r \mid r \in \bN\}$ l'ensemble de ses valeurs. Si $A$ est fini, alors il existe une sous-suite constante, donc convergente; on suppose donc $A$ infini dans la suite, sans nuire à la généralité.
	Pour tout $n\geq 0$ on se donne un entier $K(n)$ et des boules $B_1^n, \cdots, B_{K(n)}^n$ de rayon $2^{-n}$ recouvrant $X$. 
	Puis on choisit successivement $B_{k_n}^n$ tels que $A \cap B_{k_0}^0 \cap \cdots \cap B_{k_n}^n$ est infini, et donc contient un terme $y_{\sigma(n)}$ où $\sigma : \bN \to \bN$ est strictement croissante.
	On voit alors que la suite $(y_{\sigma(n)})_{n \in \bN}$ satisfait le critère de Cauchy, spécifiquement $d(y_{\sigma(p)},y_{\sigma(q)}) \leq 2 \times 2^{-N}$ si $p,q \geq N$, car ces points appartiennent à une même boule de rayon $2^{-N}$. Par complétude, on obtient une sous-suite convergente. 
	
	Preuve de (ii) $\Rightarrow$ (i). 
	Soit $(U_i)_{i \in I}$ une couverture ouverte de $X$. Observons d'abord que la propriété de Bolzano-Weierstrass implique l'existence d'un \emph{nombre de Lebesgue} $r>0$ tel que $\forall x\in X$, $\exists i \in I$, $B(x,r)\subset U_i$. 
	
	En effet, dans le cas contraire, soit $(x_n)_{n \in \bN}$ tel que $B(x_n,2^{-n})$ n'est contenu dans aucun des $U_i$, $i \in I$. Soit $x_*$ une valeur d'adhérence de cette suite, $i \in I$ tel que $x_* \in U_i$, et $r>0$ tel que $B(x_*,r) \subset U_i$. On obtient une contradiction en considérant $n\in \bN$ tel que $d(x_n,x_*) < r/2$ et $2^{-n} < r/2$.
	
	On conclut en utilisant la pré-compacité, qui découle de (ii) $\Rightarrow$ (iii) vu plus haut. 
\end{proof}

\begin{theorem}[Heine]
\label{th:Heine}
\leancheck{th:Heine}%
Soit $(X,d_X)$ métrique compact, et $(Y,d_Y)$ métrique. Alors toute fonction continue $f : X \to Y$ est en fait uniformément continue. 
\end{theorem}

\begin{proof}
	Par contraposée, supposons $f$ non-uniformément continue. Alors il existe $\ve>0$ et une suite $(x_n,y_n)$ telle que $d_X(x_n,y_n) \to 0$ et $d_Y(f(x_n),f(y_n)) \geq \ve$ pour tout $n \in \bN$. Par compacité, on a une sous-suite convergente $x_{\sigma(n)} \to x_*$ lorsque $n \to \infty$; on en déduit que $y_{\sigma(n)} \to x_*$ également.
	Comme $\max \{d_Y(f(x_n),f(x_*)),\ d_Y(f(y_n),f(x_*))\} \geq d_Y(f(x_n),\ f(y_n))/2 \geq \ve/2$, pour tout $n \in \bN$, on obtient que $f$ n'est pas continue en $x_*$, comme annoncé.
%	Enfin, l'équicontinuité (resp. uniforme) de $(f_i)_{i \in I}$ est équivalente à celle de $F : X \to \prod_{i \in I} Y : x \mapsto (f_i(x))_{i \in I}$ pour la distance de la convergence uniforme, voir 
\end{proof}

Sous les hypothèses du \cref{th:Heine}, toute famille de fonctions équicontinue $(f_i : X \to Y)_{i \in I}$ est en fait uniformément équicontinue. 
Cela découle de la reformulation de l'équicontinuité en termes de continuité usuelle, voir la \cref{rem:equicont}. 	


\begin{proposition}
\label{prop:noAdhCv}
\leancheck{prop:noAdhCv}%
	Soit $(x_n)_{n \in \bN}$ une suite à valeurs dans un espace métrique compact $X$, n'admettant pas de valeur d'adhérence excepté peut-être $x_*\in X$. Alors $x_n \to x_*$ lorsque $n \to \infty$. 
\end{proposition}

\begin{proof}
	Si par l'absurde $x_n \not \to x_*$, alors il existe $\ve>0$ et une extractrice $\vp$ tels que $d(x_{\vp(n)},x_*) \geq \ve$ pour tout $n \in \bN$. 
	La suite $(x_{\vp(n)})_{n \in \bN}$ est ainsi à valeur dans le compact $X \sm B(x_*,\ve)$, donc admet une valeur d'adhérence $y_*\in X \sm B(x_*,\ve)$ voir la \cref{rem:fermeBorelLeb}, qui est aussi valeur d'adhérence de $(x_n)_{n \in \bN}$ et satisfait $d(x_*,y_*) \geq \ve$, d'où la contradiction.
\end{proof}

\subsection{Compacité en dimension finie}
\label{subsec:compactRd}

On rappelle la compacité de la boule unité de $\bR^d$, sa conséquence l'unicité de la structure d'evn en dimension finie à équivalence de normes près, et l'obstruction en dimension infinie révélée par le théorème de Riesz. 

\begin{retenir}
La caractérisation des parties compactes de $\bR^d$, les corollaires des bornes atteintes et d'équivalence des normes, le théorème de compacité de Riesz. 
(\Cref{prop:compSsiFermeBorne,corol:bornesAtteintes,corol:normesEquivRd,th:bouleCompRiesz})
\end{retenir}

\begin{proposition}
\label{prop:compSsiFermeBorne}
\leancheck{prop:compSsiFermeBorne}%
Une partie de $\bR^d$ est compacte ssi elle est fermée et bornée.	
\end{proposition}

\begin{proof}
	On sait que $\bR^d$ est complet, et donc une partie de $\bR^d$ est complète ssi elle est fermée.
	Par ailleurs, une partie de $\bR^d$ est précompacte ssi elle est bornée. 
	(En effet, on pourra recouvrir cette partie par les boules de rayon $\ve > 0$ dont les centres sont disposés aux sommets d'un cube discrétisé $[-M,M]^d \cap (\ve/\sqrt d) \bZ^d$, en choisissant $M$ assez grand.)
	%Note : supposons $A \subset [-M,M]$ et  supposons $A \subset \cup_{i=1}^I B(x_i,\ve \sqrt d)$. Notons $y_i$ un point de $A \cap B(x_i,\ve)$ si cet ensemble est non-vide, ou un point quelconque sinon. Alors $A \subset \cup_{i=1}^I B(y_i,2\ve)$.)	
%	(Note : on peut trivialement recouvrir une partie bornée de $\bR^d$ par un nombre fini de boules de rayon $\ve>0$ disposées sur la grille cartésienne $(\ve/\sqrt d) \bZ^d$. donné non centrées en des points de cette partie, donc par un nombre fini de boules de rayon $2\ve$ centrées en des points de cette partie.)
	On conclut par le critère (iii) du \cref{th:BolzanoWeierstrass}.
\end{proof}


On déduit de la \cref{prop:compSsiFermeBorne} qu'un compact non-vide de $\bR$ contient ses bornes, ce qui a des applications immédiates en optimisation. 

\begin{corollary}[Bornes atteintes]
\label{corol:bornesAtteintes}
\leancheck{corol:bornesAtteintes}%
	Une application continue d'un compact non-vide dans $\bR$ est bornée et atteint ses bornes.
\end{corollary}

\begin{proof}
	L'image de l'application est un compact de $\bR$ par la \nref{prop:compactImage}, donc contient ses bornes par la \cref{prop:compSsiFermeBorne}.
\end{proof}

\begin{corollary}[Équivalence des normes en dimension finie]
\label{corol:normesEquivRd}
\leancheck{corol:normesEquivRd}%
	Soit $E$ un espace vectoriel de dimension finie, muni de deux normes $\|\cdot\|$ et $\|\cdot\|'$.
	Alors il existe $C,c>0$ tels que $c\|x\| \leq \|x\|' \leq C \|x\|$ pour tout $x \in E$. 
\end{corollary}

\begin{proof}
	On peut supposer, sans nuire à la généralité, que $E = \bR^d$ et $\|x\| := \sum_{i=1}^d |x_i|$. 
	Alors 
	\begin{align*}
	\|x\|' &= \Big\| \sum_{1 \leq i \leq d} x_i e_i\Big\|' \leq \sum_{1 \leq i \leq d} |x_i| \|e_i\|' \leq C \|x\|,&
	\text{où } C := \max_{1 \leq i \leq d} \|e_i\|'.	
	\end{align*}
	Par l'inégalité triangulaire on a $|  \|x\|'-\|y\|' | \leq \|x-y\|' \leq C \|x-y\|$, donc la seconde norme $\|\cdot\|'$ est $C$-Lipschitz, donc en particulier continue. La constante $c := \inf\{ \|x\|' \mid \|x\| = 1\}$ est atteinte en tant que minimum d'une fonction continue sur un compact, par le \cref{corol:bornesAtteintes}, donc $c>0$, et cette constante convient pour la borne inférieure. 
\end{proof}

Le théorème de compacité de Riesz est une forme de réciproque à la \cref{prop:compSsiFermeBorne} : il établit que la boule unité fermée n'est compacte que dans un evn de dimension finie. % que seuls les evn de dimension finie ont une boule unité fermée compacte.
\begin{lemma}[Riesz]
\label{lem:normaleRiesz}
\leancheck{lem:normaleRiesz}%
	Soit $E$ un evn, soit $F\subsetneq E$ un sous espace vectoriel strict fermé. 
	Pour tout $\rho<1$, il existe $x \in B'_E(0,1)$ tel que $d(x,F) \geq \rho$.
	Si $F$ est de dimension finie, alors il existe $x \in B'_E(0,1)$ tel que $d(x,F)=1$.  
\end{lemma}

\begin{proof}
	Choisissons $u \in E \sm F$, et observons que $d(u,F) := \inf_{v \in F} \|u-v\|$ est strictement positif par fermeture de $F$. 
	Si $F$ est de dimension finie, alors par compacité de $F \cap B'(u,d(u,F)+1)$ il existe $v \in F$ tel que $d(u,F) = \|u-v\|$, et on pose $\rho=1$. À l'inverse, si $F$ est de dimension infinie, on peut trouver $v \in F$ tel que $d(u,F) \geq \rho\|u-v\|$ par définition de l'infimum. 
	Alors $x := \frac{u-v}{\|u-v\|} \in B'_E(0,1)$ convient, puisque $d(x,F) =  d(u/\|u-v\|,F) = d(u,F)/\|u-v\| \geq \rho$. 
\end{proof}

\begin{theorem}[Compacité de Riesz]
\label{th:bouleCompRiesz}
\leancheck{th:bouleCompRiesz}%
Pour tout evn $E$, les propositions suivantes sont équiva\-lentes : (i) $E$ est de dimension finie, (ii) $B'_E(0,1)$ est compacte, (iii) $B'_E(0,1) \subset \cup_{1 \leq i \leq I} B_E(x_i,1)$ pour une famille finie $x_1,\cdots,x_I \in E$. 
\end{theorem}

\begin{proof}
	Si $E$ est de dimension finie (i), alors par le \cref{corol:normesEquivRd} sa topologie est celle de $\bR^d$, et donc $B'_E(0,1)$ est compacte car fermée et bornée, d'où (ii). L'implication (ii) $\Rightarrow$ (iii) découle de l'axiome de Borel-Lebesgue et de la couverture triviale par des boules ouvertes.
	
		Finalement, supposons (iii), et par l'absurde $E$ de dimension infinie. Soit $F := \Vect\{x_1,\cdots,x_I\}$. Par le \cref{lem:normaleRiesz} il existe $x \in B_E'(0,1)$ tel que $1 = d(x,F) \leq \min_{1 \leq i \leq I} d(x,x_i)$. C'est une contradiction, ce qui conclut. 
\end{proof}


\subsection{Produits de compacts}
\label{subsec:compactProd}

Le résultat fondamental de cette section est la compacité d'un produit de compacts, pour la topologie associée à la convergence simple. On en déduit la compacité pour les topologies adéquates de certains ensembles de fonctions : la boule unité du dual d'un evn, et la famille des fonctions $1$-Lipschitziennes entre deux compacts, entre autres. %en particulier. 

\begin{retenir}
Le théorème de Tychonov, le procédé d'extraction diagonal, le théorème de Banach-Alaoglu, le théorème d'Ascoli.
(\Cref{th:compProdTycho,th:BanachAlaoglu,th:AscoliLip}.)
\end{retenir}


\begin{theorem}[Tychonov]
\label{th:compProdTycho}
\leancheck{th:compProdTycho}%
Soit $(X_i,\bU_i)_{i \in I}$ une famille d'espaces topologiques compacts, alors le produit $X_* := \prod_{i \in I} X_i$ est compact muni de la topologie produit.
Le produit d'une famille dénombrable de compacts métriques est un compact métrique.
\end{theorem}

\begin{proof}[Preuve dans le cas métrisable dénombrable]
	Le résultat vaut dans le cadre général, mais sa preuve fait intervenir des notions (ultrafiltres) qui sortent du cadre de ce cours. On se contentera donc de l'établir dans le cas d'une famille dénombrable $(X_i,d_i)_{i \in \bN}$ de compacts métriques.
	
	La topologie produit correspond alors à la distance suivante : pour $(x_i)_{i \in \bN}, (y_i)_{i \in \bN} \in X_*$,
	\begin{equation*}
		d( (x_i)_{i \in \bN}, (y_i)_{i \in \bN}) := \max_{i \in \bN} \min\{2^{-i}, d_i(x_i,y_i)\}.
	\end{equation*}
	Considérons une suite à valeurs dans $X_*$, dont on écrit le $k$-ième terme $(x_i^k)_{i \in \bN}$. Par la propriété de Bolzano-Weierstrass, on peut trouver $\sigma_0 : \bN \to \bN$ strictement croissante telle que $(x_0^{\sigma_0(k)})_{k \in \bN}$ converge dans $X_0$. 
	De même, on construit pour tout $i \in \bN$ une extractrice $\sigma_i : \bN \to \bN$ telle que 
	$(x_i^{\sigma_0\circ \cdots \circ \sigma_i(k)})_{k \in \bN}$ converge dans $X_i$.
	On utilise alors le procédé d'extraction diagonal, en définissant 
	$\vp(k) := \sigma_0 \circ \cdots \circ \sigma_k(k)$, et en notant que $(x_i^{\vp(k)})_{k \in \bN}$ converge dans $X_i$ pour tout $i \in \bN$, de sorte que la suite dans l'espace produit est convergente. 
\end{proof}

En choisissant $X_i = X$ dans le théorème précédent, où $X$ est un espace compact, on obtient en particulier que les fonctions à valeurs dans un espace compact forment un ensemble compact pour la convergence simple.
On illustre ce résultat par une application aux formules logiques, sans rapport avec le reste du cours mais qui est un point de départ en théorie des modèles
\begin{exemple}[Satisfiabilité de formules logiques \facultatif]
\label{ex:formuleLogique}
\leancheck{ex:formuleLogique}%
Une formule logique est une application $f : X := \{0,1\}^\bN \to \{0,1\}$ qui ne dépend que d'un nombre fini de variables : 
$$
\exists N(f) \in \bN,\ \forall x=(x_0,x_1,\cdots)\in X,\
f(x) = f(x_0,\cdots,x_{N(f)},0,\cdots,0).
$$
Un ensemble $\cF$ de formules logiques est dit satisfiable si $\exists x \in X$, $\forall f \in \cF$, $f(x)=1$. 

Un ensemble de formules logiques est satisfiable ssi toutes ses parties finies sont satisfiables. 
%Si toute partie finie d'un ensemble $\cF$ de formules logiques est satisfiable, alors $\cF$ est satisfiable : 
\end{exemple}
	
\begin{proof}
	Par définition de la topologie produit sur $X$, les formules logiques sont des applications continues, et par le \cref{th:compProdTycho} l'ensemble $X$ est compact. On procède par contraposée, en supposant que $\cF$ n'est pas satisfiable, c'est-à-dire que $\cap_{f \in \cF} f^{-1}(1)= \emptyset$. Il s'agit d'une intersection vide de fermés, donc une sous famille finie est d'intersection vide, ce qui conclut.  
\end{proof}

\begin{theorem}[Banach-Alaoglu]
\label{th:BanachAlaoglu}
\leancheck{th:BanachAlaoglu}%
La boule unité fermée $B := B'_{E^*}(0,1)$ du dual $E^*$ d'un Banach est compacte pour la topologie $*$faible. 
Si $E$ est séparable, alors $B$ est un compact métrique.
\end{theorem}

\begin{proof}[Preuve dans le cas séparable]
On rappelle que $B$ est métrisable, voir la \cref{prop:bouleMetrisable}. Il suffit donc de montrer que toute suite $(f_n)_{n \in \bN}$ de $B$ admet une sous-suite convergente, par le \nref{th:BolzanoWeierstrass}. 
Soit $D$ une partie dénombrable dense de $E$. Par un procédé d'extraction diagonal, ou en utilisant la compacité de $\Pi_{x\in D} [-\|x\|,\|x\|]$ ce qui revient au même, on obtient qu'il existe une extractrice $\vp : \bN \to \bN$ telle que $(f_{\vp(n)}(x) )_{n \in \bN}$ converge pour tout $x \in D$.

On peut donc supposer, pour alléger les notations et quitte à extraire, que $(f_n)_{n \in \bN}$ converge simplement sur $D$ vers une application $f_*$.
On note que $\|x-y\| \geq |f_n(x)-f_n(y)| \to |f_*(x)-f_*(y)|$ lorsque $n \to \infty$, pour tous $x,y \in D$, donc $f_*$ est $1$-Lipschitzienne  sur $D$. Elle s'étend donc en une application $1$-Lipschitzienne sur $\overline D = E$, par la \nref{prop:prolongement_complet}, toujours notée $f_*$. 

Montrons maintenant que $f_n$ converge simplement vers $f_*$ sur $E$, c'est-à-dire que l'on a convergence en topologie *faible comme demandé, ce qui est aussi noté $f_n \rightharpoonup^* f_*$. En effet, étant donné $x\in E$, et $\ve>0$, on peut choisir $x_0 \in D$ tel que $\|x-x_0\| \leq \ve$, et alors 
\begin{equation*}
	|f_*(x)-f_n(x)|\leq 	|f_*(x)-f_*(x_0)| + |f_*(x_0)-f_n(x_0)| + |f_n(x_0)-f_n(x)| \leq 2 \ve + |f_*(x_0)-f_n(x_0)|,
\end{equation*}
ce qui est plus petit que $3 \ve$ pour $n$ assez grand. 
La linéarité de $f_*$ se déduit de la convergence simple : $f_*(\lambda x+y) = \lim_{n \to \infty} f_n(\lambda x+y) = \lim_{n \to \infty} \lambda f_n(x)+f_n(y) = \lambda f_*(x)+f_*(y)$, pour tous $x,y \in E$, $\lambda \in \bK$, ce qui conclut. 
\end{proof}

%\begin{proof}[Cas général]
%\def\lin{\mathrm{lin}}
%On note $B$ la boule unité fermée de $E$, et $X := [-1,1]^B$ qui est un produit de compact donc est compact.
%Soit également
%\begin{equation}
%	X_\lin := \{(l_x)_{x \in B} \in X \mid l_{x+y} = l_x+l_y,\ l_{\lambda x} = \lambda l_x,\ \forall x,y \in B, \forall \lambda \in \bK, \text{ s.t. } x+y \in B, \lambda x \in B\}.
%\end{equation}
%Il s'agit d'un fermé de $X$, donc d'un compact. De plus tout élément de $X_\lin$, étendu par homogénéité, définit une forme linéaire sur $E$, qui est bornée par $[-1,1]$ sur la boule unité. 
%On obtient donc une bijection $X_\lin \to B^*$. 
%
%Enfin, cette bijection est continue. En effet, 
%
%
%
%Notons $B$ la boule unité fermée de $E$. 
%Le produit $X := [-1,1]^B$ est compact, et l'application $\Phi : B^* \to X : l \mapsto (l(x))_{x \in B}$ est injective et continue. En effet, une base d'ouverts de $X$ est donnée par $\Pi_{x\in B} I_x$, où $I_x\subset [-1,1]$ est un intervalle ouvert et $B_0 := \{x\in B \mid I_x \neq [-1,1]\}$ est fini. Son image réciproque $\{l\in B^* \mid \forall x \in B_0, l(x) \in I_x\}$ est bien un ouvert de $B^*$ pour la topologie faible. 
%
%L'image de $\Phi$ se caractérise par 
%\begin{equation}
%	\{(l_x)_{x \in B} \mid l_{x+y} = l_x+l_y,\ l_{\lambda x} = \lambda l_x,\ \forall x,y \in B, \forall \lambda \in \bK, \text{ s.t. } x+y \in B, \lambda x \in B\}.
%\end{equation}
%En effet ces relations définissent une forme linéaire, en l'étendant à $E$ par homogénéité, et celle-ci est continue car bornée sur la boule unité. Ainsi $\mathrm{Im}(\Phi)$ est fermée, donc compacte. L'application $B^* \mapsto $
%
%\end{proof}


Une preuve d'un esprit similaire permet d'établir le théorème d'Ascoli. Nous présentons d'abord sa forme la plus simple, puis deux corollaires l'étendant d'une part aux familles de fonctions $K$-Lipschitz, $\alpha$-Hölder, ou liées à d'autres modules de continuité, et d'autre part aux dérivées d'ordre supérieur.
\begin{theorem}[Ascoli]
\label{th:AscoliLip}
\leancheck{th:AscoliLip}%
	Soient $X$ et $Y$ des espaces métriques compacts. Alors l'ensemble des applications $1$-Lipschitz est compact pour la topologie de la convergence uniforme: 
	\begin{align*}
		\Lip_1(X,Y) &:= \{ f:X \to Y \mid f \text{ est $1$-Lipschitz}\},&
		\text{muni de }
		d(f,g) &:= \max_{x \in X} d_Y(f(x), g(x)).
	\end{align*}
\end{theorem}

\begin{proof}
Le maximum définissant la distance est atteint, voir le \nref{corol:bornesAtteintes}.
Soit $(f_n)_{n \in \bN}$ une suite de $\Lip_1(X,Y)$. Soit également $D \subset X$ une partie dénombrable dense; celle-ci existe par séparabilité de $X$, voir la \cref{rem:precompact_separable}. 
Par le procédé d'extraction diagonal, ou par compacité de $Y^D$, il existe une sous-suite $f_{\vp(n)}$ qui converge simplement sur $D$. Notons $f_* : D \to Y$ sa limite, et observons qu'elle est $1$-Lipschitzienne par passage à la limite, donc s'étend en une fonction $1$-Lipschitz sur $\overline D = X$ par la \nref{prop:prolongement_complet}.

Justifions la convergence uniforme : par compacité de $X$, il existe pour tout $\ve>0$ un ensemble fini $D_\ve \subset D$ tel que $\cup_{x \in D_\ve} B'(x,\ve) = X$. 
Par finitude de $D_\ve$, il existe un rang $N_\ve$ tel que $d(f_{\vp(n)}(x), f_*(x)) \leq \ve$ pour tous $x \in D_\ve$ et $n \geq N_\ve$. Ainsi, étant donné $x \in X$, on peut choisir $x_\ve\in D_\ve$ tel que $d(x,x_\ve) \leq \ve$ et obtenir pour tout $n \geq N_\ve$
\begin{equation*}
	d(f_{\vp(n)}(x), f_*(x)) \leq d(f_{\vp(n)}(x), f_{\vp(n)}(x_\ve)) + d(f_{\vp(n)}(x_\ve), f_*(x_\ve)) + d(f_*(x_\ve), f_*(x)) \leq 3 \ve,
\end{equation*}
ce qui établit la convergence uniforme et conclut la preuve. 
\end{proof}

Les hypothèses du théorème d'Ascoli peuvent être affaiblies dans une certaine mesure.
Rappelons qu'une famille $\cF \subset C^0(X,Y)$ de fonctions d'un espace métrique $X$ dans un autre $Y$ est équicontinue (resp.\ uniformément équicontinue), voir la \cref{rem:equicont},  ssi il existe pour tout $x \in X$ un module de continuité $\omega_x$ (resp.\ un module de continuité $\omega$) tel que 
\begin{align*}
	\forall x,y \in X,\ \forall i \in I,\ d_Y(f_i(x),f_i(y)) &\leq \omega_x (d_X(x,y)). &
	(\text{resp. } d_Y(f_i(x),f_i(y)) \leq \omega(d_X(x,y)). )
\end{align*}

\begin{corollary}[Ascoli équicontinu]
\label{corol:AscoliEquicont}
\leancheck{corol:AscoliEquicont}%
	Soient $X$ et $Y$ des espaces métriques compacts, et soit $\cF \subset C^0(X,Y)$ une famille de fonctions équicontinue. Alors $\cF$ est uniformément équicontinue, et $\overline \cF$ est compact pour la topologie de la convergence uniforme. 
\end{corollary}

\begin{proof}
	L'uniforme équicontinuité découle du \nref{th:Heine}, et du commentaire qui le suit concernant les familles de fonctions. On note $\omega$ le module de continuité correspondant, que l'on peut supposer borné car $Y$ est compact donc borné. Par le \cref{lem:regModuleCont}, on peut donc aussi supposer que $\omega$ est croissant et sous-additif, de sorte que $d_X' := \omega\circ d_X$ est une distance sur $X$, définissant la même topologie. Avec cette nouvelle distance, $\cF$ est un sous-ensemble de $\Lip_1((X,d'_X),(Y,d_Y))$, donc $\overline \cF$ est compact pour la topologie de la convergence uniforme sur $Y$. On conclut en notant que cette dernière dépend de $d_Y$ mais pas du choix de $d_X$.
\end{proof}

\begin{corollary}[Ascoli $C^k$]
\label{corol:AscoliCk}
Soit $\Omega\subset \bR^d$ un ouvert convexe borné, et soit $k \in \bN$. La boule unité de $C^{k+1}_b(\Omega)$ est pré-compacte dans $C_b^k(\Omega)$.
\end{corollary}

\begin{proof}
	Adapter la preuve du \cref{corol:AscoliCkNonBorne} ci-dessous, similaire mais un peu plus technique. 
\end{proof}



\begin{corollary}[Ascoli $C^k$ non-borné \facultatif]
\label{corol:AscoliCkNonBorne}
Soit $\Omega\subset \bR^d$ un ouvert, et soit $k \in \bN \cup \{\infty\}$. Soit $C(K,\alpha)$ une constante arbitraire associée à tout $K \csubset \Omega$ et tout élément de $\bN^d_{k+1} := \{\alpha \in \bN^d\mid |\alpha|_1 \leq k+1\}$.
Alors 
\begin{equation*}
	\cF := \{ f \in C^{k+1}(\Omega) \mid \forall K \csubset \Omega,\, \forall \alpha \in \bN^d_{k+1},\,   \|\partial_\alpha f\|_{L^\infty(K)} \leq C(K,\alpha)\}
\end{equation*}
est une partie précompacte de $C^k(\Omega)$. %, fermée dans le cas $k = \infty$. 
%est une partie précompacte de $C^k(\Omega)$. Si $k=\infty$, alors c'est une partie compacte.
%Les constantes $C(K,\alpha)$ dépendent seulement des paramètres indiqués. 
\end{corollary}

\begin{proof}
	Soit $(B_n)_{n \in \bN}$ une famille dénombrable de boules fermées telles que $\Omega = \cup_{n \in \bN} \mathring B_n$. (Dans le cas $\Omega\neq \bR^d$, on peut par exemple considérer une suite dense $(x_n)_{n \in \bN}$ dans $\Omega$, et poser $B_n := B'(x_n,r_n)$ avec $r_n := \frac 1 2 d(x_n, \bR^d \sm \Omega)$.)
	
	Choisissons $\alpha \in \bN^d_k$, i.e.\ $\alpha \in \bN^d$ tel que $|\alpha|_1 \leq k$. Si $f \in \cF$, alors $\partial_\alpha f$ est $K(n,\alpha)$-Lipschitz sur $B_n$, par convexité de cet ensemble, avec $K(n,\alpha) := \sum_{1 \leq i \leq d} C(B_n, \alpha+e_i)$. De plus $\|\partial_\alpha f\|_{L^\infty(B_n)} \leq C(B_n,\alpha)$.
	Par le théorème d'Ascoli, $\cF_{n,\alpha} := \{\partial_\alpha f_{|B_n} \mid f \in \cF\}$ est une partie précompacte de $C^0(B_n)$ pour la topologie de la convergence uniforme. 
	
	Considérons une suite $(f_r)_{r \in \bN}$ de $\cF$.
	Par compacité de $\prod_{n \in \bN} \prod_{\alpha\in \bN^d_k} \overline {\cF_{n,\alpha}}$, ou par un procédé d'extraction diagonal, il existe une sous-suite notée $(f_{\vp(r)})_{r \in \bN}$ telle que $\partial_\alpha f_{\vp(r)}$ converge uniformément lorsque $r \to \infty$ sur $B_n$, pour tout $n \in \bN$ et tout $\alpha \in \bN^d_k$.
	On en déduit que $(f_{\vp(r)})_{r \in \bN}$ est de Cauchy dans $C^k(\Omega)$, muni des semi-normes $\|\partial_\alpha f\|_{L^\infty(K_n)}$ où $\alpha \in \bN^d_k$, où $(K_n)_{n \in \bN}$ est une suite exhaustive arbitraire de compacts de $\Omega$, voir la \cref{prop:exhaustComp}. 
	(En effet, on peut recouvrir tout compact $K_n$ par un nombre fini de boules prises dans la famille $(\mathring B_m)_{m \in \bN}$.) 
	Ainsi la suite $(f_{\vp(r)})_{r \in \bN}$ est convergente pour la topologie de $C^k(\Omega)$, ce qui donne la précompacité.
\end{proof}

Remarquons que $\cF$ du corollaire précédent est une partie fermée de $C^{k+1}(\Omega)$, par définition. Ainsi $\cF \subset C^\infty(\Omega)$, dans le cas particulier $k=\infty$, est précompacte et fermée, donc compacte.  


%\todoi{
%Compacité dans $C_b^k$ de la boule unité de $C_b^{k+1}$. 
%Versions locales : si les dérivées sont bornées sur tout compact.
%}

On rappelle la définition de l'enveloppe convexe d'une partie $K$ d'un $\bR$-ev:
\begin{align*}
	\Hull(K) &:= \{ \sum_{0 \leq i \leq I} \lambda_i x_i \mid 
	I \geq 1,\ \lambda \in \Delta^I,\ \forall 0 \leq i \leq I,\ x_i \in K\}, \\
	\text{où }\Delta^I &:= \{\lambda = (\lambda_i)_{0 \leq i \leq I} \in \bR_+^{I+1} \mid \sum_{0 \leq i \leq I} \lambda_i = 1\}.
\end{align*}
En dimension finie, si $K\subset \bR^d$, alors le \nref{th:Caratheodory} affirme que l'on peut se limiter à $I=d$ dans l'expression ci-dessus, c.-à-d. à des combinaisons convexes de $d+1$ points. On en déduit que $\Hull(K)$ est compact si $K$ est une partie compacte d'un espace de dimension finie.

En dimension infinie, on a le résultat de \emph{précompacité} suivant pour la topologie forte. Il admet une généralisation due à Krein-Smulian pour la topologie faible.

%\begin{lemma}
%\label{lem:allToZero}
%	Supposons que $\ve_k(n) \to 0$ lorsque $n \to \infty$, quel que soit $k \in \bN$. Alors il existe une suite $\vp(n) \to \infty$ telle que pour toute suite $\psi \leq \vp$ on a $\ve_{\psi(n)}(n) \to 0$.  
%\end{lemma}
%
%\begin{proof}
%	Poser par exemple $\vp(n):= \max \{k \in \llbracket 0,n\rrbracket \mid \forall r \in \llbracket 0,k \rrbracket,\ \vp_r(n) \leq 1/k\}$, avec la convention $1/0 = \infty$ assurant que cet ensemble est non-vide. 
%\end{proof}

\begin{proposition}[Compacité de l'enveloppe convexe fermée]
\label{prop:hullCompact}
\leancheck{prop:hullCompact}%
	Soit $K$ une partie d'un $\bR$-Banach. Si $\overline K$ est compact, alors $H := \overline {\Hull(K)}$ est compact. 
\end{proposition}

\begin{proof}
	Étant donné $\ve>0$, soit $K \subset \cup_{i=1}^I B(z_i,\ve)$ une couverture finie de $K$ par des boules de rayon $\ve$. Posons $H_\ve := \Hull\{z_i \mid 0 \leq i \leq I\}$, qui est un ensemble compact, en tant qu'image continue de $\Delta^I$. Par construction, il existe pour tout $x \in H$ un point $x_\ve \in H_\ve$ tel que $\|x-x_\ve\| \leq \ve$. 
	
	Soit $(x_n)_{n \in \bN}$ une suite à valeurs dans $H$, et soit $x_n^k \in H_{1/k}$ tel que $\|x_n - x_n^k\| \leq 1/k$, pour tout $k \in \bN^*$. 
	Par un argument d'extraction diagonal, ou en utilisant la compacité de $\prod_{k \in \bN^*} H_{\frac 1 k}$ ce qui revient au même, il existe une extractrice $\vp$ telle que $x_{\vp(n)}^k$ converge pour tout $k \in \bN^*$, vers une limite notée $x_*^k$ lorsque $n \to \infty$.
	Par construction on a $\| x_n^k - x_n^l \| \leq \frac 1 k + \frac 1 l$ pour tous $n,k,l>0$, donc $\|x_*^k - x_*^l\| \leq \frac 1 k + \frac 1 l$ par passage à la limite, ce qui montre que $(x_*^k)_{k \in \bN}$ est de Cauchy, donc converge vers une limite notée $x_*^* \in \overline {\Hull(K)}$ lorsque $k \to \infty$.
	En utilisant l'inégalité triangulaire, on a $\|x_{\vp(n)} - x_*^*\| \leq \|x_{\vp(n)} - x_{\vp(n)}^k\| + \|x_{\vp(n)}^k - x_*^k\| + \|x_*^k-x_*^*\|$. 
	En choisissant %via le \cref{lem:allToZero}
	une suite $k(n)\to \infty$ telle que $\|x_{\vp(n)}^{k(n)} - x_*^{k(n)}\| \to 0$ lorsque $n \to \infty$, on conclut que $x_{\vp(n)} \to x_*^* $, d'où la précompacité annoncée.	
\end{proof}

%\hiddenproof{
%Argument détaillé pour la construction de $k(n)$. %, qui est une d'extraction diagonale dans une double limite. 
%Pour tout $n \in \bN$ donné, considérons une suite $\ve_k(n) \to 0$ lorsque $n \to \infty$. (Dans la preuve précédente $\ve_k(n) = \|x^k_{\vp(n)}-x^k_*\|$.)
%Par hypothèse il existe pour tout $k \in \bN$ une extractrice $\vp_k(r)$ telle que $\ve_k(n) \leq 1/r$ pour tout $n \geq \vp_k(r)$. On peut supposer $\vp_k(r)$ croissante en $k$ également. On note que $\ve_k(\vp_k(k)) \leq 1/k$. On définit $k(n) := \max\{k>0 \mid \vp_k(k) \leq n\}$. C'est une suite croissante, qui tend vers l'infini, et qui convient puisque $\ve_{k(n)}(n) \leq 1/k(n)$.
%}

Pour être complet, présentons le cas de la dimension finie.
\begin{theorem}[Carathéodory]
\label{th:Caratheodory}
\leancheck{th:Caratheodory}%
Pour tout $A \subset \bR^d$ on a 
\begin{equation*}
	\Hull(A) = \{\sum_{0 \leq i \leq d} \lambda_i x_i\mid \lambda \in \Delta^d,\ x \in A^{d+1}\}.
\end{equation*}
On en déduit que $\Hull(K)$ est compact pour tout $K \csubset \bR^d$ compact. 
\end{theorem}

\begin{proof}
	Soit $x \in \Hull(A)$, que l'on peut donc écrire $x = \sum_{0 \leq i \leq I} \lambda_i x_i$ avec $\lambda \in \Delta^I$ et $x_0,\cdots,x_I \in A$. Supposons que l'entier $I$ est minimal dans cette représentation et, par l'absurde, que $I\geq d+1$. Alors, par un argument de dimension, les vecteurs $(x_i-x_0)_{1 \leq i \leq I}$ sont liés, et il existe donc une relation linéaire de la forme
	\begin{equation*}
		\sum_{1 \leq i \leq I} \mu_i (x_i-x_0) = 0,
	\end{equation*}
	qui est non-triviale, c'est-à-dire que l'un au moins des $(\mu_i)_{1 \leq i \leq I}$ est non-nul.  
	On pose $\mu_0 := -\sum_{1 \leq i \leq I} \mu_i$, de sorte que 
	\begin{align*}
		\sum_{0 \leq i \leq I} \mu_i x_i &= 0, &
		\sum_{0 \leq i \leq I} \mu_i &= 0,
	\end{align*}
	et que l'un au moins des $\mu_i$ est strictement positif. Soit ensuite $\rho := \min\{\lambda_i / \mu_i \mid 0 \leq i \leq I,\, \mu_i > 0\}$. 
	Alors 
	\begin{equation*}
		x = \sum_{0 \leq i \leq I} (\lambda_i - \rho \mu_i) x_i,
	\end{equation*}
	et cette combinaison convexe comporte au moins un coefficient nul. Ceci contredit la minimalité de $I$, d'où la caractérisation annoncée de l'enveloppe convexe. 

	Finalement, $\Hull(K)$ est compact en tant qu'image de l'application continue suivante, dont le domaine est un compact, voir la \cref{prop:compactImage} : \belowdisplayskip=-12pt
	\begin{align*}
		\Delta^d \times K^{d+1} &\to \bR^d, &
		( (\lambda_i)_{0 \leq i \leq d} ,(x_i)_{0\leq i \leq d}) \mapsto \sum_{0 \leq i \leq d} \lambda_i x_i.
	\end{align*}
\end{proof}


%\todoi{Frechet Kolmogorov ? $\to$ Dans $L^p$ ? A voir. Utilise la convolution. Sinon le chapitre sur les $L^p$ sera vide.}

\subsection{Point fixe de Brouwer}
\label{subsec:Brouwer}

Le \nref{th:Brouwer} est un résultat fondamental de topologie des espaces euclidiens. Ses applications couvrent de nombreux champs des mathématiques, soit par une application directe (e.g.\ existence d'un équilibre de Nash pour un jeu fini probabilisé), soit via des extensions en dimension infinie comme discuté dans la \nref{subsec:Schauder}.

Dans toute cette sous-section, on note $B := B'(0,1)$ la boule unité fermée de $\bR^d$.

\begin{retenir}
L'énoncé du théorème de Brouwer, la formule de changement de variable dans une intégrale.
(\Cref{th:Brouwer,th:chgVarIntegr}.) 
\end{retenir}



%Le but de cette sous-section est d'établir le résultat suivant. 
\begin{theorem}[Point fixe de Brouwer]
\label{th:Brouwer}
\leancheck{th:Brouwer}%
	Toute fonction $f \in C^0(B, B)$ admet un point fixe. 
\end{theorem}

Le théorème de Brouwer s'étend de manière triviale à tout ensemble homéomorphe à une boule, ainsi qu'à tout ensemble $K$ sur lequel une boule $B$ se \emph{rétracte}.

\begin{corollary}	
	Soit $K \subset \bR^d$ un ensemble sur lequel la boule $B$ se rétracte, i.e.\ $K \subset B$ et il existe $P \in C^0(B,K)$ telle que $P_{|K}=\Id$. Alors toute  $f\in C^0(K,K)$ admet un point fixe. 
	En particulier, toute application continue d'un convexe compact non-vide de $\bR^d$ dans lui même admet un point fixe.
\end{corollary}

\begin{proof}
	\emph{Premier point : }d'après le \nref{th:Brouwer}, la composée $f\circ P$ admet un point fixe $x \in B$. Ainsi $x \in \Image(f) = K$, puis $P(x) = P_{|K}(x) = x$, et donc $x=f(x)$ comme annoncé.
	\emph{Second point : }quitte à faire une homothétie on peut supposer que $K \subset B$, puis considérer la rétractation $P_{K|B}$ définie par la restriction à $B$ de la projection orthogonale sur $K$.
\end{proof}

On suit ici la preuve de Peter Lax, également présentée dans \cite{alazard2017analyse}, qui s'appuie sur une généralisation de la formule de changement de variable dans une intégrale. 
Rappelons d'abord la forme classique de cette formule, sans démonstration. La matrice (\ref{eq:chgmt_var}, droite) de taille $d \times d$ est appelée matrice jacobienne de $\vp$, et son déterminant est nommé jacobien. 
\begin{theorem}[Changement de variable dans une intégrale]
\label{th:chgVarIntegr}
\leancheck{th:chgVarIntegr}%
Soient $U,V \subset \bR^d$ ouverts, $\vp \in C^1(U, V)$ difféomorphisme, et $f : V \to \bR$ intégrable. 
Alors %, en notant $J$ la matrice jacobienne de $\vp$,
\begin{align}
\label{eq:chgmt_var}	
	\int_V f(x) \diff x &= \int_U f(\vp(x)) |\det J(x)| \diff x,&
	\text{où }
	J(x) := \Big[ \frac{\partial \vp_i}{\partial x_j}\Big]_{1 \leq i,j \leq d}.
\end{align}	
\end{theorem}

Rappelons que la comatrice de $A \in \bR^{d\times d}$  est définie par
\begin{equation*}
	[\com A]_{ij} = \frac{\partial \det A}{\partial A_{ij}} = \det [A_1,\cdots,A_{j-1}, b_i, A_{j+1},\cdots, A_d],
\end{equation*}
où $A_{ij}$ (resp.\ $A_i$) désignent les coefficients (resp.\ colonnes) de $A$, pour tous $1 \leq i, j \leq d$, et où $(b_i)_{1 \leq i \leq d}$ est la base canonique de $\bR^d$. On a $(\com A)^\top A = (\det A) \Id$, et $\det (A+H) = \det A + \Tr((\com A)^\top H) + o(\|H\|)$ lorsque $H \to 0$. 

On note $\nabla \vp = J^\top$ la transposée de la matrice Jacobienne.
Par convention, la divergence d'un champ de matrices est calculée colonne par colonne.

\begin{lemma}[Identité de Piola] \label{lem:Piola}
\leancheck{lem:Piola}%
Pour toute $\vp \in C^2(\bR^d,\bR^d)$, on a $\diver( \com \nabla \vp) = 0$ sur $\bR^d$.
\end{lemma}

\begin{proof}
	Notons $\vp = (\vp_1,\cdots,\vp_d)$, $A = \nabla \vp$, et observons que ses colonnes sont $A_i = \nabla \vp_i$, $1 \leq i \leq d$. On a
	\begin{equation*}
		\frac{\partial [\com A]_{ij}}{\partial x_i} 
		= \sum_{1 \leq k,l \leq d} 
		\frac{\partial [\com A]_{ij}}{\partial A_{kl}} 
		\frac {\partial A_{kl}}{\partial x_i} 
		= \sum_{1 \leq k,l \leq d} 
		\frac{\partial^2 \det A}{\partial A_{ij} \partial A_{kl}} 
		\frac {\partial^2 \vp_l}{\partial x_i \partial x_k},
%		\frac{\partial A_{kl}}{\partial x_i}		 
	\end{equation*}
par dérivation composée. Par ailleurs, nous avons la propriété d'anti-symétrie: si $j<l$
\begin{equation*}
	\frac{\partial^2 \det A}{\partial A_{ij} \partial A_{kl}} 
	= \det [A_1,\cdots,A_{j-1}, b_i, A_{j+1}, \cdots, A_{l-1}, b_k, A_{l+1}, \cdots, A_d] 
	= - \frac{\partial^2 \det A}{\partial A_{kj} \partial A_{il}}
\end{equation*}
de même si $j>l$, tandis que $\frac{\partial^2 \det A}{\partial A_{ij} \partial A_{kl}} = 0$ si $j=l$ par linéarité du déterminant par rapport aux colonnes. Ainsi, pour tout $1 \leq j \leq d$, par anti-symétrie et via l'identité de Schwarz
\begin{equation*}
\diver(\com \nabla \vp)_j 
= \sum_{1 \leq i \leq d}
\frac{\partial [\com A]_{ij}}{\partial x_i} 
= \sum_{1 \leq l \leq d} \Big(\sum_{1 \leq i,k \leq d} 
\frac{\partial^2 \det A}{\partial A_{ij} \partial A_{kl}} 
\frac {\partial^2 \vp_l}{\partial x_i \partial x_k}\Big) = 0.
\qedhere
\end{equation*}
\end{proof}

\begin{lemma}[Lax]
\label{lem:chgVarLax}
\leancheck{lem:chgVarLax}%
% -- DO NOT REMOVE -- La preuve Lean évite la division par t^-1, et le cas t=0, avec un peu d'algèbre supplémentaire.
Soit $\vp \in C^2(\bR^d,\bR^d)$ telle que $\vp(x) = x$ pour tout $x \notin B$, et soit $f \in C^1_c(\bR^d)$. Alors 
	\begin{equation}		
	\label{eq:chgmt_var_nondiff}
		\int_{\bR^d} f(x) \diff x = \int_{\bR^d} f(\vp(x)) \det J(x) \diff x.
	\end{equation}
\end{lemma}

\begin{proof}
	Commençons par calculer, en notant $\id$ (resp.\ $\Id$) l'application (resp.\ la matrice) identité
	\begin{align} \nonumber
		&\diver( f\circ \vp \, [\com \nabla \vp]  (\vp - \id)) \\ \nonumber
		&= \<\nabla \vp \nabla f \circ \vp, [\com \nabla \vp] (\vp-\id)\> + f\circ \vp \<\diver (\com \nabla \vp), \vp-\id\> + f\circ \vp \Tr((\com \nabla \vp)^\top  (\nabla \vp - \Id))\\ \label{eq:divLax}
		&= \<\nabla f \circ \vp, \vp-\id\> \det (\nabla \vp) + f \circ \vp \Tr((\com \nabla \vp)^\top  (\nabla \vp - \Id)),
	\end{align}
	où l'on a distribué la divergence, puis utilisé l'identité de Piola, et l'égalité $A^\top (\com A) = (\det A) \Id$ avec $A = \nabla \vp$.
	Notons maintenant $\vp_t := (1-t) \id + t \vp$, et calculons pour $t \in ]0,1]$
	\begin{align*}
		&\partial_t \int_{\bR^d} f \circ \vp_t \, \det (\nabla \vp_t) 
		= \int_{\bR^d} \<\nabla f \circ \vp_t, \partial_t \vp_t\> \det \nabla \vp_t + f \circ \vp_t \Tr((\com \nabla \vp_t)^\top \partial_t \nabla \vp_t)\\
		&= t^{-1} \int_{\bR^d} \<\nabla f \circ \vp_t, \vp_t-\id\> \det \nabla \vp_t + f \circ \vp_t \Tr((\com \nabla \vp_t)^\top (\nabla \vp_t- \Id))\\
		&= t^{-1} \int_{\bR^d} \diver( f \circ \vp_t [\com \nabla \vp_t] (\vp_t - \id)) = 0.
	\end{align*}
	On a successivement (i) dérivé sous le signe intégral, et distribué $\partial_t$ sur le produit, (ii) remarqué que $\partial_t \vp_t = \vp-\id = t^{-1} (\vp_t-\id)$ et de même $\partial_t \nabla \vp_t = t^{-1} (\nabla \vp_t - \Id)$, (iii) utilisé \eqref{eq:divLax}, et la nullité de l'intégrale de la divergence d'un champ de vecteurs à support compact. Ainsi $t\in [0,1] \mapsto \int_{\bR^d} f \circ \vp_t \, \det (\nabla \vp_t)$ est continu et de dérivée nulle sur $]0,1]$, donc constant.
\end{proof}

%\begin{lemma}[Lax]
%\label{lem:chgVarLax}
%\leancheck{lem:chgVarLax}%
%	Soit $\vp \in C^2(\bR^d,\bR^d)$ telle que $\vp(x) = x$ pour tout $x \notin B$, et soit $f \in \cD(\bR^d)$. Alors 
%	\begin{equation}		
%	\label{eq:chgmt_var_nondiff}
%		\int_{\bR^d} f(x) \diff x = \int_{\bR^d} f(\vp(x)) \det J(x) \diff x.
%	\end{equation}
%\end{lemma}
%
%\begin{proof}[Preuve en dimension $d=2$] On renvoie à \cite{alazard2017analyse} pour le cas général. Soit $c\geq 1$ tel que $\supp(f)\subset Q := [-c,c]^d$.
%On note $g(x) := \int_{-c}^{x_1} f(s,x_2) \diff s$, de sorte que $\partial_1 g = f$. 
%Par dérivation composée, on obtient l'identité
%\begin{equation*}
%	f \circ \vp\, J = f \circ \vp\, (\partial_1 \vp_1 \,\partial_2 \vp_2 - \partial_1 \vp_2 \, \partial_2 \vp_1)
%	= \partial_1 (g \circ \vp) \, \partial_2 \vp_2  -\partial_2 (g \circ \vp) \, \partial_1 \vp_2 .
%\end{equation*}
%Par la formule usuelle de dérivation d'un produit
%%Par intégration par parties 
%\begin{equation*}
%	\int_Q f(\vp(x)) \det J(x)\diff x = \int_Q \Big[ \partial_1 (g\circ \vp \, \partial_2 \vp_2) - \partial_2(g \circ \vp \, \partial_1 \vp_2)\Big] - \int_Q g \circ \vp \, (\partial_1 \partial_2 \vp_2 - \partial_2\partial_1 \vp_2).
%\end{equation*}
%Le membre de droite est somme de deux intégrales. La première ne dépend des valeurs de $g$, $\vp$ et $\partial \vp$ que sur $\partial Q$ (par intégration de la primitive), tandis que la seconde s'annule par le théorème de Schwarz d'interversion des dérivées partielles. Le terme (\ref{eq:chgmt_var_nondiff}, droite) est donc identique pour $\tilde \vp := \Id$, car cette application coïncide avec $\vp$ sur $\partial Q$, ce qui conclut.
%\end{proof}

\begin{remark}
\label{rem:chgVarDiffeo}
Noter que l'on ne suppose pas que $\vp$ est un difféomorphisme dans \eqref{eq:chgmt_var_nondiff}. Par ailleurs le jacobien intervient dans \eqref{eq:chgmt_var_nondiff}, et sa valeur absolue\footnote{L'hypothèse que $\vp$ est un difféomorphisme dans le \cref{th:chgVarIntegr} impose toutefois que le jacobien est de signe constant sur les composantes connexes de $U$} dans \eqref{eq:chgmt_var}.
	Lax a montré que le \cref{th:chgVarIntegr} se déduit du \cref{lem:chgVarLax}.
\end{remark}


\begin{lemma}
\label{lem:noretractSmooth}
\leancheck{lem:noretractSmooth}%
	Soit $\vp \in C^2(\bR^d, \bR^d)$ telle que $\vp(x) = x$ pour tout $x \notin B$. Alors $B \subset \vp(B)$.
\end{lemma}

\begin{proof}
	Par l'absurde, considérons $y_0 \in B \sm \vp(B) \neq \emptyset$. Comme $\partial B$ est fixé par $\vp$, par continuité, on a $y_0 \in \mathring B \sm \vp(B)$. Comme cet ensemble est ouvert, il existe $\ve>0$ tel que $B(y_0, \ve) \subset \mathring B \sm \vp(B)$. 
	Choisissons $f\in C^\infty(\bR^d)$ telle que $\supp(f) \subset B(y_0,\ve)$ et $\int_{\bR^d} f = 1$.
	Le \cref{lem:chgVarLax} donne alors la contradiction suivante, ce qui conclut : 
	\begin{equation*}
		1 = \int_{\bR^d} f(x) \diff x = \int_{\bR^d} f(\vp(x)) \det J(x) \diff x = 0.
		\qedhere
	\end{equation*} 
\end{proof}

\begin{corollary}[Non rétractation de la boule sur la sphère]
\label{corol:noretract}
\leancheck{corol:noretract}%
	Soit $\vp\in C^0(B,B)$ telle que $\vp(x) = x$ pour tout $x\in \partial B$. 
	Alors $\vp(B) = B$. 
\end{corollary}

\begin{proof}
L'énoncé est très similaire au \cref{lem:noretractSmooth}, à l'exception des hypothèses de régularité qui sont affaiblies. Le résultat se déduit du précédent, par un argument de régularisation et de compacité. 
	Étendons $\vp$ à $\bR^d$ en posant $\vp(x) = x$ pour tout $x \notin B$. 
	Puis, posons pour tout $\ve>0$
	\begin{align*}
		\vp_\ve &:= \eta_\ve \star \vp,&
		\text{ où }
		\eta_\ve(x) &:= \tfrac{1}{\ve^d} \eta(\tfrac{x}{\ve}).
	\end{align*}
	On a noté $\star$ l'opération de convolution, qui sera discutée plus en détail \cref{sec:Sobolev}. Le noyau de convolution $\eta$ est choisi $C^\infty$, radial, satisfaisant $\supp(\eta) \subset B$ et $\int_B \eta = 1$, voir la \nref{prop:cinfBosse} pour une construction possible. 
	Par construction, on a $\vp_\ve(x) = x$ pour tout $x \notin B'(0,1+\ve)$, $\vp_\ve \in C^\infty(\bR^d)$, et $\vp_\ve$ converge uniformément vers $\vp$ lorsque $\ve \to 0$. 
	Par le \cref{lem:noretractSmooth}, on obtient $\vp_\ve(B'(0,1+\ve)) \supset B'(0,1+\ve)$. 
	
	Concluons, par passage à la limite : soit $x \in B$ arbitraire, et pour tout $\ve>0$ soit $y_\ve \in B'(0,1+\ve)$ tel que $\vp_\ve(y_\ve) = x$. Alors il existe par compacité un point $y_* \in B$ et une suite $\ve(n) \to 0$ telle que $y_{\ve(n)} \to y_* \in B$ lorsque $n \to \infty$. Finalement $\vp(y_*) = x$ par convergence uniforme, ce qui  établit $\vp(B) \supset B$, l'inclusion réciproque découlant des hypothèses. 
\end{proof}

\begin{proof}[Preuve du théorème de Brouwer]
	Par l'absurde, considérons $f:B\to B$ sans point fixe, 
	et pour tout $x \in B$ définissons $\vp(x) \in \partial B$ par 	
	\begin{align*}
		\{\vp(x)\} &:= \partial B \cap D_x, &
		\text{où } D_x &:= \{x + \lambda(x-f(x)) \mid \lambda\geq 0\},
	\end{align*}
	Le polynôme $P_x(\lambda) := \|x+\lambda (x-f(x))\|^2-1$ a le coefficient dominant $\|x-f(x)\|^2 > 0$, et satisfait $P_x(0) = \|x\|^2-1 \leq 0$ et $P_x(-1) = \|f(x)\|^2-1 \leq 0$, donc a une unique racine positive $\lambda(x)$, dépendant continument de $x \in B$, d'où la continuité de $\vp(x) := x+\lambda(x) (x-f(x))$. De plus $\vp(x) = x$ pour tout $x \in \partial B$ par construction.
	L'application $\vp$ contredit ainsi le résultat de non-rétractation \cref{corol:noretract}, ce qui conclut.
%\todoi{Rédiger un peu plus : le polynôme $P_x(\lambda) := \|x+\lambda (x-f(x))\|^2-1$ a un coefficient dominant $\|x-f(x)\|^2 > 0$, et satisfait $P_x(0)\leq 0$ et $P_x(-1) \leq 0$, donc a une unique racine positive, dépendant continument de $x$.}
\end{proof}

\subsection{Point fixe de Schauder} 
\label{subsec:Schauder}

Le théorème de point fixe de Schauder étend le théorème de Brouwer en dimension infinie, et permet de montrer l'existence (mais pas l'unicité !) de solutions de certaines EDO et EDP. 

\begin{retenir}
Le théorème de point fixe de Schauder, et son application au théorème d'existence d'EDO de Cauchy-Arzélà-Peano, preuve comprise.
(\Cref{th:Schauder,th:CauchyPeano}.)	
\end{retenir}

\begin{theorem}[Point fixe de Schauder]
\label{th:Schauder}
\leancheck{th:Schauder}%
Soit $K$ un convexe fermé non-vide d'un Banach, et soit $f\in C^0(K,K)$. Si $\overline {f(K)}$ est compact alors $f$ admet un point fixe. 
\end{theorem}

\begin{proof}
Quitte à remplacer $K$ par le sous-ensemble $\overline {\Hull(f(K))}$, on peut supposer que $K$ est convexe et compact, voir la \cref{prop:hullCompact}.
Étant donné $\ve>0$, considérons une couverture de la forme $K \subset \cup_{i=0}^I B(x_i, \ve)$, où $I$ est un entier et $x_0,\cdots,x_I \in K$, et posons 
\begin{align*}
	g_\ve : K &\to K_\ve := \Hull\{x_i \mid 0 \leq i \leq I\}, &
	x &\mapsto \frac{\sum_{0 \leq i \leq I} \vp_i(x) x_i}{\sum_{0 \leq i \leq I} \vp_i(x)}, 
	%&\text{ où }\vp_i(x) &:= \max \{0, \ve-\|x - x_i\|\}.
\end{align*}
où $\vp_i(x) := \max \{0, \ve-\|x - x_i\|\}$.
Par construction, $g_\ve$ est bien définie et continue sur $K$ car le dénominateur y est positif, et de plus $\|g_\ve(x)-x\| \leq \ve$ pour tout $x \in K$. 
D'après le théorème de Brouwer, appliqué à $K_\ve \subset K$ qui est compact convexe et de dimension au plus $I$, la fonction continue $g_\ve \circ f : K_\ve \to K_\ve$ admet un point fixe noté $x_\ve$. On a donc $\|f(x_\ve)-x_\ve\| \leq \ve$. Par compacité de $K$, on peut trouver une suite convergente de $x_{\ve_n} \to x_*$ lorsque $n \to \infty$, et ainsi obtenir un point fixe $x_*$ de $f$. 
\end{proof}

\begin{remark}
\label{rem:SchauderEvt}
%La preuve ci-dessus s'adapte sans difficulté au cas où la topologie induite sur $K$ est métrisable (et l'addition et la multiplication scalaire continues), ce qui permet de traiter le cas des evtlcs ayant une famille dénombrable, ou de la topologie *faible sur le dual d'un evn séparable.
Le théorème de Schauder peut se démontrer sous l'hypothèse plus faible que l'espace ambiant est un espace vectoriel topologique (plutôt qu'un Banach); en particulier, il s'applique aux evtlcs, donc aux topologies faibles. 
Attention, il faut vérifier que $f$ est continue par rapport à la topologie utilisée. 
\end{remark}


Le théorème de Schauder permet de prouver l'existence de solutions d'EDOs pour des lois d'évolution seulement continues, à comparer avec le \nref{th:Cauchy_Lipschitz} qui établit l'unicité et la stabilité, au prix d'hypothèses de régularité plus fortes.
\begin{theorem}[Cauchy-Arzelà-Peano]
\label{th:CauchyPeano}
\leancheck{th:CauchyPeano}%
	Soient $\Omega \subset \bR^d$ un ouvert, et $f\in C^0(\bR_+ \times \Omega, \bR^d)$. Alors pour tout $x_0 \in \Omega$ il existe $t_0 > 0$ et $u \in C^1([0,t_0],\Omega)$ solution de 
	\begin{align}
	\label{eq:PeanoEDO}
		u(0) &= x_0,&
		\forall t \in [0,t_0],\, u'(t) &= f(t,u(t)).
	\end{align}
\end{theorem}

\begin{proof}
	Soit $r_0>0$ tel que $B'(x_0,r_0) \subset \Omega$, et soit $t_1>0$. Soit $C_\infty$ une borne supérieure de $\|f\|$ sur le compact $[0,t_1] \times B'(x_0,r_0)$. Soit $0<t_0\leq t_1$ tel que $t_0 C_\infty \leq r_0$, et soit $\omega$ un module d'uniforme continuité de $f$ sur $[0,t_0] \times B'(x_0,r_0)$.
	
	Posons $K := C^0([0,t_0],B'(x_0,r_0))$,  qui est un convexe fermé de $(C^0([0,t_0],\bR^d), \|\cdot\|_\infty)$, et 
	\begin{equation*}
	F u(t) := x_0 + \int_0^t f(s,u(s)) \diff s	
	\end{equation*}
	pour tout $u \in K$, $0 \leq t \leq t_0$. 
	Alors $F u$ est $C_\infty$-Lipschitz comme primitive d'une fonction bornée par $C_\infty$. Par ailleurs $\|Fu(t)-x_0\| \leq C_\infty t_0 \leq r_0$, de sorte que $F$ stabilise $K$. Notons aussi que l'uniforme continuité de $F$ se déduit de celle de $f$: pour  tout $t \in [0,t_0]$
	\begin{equation*}
		\|Fu(t) - Fv(t)\| \leq \int_0^t \|f(s,u(s)) - f(s,v(s))\|\diff s \leq t_0 \omega( \|u-v\|_\infty),
	\end{equation*}
	et que l'image de $F$ est contenue dans l'ensemble $\{u \in K\mid u \text{ est $C_\infty$-Lipschitz}\}$ qui est compact par le \nref{th:AscoliLip}. Ainsi le \nref{th:Schauder} s'applique à $F$, et tout point fixe est solution de \eqref{eq:PeanoEDO}, ce qui conclut. 
\end{proof}

\begin{remark}[Preuve alternative] % fondée sur le théorème d'Ascoli
\label{rem:preuveCauchyPeano}
	Il existe une preuve plus élémentaire du théorème de Cauchy-Arzelà-Peano, fondée sur une régularisation par convolution de $f$ et sur le théorème d'Ascoli.
\end{remark}


\subsection{Théorème de Stone-Weierstrass}
\label{subsec:StoneWeierstrass}

Le théorème de Stone-Weierstrass assure la densité de certaines algèbres de fonctions, et en particulier des fonctions polynomiales réelles sur un domaine compact. En analyse fonctionnelle, il est courant d'établir des résultats sur certains ensembles de fonctions régulières, puis de les étendre par des arguments de continuité et de densité à des espaces plus grands, voir la \cref{sec:Sobolev}. Stone-Weierstrass se présente alors comme une alternative aux arguments fondés sur des convolutions, déjà entrevus dans le \cref{corol:noretract}. 

\begin{retenir}
Le théorème de Stone-Weierstrass, sa variante complexe, et leur application aux fonctions polynomiales.
(\Cref{th:StoneWeierstrass,corol:StoneComplexe,rem:polDens}.)
\end{retenir}


\begin{theorem}[Stone-Weierstrass]
\label{th:StoneWeierstrass}
\leancheck{th:StoneWeierstrass}%
	Soit $X$ un espace topologique compact, et soit $A \subset C^0(X,\bR)$ une algèbre\footnote{I.e.\ $A$ est stable par addition, multiplication interne, et externe (par les scalaires).} (pas forcément unitaire a priori) qui : 
	\begin{enumerate}[(i)]
		\item sépare les points : $\forall x, y \in X$ distincts, $\exists f \in A$, $f(x) \neq f(y)$.
		\item ne s'annule pas simultanément en un point : $\forall x \in X$, $\exists f \in A$, $f(x)\neq 0$. 
	\end{enumerate} 
	Alors $\overline A = C^0(X,\bR)$, pour la topologie de la convergence uniforme donnée par $\|\cdot\|_\infty$. 
\end{theorem}

Les deux premiers lemmes construisent des éléments bien choisis de l'algèbre fermée $\overline A$ sous la forme de séries.
\leannote{Le théorème de Stone-Weierstrass étant déjà présent dans Mathlib, les résultats intermédiaires suivants ont été trivialement déduits de celui-ci plutôt que redémontrés en Lean.} 

\begin{lemma} 
\label{lem:StoneUnite}
\leancheck{lem:StoneUnite}%
Sous les hypothèses du \cref{th:StoneWeierstrass}, $\overline A$ est une algèbre unitaire complète. 
\end{lemma}

\begin{proof}
	Les propriétés d'algèbre de $\overline A$ sont claires, et la complétude se déduit de sa fermeture dans le Banach $C^0(X,\bR)$. Il reste à montrer que $\overline A$ contient la fonction constante égale à $1$, notée $\bbone$, qui est son unité. Par l'axiome (ii), et la stabilité par la multiplication par un scalaire, pour tout $x \in X$ il existe $f \in A$ telle que $f(x)=1$. 
	Par compacité, on peut trouver $f_1,\cdots,f_I$ telles que $f := \sum_{1 \leq i \leq I} f_i^2 \in A$ est strictement positive sur $X$. Quitte à multiplier $f$ par un scalaire, on peut supposer que $0< f\leq 1$ sur $X$; par compacité, $f$ est bornée inférieurement par une constante strictement positive.
	On écrit alors 
	\begin{equation}
		\bbone = f \times \frac 1 {1-(1-f)} = \sum_{n \in \bN} f (1-f)^n.
	\end{equation}
	Le terme général $f (1-f)^n$ est un polynôme en $f$ sans terme constant, donc un élément de l'algèbre $A$. De plus la série converge uniformément, donc sa somme appartient à $\overline A$, ce qui conclut.
\end{proof}

\begin{lemma}
\label{lem:StoneSqrt}
\leancheck{lem:StoneSqrt}%
	Sous les hypothèses du \cref{th:StoneWeierstrass}, soit $f \in \overline A$. Si $f\geq 0$ sur $X$, alors $\sqrt f \in \overline A$. On a toujours $|f|\in \overline A$. Si $f(x) \neq 0$ pour tout $x\in X$, alors $1/f \in \overline A$.
\end{lemma}

\begin{proof}
	\emph{Preuve du premier point.}
	Quitte à multiplier $f$ par une constante, on peut supposer que $0\leq f \leq 1$ sur $X$. 
	Pour tout $0 < \ve < 1$ on a 
	\begin{equation}
		\sqrt{f+\ve} = \sqrt{1+(f+\ve-1)} = \sum_{n \in \bN} \binom{1/2} n (f+\ve-1)^n,
	\end{equation}
	où $\binom \alpha n := \alpha (\alpha-1) \cdots (\alpha-n+1) / n !$. 
	Comme $-1 < f+\ve -1 < 1$ sur $X$, la série converge uniformément. Noter que le premier terme de cette série est une constante, d'où l'utilité du \cref{lem:StoneUnite}. Ainsi $\sqrt{f+\ve} \in \overline A$ pour tout $\ve>0$, et par fermeture $\sqrt f \in \overline A$.  
	
	\emph{Preuve du deuxième point.} Observer que $|f| = \sqrt{f^2} \in \overline A$.  
	
	\emph{Preuve du troisième point.}
	Supposons d'abord $f>0$. Quitte à multiplier $f$ par une constante, on peut supposer que $0<f\leq 1$, et alors la série de $1/f = 1/(1-(1-f))$ converge normalement. 
	Dans le cas général où on a seulement $f(x) \neq 0$ en tout $x\in X$, on note que $|f| \in \overline A$ est strictement positive sur $X$, puis que $1/f = f\times 1/|f| \times 1/|f| \in \overline A$.
\end{proof}

\begin{corollary}
\label{corol:StoneMinMax}
\leancheck{corol:StoneMinMax}%
	Sous les hypothèses du \cref{th:StoneWeierstrass}, soient $f,g \in \overline A$. Alors $\min\{f,g\}$, $\max\{f,g\} \in \overline A$.
\end{corollary}

\begin{proof}
	Les fonctions $\max \{f,g\} = (f+g+|f-g|)/2$ et $\min \{f,g\} = (f+g-|f-g|)/2$ appartiennent à $\overline A$, par linéarité et par le \cref{lem:StoneSqrt}. 
\end{proof}

\begin{lemma}
\label{lem:StoneDirac}
\leancheck{lem:StoneDirac}%
	Sous les hypothèses du \cref{th:StoneWeierstrass}, 
	pour tout voisinage ouvert $V$ de tout $x \in X$, et tout $\ve>0$, il existe $f \in A$ telle que $0 \leq f \leq 1$, $f(x)>1-\ve$ et $f< \ve$ sur $X \sm V$. 	
\end{lemma}

\begin{proof}
	Étant donné $y \in X$, on dispose par hypothèse (i) de $f_y\in \overline A$ telle que $f_y(x) \neq f_y(y)$. Quitte à considérer $\alpha f_y + \beta\bbone$, où $\alpha,\beta \in \bR$, on peut supposer que $f_y(x)=1$ et $f_y(y) = 0$, et poser $V_y := \{f_y < \ve\}$. 
	Soit $X = V \cup (\cup_{i=1}^I V_{y_i})$ une couverture finie de $X$, par compacité, et soit $f_* := \min_{1 \leq i \leq I} f_{y_i}$. Alors $f_*(x)=1$ et $f_* < \ve$ sur $\cup_{i=1}^I V_{y_i}$, donc $f_* < \ve$ dans le compact $X \sm V$. 
	On conclut en choisissant $f \in A$ assez proche uniformément de $\max\{\ve/2, \min\{1-\ve/2, f_*\}\} \in \overline A$. 
\end{proof}

\begin{proof}[Preuve du \cref{th:StoneWeierstrass}]
Soit $f \in C^0(X, \bR)$, et soit $\ve>0$; on souhaite montrer qu'il existe $f_* \in \overline A$ telle que $\|f-f_*\| \leq \ve$.  Quitte à considérer $\alpha f+ \beta \bbone$, où $\alpha,\beta \in \bR$ et comme $\bbone \in \overline A$, on peut supposer que $f \in C^0(X,[0,1])$.
Pour tout $x \in X$, soit $V_x := \{y \in X \mid |f(y)-f(x)|< \ve\}$, qui est un voisinage ouvert de $x$. Par compacité, soit $\bX\subset X$ fini tel que $X = \cup_{x\in \bX} V_x$.


Par le \cref{lem:StoneDirac}, pour tout $y \in X$ soit $x_y \in \bX$ tel que $y \in V_{x_y}$ et soit $\vp_y \in A$ telle que $0 \leq \vp_y \leq 1$, $\vp_y(y)>1-\ve$, et $\vp_y < \ve$ hors de $V_{x_y}$. Notons $W_y := \{\vp_y>1-\ve\} \subset V_{x_y}$, et par compacité considérons une couverture finie $X = \cup_{y \in \bY} W_y$ où $\bY \subset X$ est fini.
Alors $f_*(x) := \max_{y \in \bY} f(y) \vp_{y}(x)$ convient. 

En effet on a $f(y) \vp_{y} \leq f+2\ve$ sur $V_{x_y}$, et $f(y) \vp_y\leq \ve$ hors de $V_{x_y}$, donc $f_* \leq f+2 \ve$ sur $X$. 
Par ailleurs étant donné $x\in X$, soit $y \in \bY$ tel que $x \in W_y$, et soit $x':= x_y \in \bX$ de sorte que $W_y \subset V_{x'}$. Alors $|f(x)-f(y)| \leq |f(x)-f(x')| + |f(x')-f(y)| \leq 2 \ve$ et $\vp_y(x) \geq 1-\ve$, donc $f(y) \vp_{y}(x) \geq (f(x)-2 \ve)_+(1-\ve) \geq f(x)-3 \ve$, ce qui conclut. 
\end{proof}

La généralisation du \nref{th:StoneWeierstrass} aux fonctions complexes requiert une hypothèse supplémentaire de stabilité par conjugaison, dont nous illustrons la nécessité par un contre-exemple, voir les \cref{corol:StoneComplexe,prop:StoneCtrEx}.

\begin{corollary}[Stone-Weierstrass complexe]
\label{corol:StoneComplexe}
\leancheck{corol:StoneComplexe}%
	Soit $X$ un espace topologique compact, et soit $A \subset C^0(X,\bC)$ une algèbre qui
	\begin{enumerate}[(i)]
		\item sépare les points : $\forall x,y \in X$ distincts, $\exists f \in A$, $f(x) \neq f(y)$.
		\item ne s'annule pas simultanément en un point : $\forall x \in X$, $\exists f \in A$, $f(x)\neq 0$.
		\item est stable par conjugaison complexe : $\forall f \in A$, $\overline f \in A$.
	\end{enumerate}
	Alors $\overline A = C^0(X,\bC)$, pour la topologie de la convergence uniforme donnée par $\|\cdot\|_\infty$.
\end{corollary}

\begin{proof}
	Par stabilité par conjugaison, si $f \in A$ alors $\Re(f) := (f+\overline f)/2\in A$ et $\Im(f) = (f-\overline f)/(2i)\in A$. On en déduit que la $\bR$-algèbre $A \cap C^0(X,\bR)$ sépare les points et ne s'annule pas simultanément en un point, donc est dense dans $C^0(X,\bR)$. On conclut en approchant indépendamment la partie réelle et imaginaire d'un élément de $C^0(X,\bC)$ donné. 
\end{proof}

\begin{remark}[Application aux fonctions polynomiales]
\label{rem:polDens}
\leancheck{rem:polDens}%
Les polynômes réels $\bR[X_1,\cdots,\allowbreak X_d]$ sont denses dans $C^0(K,\bR)$,  pour la convergence uniforme, pour tout $K \csubset \bR^d$. 
De même, $\bC[X_1,\overline X_1,\cdots,X_d,\overline X_d]$ est dense dans $C^0(K',\bC)$, pour tout $K'\csubset \bC^d$. En particulier, $C^0(K,\bR)$ et $C^0(K',\bC)$ sont séparables.
\end{remark}

\begin{proof}
	Application directe du \cref{th:StoneWeierstrass}. La séparabilité découle de la densité de $\bQ[X_1,\cdots,\allowbreak X_d]$ dans $\bR[X_1,\cdots,X_d]$ (resp.\ $\bQ[i][X_1,\overline X_1,\cdots,X_d,\overline X_d]$ dans $\bC[X_1,\overline X_1,\cdots,X_d,\overline X_d]$).
\end{proof}

Attention : l'hypothèse (iii) de stabilité par conjugaison est essentielle dans le \nref{corol:StoneComplexe}, comme l'illustre le résultat élémentaire suivant d'analyse complexe. Celui-ci implique en effet que $\bC[X]$ \emph{n'est pas dense} dans $C^0(B,\bC)$, où $B:=B'_\bC(0,1)$, compte tenu de \eqref{eq:f_int_dB}. 
%Rappelons que le rayon de convergence d'une suite $(a_n)_{n \in \bN}$ est défini par $R:= \sup\{r \geq 0 \mid a_n r^n \text{ est bornée}\}$.

\begin{proposition}[Contre exemple à Stone-Weierstrass sans la stabilité par conjugaison \facultatif]
\label{prop:StoneCtrEx}
Soit $B := B'_\bC(0,1)$ la boule unité fermée du plan complexe, soit $\mathring B$ la boule ouverte, et soit $f \in C^0(B, \bC)$. Les propriétés suivantes sont équivalentes:
\begin{enumerate}[(i)]
	\item $f \in \overline{\bC[X]}$
	\item (Développabilité en série entière.) Il existe une suite $(a_n)_{n \geq 0} \in \bC^\bN$ bornée et telle que pour tout $z \in \mathring B$ on a %, de rayon de convergence au moins $1$, 
	\begin{equation*}
		f(z) = \sum_{n \geq 0} a_n z^n.
	\end{equation*}
	\item (Expression intégrale le long de $\partial B$.) Pour tout $z \in \mathring B$ on a 
	\begin{equation}
	\label{eq:f_int_dB}
		f(z) = \frac{1}{2i\pi}\int_0^{2 \pi} \frac{f(e^{it})}{e^{it}-z} \, i e^{it} \diff t.		
	\end{equation}
\end{enumerate}

\begin{proof}
	Preuve de (i) $\Rightarrow$ (iii). Supposons que $f(z) = \sum_{k=0}^K a_k z^k$ est un polynôme, de degré noté $K$. 
	Alors pour tout $z \in \mathring B$ on a 
	\begin{align}
	\nonumber
		\int_0^{2 \pi} \frac{f(e^{it})}{e^{it}-z} e^{it} \diff t 
		&= \int_0^{2 \pi} \frac{f(e^{it})}{1-z e^{-it}} \diff t \\
	\label{eq:f_dB_serie}
		&= \sum_{n \geq 0} \int_0^{2 \pi} f(e^{it}) (z e^{-it})^n \diff t \\
	\nonumber
		&= \sum_{0 \leq k \leq K} \sum_{n \geq 0}  z^n \int_0^{2 \pi} a_k e^{ikt}  e^{-int} \diff t \\
	\nonumber
		&= 2 \pi \sum_{0 \leq k \leq K} a_k z^k. 
	\end{align}
	Noter que la série dans \eqref{eq:f_dB_serie} est normalement convergente, ce qui justifie l'interversion avec l'intégrale.
	Ainsi \eqref{eq:f_int_dB} est satisfaite pour tout $|z|<1$ et pour toute $f \in \bC[X]$, donc aussi pour toute $f \in \overline {\bC[X]}$ car cette propriété est fermée sous convergence uniforme.
	
	Preuve de (iii) $\Rightarrow$ (ii). Poser $a_n := \int_0^{2 \pi} f(e^{it}) e^{-int} \diff t / (2 \pi)$, puis utiliser l'expression \eqref{eq:f_dB_serie}.
	
	Preuve de (ii) $\Rightarrow$ (i). On note que la série dans (ii) est uniformément convergente sur $B'(0,1-\ve)$ pour tout $\ve>0$. On en déduit que $f_\ve \in C^0(B)$, définie par $f_\ve(z) := f( (1-\ve) z)$, est limite uniforme de polynômes. 
	Finalement, on note que $f$ est uniformément continue sur $B$ par le \cref{th:Heine}, donc $f_\ve \to f$ uniformément lorsque $\ve\to 0$, et ainsi $f \in  \overline{\bC[X]}$ par fermeture, ce qui conclut.
\end{proof}
	
\end{proposition}

\subsection{Opérateurs linéaires compacts}
\label{subsec:linCompact}

Les opérateurs dits compacts, c.-à-d.\ dont l'image de la boule unité est précompacte, jouent un rôle important en analyse fonctionnelle et en particulier dans la théorie des espaces de Sobolev,  via le \nref{th:Rellich} d'injection compacte. 
Nous présentons ici quelques propriétés élémentaires, ainsi que le théorème de Fredholm qui est un point de départ en théorie spectrale. 
%, ainsi que leur lien avec les opérateurs de rang fini. 


\begin{retenir}
	La complète continuité des opérateurs compacts, 
	le théorème de Fredholm.
	(\Cref{prop:opCompactComplCont,th:Fredholm})
%	, et la continuité des opérateurs compacts de la topologie faible dans la forte (sur la boule unité et en supposant la \emph{propriété d'approximation}). 
%	(\Cref{th:Fredholm,prop:limRangFini})
\end{retenir}

\begin{definition} %[Opérateur linéaire compact]
\label{def:linCompact}
\leancheck{def:linCompact}%
	Un opérateur linéaire continu $T \in L(E,F)$ entre evn $E,F$ est dit compact si $T(B_E(0,1))$ est relativement compact (i.e.\ l'adhérence $\overline{T(B_E(0,1))}$ est compacte).
	%(i.e.\ $\overline{T(B_E(0,1))}$ est compact). 
	On note $L_c(E,F)$ leur ensemble.
\end{definition}

Présentons deux propriétés élémentaires mais fondamentales des opérateurs compacts.

\begin{proposition}[Fermeture par composition et limite des opérateurs compacts]
\label{prop:opCompactComposition}
\leancheck{prop:opCompactComposition}%
	La composition d'un opérateur compact avec un opérateur continu, à droite ou à gauche, est un opérateur compact.
\end{proposition}

\begin{proof}
	On désigne par $E,F,G$ des evn, et on note $B := B_E(0,1)$.
	Si $T \in L_c(E,F)$ et $A \in L(F,G)$, alors $A (\overline  {T B})$ est compact en tant qu'image continue d'un compact, et contient $A T B$ qui est donc relativement compact comme annoncé.
	Si $A \in L(E,F)$ et $T \in L_c(F,G)$,  alors $T(A(B)) \subset T(\vertiii{A} B) = \vertiii{A} T(B)$ qui est homothétique à $T(B)$ donc relativement compact.
\end{proof}

\begin{proposition}
\label{prop:opCompactLimite}
\leancheck{prop:opCompactLimite}%
	Soient $E,F$ des evn.
	Si $F$ est un Banach, alors l'ensemble $L_c(E,F)$ est fermé dans $L(E,F)$. 
	Si $F$ est un Hilbert, alors $L_c(E,F) = \overline{ \{u \in L(E,F)\mid u \text{ est de rang fini}\}}$.
\end{proposition}

\begin{proof}
Notons $B := B_E(0,1)$, et soit $T \in \overline {L_c(E,F)}$, où $E,F$ sont des evn. Soit aussi $\ve>0$. 

\emph{Premier point}.
Supposons que  $F$ est un Banach. Soit $T_0 \in L_c(E,F)$ tel que $\|T-T_0\|<\ve$, et soit $K_0 := \overline{T_0(B)}$. Par compacité, il existe une couverture finie $K_0 \subset \cup_{1 \leq i \leq I} B(y_i,\ve)$ avec $y_1,\cdots,y_I \in K_0$. Ainsi $T(B) \subset \cup_{1 \leq i \leq I} B(y_i,2\ve)$, ce qui établit la précompacité de cet ensemble puisque $\ve>0$ est arbitraire, comme annoncé. 

\emph{Second point.}
Soit $T \in L_c(E,F)$, où $F$ est un Hilbert. Soient $x_1,\cdots,x_I \in F$ tels que $T(B) \subset \cup_{1 \leq i \leq I} B(x_i,\ve)$. Soit $P$ la projection orthogonale sur $\Vect\{x_1,\cdots,x_I\}$. Alors $P \circ T$ est de rang fini, et pour tout $x \in B$ on a $\|T(x)-P(T(x))\| \leq \|T(x)-x_i\| + \|P(x_i - T(x))\| \leq 2 \ve$, où l'indice $1 \leq i \leq I$ est choisi tel que $T(x) \in B(x_i,\ve)$. Le résultat s'ensuit puisque $\ve>0$ est arbitraire.
\end{proof}	

\begin{proposition}[Complète continuité des opérateurs compacts]
\label{prop:opCompactComplCont}
\leancheck{prop:opCompactComplCont}%
	Soit $T \in L_c(E,F)$ un opé\-rateur linéaire compact entre des evn $E,F$, et soit $x_n \weakto x_*$ une suite à valeurs dans $E$ convergeant faiblement lorsque $n \to \infty$. 
	Alors $T x_n \to T x_*$ lorsque $n \to \infty$. 
	Si $E^*$ est séparable, alors la restriction $T : B'_E(0,1) \to F$ est continue de la topologie faible dans la topologie forte. 
\end{proposition}

\begin{proof}
	La preuve fait intervenir deux conséquences du théorème de Hahn-Banach: \cref{corol:faibleSepare} (la topologie faible est séparée), et surtout \cref{corol:faibleFortBorne} (une famille faiblement bornée est bornée); 
	il s'agit d'un analogue pour la topologie faible du \nref{th:BanachSteinhaus} précédemment établi pour la topologie $*$-faible.
	Comme la suite $x_n$ est faiblement convergente, elle est faiblement bornée, et donc bornée par le \cref{corol:faibleFortBorne}. % (qui découle lui-même du \nref{corol:BanachSteinhausK}).
	Soit $R := \sup_{n \in \bN} \|x_n\|$ et $K := \overline {T(B_E(0,1))}$ qui est compact par hypothèse. Considérons une valeur d'adhérence $T x_{\sigma(n)} \to y_*$ de la suite image, car celle-ci est à valeurs dans le compact $R K$. 
	Alors pour tout $\vp \in F^*$, on a $\<\vp, T x_{\sigma(n)}\> = \<T^* \vp, x_{\sigma(n)}\> \to \<T^* \vp, x_*\> = \<\vp, T x_*\>$ et par ailleurs $\<\vp, T x_{\sigma(n)}\> \to \< \vp, y_*\>$, ainsi $y_* = T x_*$ par le \cref{corol:faibleSepare}.
	Finalement, $y_*$ est l'unique valeur d'adhérence de la suite $(T x_n)_{n \in \bN}$ à valeurs dans le compact $R K$, donc $T x_n \to y_*$ par la \cref{prop:noAdhCv} comme annoncé.	
	
	Si $E^*$ est séparable, alors la topologie faible sur $B'_E(0,1)$ est métrisable par une adaptation immédiate de la \cref{prop:bouleMetrisable}, de même que la topologie forte sur l'evn $F$, ce qui permet d'utiliser la caractérisation séquentielle de la continuité \nref{lem:contDistance}.
\end{proof}

%Si la topologie faible sur $E$ est séparable, et $T \in L_c(E,F)$, alors la restriction $T : B'_E(0,1) \to F$ est continue de la topologie faible dans la topologie forte. 

La suite de cette section est une introduction à l'analyse du spectre des opérateurs compacts.

\begin{definition}
\label{def:linSpectre}
\leancheck{def:linSpectre}%
Le spectre d'un endomorphisme $T \in L(E)$ d'un Banach $E$ est l'ensemble 	
\begin{equation*}
	\sigma(T) := \{\lambda \in \bK \mid \lambda\Id - T \text{ n'a pas d'inverse continu}\},
\end{equation*}
dont les éléments sont appelés valeurs spectrales. Une valeur propre est une valeur spectrale $\lambda\in \sigma(T)$ pour laquelle l'espace propre $E_\lambda := \Ker(\lambda \Id-T)$ est non réduit à $\{0\}$.
\end{definition}


%Comme annoncé, nous évoquons d'abord les propriétés spectrales des opérateurs compacts - 
Si $T\in L_c(E,E)$ est compact, alors $E_1 := \Ker(\Id-T)$ est de dimension finie par le \nref{th:bouleCompRiesz}, puisque $B_{E_1}(0,1) \subset T(B_E(0,1))$ est relativement compacte. De même, tout sous-espace propre de $T$ associé à une valeur propre non-nulle est de dimension finie. Le théorème de Fredholm, présenté ci-dessous et précisé dans le \cref{corol:FredholmPrecise}, implique que les seules valeurs spectrales possibles pour un opérateur compact sont $0$ et ses valeurs propres.

\begin{theorem}[Fredholm]
\label{th:Fredholm}
\leancheck{th:Fredholm}%
	Soit $E$ un Banach, et soit $T \in L_c(E)$ un opérateur compact tel que $\Ker(\Id-T) = \{0\}$. Alors $\Image(\Id-T) = E$, et $\Id-T$ admet un inverse continu.
\end{theorem}

\begin{proof}
	Montrons d'abord que $\Image(\Id-T)$ est fermé. 
	Pour cela, considérons $f \in \overline {\Image(\Id-T)}$, qui est limite d'une suite de terme général noté $f_n = u_n - T u_n \in \Image(\Id-T)$. 
	Si par l'absurde $\|u_n\| \to \infty$, alors posons $v_n := u_n/\|u_n\|$. Par précompacité de $T(B_E(0,1))$, on dispose d'une sous-suite convergente de la forme $T v_{\vp(n)} \to w$, donc $v_n = T v_n + f_n/\|u_n\| \to w$. On obtient $T w = w$ par continuité de $T$, et $\|w\|=1$ par continuité de la norme, ce qui contredit l'hypothèse $\Ker(\Id-T) = \{0\}$. 
	Ainsi, $u_n$ reste bornée (quitte à extraire une sous-suite), disons $\|u_n\| \leq R$ pour tout $n \in \bN$, et donc son image par $T$ admet une sous-suite convergente notée $T u_{\vp(n)} \to g$, par précompacité de $T(B_E(0,R))$. 
	Finalement $u_{\vp(n)} = f_{\vp(n)} + T u_{\vp(n)} \to u_* := f+g$, et donc $f = u_*-T u_*$ par continuité de $T$, ce qui établit la fermeture de $\Image(\Id-T)$.

	Ainsi l'espace $F_1 := \Image(\Id-T)$ est un Banach, et une application du raisonnement précédent à la restriction $T_{|F_1}$ (qui est compact et laisse stable $F_1$) montre que $(\Id-T)F_1$ est fermé dans $F_1$, donc complet puisque $F_1$ est un Banach, donc fermé dans $E$. De même, par une récurrence immédiate, $F_n := (\Id-T)^n E$ est fermé pour tout $n \in \bN$.
	
%	Par récurrence, on obtient que $F_n := (\Id-T)^n(E)$ est fermé, donc est un Banach, pour tout $n \in \bN$. 
	Si par l'absurde $(\Id-T)E \neq E$, considérons $u_0 \in E\sm (\Id-T)E$, et notons que $u_n := (\Id-T)^n(u_0) \in  (\Id-T)^n E\sm (\Id-T)^{n+1}E$ pour tout $n \in \bN$ par injectivité de $\Id-T$, c.f.\ l'hypothèse $\Ker(\Id-T)=\{0\}$. Ainsi, on a une suite strictement décroissante pour l'inclusion de sous-espaces fermés $(F_n)_{n \in \bN}$. 
	Par le \nref{lem:normaleRiesz} il existe $x_n \in F_n \cap B'_E(0,1)$ tel que $d(x_n,F_{n+1}) \geq 1/2$, pour tout $n \in \bN$. Alors $\|T x_p - T x_q\| = \|x_p - [(\Id-T) x_p + Tx_q]\| \geq d(x_p,F_{p+1}) \geq 1/2$, pour tous $p<q \in \bN$. Ceci contredit la compacité de $\overline {T(B_E(0,1))}$, d'où l'on conclut que $(\Id-T)E = E$ comme annoncé.
	Finalement, $\Id-T$ est une bijection linéaire continue entre Banach, donc admet un inverse continu par le \nref{corol:isomBanach}. 
\end{proof}


Illustrons maintenant la notion d'opérateur compact, et l'alternative de Fredholm avec une équation fonctionnelle, la compacité de l'opérateur découlant du \nref{corol:AscoliEquicont}. 
Un analogue dans le cadre $L^2$ est présenté au \cref{ex:operCompactL2}, la compacité découlant dans ce cas du \nref{th:FrechetKolmogorov}.

\begin{lemma}[Compacité d'un opérateur intégral]
\label{lem:convK_compact}
	Soit $K \in C^0([0,1]^2,\bR)$, et soit 
	\begin{equation*}
		\cK(f)(x) := \int_0^1 K(x,y) f(y) \diff y.
	\end{equation*}
	pour tous $f \in C^0([0,1],\bR)$ et $x \in [0,1]$. 
	Alors $\cK$ est un endomorphisme compact de $C^0([0,1],\bR)$. 
\end{lemma}

\begin{proof}
	Remarquons d'abord que $\|\cK(f)\|_\infty \leq \|K\|_\infty \|f\|_\infty$. 
	Soit $\omega$ un module d'uniforme continuité de $K$ sur $[0,1]^2$. Alors pour tous $x,y \in [0,1]$ on a 
	\begin{equation*}
		|\cK(f)(x)-\cK(f)(y)| 
		%= | \int_0^1 (K(x,z)-K(y,z)) f(z) \diff z| 
		\leq \int_0^1 |K(x,z)-K(y,z)| |f(z)| \diff z \leq \omega(|x-y|) \|f\|_\infty.
	\end{equation*}
	Ainsi l'ensemble $\{\cK(f) \mid \|f\|_\infty \leq 1\}$ est uniformément borné et uniformément équicontinu, donc relativement compact par le \nref{corol:AscoliEquicont}.
\end{proof}

\begin{corollary}
\label{corol:eqFuncNoyau}
	Soit $K \in C^0([0,1]^2,\bR_+)$ tel que $K(x,y) = K(y,x)$ pour tous $x,y \in [0,1]$. Alors pour toute $g \in C^0([0,1])$, il existe $f \in C^0([0,1])$ telle que 
	\begin{equation*}
		\forall x \in [0,1],\ \int_0^1 K(x,y) (f(y)-f(x)) \diff y - f(x) = g(x). 
	\end{equation*}
\end{corollary}

\begin{proof}
	Posons $k(x) := \int_0^1 K(x,y) \diff y$, pour tout $x\in [0,1]$.
	On peut réécrire l'égalité souhaitée sous la forme 
	\begin{equation*}
		\cK(f) - (1+k) f = g,
	\end{equation*}
	avec les notations du \cref{lem:convK_compact}. Soit encore 
	\begin{align*}
		\hat \cK(f) - f &= \hat g, &
		\text{avec } \hat \cK(f) &:= (1+k)^{-1}\cK(f), &
		\text{et } \hat g := (1+k)^{-1} g.
	\end{align*}
	L'opérateur $\hat \cK$ est compact en tant que composition d'un opérateur compact, voir le \cref{lem:convK_compact}, et d'un opérateur continu, et on a bien $\hat g \in C^0([0,1])$.
	D'après l'alternative de Fredholm, il suffit donc de vérifier que $\hat \cK - \Id$ est injectif, pour obtenir sa surjectivité, et donc l'existence annoncée. Or pour tout $f \in C^0([0,1])$,
	\begin{align*}
		\hat \cK(f) = f \ 
		& \Rightarrow \ \forall x\in [0,1],\, \int_{y \in [0,1]} K(x,y) (f(y)-f(x)) \diff y = f(x)\\
		& \Rightarrow \int_{x\in [0,1]} \int_{y\in [0,1]}  K(x,y) (f(y)-f(x)) f(x) \, \diff y \, \diff x = \int_0^1 f(x)^2 \diff x.
	\end{align*}
	Or la première intégrale se réécrit, par symétrie de $K$, sous la forme suivante, négative ou nulle
	\begin{equation*}
		-\frac 1 2 \int_{x\in [0,1]} \int_{y\in [0,1]} K(x,y) |f(y)-f(x)|^2 \diff y \diff x.
		\qedhere
	\end{equation*}	
\end{proof}

Comme annoncé, nous précisons maintenant le \nref{th:Fredholm}. Pour inspiration, on peut  considérer un autre exemple d'opérateur compact, plus élémentaire : l'opérateur $L^\infty(\bN) \to L^\infty(\bN)$, $(x_n)_{n \geq 0} \mapsto (\lambda_n x_n)_{n \geq 0}$, où $(\lambda_n)_{n \in \bN}$ est une suite convergeant vers $0$. Son spectre est l'ensemble $\{\lambda_n \mid n \in \bN\} \cup \{0\}$. 
%Encore un autre exemple : l'opérateur
%$(x_n,y_n)_{n \geq 0} \mapsto (\ve_n (y_n,-x_n))_{n \geq 0}$ avec $\ve_n \to 0$, $\ve_n \neq 0$, défini sur $L^\infty(\bN\times \bN)$.
%Si $\bK = \bR$, alors $T$ n'a pas de valeur propre, et $0$ est sa seule valeur spectrale. \\


\begin{corollary}[Fredholm précisé]
\label{corol:FredholmPrecise}
\leancheck{corol:FredholmPrecise}%
	Soit $E$ un Banach, et $T \in L_c(E)$. Alors % un opérateur compact. Alors :
	\begin{enumerate}[(i)]
		\item Si $\dim(E)=\infty$ alors $0 \in \sigma(T)$. (Mais $0$ n'est pas forcément valeur propre.)
		\item Pour tout $\lambda \in \sigma(T)\sm \{0\}$, il existe $x \in E\sm \{0\}$ tel que $T x = \lambda x$. (i.e.\ $\lambda$ est valeur propre)
		\item Pour tout $\lambda \in \sigma(T)\sm \{0\}$, il existe $m = m(\lambda)$, tel que $\Ker( (\lambda\Id-T)^m ) = \Ker((\lambda\Id-T)^{m+1})$. De plus $\dim \Ker( (\lambda\Id-T)^m ) < \infty$. (i.e.\ la multiplicité de $\lambda$ est finie)
		\item L'ensemble $\sigma(T)\subset B'(0,\|T\|)$ est dénombrable, et $0$ est son seul point d'accumulation possible.
	\end{enumerate}
\end{corollary}

\begin{proof}
	Preuve de (i). Par contraposée, supposons $0 \notin \sigma(T)$. Alors $T^{-1}$ existe et est continu, par définition du spectre, donc $T^{-1}(\overline {T(B'(0,1))}) \supset B'(0,1)$ est compact comme image directe d'un compact par une application continue. Par le \nref{th:bouleCompRiesz}, $E$ est donc de dimension finie.

	Preuve de (ii). Appliquer le \nref{th:Fredholm} à $T/\lambda$.
	
	Preuve de (iii). 
	Posons $F_n := \Ker((\lambda\Id-T)^n)$ et montrons d'abord qu'il existe $n\in \bN$ tel que $F_n = F_{n+1}$ (ce qui implique $F_n = F_m$ pour tout $m \geq n$).
	Sinon, par l'absurde et en utilisant le \nref{lem:normaleRiesz}, choisissons $x_n \in F_n$ tel que $\|x_n\| = 1$ et $d(x_n,F_{n-1}) \geq 1/2$, avec la convention $F_{-1} = \{0\}$.
	Par construction, $(\lambda \Id-T) x_n \in F_{n-1}$ pour tout $n \geq 0$, et de plus $F_n$ est stable par $T$. 
	Ainsi, pour tous $n>m$ on a $\|T x_n - T x_m\| = 
	\|\lambda x_n - [(\lambda \Id-T)x_n + T x_m]\| \geq |\lambda|/2$.
	Ceci contredit la compacité de $T$, d'où l'existence de l'entier $m(\lambda)$ annoncée.
	
	Par ailleurs, si $x \in \Ker( (\lambda\Id-T)^n)$ alors en développant le binôme de Newton dans l'identité $(\lambda\Id - T)^n x = 0$, on obtient
	\begin{align*}
		x &= T P(T) x, & 
		\text{avec }
		P(T) &:= \lambda^{-n} \sum_{0 \leq k\leq n-1} {\textstyle \binom n k} \lambda^k (-T)^{n-1-k}.
	\end{align*}
	 Or $T P(T)$ est un opérateur compact, comme composée d'un opérateur continu et d'un opérateur compact. Ainsi $\Ker( (\lambda\Id-T)^n) \cap B'(0,1) \subset T P(T) B'(0,1)$ qui est un ensemble précompact, donc $\Ker( (\lambda\Id-T)^n)$ est de dimension finie par le \nref{th:bouleCompRiesz}.
	
	Preuve de (iv).
	Considérons par l'absurde un point d'accumulation non-nul de $\sigma(T)$, noté $\lambda_* \in \bK\sm \{0\}$. Par hypothèse, il existe donc une suite de valeurs spectrales deux à deux distinctes de $T$, notées $(\lambda_n)_{n \in \bN}$ et que l'on peut supposer non-nulles, convergeant vers $\lambda_*$.
	Par le point (ii) il existe un vecteur propre $x_n \in E\sm \{0\}$ pour $\lambda_n$. Notons $F_n := \Vect\{x_1,\cdots,x_n\}$, et choisissons $y_n\in F_n$ tel que $d(y_n,F_{n-1}) \geq 1/2$ et $\|y_n\|=1$, par le \nref{lem:normaleRiesz}. 
	Alors $(\lambda_n \Id-T) y_n  \in F_{n-1}$. Le raisonnement précédent s'applique : on a 
	$\|T y_n - T y_m\| = 
	\| \lambda_n y_n - [(\lambda_n \Id - T) y_n + T y_m]\| \geq |\lambda_n|/2$ pour tous $n>m$, ce qui contredit la compacité de $T$. 
	
	Pour tout $\lambda \in \sigma(T)$, on a $|\lambda| \leq \vertiii{T}$; en effet, dans le cas contraire, la série $\lambda^{-1} \sum_{n \geq 0} (T/\lambda)^n$ est sommable, et fournit un inverse à $\lambda \Id - T$. 
	Comme $0$ est le seul point d'accumu\-lation de $\sigma(T)$, l'ensemble $\sigma(T) \sm B(0,\ve)$ est fini pour tout $\ve>0$. On en déduit que $\sigma(T)$ est dénombrable.
\end{proof}



\newpage
\section{Dualité et topologie faible}
\label{sec:dual}
Cette section est consacrée à la dualité, et à la topologie faible dans les evn. 
Le point de départ est le théorème de Hahn-Banach, qui établit que le dual continu $E^*$ d'un evn $E$ contient des éléments non triviaux. 
La sous-section suivante est dédiée à la propriété de réflexivité qui, lorsqu'elle est satisfaite, permet d'identifier $E$ avec son bidual $E^{**}$. Finalement, on présente les formes géométriques du théorème de Hahn-Banach.

Deux annexes complètent cette discussion.
L'\cref{sec:Hilbert} présente le cadre hilbertien, dans lequel la discussion précédente devient essentiellement triviale grâce au \nref{th:reprRiesz} et au  \nref{th:projConvexe}. 
L'\cref{sec:LegendreFenchel} est une introduction à l'optimisation convexe, via le \nref{th:LegendreFenchel} et le \nref{th:FenchelRockafellar}.

Pour la simplicité, on considère seulement des $\bR$-evn dans cette section.

\inventaire{
\Nref{subsec:HahnBanach}\\
\Nref{def:inductif}\\
\Nref{lem:Zorn}\\
\Nref{def:axiomChoix}\\
\Nref{def:sousLin}\\
\Nref{th:HahnBanach}\\
\Nref{corol:prol_forme_evn}\\
\Nref{corol:dens_banach}\\
\Nref{ex:dualSeparable} (Si le dual est séparable, alors le primal aussi.)\\
\Nref{corol:projDimFinie}\\

\Nref{subsec:reflexivite}\\
\Nref{def:reflexif}\\
\Nref{prop:conjugueDual}\\
\Nref{corol:faibleSepare} (La topologie faible sur un evn est séparée)\\
\Nref{corol:faibleFortBorne} (Toute partie faiblement bornée d'un evn est bornée)\\
\Nref{corol:bidualComplete} ($\overline{\Phi(E)}$ s'identifie au complété de $E$)\\
\Nref{th:unifCvxReflex} (Uniforme convexité implique réflexivité)\\
\Nref{rem:unifCvxHilbert}\\
\Nref{th:HolderIneq}\\
\Nref{prop:ClarksonIneq}\\
\Nref{th:lpDualLq}\\
\Nref{ex:nonReflex}\\

\Nref{subsec:HahnBanachGeom}\\
\Nref{prop:jaugeCvx}\\
\Nref{lem:jaugeOuvertCvx} (Jauge d'un ouvert et d'un fermé)\\
\Nref{th:separOuvertPt}\\
\Nref{corol:separOuvertPt} (Séparation de $0$. Séparation d'un convexe et d'un ouvert convexe.)\\
\Nref{th:separFermeComp}\\
\Nref{def:convexePolaire}\\
\Nref{prop:cvxFaibleFerme}\\
\Nref{corol:cvxFermeBdComp} (Un convexe fermé borné d'un réflexif est faiblement compact)\\
\Nref{corol:minFctCvxRefl}
}


\subsection{Théorème de Hahn-Banach}
\label{subsec:HahnBanach}

Le théorème de Hahn-Banach permet d'établir l'existence de prolongements continus de formes linéaires, dans un espace vectoriel normé de dimension infinie. Il montre que ces formes linéaires sont suffisamment nombreuses pour rendre la topologie faible intéressante. La preuve du théorème de Hahn-Banach fait intervenir l'axiome du choix, via le Lemme de Zorn, dont l'utilisation est assez courante en analyse.


\begin{retenir}
Le Lemme de Zorn et le théorème de Hahn-Banach, avec les définitions associées. Le prolongement de même norme d'une forme linéaire, et le critère de densité associé.
(\Cref{def:inductif,lem:Zorn,def:sousLin,th:HahnBanach,corol:prol_forme_evn,corol:dens_banach}.)
\end{retenir}

\begin{definition}[Ensemble inductif]
\label{def:inductif}
	Un ensemble (partiellement) ordonné $(E, \leq)$ est dit inductif si toute partie \emph{totalement} ordonnée $F\subset E$ admet un majorant dans $E$ : 
	\begin{equation}
		\forall F \subset E,\ (\forall y,z \in F,\ y \leq z \text{ ou } z \leq y) \Rightarrow (\exists x \in E,\ \forall y \in F,\ x \geq y). 
	\end{equation} 
\end{definition}

``Ce n'est pas un Lemme, et il n'est pas de moi'' aurait déclaré Max August Zorn au sujet du résultat suivant, désormais associé à son nom.
\begin{lemma}[Zorn] 
\label{lem:Zorn}
	Tout ensemble non-vide et inductif admet un élément maximal. 
\end{lemma}

Le Lemme de Zorn est équivalent à l'axiome du choix, moyennant les autres axiomes de la théorie de Zermelo-Fraenkel des ensembles, voir \cite{krivine2007theorie}. On présente ci-dessous la preuve de l'axiome du choix à partir du lemme de Zorn, qui est élémentaire et assez représentative de l'utilisation de ce dernier en analyse.



\begin{definition}[Axiome du choix \facultatif]
\label{def:axiomChoix}
Soit $\cA$ un ensemble d'ensembles non vides. Alors il existe $f : \cA \to \cup\cA$, appelée \emph{fonction de choix}, telle que $f(a) \in a$ pour tout $a \in \cA$. 	
\end{definition}


\begin{proof}[Preuve de l'axiome du choix à partir du Lemme de Zorn]
	Soit $\cA$ un ensemble d'ensembles non-vides, et soit $\cB := \cup \cA$ sa réunion. On définit $E := \{f : A \to \cB \mid A \subset \cA \text{ et } \forall a \in A, f(a) \in a\}$ l'ensemble des fonctions de choix partielles sur $\cA$, ordonné par la restriction : $(f_0 : A_0 \to \cB) \leq (f_1 : A_1 \to \cB)$ ssi $A_0 \subset A_1$ et $f_{1|A_0} = f_0$. 
	Cet ensemble est inductif : si $(f_i : A_i \to \cB)_{i \in I}$ est une famille totalement ordonnée, alors on peut définir $A_* := \cup_{i \in I} A_i$ et $f_* : A_* \to \cB$ par ses restrictions $f_{*|A_i} = f_i$ pour tout $i \in I$, qui sont consistantes les unes avec les autres par définition de l'ordre, ce qui donne un majorant de la famille. 
	
	Par le Lemme de Zorn, l'ensemble $E$ admet un élément maximal noté $f_0 : A_0 \to \cB$. Si par l'absurde $A_0 \neq \cA$, alors on peut choisir $a \in \cA\sm A_0$ et un élément $\alpha \in a$, puis considérer l'extension définie par $A_1 := A_0 \cup\{a\}$ et $f_1(a) := \alpha$, $f_{1|A_0} = f_0$. Cela contredit la maximalité de $f_0$, et achève ce raisonnement par l'absurde. 
	Ainsi $A_0 = \cA$ et $f_0$ est une fonction de choix sur $\cA$, ce qui conclut. % établit l'axiome du choix, et conclut la preuve. 
\end{proof}


\begin{definition}[Fonction sous-linéaire]
\label{def:sousLin}
Une fonction $\rho : E \to \bR$, sur un $\bR$-espace vectoriel $E$, est dite sous-linéaire si elle est positivement $1$-homogène et convexe. De manière équivalente
\begin{align*}
	\rho(\lambda x) &= \lambda \rho(x),&
	\rho(x+y) &\leq \rho(x)+\rho(y), &
	\forall \lambda \in \bR_+,\ x, y \in E.
\end{align*}
\end{definition}

Noter qu'une semi-norme est une application sous-linéaire qui de plus est \emph{positive et symétrique}: $\forall x\in E, \rho(-x) = \rho(x) \geq 0$.

\begin{theorem}[Hahn-Banach]
\label{th:HahnBanach}
	Soit $E$ un $\bR$-espace vectoriel, $F$ un sous espace vectoriel, et $\rho : E \to \bR$ sous-linéaire. 
	Soit $\vp_F : F \to \bR$ une forme linéaire telle que $\vp_F \leq \rho$ sur $F$.
	Alors il existe $\vp : E \to \bR$ linéaire et telle que $\vp \leq \rho$ sur $E$ et $\vp_{|F} = \vp_F$. 
\end{theorem}

\begin{proof}
Introduisons l'ensemble des extensions partielles satisfaisant aux hypothèses : 
\begin{equation*}
	\cE := \{\vp : G \to \bR \mid G \text{ s.e.v.\ de } E  \text{ t.q. }  F \subset G,\,  \vp \leq \rho \text{ sur } G,\, \vp \text{ linéaire et } \vp_{|F} = \vp_F\}
\end{equation*}
L'ensemble $\cE$ est ordonné par la restriction : $(\vp : G\to \bR) \leq (\psi : H \to \bR)$ ssi $G \subset H$ et $\psi_{|G} = \vp$. 
Cet ensemble est inductif : si $(\vp_i : F_i \to \bR)_{i \in I}$ est totalement ordonnée, alors on peut définir $F_* := \cup_{i \in I} F_i$ et $\vp_* : F_* \to \bR$ par ses restrictions $\vp_{*|F_i} = \vp_i$ qui sont consistantes, ce qui donne un majorant. Il contient $\vp_F : F \to \bR$, donc est non-vide.

Par le \nref{lem:Zorn}, $\cE$ admet un élément maximal, noté $\vp : G \to \bR$. Supposons par l'absurde que $G \neq E$, autrement le résultat est prouvé. Soit $x \in E \sm G$ et soit $H := \{\lambda x+e \mid \lambda \in \bR,\, e \in G\}$ l'espace vectoriel engendré. On souhaite définir une extension $\psi : H \to \bR$ de la forme $\psi(\lambda x + e) = \lambda \alpha + \vp(e)$ pour tous $\lambda \in \bR$, $e \in G$, ce qui contredira la maximalité de $\vp$.

Il faut pour cela montrer l'existence de $\alpha \in \bR$ tel que $\lambda \alpha + \vp(e) \leq \rho(\lambda x + e)$ pour tous $\lambda \in \bR$ et $e \in G$. 
On peut se ramener à $\lambda = \pm 1$, en utilisant la positive homogénéité de $\rho$, d'où la condition équivalente suivante : %et cette condition est donc équivalente à : 
\begin{equation*}
	\sup_{e \in G} \vp(e) - \rho(e-x) \leq \alpha \leq \inf_{f \in G} \rho(f+x)-\vp(f). 
\end{equation*}
Or pour tous $e,f \in G$ on a $\vp(e)+\vp(f) = \vp(e+f) \leq \rho(e+f) \leq \rho(e-x) + \rho(f+x)$, par sous-linéarité de $\rho$, donc le $\sup$ est majoré par l'$\inf$, et le choix de $\alpha$ est possible, ce qui conclut. 
\end{proof}

Le \nref{th:HahnBanach} s'applique aux evn en choisissant typiquement $\rho$ proportionnel à la norme, ce qui permet de prolonger une forme linéaire continue, et d'obtenir un critère de densité, voir ci-dessous. 
%(On pourrait également prolonger une forme linéaire continue sur un sous-espace d'un evtlc.)
\begin{corollary}[Prolongement de même norme d'une forme linéaire]
\label{corol:prol_forme_evn}
	Soit $E$ un evn, et soit $\vp_F \in F^*$ où $F\subset E$ est un s.e.v.
	Alors il existe $\vp\in E^*$ prolongeant $\vp_F$ et de même norme. 
\end{corollary}

\begin{proof}
	Poser $\rho(x) := \|\vp_F\|_{F^*} \|x\|_E$, qui est bien sous-linéaire et domine $\vp_F$, puis appliquer le \cref{th:HahnBanach} de Hahn-Banach.
\end{proof}

\begin{corollary}[Critère de densité]
\label{corol:dens_banach}
	Soit $E$ un $\bR$-evn, et $F \subset E$ un sev. Alors $F$ est dense ssi la seule forme linéaire continue s'annulant sur $F$ est $0$.
\end{corollary}

\begin{proof}
	Si $F$ est dense et $\vp \in E^*$ est nulle sur $F$, alors par continuité $\vp = 0$ sur $E$, comme annoncé. 
	Réciproquement, supposons $\overline F \neq E$, et par le \nref{lem:normaleRiesz} choisissons $x_0\in E$ tel que $\|x_0\|_E = 1$ et $d(x_0,F)\geq 1/2$. Posons $\vp : \overline F \oplus \bR x_0 \to \bR$, définie par $\vp(x + \lambda x_0) := \lambda$, et notons que $\|x+\lambda x_0\| \geq d(x+\lambda x_0,F) = |\lambda| d(x_0,F) \geq |\lambda|/2$, de sorte que $\vp$ est continue. 
	Par le \cref{corol:prol_forme_evn} de Hahn-Banach, $\vp$ admet un prolongement continu et de même norme, qui est donc une forme linéaire continue s'annulant sur $F$ mais pas sur $E$. 
\end{proof}

\begin{exemple}
\label{ex:dualSeparable}
	Soit $E$ un evn. Si $E^*$ est séparable, alors $E$ est séparable. 
\end{exemple}

\begin{proof}
	Soit $(\vp_n)_{n \in \bN}$ une famille dense dans $E^*$, et soit $(x_n)_{n \in \bN}$ une famille de $E$ telle que $\vp_n(x_n) \geq \|\vp_n\|/2$ et $\|x_n\|=1$ pour tout $n \in \bN$ (celle-ci existe par définition de la norme duale).
	Posons $F := \Vect\{x_n \mid n \in \bN\}$, qui est un espace vectoriel séparable (considérer toutes les combinaisons linéaires finies à coefficients rationnels), et considérons $\vp \in E^*$ s'annulant sur $F$. Par l'absurde supposons $\vp \neq 0$; on peut donc supposer $\|\vp\|=1$ par normalisation. Choisissons $n\in \bN$ tel que $\|\vp -\vp_n\| \leq \frac 1 4$. 
	On obtient alors la contradiction suivante : 
	\begin{equation*}
		\vp(x_n) \geq \vp_n(x_n) - \|\vp-\vp_n\| \geq \frac 1 2 \|\vp_n\| - \|\vp-\vp_n\| \geq \frac 1 2 \|\vp\| - \frac 3 2 \|\vp-\vp_n\| \geq \frac 1 8 > 0.
	\end{equation*} 
	Ainsi $\vp$ est nécessairement nulle, donc $F$ est dense dans $E$ par le \cref{corol:dens_banach}, ce qui conclut. 
\end{proof}


Dans un evn général, et contrairement au cas d'un Hilbert discuté dans le \cref{th:projConvexe}, on ne dispose pas de projection continue sur tout sous espace fermé. %En particulier, un sous-espace fermé n'a pas forcément de supplémentaire fermé. 
Le résultat suivant permet toutefois, via le \nref{th:HahnBanach}, de construire une projection continue sur tout sous-espace de dimension finie. Noter qu'un théorème de Kadets-Snobar montre plus précisément l'existence d'une projection $P_F$ de norme au plus $\sqrt{\dim(F)}$. 

\begin{corollary}[Projection sur un sous espace de dimension finie]
\label{corol:projDimFinie}
Pour tout sous-espace de dimension finie $F$ d'un evn $E$, il existe une projection linéaire continue $P_F : E \to F$. 
\end{corollary}

\begin{proof}
	Soit $(e_1,\cdots,e_n)$ une base de $F$, et soit $(\vp_1,\cdots,\vp_n) \in F^*$ la base duale, c'est-à-dire que $\vp_i(e_j) = \delta_{ij}$ pour tous $1 \leq i,j \leq n$. Le projecteur s'obtient en prolongeant chacune des $\vp_i$ de manière continue à $E$ par le \cref{corol:prol_forme_evn}, et en posant $P_F(x) := \sum_{1 \leq i \leq n} \vp_i(x) e_i$. 
\end{proof}

\subsection{Réflexivité}
\label{subsec:reflexivite}

La propriété de réflexivité, lorsqu'elle est satisfaite, permet d'identifier un evn $E$ avec son bidual $E^{**}$, et la topologie faible sur $E$ avec la topologie *faible sur $E^{**}$. Les résultats fondés sur cette dernière topologie, comme le \nref{th:BanachAlaoglu}, s'appliquent donc aux espaces réflexifs.

\begin{retenir}
La définition d'un espace réflexif, l'inégalité de Hölder et la dualité dans les espaces $L^p$.
(\Cref{def:reflexif,th:HolderIneq,th:lpDualLq}.)	
\end{retenir}


\begin{definition}[Réflexivité]
\label{def:reflexif}
	La fonction d'évaluation $\Phi : E \to E^{**}$, d'un evn dans son bidual, est définie par
	\begin{equation*}
			\forall x \in E,\, \forall \vp \in E^*,\,  \Phi(x)(\vp) := \vp(x).
	\end{equation*}
	L'evn $E$ est dit réflexif si $\Phi$ est bijective.
\end{definition}


Pour étudier les propriétés de la fonction d'évaluation $\Phi$ nous introduisons la notion de \emph{conjugué dual}, dont l'existence est assurée par le \nref{th:HahnBanach}. Comme illustré dans le \cref{th:HolderIneq} consacré aux espaces de Lebesgue $L^p$, il peut exister plusieurs conjugués duaux, et plusieurs éléments peuvent avoir le même conjugué dual.

\begin{proposition}[Élément conjugué dual]
\label{prop:conjugueDual}
	Pour tout élément  $x \in E$ d'un evn, il existe $\vp\in E^*$ telle que 
	\begin{align*}
		\vp(x) &= \|x\|_E \|\vp\|_{E^*}, &
		\|\vp\|_{E^*} &= \|x\|_E.
	\end{align*}
	On en déduit que la fonction d'évaluation $\Phi : E \to E^{**}$ est isométrique, donc aussi injective.
\end{proposition}

\begin{proof}
	On obtient un conjugué dual $\vp$ en appliquant le \nref{corol:prol_forme_evn} au sous espace unidimensionnel $F := \bR x$ et à la forme linéaire $\vp_F(\lambda x) := \lambda \|x\|^2_E$.
	
	On a $\|\Phi(x)\|_{E^{**}} = \sup \{|\vp(x)| \mid \vp \in E^*,\, \|\vp\|_{E^*} = 1\} \leq \|x\|_E$. De plus, en choisissant pour $\vp$ un conjugué dual de $x$, on obtient $\|\Phi(x)\|_{E^{**}} \geq |\vp(x)|/\|\vp\|_{E^*} = \|x\|_E$, ce qui prouve l'isométrie et donc l'injectivité. 
\end{proof}

À titre illustratif, mentionnons quelques conséquences immédiates des propriétés de $\Phi$.
%L'application $\Phi$ permet d'importer sur un evn $E$ des résultats montr \nref{corol:BanachSteinhausK} 


\begin{corollary}
\label{corol:faibleSepare}
	La topologie faible sur un evn $E$ est séparée.
\end{corollary}

\begin{proof}
%Si $x\in E$ est tel que $\vp(x)=0$ pour tout $\vp \in E^*$, alors $x=0$ par le \cref{corol:prol_forme_evn}. 
D'après le \nref{lem:evtlc_separe}, il faut montrer que $(\forall \vp \in E^*,\ \vp(x)=0) \Rightarrow x=0$, pour tout $x \in E$. On conclut en choisissant pour $\vp$ un élément conjugué dual. 
\end{proof}

\begin{corollary}
\label{corol:faibleFortBorne}
	Toute partie $A$ d'un evn $E$ qui est faiblement bornée (i.e.\ $\{\vp(x) \mid x \in A\}$ est bornée pour toute $\vp \in E^*$), est bornée.
\end{corollary}

\begin{proof}
	%Une partie $A \subset E$ est faiblement bornée ssi $\{\vp(a) \mid a \in A\}$ est borné pour tout $\vp \in E^*$; de manière équivalente $\Phi(A) \subset E^{**}$ est *faiblement bornée. 
	Remarquer que $\Phi(A) \subset E^{**}$ est *faiblement bornée, et que $E^* = L(E,\bR)$ est complet comme tout espace dual, voir la \cref{prop:LEFcomplet}.
	Par le \nref{corol:BanachSteinhausK}, $\Phi(A)$ est donc fortement bornée, donc $A\subset E$ aussi par isométrie.
\end{proof}


\begin{corollary}
\label{corol:bidualComplete}
L'image fermée $\overline {\Phi(E)} \subset E^{**}$ d'un evn $E$ s'identifie à son complété. % s'identifie au complété de $E$, voir la \cref{prop:espaceComplete}.  	
\end{corollary}

\begin{proof}
	Pour tout evn $E$, le dual $E^*$ est complet par la \cref{prop:LEFcomplet}, et donc $E^{**}$ est également complet. Ainsi $\overline {\Phi(E)}$ est un complet, en tant que partie fermée d'un complet, et contient le sous-ensemble dense $\Phi(E)$ qui est en bijection isométrique avec $E$. Par la \nref{prop:espaceComplete}, $\overline{\Phi(E)}$ s'identifie au complété de $E$.
\end{proof}


Pour établir la propriété de réflexivité, on dispose de la condition géométrique suffisante suivante.
Sa preuve, élémentaire mais assez technique, est admise. Il existe également un critère nécessaire et suffisant, donné par le théorème de James. 

%\begin{theorem}[Théorème de James]
%Soit $E$ un Banach et $B := B'_E(0,1)$ sa boule unité fermée. Sont équivalents:
%\begin{enumerate}[(i)]
%	\item $E$ est réflexif.
%	\item $E^*$ est réflexif.
%	\item La boule unité $B$ est faiblement compacte.
%	\item Pour tout $\vp\in E^*$, $\exists x \in B$, $\vp(x) = \sup_B \vp$. 
%\end{enumerate}
%\end{theorem}
%
%\begin{proof}[Preuve des implications faciles]
%	On a (i) $\Leftrightarrow$ (ii). En effet de l'identification $E \approx E^{**}$ on obtient par dualité $E^* \approx E^{***}$. 
%	Réciproquement supposons $E\subsetneq E^{**}$, et remarquons que $E$ est un Banach, donc un sev fermé, donc n'est pas dense, et donc il existe une forme linéaire continue s'annulant sur $E$ mais pas sur $E^{**}$, de sorte que $E^* \subsetneq E^{***}$.
%
%	On a (i) $\Rightarrow$ (iii) car la topologie faible sur $E$ s'identifie à la topologie *faible sur $E^**$ et par le \cref{th:BanachAlaoglu} de Banach Alaoglu. 
%	On a (iii) $\Rightarrow$ (iv) car les formes linéaires sont faiblement continues, et car $B$ est faiblement compacte. L'implication (iv) $\Rightarrow$ (i), plus difficile, est l'objet principal du Théorème de James et admise.
%\end{proof}


\begin{theorem}[Milman-Pettis \facultatif]
\label{th:unifCvxReflex}
	Un evn $E$ est dit uniformément convexe si $\forall \ve>0$, $\exists \delta>0$, $\forall x,y \in E$,
	\begin{align}
	\label{eq:unif_cvx_ball}
		(\|x\|=\|y\|=1 \text{ et } \|x-y\| \geq \ve) \Rightarrow \big\|\frac{x+y} 2\big\| \leq 1-\delta.
	\end{align}
	Tout espace de Banach uniformément convexe est réflexif. 
	%es espaces uniformément convexes sont parfois appelés super-réflexifs
\end{theorem}

%L'uniforme convexité implique la réflexivité, comme le montre la proposition ci-dessous. 
%Noter que les espaces uniformément convexes sont parfois appelés super-réflexifs, et que leur réflexivité peut être établie de manière élémentaire (sans passer par le théorème de James).

%\todoi{Peut-être préférer l'approche élémentaire pour l'année prochaine ? Éventuellement passer la preuve, ou la proposer en DM ?} 
%
%\begin{proposition}
%\label{prop:unif_cvx_lin}
%Soit $E$ un Banach uniformément convexe. Alors toute $\vp \in E^*$ atteint son maximum en un unique point de la boule unité fermée $B$, et donc $E$ est réflexif. 
%%L'application (non-linéaire) $J_E : E \to E^*$ associant à un point son conjugué dual est uniquement définie, bijective, de même que $J_{E^*} : E^* \to E^{**}$, et $J_{E^*} \circ J_E = \Phi : E \to E^{**}$ est l'identification naturelle.
%\end{proposition}
%
%\begin{proof}
%	Soit $\vp \in E^*$, avec $\|\vp\|=1$ sans réduire la généralité.
%	Considérons $(x_n)_{n \in \bN} \in B^\bN$ telle que $\vp(x_n) = 1-\eta_n$ avec $\eta_n \to 0$. On peut supposer $\|x_n\|=1$ quitte à considérer $x_n/\|x_n\|$. Pour tous $m,n\in \bN$ on a 
%	\begin{equation}
%		1-\frac{\eta_n+\eta_m} 2 \leq \vp\big(\frac{x_m+x_n} 2\big) \leq \|\frac{x_m+x_n} 2\| \leq 1-\delta(\|x_n-x_m\|),
%	\end{equation}
%	ainsi $\delta(\|x_n-x_m\|) \leq \frac{\eta_n+\eta_m} 2$, en notant $\delta(\ve)>0$ un paramètre satisfaisant \eqref{eq:unif_cvx_ball}, donc $ (x_n)_{n \in \bN}$ est de Cauchy, donc convergente vers $x_*$. On a $x_* \in B$ par fermeture, et $\vp(x_*) = 1$ par continuité. 
%	L'unicité découle de  l'uniforme convexité, et la réflexivité du théorème de James.
%\end{proof}

\begin{remark}[Cas Hilbertien]
\label{rem:unifCvxHilbert}
Un espace de Hilbert $H$ satisfait l'identité du parallélogramme $\|x+y\|^2 + \|x-y\|^2 = 2(\|x\|^2+\|y\|^2)$, qui implique l'uniforme convexité. L'identification linéaire (ou antilinéaire dans le cas complexe) $H \to H^* : x \mapsto \<x,\cdot\>$ correspond au conjugué dual.
\end{remark}

\todoi{L'uniforme convexité de l'espace dual implique l'unicité du conjugué dual.}

\paragraph{Dualité dans les espaces de Lebesgue.}
Deux exposants $p,q \in [1,\infty]$ sont dits conjugués si 
%Commençons par l'inégalité de Holder, qui montre que $L^p$ s'injecte isométriquement dans $(L^q)^*$ si les exposants $p$ et $q$ sont conjugués, c'est-à-dire si 
\begin{equation*}
	\frac 1 p + \frac 1 q = 1.
\end{equation*}
L'objectif de cette fin de section est d'établir l'identification $L^p(X,\mu)^* = L^q(X,\mu)$, lorsque $p$ et $q$ sont conjugués et $p<\infty$, et où $(X,\mu)$ désigne un espace mesuré. 
%. Le crochet de dualité est donné par l'intégrale $\int_X f g\, \diff \mu$,%utilisant l'intégrale comme crochet de dualité,
%\begin{align*}
%	L^p(X,\mu)^* &= L^q(X,\mu), 
%	\<f,g\> := \int f g
%\end{align*}
%
%et qui donne par son cas d'égalité le conjugué dual d'un élément. %On montre ensuite dans le \cref{th:lpDualLq} que cette injection est bijective, pour $p<\infty$. 
%Dans toute la suite, $(X,\mu)$ désigne un espace mesuré. 

\begin{theorem}[Inégalité de Hölder]	
\label{th:HolderIneq}
Soient $p,q\in [1,\infty]$ satisfaisant $\frac 1 p + \frac 1 q = 1$. Pour tous $f\in L^p(X,\mu)$ et $g\in L^q(X,\mu)$ on a $fg \in L^1$ et 
\begin{equation}
\label{eq:Holder}
	\int_X f g \leq \|f\|_p \|g\|_q.
\end{equation}
%\begin{align*}
%	\int_X f g &\leq \|f\|_p \|g\|_q,&
%	\text{si}\quad \tfrac 1 p + \tfrac 1 q &= 1.
%\end{align*}
Supposons $\|f\|_p \|g\|_q \neq 0$. Dans le cas $1< p < \infty$, on a égalité dans \eqref{eq:Holder} ssi $f = \lambda \sign(g) |g|^{q/p}$ $\mu$-presque partout, avec $\lambda > 0$. (Dans le cas $p=1$, on a égalité dans \eqref{eq:Holder} ssi $g = \lambda \sign(f)$ $\mu$-presque partout où $|f|>0$, et $\|g\|_\infty = \lambda$.)
\end{theorem}
\begin{proof}
Supposons $1 < p < \infty$, le cas $p \in \{1,\infty\}$ étant trivial. 
L'inégalité suivante, due à Young, se déduit de la convexité de l'exponentielle : $\forall a,b>0$,
\begin{equation*}
	a b = \exp(\tfrac 1 p \ln(a^p) + \tfrac 1 q \ln(b^q)) \leq \tfrac 1 p \exp(\ln(a^p)) + \tfrac 1 q \exp(\ln(b^q)) = \tfrac 1 p a^p + \tfrac 1 q b^q.
\end{equation*}
Cette inégalité reste vraie si $a,b\geq 0$, et par stricte convexité de l'exponentielle l'égalité a lieu ssi $a^p = b^q$.
Par homogénéité de l'inégalité \eqref{eq:Holder} considérée, on peut supposer $\|f\|_p = \|g\|_q = 1$. Ainsi 
\begin{equation*}
	\int_X |fg| \leq \frac 1 p \int_X |f|^p + \frac 1 q \int_X |g|^q = 1 = \|f\|_p  \|g\|_q,
\end{equation*}
comme annoncé. On en déduit $fg \in L^1$ et l'inégalité demandée. Le cas d'égalité dans \eqref{eq:Holder} se déduit de celui de l'inégalité de Young. 
\end{proof}

L'uniforme convexité des espaces de Lebesgue $L^p(X,\mu)$, $1 < p < \infty$ découle des inégalités de Clarkson.
Rappelons la première d'entre elles\footnote{
Pour être complet, la seconde énonce que 
$
\big\|\frac{f+g} 2\big\|_p^q + \big\|\frac{f-g} 2\big\|_p^q \leq \big(\frac 1 2 \|f\|_p^p + \frac 1 2 \|g\|_p^p\big)^\frac q p
$,
si $1<p \leq 2$ et $\frac 1 p + \frac 1 q =1$.
}, correspondant au cas $p \geq 2$, qui suffit à la preuve du \cref{th:lpDualLq} ci-dessous et qui se montre simplement.

\begin{proposition}[Première inégalité de Clarkson \facultatif]
\label{prop:ClarksonIneq}
Soient $2 \leq p < \infty$ et $f,g\in L^p(X,\mu)$. Alors
	\begin{equation*}
		\Big\|\frac{f+g} 2\Big\|_p^p + \Big\|\frac{f-g} 2\Big\|_p^p \leq \frac 1 2 \|f\|_p^p + \frac 1 2 \|g\|_p^p.
	\end{equation*}
\end{proposition}

\begin{proof}%[Preuve de la première inégalité (la plus facile)]
Étant donné $a,b\geq 0$ et $s\geq 1$, on a $a^s+b^s\leq (a+b)^s$; en effet, on peut supposer $a+b=1$ par homogénéité, et remarquer alors que $a^s \leq a$ et $b^s\leq b$. On en déduit que, ponctuellement
\begin{equation*}
	\big|\frac{f+g} 2\big|^p + \big|\frac{f-g} 2\big|^p \leq \Big(\big|\frac{f+g} 2\big|^2 + \big|\frac{f-g} 2\big|^2\Big)^\frac p 2 = \big(\frac{f^2+g^2} 2\big)^\frac p 2 \leq \frac 1 2 |f|^p + \frac 1 2 |g|^p.
\end{equation*}
	On a utilisé $s=p/2$ pour la première inégalité, et la convexité de $x\in \bR_+ \mapsto x^\frac p 2$ pour la seconde. 
	Par intégration sur le domaine $X$, on obtient le résultat annoncé.
\end{proof}

%Il existe une seconde inégalité de Clarkson, s'appliquant aux exposants $1 < p \leq 2$, et un peu plus difficile à établir. Elle énonce que 
%	\begin{align*}
%		\big\|\frac{f+g} 2\big\|_p^q + \big\|\frac{f-g} 2\big\|_p^q &\leq \big(\frac 1 2 \|f\|_p^p + \frac 1 2 \|g\|_p^p\big)^\frac q p, 
%		& & \text{si } 1<p \leq 2 \text{ et } \frac 1 p + \frac 1 q =1.
%	\end{align*}

%Les inégalités de Clarkson impliquent que $L^p$ est uniformément convexe pour tout $1 < p < \infty$. 
%En effet si $\|f\|_p=\|g\|_p=1$ et $\|f-g\|_p=\delta$, alors $\|\frac {f+g} 2\|_p \leq (1-(\delta/2)^p)^\frac 1 p$ si $p\geq 2$ en utilisant la première inégalité, et par ailleurs $\|\frac {f+g} 2\|_p \leq (1-(\delta/2)^q)^\frac 1 q$ si $p<2$ en utilisant la seconde. 



%Noter que la preuve de $(L^p(X,\mu) )^* = L^q(X,\mu)$, pour tout $1 < p < \infty$, peut s'établir en utilisant seulement la première inégalité de Clarkson. 

\begin{theorem}[Dualité dans les espaces de Lebesgue]
\label{th:lpDualLq}
Soit $1 \leq p \leq \infty$, et $q$ l'exposant conjugué, i.e. $\frac 1 p + \frac 1 q =1$. Posons pour tous $f\in L^p(X,\mu)$ et $g \in L^q(X,\mu)$
\begin{equation*}
	\vp_g(f) := \int_X fg.
\end{equation*}
Alors $g\in L^q \mapsto \vp_g \in (L^p)^*$ est une injection isométrique, bijective si $p<\infty$. 
\end{theorem}

\begin{proof}[Preuve, seulement dans le cas $1 < p < \infty$ pour la bijectivité]
D'après l'inégalité de Hölder, 
la for\-me linéaire $\vp_g \in (L^p)^*$ est continue et de norme au plus $\|g\|_q$. Par le cas d'égalité, on a l'injection et l'isométrie. Il reste à établir la surjectivité.

Supposons d'abord que $2 \leq p < \infty$. Notons $E := (L^p)^*$ et $F = \{\vp_g \mid g \in L^q\} \subset E$, qui est un sous espace complet car isométrique à $L^q$, donc fermé. 
Supposons par l'absurde que $F \neq E$, de sorte qu'il existe $\vp \in E^*\sm \{0\}$ telle que $F \subset \Ker(\vp)$, par le \nref{corol:dens_banach}.

Par la \nref{prop:ClarksonIneq}, l'espace $L^p$ est uniformément convexe donc réflexif. L'élément du bidual $\vp \in E^* = (L^p)^{**}$ s'écrit donc $\vp = \Phi(f)$, c'est-à-dire $\vp(\psi) =  \psi(f)$ pour toute $\psi \in E$, pour un certain $f \in L^p$, non-nul car $\vp$ l'est. Par le cas d'égalité dans l'inégalité de Hölder, il existe $g \in L^q$ telle que $\Phi(f)(g) := \int_X f g \neq 0$, ce qui contredit l'hypothèse que $\vp$ s'annule sur $F$. 

Nous avons établi que $F=E$, c'est-à-dire l'identification isométrique bijective $L^p(X,\mu)^* = L^q (X,\mu)$. Par ailleurs $L^q(X,\mu)^* = L^p(X,\mu)^{**} = L^p(X,\mu)$ par réflexivité de $L^p$, ce qui conclut.
\end{proof}

%Les normes $\|x\|_1 := \sum_{i=1}^d |x_i|$ et $\|x\|_\infty := \max_{i=1}^d |x_i|$ sont bien duales sur l'espace de dimension finie $\bR^d$. En revanche
%L'espace $\bR^d$ de dimension finie $d$, 

Pour terminer, nous montrons que $\ell^1\subsetneq (\ell^\infty)^*$, où $\ell^p := L^p(\bN)$ muni de la mesure de comptage, c'est-à-dire la non-bijectivité de l'inclusion du \cref{th:lpDualLq} dans le cas $p=\infty$, et la non-réflexivité de $\ell^1$. En revanche, les evn de dimension finie $(\bR^d,\|\cdot\|_1)$ et $(\bR^d,\|\cdot\|_\infty)$, correspondant à un domaine $X$ fini, sont bien duaux isométriques l'un de l'autre. 

\begin{exemple}[Non réflexivité]	
\label{ex:nonReflex}
On a $(\ell^\infty_0)^* = \ell^1$ et $(\ell^1)^* = \ell^\infty$ (identification via des isométries bijectives), où $\ell^\infty_0$ désigne l'espace des suites convergeant vers $0$, le crochet de dualité étant 
\begin{equation*}
	\<(x_n)_{n \in \bN}, (y_n)_{n\in \bN}\> := \sum_{n \in \bN} x_n y_n.
\end{equation*}
De plus il existe $\vp \in (\ell^\infty)^*$ telle que $\vp( (x_n)_{n \in \bN} ) = \lim_{n \to \infty} x_n$ pour toute suite $(x_n)_{n\in \bN}$ convergente, et $\|\vp\|_{(\ell^\infty)^*}=1$. En particulier $(\ell^\infty)^* \neq \ell^1$. 
\end{exemple}

\begin{proof}
Considérons $\vp \in (\ell_0^\infty)^*$, posons $x_n := \vp(\delta_n)$ où $\delta$ est le symbole de Kronecker, et notons $\sigma_n$ le signe de $x_n$.
La continuité de $\vp\in (\ell^\infty_0)^*$ donne $\sum_{0 \leq n \leq N} |x_n| = \vp(\sigma_0 \delta_0+\cdots+ \sigma_N\delta_N) \leq \|\vp\|_{(\ell^\infty_0)^*}$ pour tout $N\in \bN$, puisque $\|\sigma_0 \delta_0+\cdots+ \sigma_N\delta_N\|_\infty = 1$, donc $(x_n)_{n\in \bN} \in \ell^1$. 
Ainsi  $\tilde \vp(y) := \sum_{n\in \bN} x_n y_n$ est continue sur $\ell_0^\infty$, et donc $\tilde \vp = \vp$ par densité des suites à support fini. On a montré que $(\ell_0^\infty)^* = \ell^1$, comme annoncé.

Supposons maintenant $\vp\in (\ell^1)^*$, définissons de même $x_n := \vp(\delta_n)$ et $\sigma_n$ son signe, et notons que $|x_n| = \vp(\sigma_n \delta_n) \leq \|\vp\|_{(\ell^1)^*}$ puisque $\|\delta_n\|_1 = 1$. Ainsi la suite $(x_n)_{n \in \bN} \in \ell^\infty$, et on conclut comme précédemment que $(\ell^1)^* = \ell^\infty$, comme annoncé. 

Soit $F \subset \ell^\infty$ le sous espace vectoriel des suites convergentes. On définit $\vp_F((x_n)_{n \in \bN}) := \lim_{n \to \infty} x_n$ sur cet espace, qui est une forme linéaire continue de norme $1$. Puis par le \cref{corol:prol_forme_evn} on étend $\vp_F$ à $\ell^\infty$ en conservant sa norme, comme annoncé.
Si par l'absurde on pouvait écrire $\vp( (x_n)_{n\in \bN}) = \sum_{n\in \bN} x_n y_n$, pour une certaine suite $(y_n)_{n \in \bN} \in \ell^1$, alors on obtiendrait $\sum_{n \geq N} y_n = 1$ pour tout $N\in \bN$ (en considérant la suite constante égale à $1$ à partir du rang $N$, qui est bien convergente), ce qui est impossible. Ainsi $(\ell^\infty)^* \neq \ell^1$, comme annoncé, ce qui conclut.
\end{proof}


\subsection{Formes géométriques de Hahn-Banach}
\label{subsec:HahnBanachGeom}

Le théorème de Hahn-Banach possède des conséquences géométriques, affirmant l'existence d'hyperplans séparant des ensembles convexes. Elles complètent l'interprétation analytique, de prolongement continu de formes linéaires, et ouvrent de nouvelles applications.

\begin{retenir}
Le théorème de séparation d'un ouvert convexe et d'un point, et de séparation forte d'un fermé et d'un compact convexes.
(\Cref{th:separOuvertPt,th:separFermeComp}.)
\end{retenir}

Commençons par décrire la jauge d'un convexe, qui jouera le rôle de la fonction sous-linéaire dans le \nref{th:HahnBanach}. Noter que la jauge de la boule unité d'un evn est sa norme. 

Pour la lisibilité, et en l'absence d'ambiguïté, on se permet dans la suite de noter 
$\{f < 1\} := \{x \in E \mid f(x) < 1\} = f^{-1}(]-\infty,1[)$, et de même pour $\{f\leq 1\}$, pour toute $f : E \to \bR$.

\begin{proposition}[Jauge d'un convexe]
\label{prop:jaugeCvx}
Soit $C$ un convexe contenant $0$ d'un $\bR$-ev $E$. Définissons 
\begin{align*}
	p_C &: E \to [0,\infty], & p_C(x) &:= \inf\{t > 0 \mid x/t \in C\},
\end{align*} 
avec $\inf(\emptyset) := \infty$ par convention.
Alors $p_C(\lambda x) = \lambda p_C(x)$, et $p_C(x+y) \leq p_C(x)+p_C(y)$, pour tous $\lambda \geq 0$, $x,y \in E$ (avec la convention $0 \times \infty = 0$). 
Si $p_C$ prend des valeurs finies, alors c'est une fonction sous-linéaire; si de plus elle est symétrique, alors c'est une semi-norme. 
%On a $\{p_C < 1\} \subset C \subset \{p_C \leq 1\}$.
On a 
\begin{equation*}
%\{x \in E \mid p_C(x) < 1\} \subset C \subset \{x \in E \mid p_C(x) \leq 1\}.
\{p_C < 1\} \subset C \subset \{p_C \leq 1\}.	
\end{equation*}
\end{proposition}

\begin{proof}
	\emph{Preuve de l'homogénéité positive de $p_C$}~:
	si $x/t \in C$, alors $(\lambda x)/(\lambda t) = x/t \in C$ pour tout $\lambda>0$, donc $p_C(\lambda x) = \lambda p_C(x)$; de plus $(0\times x) / t = 0 \in C$ pour tout $t>0$, par hypothèse, donc $p_C(0\times x) = 0$ comme annoncé. 
	\emph{Preuve de la sous-additivité}~: 
	si $x/s \in C$ et $y/t \in C$, alors par convexité $(x+y)/(s+t) = (1-\theta)(x/s) + \theta (y/t) \in C$ avec $\theta := t/(s+t)$. 
	Les autres propriétés sont immédiates. 
\end{proof}


\begin{lemma}
\label{lem:jaugeOuvertCvx}
	Soit $C$ un convexe contenant $0$ d'un $\bR$-evn. Si $C$ est ouvert, alors $p_C : E \to \bR$ est Lipschitz et $C = \{p_C < 1\}$. Si $C$ est fermé, alors $p_C : E \to [0,\infty]$ est semi-continue inférieurement, et $C = \{p_C \leq 1\}$.
\end{lemma}

\begin{proof}
	Supposons $C$ ouvert, et donc $B(0,r) \subset C$ pour un certain $r>0$. Alors $p_C(x) \leq \|x\|/r$ et donc $p_C(x) - \|h\|/r \leq p_C(x) - p_C(-h) \leq p_C(x+h) \leq p_C(x) + p_C(h) \leq p_C(x) + \|h\|/r$ d'où la $r^{-1}$-Lipschitzianité. Enfin $\{p_C < 1\}$ est l'intérieur de $\{p_C\leq 1\}$, donc est égal à l'ouvert $C$. 
	
	Supposons $C$ fermé, et notons que $x/p_C(x) \in C$ par fermeture dès que $p_C(x) >0$, soit encore $C = \{p_C \leq 1\}$ par homogénéité comme annoncé. 
	Les ensembles $\{p_C \leq \lambda\}$ pour $\lambda \in ]0,\infty[$ sont homothétiques à $C$ donc fermés, et le cône asymptotique $\{p_C \leq 0\} = \cap_{t>0} t C$ est fermé comme intersection de fermés, d'où la semi-continuité inférieure. 
\end{proof}

Nous déduisons du \nref{th:HahnBanach} deux résultats de séparation géométrique de convexes par un hyperplan. Dans le premier, le \cref{th:separOuvertPt} et son corollaire, l'hyperplan peut être adhérent aux ensembles qui sont séparés. Une séparation dite stricte est en revanche établie dans le \cref{th:separFermeComp}.

\begin{theorem}[Séparation d'un ouvert convexe et d'un point]
\label{th:separOuvertPt}
	Soit $E$ un evn, soit $C \subset E$ un ouvert convexe contenant $0$, et soit $x \in E\sm C$. Alors il existe $\vp \in E^*$ telle que $\vp(x)=1$ et $\vp<1$ sur $C$. 
	\end{theorem}

\begin{proof}
Définissons $F := \bR x$, et $\vp_F : F \to \bR$ par $\vp_F(\lambda x) := \lambda$.
Cette forme linéaire est bornée sur $F$ par la jauge $p_C$ :  en effet, $p_C(x) \geq 1$ puisque $x \notin C$, et donc $p_C(\lambda x) \geq \lambda$ pour tout $\lambda \geq 0$; par ailleurs $p_C(-\lambda x) \geq 0 \geq -\lambda$ pour tout $\lambda \geq 0$. 
	D'après le \nref{th:HahnBanach}, $\vp_F$ admet une extension à $E$ notée $\vp$, satisfaisant toujours $\vp \leq p_C$. 
	
	Reste à montrer la continuité de $\vp$. 
	Par hypothèse, $C$ est ouvert, donc il existe $r>0$ tel que $B(0,r) \subset C$.
	D'après le \cref{lem:jaugeOuvertCvx} on a $\vp \leq p_C < 1$ sur $C$, donc aussi sur $B(0,r)$, ce qui montre que $\|\vp\|_{E^*} \leq r^{-1}$ et conclut. 
\end{proof}

La preuve des résultats suivants fait intervenir la \emph{somme de Minkowski} de deux parties $A$ et $B$ d'un espace vectoriel, définie par $A+B := \{a+b \mid a\in A,\, b\in B\}$. 

\begin{corollary}
\label{corol:separOuvertPt}
	Soit $C$ un ouvert convexe d'un evn $E$ ne contenant pas $0$, alors il existe $\vp \in E^*$ telle que $\vp>0$ sur $C$.
	Soient $A$ un convexe et $B$ un ouvert convexe d'un evn $E$, non-vides et disjoints, alors il existe $\vp \in E^*$ et $\alpha \in \bR$ tels que $\vp(a) \leq \alpha < \vp (b)$ pour tous $a \in A$ et $b \in B$. 
\end{corollary}

\begin{proof} 
	On peut supposer $C \neq \emptyset$.
	\emph{Premier point.} Choisissons $x_0 \in C$, et notons que $0 \in C-x_0$, et que $-x_0 \notin C-x_0$.
	Par le \cref{th:separOuvertPt} appliqué à $C-x_0$ et $-x_0$, il existe $\vp \in E^*$ telle que $\vp(-x_0)=1$ et $\vp<1$ sur $C-x_0$; ainsi $-\vp>0$ sur $C$, comme annoncé. 
	
	\emph{Second point.} Posons $C:=B-A = \{b-a\mid a \in A, b \in B\}$ qui est un ouvert convexe ne contenant pas $0$. Soit $\vp\in E^*$ telle que $\vp>0$ sur $C$, donc $\vp(a) \leq \vp(b)$ pour tous $(a,b)\in A \times B$. On pose $\alpha = \inf_B \vp$, et on note que $\vp(B)$ est ouvert car $\vp\neq 0$ et $B$ est ouvert, donc $\vp(A) \subset ]-\infty,\alpha]$ et $\vp(B)\subset ]\alpha,\infty[$ comme annoncé.
\end{proof}


\begin{theorem}[Séparation stricte d'un fermé et d'un compact convexes]
\label{th:separFermeComp}
	Soit $E$ un evn, et soient $A,B \subset E$ convexes, non-vides, disjoints, avec $A$ fermé et $B$ compact. Alors il existe $\vp \in E^*$ telle que $\sup_A \vp < \inf_B \vp$. 
	On dit que l'hyperplan affine $\{x \in E \mid \vp(x) = \tfrac 1 2 (\sup_A \vp + \inf_B \vp)\}$ sépare $A$ et $B$ au sens strict. Si $0 \in A$ alors on peut choisir $\vp$ telle que $\sup_A \vp < 1 < \inf_B \vp$. 
\end{theorem}

\begin{proof}
	Posons $F := B-A = \{b-a \mid a\in A, b \in B\}$ qui est convexe et non-vide. \emph{Montrons que $F$ est fermé} : considérons une suite convergente $b_n-a_n \to c_*$ lorsque $n \to \infty$, avec $a_n \in A$ et $b_n\in B$. Extrayons une sous-suite convergente $b_{\sigma(n)} \to b_*\in B$ par compacité. Par différence $a_{\sigma(n)} \to a_* := b_* - c_*$  lorsque $n \to \infty$, et par fermeture $a_* \in A$, donc $c_* = b_*-a_* \in B-A$ comme souhaité. 
	
	Notons que $0 \notin F$, et donc par fermeture il existe $r>0$ tel que $B(0,r) \cap F =\emptyset$. Le convexe ouvert $F-B(0,r)$ ne contient pas $0$, et donc par le \cref{corol:separOuvertPt} il existe $\vp \in E^*$ telle que $\vp>0$ sur $F-B(0,r)$. On en déduit que $\vp(b)-\vp(a) \geq r\|\vp\|_{E^*} > 0$ pour tous $a\in A$ et $b \in B$, comme annoncé.
	Pour le dernier point, considérer la forme linéaire $\psi := 2\vp/(\sup_A \vp+\inf_B\vp)$, le dénominateur étant strictement positif puisque $0 \leq \sup_A \vp < \inf_B \vp$.
\end{proof}

%\subsection{Compacité faible et optimisation} % des convexes fermés bornés} % et optimisation} %Dualité polaire des ensembles convexes} 
%\label{subsec:CompaciteConvexes}

%Nous montrons également comment le \cref{th:HahnBanach} permet de décrire la géométrie des convexes, en les exprimant comme intersections de demi-espaces, ou en montrant l'existence de plans tangents, sous des hypothèses adéquates. 

%Le théorème de Hahn-Banach permet de décrire des ensembles convexes comme intersections de demi-espaces.


La dualité polaire consiste à exprimer un ensemble convexe fermé comme une intersection de demi-espaces, en utilisant les théorèmes de séparation précédents.

\begin{definition}[Ensemble polaire \facultatif]
\label{def:convexePolaire}
L'ensemble polaire d'une partie $C$ d'un evn $E$ est défini par 
%	Soit $C \subset E$ une partie non-vide d'un evn. On définit l'ensemble polaire par 
	\begin{equation*}
		C^\circ := \{\vp \in E^* \mid \forall x \in C,\, \vp(x) \leq 1\}.
	\end{equation*}
	On note $\kC(E)$ l'ensemble des \emph{convexes fermés contenant $0$} d'un evn $E$ .
\end{definition}
Observons que $C^\circ \in \kC(E^*)$, et que $C^\circ$ est *faiblement fermé.
Si $C$ est un cône, c'est-à-dire s'il est stable par multiplication par des scalaires positifs, alors c'est aussi le cas de l'ensemble polaire, que l'on peut écrire sous la forme $C^\circ = \{\vp \in E^*\mid \forall x \in C, \vp(x) \leq 0\}$.
Par construction $B'_E(0,1)^\circ = B'_{E^*}(0,1)$. 

\begin{proposition}[Dualité polaire \facultatif]
\label{prop:cvxFaibleFerme}
Pour tout evn $E$, et tout $C \in \kC(E)$, on a 
\begin{equation}
\label{eq:cvxFaibleFerme}
	C = \{x \in E \mid \forall \vp \in C^\circ, \vp(x) \leq 1\} 
\end{equation}
\end{proposition}

\begin{proof}
Notons $\tilde C$ le membre de droite, et remarquons que l'inclusion $C \subset \tilde C$ est évidente par définition. Si $x\notin C$, alors le \nref{th:separFermeComp} donne $\vp \in C^\circ$ les séparant strictement, d'où l'inclusion réciproque.
\end{proof}


Si $E$ est un espace réflexif, alors on se permet de l'identifier à $E^{**}$, et la \cref{prop:cvxFaibleFerme} se réécrit donc : pour tout $C \in \kC(E)$
\begin{equation*}
	C = C^{\circ\circ}.
\end{equation*}
Ainsi la dualité polaire $C \in \kC(E) \mapsto C^\circ \in \kC(E^*)$ définit une bijection de réciproque $C \in \kC(E^*) \mapsto C^\circ \in \kC(E^{**}) = \kC(E)$.


\begin{corollary} %\facultatif
\label{corol:cvxFermeBdComp}
Tout convexe fermé $C$ d'un evn $E$ est faiblement fermé.
Si de plus $C$ est borné et $E$ est réflexif, alors $C$ est faiblement compact.
%	Tout convexe fermé borné $C$ d'un espace réflexif $E$ est faiblement compact.
\end{corollary}

\begin{proof}
\emph{Premier point}. Si $C=\emptyset$ alors il est clairement fermé, et sinon on peut supposer $0 \in C$ quitte à faire une translation. Alors \eqref{eq:cvxFaibleFerme} montre que $C$ est une intersection de demi-espaces fermés, qui sont faiblement fermés.


\emph{Second point.} Le convexe $C$ est fermé, et contenu dans une boule fermée qui est compacte par le \nref{th:BanachAlaoglu}, puisque la topologie faible sur $E$ correspond à la topologie *faible sur $E^{**}$.
\end{proof}

On déduit du \cref{corol:cvxFermeBdComp} qu'une fonction $f$ convexe et s.c.i.\ (semi continue inférieurement, voir la \cref{def:semiContinuite}) est également s.c.i.\ pour la topologie faible, puisque les ensembles de sous-niveau $\{f \leq \lambda\}$ sont convexes et fermés, donc faiblement fermés. 
Cette observation, et la propriété de compacité faible établie dans le \cref{corol:cvxFermeBdComp}, permettent dans le résultat suivant de montrer l'existence de minimiseurs à des problèmes d'optimisation convexes, et à certaines de leurs perturbations. 

L'\nref{ex:nonLinMinH1} en est une application représentative. Dans celui-ci, la fonction $g$ est la composition d'une application linéaire compacte et d'une fonction non-linéaire continue. 
Par la \nref{prop:opCompactComplCont}, et comme l'espace de départ $H^1$ est un Hilbert séparable, $g$ est faiblement continue sur les bornés comme nécessaire ici.

\begin{corollary}[Minimisation par convexité et compacité \facultatif]
\label{corol:minFctCvxRefl}
	Soit $f : E \to ]-\infty,\infty]$ convexe s.c.i. sur un evn $E$ réflexif, soit $g : E \to \bR$ faiblement continue sur les bornés, et soit $M \in \bR$. %soient $M,N,R \in \bR$.
	\begin{itemize}
		\item Si $\{x\in E\mid f(x) \leq M\}$ est non-vide et borné, alors $f$ admet un minimiseur. 
		\item Si $\{x \in E\mid f(x)+g(x)\leq M\}$ est non-vide et borné, alors $f+g$ admet un minimiseur. 
		\end{itemize}	
\end{corollary}

\begin{proof}
\emph{Premier point.} Comme $f$ est convexe et s.c.i., l'ensemble $E_M := \{x\in E \mid f(x) \leq M\}$ est convexe et fermé, donc faiblement fermé. Ainsi $f$ est s.c.i.\ pour la topologie faible. 
Ainsi $(E_m)_{\inf(f)<m\leq M}$ est une famille emboîtée de fermés non-vides, contenus par hypothèse dans une boule $B'(0,R)$ qui est compacte par le \nref{th:BanachAlaoglu}. L'intersection $\cap_{\inf(f) < m\leq M} E_m = \{x \in E \mid f(x) = \inf(f)\}$ est donc non-vide par la \nref{rem:fermeBorelLeb}, comme annoncé.

\emph{Second point.} 
Soit $B$ une boule fermée contenant l'ensemble d'intérêt $\{x \in E\mid f(x)+g(x)\leq M\}$ dans son intérieur.
La fonction $g$ est continue sur $B$ pour la topologie faible, donc s.c.i.\ pour cette topologie, tout comme $f$. Ainsi $f+g$ est faiblement s.c.i.\ sur $B$ comme somme de fonctions faiblement s.c.i. \eqref{eq:somme_sci}. Les ensembles $E'_m := \{x \in E\mid f(x)+g(x)\leq m\}$ pour $\inf(f+g) < m\leq M$ sont donc fermés pour la topologie faible, inclus dans un compact, et de plus non-vides et emboîtés. Comme précédemment, l'ensemble des minimiseurs $\cap_{\inf(f+g) < m\leq M} E'_m = \{x \in E \mid f(x)+g(x) = \inf(f+g)\}$ est non-vide, ce qui conclut. 
%	On note que $\{f\leq M\}$ est convexe fermé borné, donc faiblement compact par le \cref{corol:cvxFermeBdComp}. 
%	Ainsi les ensembles $\{f \leq M'\}$ pour $\inf(f)<M'\leq M$ sont fermés dans $\{f\leq M\}$ donc faiblement compacts, et de plus non non-vides et emboités, donc d'intersection non-vide par la \cref{rem:fermeBorelLeb}, comme annoncé. 
%	
%	Dans le second cas, $f+g$ est faiblement s.c.i.\ comme somme de fonctions faiblement s.c.i. \eqref{eq:somme_sci}, donc les ensembles $\{f+g\leq N'\}$ pour $\inf(f+g) < N'\leq N$ sont faiblement fermés dans $\{f\leq M\}$, donc compacts, et de nouveau non-vides et emboités, donc d'intersection non-vide.
\end{proof}



\newpage
\section{Espaces de Sobolev}
\label{sec:Sobolev}

Les espaces de Sobolev, qui expriment que les dérivées d'une fonction jusqu'à un ordre $n$ donné sont intégrables à une puissance $p$ donnée, constituent un cadre naturel pour l'étude des problèmes variationnels et des EDPs, par exemple. Cette section décrit les injections entre ces espaces, et notamment le théorème d'injection compacte de Rellich, ainsi que les techniques liées que sont la régularisation par convolution et le théorème de compacité de Fréchet-Kolmogorov.\\

Dans toute cette partie, $\Omega$ désigne un ouvert de $\bR^d$. %, sauf mention contraire.
En préliminaire de cette section, nous énonçons deux résultats de densité, l'un dans $L^p$ et l'autre dans l'espace des fonctions test $\cD(\Omega)$.

\inventaire{
\Nref{prop:densC0Lp}\\
\Nref{lem:intDerivIter} (Intégration et dérivation itérée)\\
\Nref{prop:testFunSeparab} (Séparabilité de $\cD(\Omega)$)\\

\Nref{subsec:convLp} \\
\Nref{prop:Jensen}\\
\Nref{corol:lpInclusLq} (Injection $L^p$)\\
\Nref{prop:convolLp}\\
\Nref{lem:convolTranslat} (Convolution et translation)\\
\Nref{prop:convolApprox}\\
\Nref{prop:convolDeriv}\\
\Nref{corol:testDenseLp} (Densité de $\cD(\Omega)$ dans $L^p$)\\
\Nref{corol:convolTestFun} (La mollification approche l'identité dans $\cD(\Omega)$)\\
\Nref{th:FrechetKolmogorov}\\
\Nref{ex:operCompactL2}\\


\Nref{subsec:convolSobolev}\\
\Nref{prop:distribL1Loc}\\
\Nref{corol:distribLp}\\
\Nref{def:Sobolev}\\
\Nref{prop:SobolevBanach} ($W^{s,p}$ est un Banach)\\
\Nref{def:convolDistrib}\\
\Nref{prop:convolDistrib} (Propriétés de la convolution des distributions)\\
\Nref{prop:testDenseWsp} (Densité de $\cD(\bR^d)$ dans $W^{s,p}$)\\
\Nref{def:ouvertReg}\\
\Nref{th:prolSobolev}\\
\Nref{corol:testDenseWspO} (Densité des restrictions de $\cD(\bR^d)$ dans $W^{s,p}(\Omega)$)\\
\Nref{rem:bcDirichlet}\\
\Nref{th:pbPoisson}\\

\Nref{subsec:injSousCrit}\\
\Nref{th:YoungFaible}\\
\Nref{prop:convolPow} ($\diver V = \delta_0$)\\
\Nref{rem:homogSobolev}\\
\Cref{th:injSousCrit} (Injection sous-critique $W^{s,p} \subset L^{p^*}$, $sp<d$)\\
\Nref{corol:injSousCritOm} (Injection sous-critique, domaine régulier)\\
\Nref{th:Young}\\
\Nref{corol:noyauGreen}\\

\Nref{subsec:injSurCrit}\\
\Nref{prop:convolPow2} (Lemme type convolution, sur boule bornée)\\
\Nref{th:injSurCrit}\\
\Nref{corol:ingSurCritOm} (Injection sur-critique, dans un domaine)\\

\Nref{subsec:injCompacte}\\
\Nref{prop:interpLp}\\
\Nref{th:Rellich}\\
\Nref{corol:Rellich} (Rellich, séquentiel)\\
\Nref{ex:nonLinMinH1}\\
\Nref{rem:bouleW1pFerme}
}


\begin{proposition}[Densité dans $L^p$ des fonctions continues à support compact]	
\label{prop:densC0Lp}
Soit $(X,d)$ un espace métrique localement compact muni d'une mesure borélienne régulière, c'est-à-dire que pour toute partie mesurable $A \subset X$
\begin{align}
\label{eqdef:mes_reguliere}
	\mu(A) &= \sup\{\mu(K) \mid K \csubset A, K \text{ compact}\}, &
	\mu(A) &= \inf\{\mu(U) \mid U \supset A, U \text{ ouvert}\}.
\end{align}
Alors $C^0_c(X)$, l'ensemble des fonctions continues à support compact, est dense dans $L^p(X,\mu)$ pour tout $1 \leq p < \infty$; 
en particulier, $C^0_c(\Omega)$ est dense dans $L^p(\Omega)$, pour tout ouvert $\Omega \subset \bR^d$.
\end{proposition}


\begin{proof}[Preuve \facultatif] 
	La preuve est fondée sur des arguments de théorie de la mesure, et n'est donc pas exigible pour ce cours.
	On note $\bbone_A : X\to \{0,1\}$ la fonction caractéristique d'un ensemble $A\subset X$. 
	L'espace des fonctions étagées est noté $E := \Vect\{\bbone_A \mid A \subset X \text{ mesurable},\, \mu(A)<\infty \}$, et on rappelle que $L^p(X,\mu)$ s'identifie à son  séparé complété pour la norme $\|\cdot\|_p$. % séparé car on identifie les fonctions égales presque partout
	
	Il suffit donc d'approcher les fonctions étagées.
	Soit $A \subset X$ mesurable avec $\mu(A) < \infty$, et soit $\ve>0$. Par régularité de la mesure il existe $K \subset A$ compact et $U \supset A$ ouvert tels que $\mu(U\sm K) \leq \ve$. 
	Par compacité locale de $X$, et compacité de $K$, il existe une couverture finie de $K$ par des boules $K \subset \cup_{i=1}^I B(x_i,r_i)$, telles que $B'(x_i,2r_i)\csubset U$ pour tout $1 \leq i \leq I$. Posons $V := \cup_{i=1}^I B(x_i,2r_i)$, de sorte que $K \csubset V$ et $\overline V \csubset U$.  
	
%	Comme $X$ est localement compact, et $K$ est compact, il existe $U'$ ouvert d'adhérence compacte tel que $K \csubset U' \subset \overline {U'} \csubset U$.
	
	Posons $f(x) := d(x,X\sm V)/(d(x,K) + d(x,X\sm V))$, qui est continue et à support $\supp(f) \subset \overline V$ compact. Alors $f-\bbone_A : X \to [-1,1]$ est nulle hors de $V\sm K$, donc $\|f-\bbone_A\|_p^p \leq \mu(U\sm K) \leq \ve$, comme souhaité.
	La dernière affirmation découle de la régularité de la mesure de Lebesgue sur un ouvert de $\bR^d$, ce qui conclut.
\end{proof}



Nous établissons dans la suite que l'espace $\cD(\Omega)$ des fonctions $C^\infty$ à support compact est dense dans de nombreux espaces fonctionnels, voir par exemple la \cref{prop:testDenseWsp}.
La proposition suivante implique que ces espaces sont tous séparables. 
On note $\alpha \preceq \beta$ l'ordre partiel coordonnée par coordonnée sur $\bN^d$. Le lemme suivant, présenté sans preuve, rappelle que l'intégration et la dérivation sont des opérations réciproques.

\begin{lemma}
\label{lem:intDerivIter}
	Posons pour $1 \leq i \leq d$, $\alpha = (\alpha_1,\cdots,\alpha_d) \in \bN^d$, et toute $f \in C^\infty([0,1]^d)$ % élément de l'ensemble $E := \{f \in C^\infty([0,1]^d) \mid \supp(f) \subset ]0,1]^d\}$, 
%	$f \in C^0([0,1]^d)$ nulle sur $\{(x_1,\cdots,x_d) \in [0,1]^d \mid x_1 \cdots x_d = 0\}$, 
	\begin{align*}
		I_i f(x_1,\cdots,x_d) &:= \int_0^{x_i} f(x_1,\cdots,s,\cdots, x_d) \diff s, &
		I_\alpha f := (I_1)^{\alpha_1} \cdots (I_d)^{\alpha_d} f.
	\end{align*}
	Alors $\|I_\alpha f\|_\infty \leq \|f\|_\infty$, $\partial_\alpha I_\beta f = \partial_{\alpha-\beta} f$ si $\alpha \succeq \beta$, et $\partial_\beta I_\alpha f = I_{\alpha-\beta} f$ si $\beta\preceq\alpha$. 
	De plus, $f = I_\alpha \partial_\alpha f$ pour toute $f \in C^\infty_c(]0,1]^d)$, $\alpha \in \bN^d$.
\end{lemma}


\begin{proposition}[Séparabilité de $\cD(\Omega)$]
\label{prop:testFunSeparab}	
	L'espace des fonctions test $\cD(\Omega)$ est séparable.
\end{proposition}

\begin{proof}
Il suffit de traiter le cas où $\Omega$ est borné, car les fonctions test sont à support compact; quitte à faire une homothétie, on peut donc supposer que $\Omega \subset ]0,1[^d$.
Soit $K \csubset \Omega$, et soit $\rho_K \in \cD(\Omega)$ égale à $1$ sur $K$, voir la \cref{prop:cinfSupport}.
Soit $\vp \in \cD_K(\Omega)$ et $\alpha \in \bN^d$. Soit $(P_n)_{n \in \bN}$ une suite de polynômes à coefficients rationnels convergeant uniformément vers $\partial_\alpha \vp$ sur $[0,1]^d$. Une telle suite existe d'après la \cref{rem:polDens}, issue du \nref{th:StoneWeierstrass}. On a alors, pour tout $\beta \preceq \alpha$, %en utilisant le \cref{lem:intDerivIter},
\begin{equation*}
	\partial_\beta (\vp - \rho_K I_\alpha P_n) 
	= \partial_\beta(\rho_K (\vp-I_\alpha P_n)) 
	= \partial_\beta (\rho_K I_\alpha (\partial_\alpha \vp - P_n)) 
	= \sum_{\gamma \preceq \beta} \binom \beta \gamma \partial_\gamma \rho_K I_{\alpha-\beta+\gamma} (\partial_\alpha \vp-P_n),
\end{equation*}
où l'on a utilisé successivement (i) $\vp = \rho_K \vp$, (ii) $I_\alpha \partial_\alpha \vp = \vp$ par le \cref{lem:intDerivIter}, et (iii) la règle de Leibniz de dérivation d'un produit, et $\partial_{\beta-\gamma} I_\alpha = I_{\alpha-\beta+\gamma}$ d'après le \cref{lem:intDerivIter} encore.
Ainsi, pour tout $\beta \preceq \alpha$
\begin{align*}
	\|\partial_\beta (\vp - \rho_K I_\alpha P_n)\|_\infty &\leq C(K,\beta) \|\partial_\alpha \vp-P_n\|_\infty, &
	C(K,\beta) := \sum_{\gamma \preceq \beta} \binom \beta \gamma \|\partial_\gamma \rho_K\|_\infty.
\end{align*}
Comme $\|\partial_\alpha \vp -  P_n\|_\infty \to 0$  lorsque $n \to \infty$, on obtient $\rho_K I_\alpha P_n \to \vp$ dans la norme $|f|_\alpha := \max_{\beta\preceq \alpha} \|\partial_\beta f\|_\infty$. 
En notant que $I_\alpha P$ est un polynôme, et en introduisant une suite exhaustive $(K_r)_{r \in \bN}$ de compacts de $\Omega$, voir la \cref{prop:exhaustComp}, on conclut que $\{\rho_{K_r} Q \mid r \in \bN, Q \in \bQ[X_1,\cdots,X_d]\}$ est une partie dénombrable dense de $\cD(\Omega)$.
\end{proof}

\subsection{Convolution dans les espaces $L^p$} %Densité et compacité d
\label{subsec:convLp}

On décrit dans cette sous-section le procédé de régularisation par convolution dans les espaces de Lebesgue $L^p(\bR^d)$, et ses premières conséquences.


\begin{retenir}
	L'inégalité de Jensen, la continuité de la convolution, le principe d'approximation et de régularisation par convolution, le théorème de compacité de Fréchet-Kolmogorov.
	(\Cref{prop:Jensen,prop:convolLp,prop:convolApprox,prop:convolDeriv,th:FrechetKolmogorov}.)
\end{retenir}

\notei{Suppression d'une partie de la discussion sur les espaces locaux, qui n'est pas particulièrement utile ni éclairante}



Rappelons d'abord l'inégalité de Jensen, ensuite illustrée par les espaces de Lebesgue.

\begin{proposition}[Inégalité de Jensen]
\label{prop:Jensen}
Soit $(X,\mu)$ un espace probabilisé, soit $f : X \to \bR^d$ intégrable, et soit $g : \bR^d \to ]-\infty,\infty]$ convexe et s.c.i. Alors
\begin{equation*}
	g\Big(\int_X f(x) \diff \mu(x)\Big) \leq \int_X g(f(x)) \diff \mu(x).
\end{equation*}
\end{proposition}

\begin{proof}
	Soit $g_0 := \alpha + \<\vp,\cdot\>$, où $\alpha \in \bR$ et $\vp \in \bR^d$, une fonction affine minorant $g$. Alors 
	\begin{equation*}
		\int_X g(f(x)) \diff \mu(x) \geq \int_X (\alpha + \<\vp,f(x)\>) \diff \mu(x) = \alpha + \Big\<\vp, \int_X f(x) \diff \mu(x)\Big\> = g_0\Big(\int_X f(x) \diff \mu(x)\Big),
	\end{equation*}
	par linéarité de l'intégrale.
	On conclut en remarquant que $g$ est le supremum d'une famille de fonctions affines, par le \nref{th:LegendreFenchel}. %, d'où le résultat.
\end{proof}

%On rappelle que l'espace local $L^p_\loc(\Omega)$ est un Fréchet, muni de la famille de semi-normes $|f|_n := \|f_{|K_n}\|_{L^p(K_n)}$ où $(K_n)_{n \in \bN}$ est une suite exhaustive de compacts, voir \eqref{eqdef:exhaustive_compacts}. 

\begin{corollary}
\label{corol:lpInclusLq}
	Soit $(X,\mu)$ un espace mesuré de mesure finie (non-nulle), et soient $1 \leq p \leq q \leq \infty$. Alors $L^p(X,\mu) \supset L^q(X,\mu)$, et pour toute $f \in L^q(X,\mu)$ 
\begin{equation*}
	\mu(X)^{-\frac 1 p} \|f\|_p \leq \mu(X)^{-\frac 1 q} \|f\|_q.
\end{equation*}
%En particulier, on a les inclusions continues $L^q(\Omega) \subset L^q_\loc(\Omega) \subset L^p_\loc(\Omega)$.
\end{corollary}

\begin{proof}
On suppose $q<\infty$, le cas $q=\infty$ étant trivial.
Il s'agit de montrer que 
\begin{equation*}
	\Big(\frac 1 {\mu(X)} \int_X |f(x)|^p \diff \mu(x)\Big)^{\frac q p} \leq \frac 1 {\mu(X)}\int_X |f(x)|^q \diff \mu(x).
\end{equation*}
Cela correspond à la \nref{prop:Jensen} appliquée à $\tilde f := |f|^p$, la fonction $g(s) := |s|^{\frac q p}$ qui est continue et convexe car $q\geq p$, et à la mesure de probabilité $\mu/\mu(X)$.
Noter que le \nref{th:HolderIneq} aurait pu être utilisé ici à la place de Jensen. 
% (évident si $\mu(X)=0$).
%L'inclusion s'en déduit car toute partie compacte de $\Omega$ est de mesure de Lebesgue finie.
\end{proof}

L'inégalité de Jensen est également au coeur de la définition de la convolution.

\begin{proposition}[Convolution dans $L^p$]
\label{prop:convolLp}
	Soient $f\in L^1(\bR^d)$ et $g\in L^p(\bR^d)$, où $p \in [1,\infty]$. Posons
	\begin{equation*}
		(f\star g) (x) := \int_{\bR^d} f(h) g(x-h) \diff h,
	\end{equation*}
	Alors cette intégrale converge pour presque tout $x\in \bR^d$, $f \star g \in L^p$ et $\|f\star g\|_p \leq \|f\|_1 \|g\|_p$.
\end{proposition}

\begin{proof}
Notons $X := \bR^d$, et supposons d'abord que $f$ et $g$ sont positives et que $\int_X f = 1$. 
% Par l'inégalité de Jensen
%	\begin{equation*}
%		(f\star g)(x)^p \leq \int_X f(h) g(x-h)^p \diff h.
%	\end{equation*}
	Alors $f\star g : X \to [0,\infty]$ est bien définie, mesurable et 
	\begin{align*}
		\int_X (f\star g)(x)^p \diff x & 
		= \int_{x \in X} \Big(\int_{h \in X} f(h) g(x-h) \, \diff h\Big)^p \, \diff x\\
		& \leq \int_{x \in X} \int_{h \in X} f(h) g(x-h)^p \diff h \, \diff x\\
%		& = \int_{h\in X} \int_{x \in X} f(h) g(x-h)^p \diff x \, \diff h\\
		& = \int_{h \in X} f(h) \Big(\int_{x \in X} g(x-h)^p \diff x\Big) \diff h  \\
		& = \Big(\int_{h \in X} f(h) \diff h\Big) \Big(\int_{x \in X} g(x)^p \diff x\Big).
	\end{align*}
	On a utilisé, successivement, (i) la définition de $f\star g$, (ii) l'inégalité de Jensen appliquée à la mesure de probabilité $f(h)\diff h$, (iii) le théorème de Fubini, (iv) l'invariance de la mesure de Lebesgue par translation.
	On a donc $\|f\star g\|_p \leq \|f\|_1 \|g\|_p$, et par homogénéité cette inégalité reste valable sans l'hypothèse $\int_X f = 1$. 
	Ainsi la première intégrale a une valeur finie, et donc $f\star g < \infty$ presque partout sur $X$. 
	
	Dans le cas de $f\in L^1$ et $g \in L^p$ non supposées positives, le résultat précédent appliqué à $|f|$ et $|g|$ montre que $h \mapsto f(h) g(x-h)$ est intégrable pour presque tout $x$, donc $f \star g$ est bien définie. 
	Finalement, on obtient l'estimation annoncée en notant que $|f\star g| \leq |f| \star |g|$. 
\end{proof}

La convolution $(f,g)\in L^1 \times L^p \mapsto f \star g \in L^p$ est clairement bilinéaire, continue par la \cref{prop:convolLp}, et commutative (dans le cas $p=1$). 
On note également que %, pour toutes $f\in L^1$ et $g \in L^p$
\begin{equation*}
	\supp(f\star g) \subset \overline {\supp(f) + \supp(g)} 
	:= \overline{\{x+y \mid x\in \supp(f),\, y\in \supp(g)\}}. %f(x) \neq 0, g(y) \neq 0\}}.
\end{equation*}
Enfin, on a pour $f,g \in L^1$ et $h \in L^p$ l'associativité 
\begin{align*}
	f\star (g\star h) &= (f\star g) \star h, &
	(f\star g \star h)(x) = \int_{u,v \in \bR^d} f(x-u-v) g(u) h(v) \, \diff u \, \diff v.
\end{align*}

La convolution avec une fonction d'intégrale unité et de support restreint permet d'approcher une fonction en la régularisant, comme le montrent les \cref{prop:convolApprox,prop:convolDeriv} suivantes.
%Le premier lemme est fondé sur un calcul proche de celui de la \cref{prop:convolLp}. 
On note $\tau_h g(x) := g(x-h)$ l'opérateur de translation, et $C^0_0(\bR^d)$ les fonctions continues tendant vers $0$ à l'infini.
%Une variante du calcul de la \cref{prop:convolLp}
\begin{lemma}
\label{lem:convolTranslat}
Soit $f \in L^1(\bR^d,\bR_+)$ d'intégrale unité, $g \in L^p(\bR^d)$, et $A \subset \bR^d$ mesurable. On a
%Pour toutes $f \in L^1(\bR^d,\bR_+)$ satisfaisant $\int f = 1$, et $g \in L^p(\bR^d)$, on a pour tout $A \subset \bR^d$ mesurable
\begin{equation*}
	\|f \star g - g\|^p_{L^p(A)} \leq \int_{\bR^d} f(h) \, \|\tau_h g - g\|^p_{L^p(A)} \, \diff h.
\end{equation*}	
\end{lemma}
\begin{proof}
	Le calcul est similaire à celui de la preuve de la \cref{prop:convolLp}. En effet
	%est proche de celui utilisé dans la \cref{prop:convolLp} définissant la convolution: 
	\begin{equation*}
		|f\star g(x) - g(x)|^p = \Big|\int_{\bR^d} f(h) (g(x-h)-g(x)) \diff h\Big|^p \leq \int_{\bR^d} f(h) |g(x-h)-g(x)|^p \diff h,
	\end{equation*}
	par la \nref{prop:Jensen}.
	Le résultat s'ensuit en intégrant par rapport à $x$, et en intervertissant les intégrales par Fubini.
\end{proof}

Dans toute la suite, étant donnée $\rho \in L^1(\bR^d)$ supportée dans $B'(0,1)$, on définit 
\begin{align*}
	\rho_\ve(x) &:= \tfrac 1 {\ve^d} \rho(\tfrac x \ve), &
	\text{donc } \int_{\bR^d} \rho_\ve &= \int_{\bR^d} \rho, &
	\text{et } \supp(\rho_\ve) =  \ve \supp(\rho) \subset B'(0,\ve).
\end{align*}
pour tous $\ve>0$ et $x\in \bR^d$. En particulier, $(f\star \rho_\ve) (x)$ ne dépend des valeurs de $f$ que sur $B(x,\ve)$. 

\begin{proposition}[Approximation par convolution]
\label{prop:convolApprox}
Soit $\rho \in L^1(\bR^d,\bR_+)$ supportée dans $B'(0,1)$ et d'intégrale unité. 
Soit $p\in [1,\infty[$ et $f \in L^p(\bR^d)$ (resp.\ $p=\infty$ et $f\in C^0_0(\bR^d)$). Alors 
\begin{align*}
\|f-\rho_\ve\star f\|_{p} &\leq \omega_p(\ve),& \text{où}\quad
\omega_p(\ve) &:= \sup_{|h| \leq \ve} \|f-\tau_h f\|_{p},
\end{align*}
pour tout $\ve>0$. De plus $\omega_p(\ve)\to 0$ lorsque $\ve\to 0$.
\end{proposition}

\begin{proof} 
L'estimation de $\|f-\rho_\ve \star f\|_p$ découle directement du \cref{lem:convolTranslat}, et il suffit donc de justifier que $\omega_p(\ve)\to 0$ lorsque $\ve\to 0$.
Dans le cas $p=\infty$, le \nref{th:Heine} appliqué sur une boule $B'(0,R)$ de grand rayon $R$, et la convergence vers $0$ à l'infini, montrent que la fonction $f$ est uniformément continue, comme désiré.
Dans le cas $p<\infty$, cela découle de la densité des fonctions continues à support compact dans $L^p(\bR^d)$, établie dans la \cref{prop:densC0Lp}.
\end{proof}

\begin{proposition}[Dérivation d'une convolution]
\label{prop:convolDeriv}
Soit $f \in L^1(\bR^d)$, et $g \in C^k_c(\bR^d)$. Alors $f\star g \in C^k(\bR^d)$ et $\partial_\alpha (f\star g) = f \star (\partial_\alpha g)$ pour tout $\alpha \in \bN^d$ tel que $|\alpha|_1 \leq k$. 
\end{proposition}

\begin{proof}
	Il suffit de montrer que $\partial_\alpha (f\star g) = f \star (\partial_\alpha g)$ lorsque $|\alpha|_1 = 1$ et $g\in C^1_c(\bR^d)$, puis d'appliquer une récurrence immédiate. 
	On calcule
\begin{equation*}
	\frac {(f\star g)(x+\delta e_i) - (f\star g) (x)} \delta - (f\star \partial_i g) (x) = \int_{\bR^d} f(h) \Big[\frac{g(x+\delta e_i-h) - g(x-h)} \delta - \partial_i g(x-h)\Big] \diff h,
\end{equation*}
et le terme en $g$ converge uniformément vers $0$ lorsque $\delta\to 0$. 
\end{proof}

\begin{corollary}
\label{corol:testDenseLp}
$\cD(\Omega)$ est dense dans $L^p(\Omega)$ pour tout $1 \leq p < \infty$. 
%En particulier, ces espaces sont séparables.
%	$\cD(\Omega)$ est dense dans $L^p(\Omega)$, et dans $L^p_\loc(\Omega)$, pour tout $1 \leq p < \infty$. En particulier, ces espaces sont séparables.
\end{corollary}

\begin{corollary}
\label{corol:convolTestFun}
	Soient $\rho\in \cD(\bR^d)$ et $\phi \in \cD(\Omega)$. 
	Alors $\rho_\ve \star \phi \to  (\int_{\bR^d} \rho)\, \phi$ lorsque $\ve \to 0$ dans $\cD(\Omega)$.
\end{corollary}

\begin{proof}
Notons que $\supp(\rho_\ve \star \phi) \subset \ve\supp(\rho) + \supp(\phi)$, donc il existe $K \csubset \Omega$ et $\ve_*>0$ tels que $\supp(\rho_\ve \star \phi)\subset K$ pour tout $0 <\ve \leq \ve_*$. 
	On obtient la convergence uniforme $\|\rho_\ve \star \phi - m \phi\|_\infty\to 0$ lorsque $\ve \to 0$, où $m:=\int_{\bR^d} \rho$, en appliquant la \cref{prop:convolApprox} aux parties positive et négative de $\rho$, divisées par leurs intégrales respectives (avec des adaptations évidentes si l'une ou l'autre est nulle, et si le compact $\supp(\rho)$ n'est pas contenu dans $B'(0,1)$). Pour les dérivées d'ordre supérieur utiliser l'identité $\partial_\alpha (\rho_\ve \star \phi) = \rho_\ve \star (\partial_\alpha \phi)$, puis le même argument donnant la convergence uniforme. 
\end{proof}


La convolution permet, outre les propriétés de densité, d'établir un résultat de compacité.
%Pour cela, nous rappelons d'abord certaines conséquences immédiates de l'inégalité de Jensen.

\begin{theorem}[Fréchet-Kolmogorov]
\label{th:FrechetKolmogorov}
	Soit $p \in [1,\infty[$ et $\cF$ une partie bornée de $L^p(\Omega)$ t.q.\ %satisfaisant aux hypothèses suivantes :
	\begin{enumerate}[(i)]
		\item (Régularité sous translation) Pour tout $K \csubset \Omega$, il existe un module de continuité $\omega_K$ t.q.\
		\begin{equation*}
			\|f - f(\cdot-h)\|_{L^p(K)} \leq \omega_K(|h|),\, \quad \forall f \in \cF,\, \forall h \in \bR^d \text{ t.q. } |h| \leq d(K, \bR^d \sm \Omega).
		\end{equation*}
		\item (Masse évanescente au bord) Pour tout $\ve>0$, il existe $K \csubset \Omega$ tel que 
		\begin{equation*}
			\|f\|_{L^p(\Omega \sm K)} \leq \ve, \quad \forall f \in \cF.
		\end{equation*}
	\end{enumerate}
	Alors $\cF$ est une partie précompacte de $L^p(\Omega)$. 
\end{theorem}

\begin{proof}
On note $\|\cdot\|_p := \|\cdot\|_{L^p(\Omega)}$.
Dans la suite, on étend implicitement toute fonction $f \in \cF$ sur $\bR^d\sm \Omega$, par exemple par $0$, pour que toutes les quantités soient bien définies; on peut vérifier que les résultats ne dépendent pas de cette extension, grâce spécifiquement à l'utilisation de la fonction indicatrice lissée $\phi$ et à la petitesse du vecteur $h\in \bR^d$ ci-dessous.

	Commençons par introduire un certain nombre d'objets.
	Soit $\ve>0$, arbitraire mais fixé dans la preuve.
	Soit $K_0 \csubset \Omega$ donné par la propriété (ii) pour la valeur $\ve$. 
	Soit $K_1 \csubset \Omega$ tel que $K_0 \csubset \mathring K_1$ et soit $\vp \in C^1_c(\Omega,[0,1])$ telle que $\vp=1$ sur $K_0$ et $\supp(\vp) \subset \mathring K_1$.
	On a pour toute $f \in \cF$ et tout $|h| \leq d(K_1, \bR^d \sm \Omega)$
	\begin{equation}
	\label{eq:fphfp}
		\|f\vp - \tau_h(f\vp)\|_p \leq \|(f-\tau_h f)\vp\|_p + \|\tau_h f(\vp-\tau_h\vp)\|_p \leq \omega(|h|),
	\end{equation}
	pour un certain module de continuité $\omega$, par la propriété (i) appliquée à $\supp(\vp)$, et la régularité de $\vp$.
	Soit $\eta>0$ tel que $\omega(\eta)<\ve$ et $\supp(\vp) + B'(0,\eta) \subset K_1$.
	Soit enfin $\rho \in C^1(\bR^d,\bR_+)$ tel que $\supp(\rho) \subset B'(0,\eta)$ et $\int_{\bR^d} \rho=1$. 
	Posons 
	\begin{equation*}
		\cG := \{(f \vp) \star \rho  \mid f \in \cF\} \subset L^p(\Omega).
	\end{equation*}
	Étant donnée $g = (f\vp) \star \rho \in \cG$, on a par construction $g \in C^1$ et $\supp(g) \subset K_1$, de plus
	\begin{align*}
		\|g\|_\infty &= \|(f\vp) \star\rho\|_\infty \leq \|f \vp\|_1 \|\rho\|_\infty \leq |K_1|^{1-\frac 1 p} \|f\|_p \|\rho \|_\infty,\\
		\|\nabla g\|_\infty &= \| (f\vp) \star \nabla \rho\|_\infty \leq \|f\vp\|_1 \|\nabla \rho\|_\infty \leq |K_1|^{1-\frac 1 p} \|f\|_p \|\nabla \rho\|_\infty.
	\end{align*}
	Comme $\cF$ est borné dans $L^p$, l'ensemble $\cG$ est constitué de fonctions uniformément bornées ainsi que leur gradient, et nulles en dehors d'un compact. Par le \nref{th:AscoliLip}, l'ensemble $\cG$ est donc précompact pour la topologie $L^\infty$. 
	Ainsi, il existe $I \geq 1$ un entier et $f_1, \cdots, f_I \in \cF$ telles que $\forall f \in \cF$, $\exists i\in \braii{1,I}$, $\|[(f-f_i)\vp] \star \rho \|_\infty \leq \ve |K_1|^{-\frac 1 p}$, 
	 et donc $\|[(f-f_i)\vp] \star \rho \|_p\leq \ve$ par le \cref{corol:lpInclusLq}. 
	 
	En utilisant \cref{lem:convolTranslat} et l'hypothèse (ii) on obtient par ailleurs pour toute $f \in \cF$
	\begin{equation*}
		\|f-(f\vp)\star \rho\|_p \leq \|f-f\vp\|_p + \|f\vp - (f\vp) \star\rho\|_p \leq 2 \ve
	\end{equation*}
	de sorte que $\cF \subset \cup_{1 \leq i \leq I} B_{L^p} ((f_i\vp)\star \rho,3 \ve)$, ce qui établit la précompacité et conclut.
\end{proof}

Le \nref{th:FrechetKolmogorov} permet de montrer la compacité de l'opérateur intégral suivant, analogue au \cref{lem:convK_compact}.

\begin{exemple}[Compacité d'un opérateur intégral, version $L^2$]
\label{ex:operCompactL2}
Pour tout $K \in L^2([0,1]^2,\bR)$, l'endomorphisme de $L^2([0,1])$ défini par $\cK(f)(x) := \int_0^1 K(x,y) f(y) \diff y$ est compact.
\end{exemple}

\begin{proof}
La bonne définition et la continuité de $\cK$ découlent de l'inégalité de Cauchy-Schwarz : 
	\begin{align*}
		|\cK(f)(x)|^2 &\leq \Big(\int_0^1 K(x,y)^2 \diff y\Big) \Big( \int_0^1 f(y)^2 \diff y\Big), &
		\text{ donc } 
		\| \cK(f)\|^2_{L^2([0,1])} \leq \|K\|_{L^2([0,1]^2)}^2 \|f\|^2_{L^2([0,1])}, 
	\end{align*}
	la seconde inégalité s'obtenant à partir de la première par intégration par rapport à $x \in [0,1]$. 
	Par ailleurs en notant $I_\ve := [\ve,1-\ve] \subset I := [0,1]$, on obtient %de sorte que $d(x,\partial I) \leq \ve$ ssi $x \in [0,\ve] \cup [1-\ve,1]$, on a 
	\begin{equation*}
		\|\cK(f)\|_{L^2(I \sm I_\ve)}^2 = 
		\int_{x \notin I_\ve} \Big(\int_{y\in I} K(x,y) f(y) \diff y\Big)^2 \diff x \leq 
		\|f\|^2_2 \int_{x \notin I_\ve, y \in I} K(x,y)^2 \diff x \, \diff y.
	\end{equation*}
	D'autre part, pour tout $h\in \bR$ tel que $|h| \leq \ve$, on a 
	\begin{align*}
		\|\cK(f)-\tau_h \cK(f) \|_{L^2(I_\ve)}^2 &= 
		\int_{x \in I_\ve} \Big| \int_{y\in I} K(x,y) f(y) \diff y - \int_{y\in I} K(x-h,y) f(y) \diff y \Big|^2 \diff x \\
		&\leq \|f\|^2_2 \int_{x \in I_\ve, y \in I} \big(K(x,y)-K(x-h,y)\big)^2 \diff x \diff y.
	\end{align*}
	En utilisant la densité des fonctions continues dans $L^2([0,1]^2)$, on obtient que les intégrales de $K$ apparaissant dans les membres de droite des deux dernières estimations tendent vers $0$ lorsque $\ve\to 0$. L'ensemble $\{\cK(f)\mid \|f\|_2 \leq 1\}$ satisfait donc les hypothèses du \cref{th:FrechetKolmogorov}; il est donc précompact, et ainsi l'opérateur $\cK$ est compact, comme annoncé.
\end{proof}


\subsection{Convolution de distributions et espaces de Sobolev}
\label{subsec:convolSobolev}


Nous définissons dans cette sous-section les espaces de Sobolev, qui sont constitués des fonctions dont les dérivées (au sens des distributions) jusqu'à un certain ordre admettent des représentants dans $L^p$. La convolution est de nouveau un outil fondamental, entre autres pour  montrer les résultats de densité. Nous concluons par le théorème de prolongement dans les espaces de Sobolev, admis, et une application au problème de Poisson.

\begin{retenir}
La définition des espaces de Sobolev, la densité de $\cD(\Omega)$ dans ceux-ci, le théorème de prolongement, le problème de Poisson.
(\Cref{def:Sobolev,prop:testDenseWsp,th:prolSobolev,th:pbPoisson}.)	
\end{retenir}

Rappelons que $\Omega\subset \bR^d$ désigne un ouvert quelconque, sauf mention contraire.
On note $\<\vp,f\> := \vp(f)$ le crochet de dualité entre une distribution $\vp \in \cD(\Omega)^*$ et une fonction test $f \in \cD(\Omega)$, 
et $\ordre(\vp) \in \bN\cup \{\infty\}$ l'ordre de cette distribution.


L'identité de convolution suivante sera utile : pour tous $\rho, f \in L^1(\bR^d)$, $g \in L^\infty( \bR^d)$, on a par l'identité de Fubini en notant $\overline \rho (x) = \rho(-x)$ la fonction réfléchie
\begin{equation*}
	\int_{\bR^d} (\rho\star f) (x) \, g(x) \, \diff x = 
	\int_{\bR^d} f(x) \, (\overline \rho \star g)(x)\, \diff x = 
	\iint_{\bR^d\times \bR^d} f(y) \rho(x-y) g(x) \diff x \diff y
\end{equation*}


%\begin{lemma}
%Si $f \in L^1_\loc(\Omega)$ est telle que $\int f\vp = 0$ pour tout $\vp \in \cD(\Omega)$, alors $f=0$. 
%\end{lemma}
%
%\begin{proof}
%	Soit $\rho \in C^\infty(\bR^d,\bR_+)$ tel que $\supp(\rho) \subset B(0,1)$ et $\int \rho=1$, et soit $\rho_\ve(x) := \frac 1 {\ve^d} \rho(\frac x \ve)$ pour tout $\ve>0$. 
%	Soit $K \csubset \mathring K'$ avec $K' \subset  \Omega$, et soit $\delta := d(K',\bR^d \sm \Omega)$. Alors pour tout $0 < \ve < \delta$ et toute $\phi \in C^0_c(K')$, on a 
%	\begin{equation*}
%		0 = \int_\Omega f(x) (\overline \rho_\ve \star \phi) (x) \diff x  = \int_{K'} (\rho_\ve \star f)(x) \phi(x) \diff x.
%	\end{equation*}
%	On en déduit que la fonction continue $\rho_\ve \star f = 0$ est nulle sur $K$. Comme $\|\rho_\ve \star f - f\|_{L^1(K)} \to 0$ lorsque $\ve\to 0$ d'après le \cref{prop:convolApprox}, $f$ est nulle sur $K$. Comme $K\csubset \Omega$ est arbitraire, $f$ est nulle sur $\Omega$. 
%\end{proof}

\begin{proposition}[Distribution associée à une fonction]
\label{prop:distribL1Loc}
Soit $f \in L^1_\loc(\Omega)$. On définit 
\begin{equation*}
	\<U_f,\phi\> := \int_\Omega f \phi
\end{equation*}
pour toute $\phi \in \cD(\Omega)$.
Alors $U_f$ est une distribution d'ordre $0$. De plus $f \in L^1_\loc \mapsto U_f \in \cD(\Omega)^*$ est continue et injective, et pour cette raison on identifie souvent $f$ et $U_f$. 
\end{proposition}

\begin{proof}
L'intégrale est bien définie, car $\phi$ est à support compact, et linéaire.
De plus $\<U_f,\phi\> \leq \|f\|_{L^1(K)} \|\phi\|_\infty$ pour tout $K \csubset \Omega$ et toute $\vp \in \cD_K(\Omega)$, ce qui montre qu'elle est d'ordre $0$. 

Pour la preuve de l'injectivité, supposons $U_f = 0$. 
%La preuve de l'injectivité demande un peu de travail, et procède par convolution.
	Soit $\rho \in C^\infty(\bR^d,\bR_+)$ tel que $\supp(\rho) \subset B(0,1)$ et $\int \rho=1$, et soit $\rho_\ve(x) := \frac 1 {\ve^d} \rho(\frac x \ve)$ pour tout $\ve>0$. 
	Soient $K \csubset \Omega$ et soit $\ve >0$ tels que $K + B'(0,\ve) \csubset \Omega$. Alors pour toute $\vp \in \cD_K(\Omega)$ on a 
%Étant donnée $\vp \in \cD(\Omega)$, on a $\supp(\overline \rho_\ve \star \vp) \subset \supp(\vp) + B'(0,\ve) \csubset \Omega$ pour tout $\ve>0$ assez petit, d'où l'on tire	
%	Soit $K \csubset \Omega$, et soit $\delta := d(K,\bR^d \sm \Omega)$.
%	Soit $K \csubset \mathring K'$ avec $K' \csubset  \Omega$, et soit $\delta := d(K',\bR^d \sm \Omega)$. 
%Alors pour tout $0 < \ve < \delta$ et toute $\phi \in \cD_K(\Omega)$, on a en supposant $U_f = 0$ et comme $\overline \rho_\ve \star \phi \in \cD(\Omega)$
	\begin{equation*}
		0 = \int_\Omega f(x) (\overline \rho_\ve \star \phi) (x) \diff x  = \int_K (\rho_\ve \star f)(x) \phi(x) \diff x,
	\end{equation*}
	d'où l'on déduit que $\rho_\ve \star f$ est nulle sur $\mathring K$ (par densité de $\cD(\mathring K)$ dans $L^1(\mathring K)$).
%	On en déduit que la fonction $C^\infty$ donc continue $\rho_\ve \star f = 0$ est nulle sur $\mathring K$ (approcher $\rho_\ve \star f$ dans $L^1$ avec $\phi$). 
Comme $\|\rho_\ve \star f - f\|_{L^1(K)} \to 0$ lorsque $\ve\to 0$ d'après la \cref{prop:convolApprox}, $f$ est nulle sur $\mathring K$. Comme $K\csubset \Omega$ est arbitraire, $f$ est nulle sur $\Omega$. 
\end{proof}

La \cref{prop:distribL1Loc} permet de concevoir les distributions comme des fonctions généralisées.
En particulier, on a pour tout $p \in [1,\infty]$ l'injection continue $L^p(\Omega)\subset L^1_{\loc}(\Omega) \subset \cD^*(\Omega)$, par le \cref{corol:lpInclusLq}. Réciproquement, le résultat suivant donne un critère pour identifier les distributions issues des espaces de Lebesgue $L^p$, $p>1$. Noter que le critère ne s'étend pas à l'exposant $p=1$; en fait, les distributions continues pour la norme $L^\infty$ s'identifient aux mesures de Radon de masse finie sur $\Omega$.

\begin{corollary}[Distributions image de $L^p$]
\label{corol:distribLp}
Soit $T \in \cD(\Omega)^*$, $p\in ]1,\infty]$ et $q$ t.q.\ $\frac 1 p + \frac 1 q = 1$. On a (a) $|T(\phi)| \leq C\|\phi\|_q$ pour toute $\phi \in \cD(\Omega)$, ssi (b) il existe $f \in L^p(\Omega)$, telle que $T = U_f$. 
\end{corollary}

\begin{proof} 
Si (a) est satisfaite, alors $T$ définit une forme linéaire continue sur $L^q(\Omega)$, car $\cD(\Omega)$ est dense dans cet espace par le \cref{corol:testDenseLp}, d'où l'existence de $f\in L^p$ par le \cref{th:lpDualLq}.

Réciproquement, supposons (b). On a $L^p(\Omega) \subset L^1_\loc(\Omega)$ par le \cref{corol:lpInclusLq}, donc $U_f$ est bien définie par \cref{prop:distribL1Loc}, et on conclut par le \nref{th:HolderIneq}.
\end{proof}

 De nombreuses opérations sur les fonctions s'étendent aux distributions par dualité, comme la dérivation dans la proposition suivante, ou encore la multiplication par une fonction lisse définie par $\<\phi T,\psi\> := \<T, \phi \psi\>$ pour toutes $\vp,\psi \in \cD(\Omega)$, $T \in \cD^*(\Omega)$.

\begin{proposition}[Dérivation d'une distribution]
\label{prop:deriv_distrib}
Soit $T \in \cD(\Omega)^*$, $\alpha \in \bN^d$ et $s:=|\alpha|_1$. On définit
\begin{equation*}
	\<\partial_\alpha T, \phi\> := (-1)^s \<T,\partial_\alpha \phi\>, %{|\alpha|_1}
\end{equation*}
pour toute $\phi \in \cD(\Omega)$. Alors $\partial_\alpha T \in \cD(\Omega)^*$, et $\ordre(\partial_\alpha T) \leq \ordre(T) + s$. De plus $\partial_\alpha U_f = U_{\partial_\alpha f}$ pour toute $f \in C^s(\Omega)$. Enfin $\partial_\alpha : \cD^*(\Omega) \mapsto \cD^*(\Omega)$ est continue.
\end{proposition}

\begin{proof}
	Découle des définitions de l'ordre et de la continuité des distributions.
\end{proof}

Nous arrivons à la définition des espaces de Sobolev.
La condition $\partial_\alpha f \in L^p(\Omega)$ apparaissant dans celle-ci signifie que la dérivée distributionnelle $\partial_\alpha f \in \cD^*(\Omega)$ appartient au sous-ensemble $L^p(\Omega) \subset \cD^*(\Omega)$ défini par l'identification de la \cref{prop:distribL1Loc}. 

%La condition ``$\partial_\alpha f \in L^p(\Omega)$'' apparaissant dans celle-ci est un abus de notation, devant être compris comme suit : ``il existe $g \in L^p(\Omega)$ telle que $\partial_\alpha U_f = U_g$, en tant qu'éléments de $\cD^*(\Omega)$.''
%busive et signifie en fait que la distribution $\partial_\alpha U_f$ est l'image $U_g$ d'un élément de $L^p$, nécessairement unique, que l'on note $g = \partial_\alpha f$.


\begin{definition}[Espace de Sobolev]
\label{def:Sobolev}
	Pour tout $s \in \bN$, $p \in [1,\infty]$, on définit l'espace de Sobolev $W^{s,p}(\Omega) \subset L^p(\Omega)$ et sa norme comme suit : 
	\begin{align*}
		W^{s,p}(\Omega) &:= \{f \in L^p(\Omega) \mid \forall |\alpha|_1 \leq s,\, \partial_\alpha f \in L^p(\Omega)\}, &
		\|f\|_{W^{s,p}(\Omega)}^p &:= \sum_{|\alpha|_1\leq s} \|\partial_\alpha f\|_{L^p(\Omega)}^p,
	\end{align*}
	avec la convention $\|f\|_{W^{s,\infty}(\Omega)} := \max_{|\alpha|_1 \leq s} \|\partial_\alpha f\|_\infty$.
\end{definition}


\begin{proposition}
\label{prop:SobolevBanach}
	Pour tous $s \in \bN$ et $p \in [1,\infty]$, l'espace $W^{s,p}(\Omega)$ est un Banach. Si $p=2$, c'est un espace de Hilbert, aussi noté $H^s(\Omega) := W^{s,2}(\Omega)$.
\end{proposition}

\begin{proof}
	Considérons une suite de Cauchy $(f_n)_{n \in \bN}$ de $W^{s,p}(\Omega)$. Alors $\partial_\alpha f_n$ est de Cauchy dans $L^p(\Omega)$ pour tout $|\alpha|_1 \leq s$, et donc converge vers une limite notée $f^\alpha \in L^p$, par complétude. 
	Par ailleurs, $\partial_\alpha f_n$ converge vers $\partial_\alpha f^0$ au sens des distributions, par continuité de l'inclusion $L^p(\Omega) \subset \cD^*(\Omega)$ et continuité de la dérivation pour la topologie de $\cD^*(\Omega)$. Ainsi $\partial_\alpha f^0 = f^\alpha$ pour tout $|\alpha|_1 \leq s$, d'où l'on déduit que $f_n \to f^0$ dans $W^{s,p}$, ce qui établit la complétude. 
	
	Dans le cas $p=2$, la norme dérive du produit scalaire $\<f,g\>_{W^{s,2}} := \sum_{|\alpha|_1 \leq s} \<\partial_\alpha f,\partial_\alpha g\>_{L^2}$. 
\end{proof}

\notei{Suppression de la remarque sur les espaces locaux, dont la definition est évidente et qui ne sont pas utiles ici.}
%\begin{remark}
%	Les espaces locaux $W^{s,p}_\loc(\Omega) \subset L^p_\loc(\Omega)$ sont les espaces de Fréchet munis de la famille de semi-normes $\|f\|_{W^{s,p}(\mathring K_n)}$ associée à une suite exhaustive $(K_n)_{n \in \bN}$ de compacts \eqref{eqdef:exhaustive_compacts}. 
%\end{remark}

La convolution des distributions prolonge clairement celle des fonctions. 
Elle permet de régulariser, et de montrer des résultats de densité, au sein des distributions et des espaces de Sobolev.


\begin{definition}[Convolution d'une distribution]	
\label{def:convolDistrib}
	Soient $T \in \cD(\bR^d)^*$ et $\rho \in \cD(\bR^d)$. On définit la convolée $T \star \rho : \bR^d \to \bR$ par 
	$
		(T \star \rho) (x) := \<T, \rho(x - \cdot)\>.
	$
\end{definition}

\begin{proposition}
\label{prop:convolDistrib}
	Soient $T \in \cD(\bR^d)^*$ et $\rho,\phi \in \cD(\bR^d)$. 
	Alors $\<T \star \rho, \phi\> = \<T,\overline \rho \star \phi\>$, d'où l'on déduit que 
	(i) $T \star \rho \in C^\infty (\bR^d)$, 
	(ii) $\partial_\alpha (T\star \rho) = (\partial_\alpha T) \star \rho = T \star \partial_\alpha \rho$, pour tout $\alpha \in \bN^d$, 
	(iii) $T \star \rho_\ve \weakto (\int_{\bR^d} \rho) T $ lorsque $\ve \to 0$ au sens des distributions. En particulier $C^\infty(\bR^d)$ est dense dans $\cD(\bR^d)^*$. 
\end{proposition}

\begin{proof}
Définissons $T*\rho \in \cD^*(\bR^d)$ par $\<T * \rho,\phi\> := \<T,\overline \rho \star \phi\>$.
Soit $K$ un compact, $K' := K + \supp(\rho)$ qui est également compact, et $s\in \bN$ l'ordre de la distribution $T$ restreinte à $K'$. On a pour toute $\phi \in \cD_K(\bR^d)$
\begin{equation*}
	|\<T*\rho,\phi\>| 
	= |\<T,\overline \rho \star \phi\>| 
	\leq C \sum_{|\alpha|_1 \leq s} \|\partial_\alpha(\overline \rho \star \phi)\|_\infty 
	\leq \Big(C \sum_{|\alpha|_1 \leq s} \|\partial_\alpha \overline \rho\|_\infty\Big) \|\phi\|_1,
\end{equation*}
compte tenu de l'identité $\partial_\alpha(\overline \rho \star \phi) = (\partial_\alpha\overline \rho) \star \phi$ et de la continuité de la convolution \cref{prop:convolLp}.
Il s'ensuit que $T * \rho \in L^\infty_\loc(\bR^d)$, par le \nref{corol:distribLp}. Comme le même raisonnement montre que $\partial_\beta (T * \rho) \in L^\infty_\loc(\bR^d)$ pour $\beta \in \bN^d$, on obtient que $T * \rho \in C^\infty(\bR^d)$.

Considérons  $\eta \in \cD(\bR^d)$ avec $\int \eta = 1$. Alors $\<T * \rho, \eta_\ve \> = \<T,\overline \rho \star \eta_\ve\>$ converge à la fois vers $T*\rho(0)$ et $T\star \rho(0)$ lorsque $\ve\to 0$, puisque $\overline \rho \star \eta_\ve \to \overline \rho$ dans $\cD(\bR^d)$ par le  \cref{corol:convolTestFun}.  De même $T*\rho(x)=T\star \rho(x)$ pour tout $x \in \bR^d$.
On en déduit que les définitions notées $*$ et $\star$ de la convolution d'une distribution coïncident, ce qui établit (i). 
Finalement, les propriétés (ii) et (iii) annoncées découlent immédiatement de la définition notée $*$ de la convolution.
\end{proof}

\begin{proposition}
\label{prop:testDenseWsp}
	Soient $s\in \bN$, $p \in [1,\infty[$, $f \in W^{s,p}( \bR^d)$ et $\rho \in \cD(\bR^d)$ telle que $\int \rho=1$. Alors 
	$f \star \rho_\ve \to f$ dans $W^{s,p}(\bR^d)$, et de plus $\cD(\bR^d)$ est dense dans $W^{s,p}(\bR^d)$. % qui est donc séparable.  
\end{proposition}

\begin{proof}
	On a $f \star \rho_\ve \to f$ lorsque $\ve \to 0$ pour la topologie de $L^p$ par la \cref{prop:convolApprox}, et de même $\partial_\alpha (f \star \rho_\ve) = (\partial_\alpha f) \star \rho_\ve \to \partial_\alpha f$ pour la topologie de $L^p$, d'où la convergence annoncée. On a utilisé implicitement le fait que la convolution au sens des distributions et au sens de $L^p$ coïncident sur $L^p(\bR^d) \times \cD(\bR^d)$, voir les \cref{prop:convolLp,def:convolDistrib} (ce qui découle de leur coïncidence sur $\cD(\bR^d)\times \cD(\bR^d)$ et de leur continuité pour la topologie de $L^p(\bR^d) \times \cD(\bR^d)$).
	
	Pour établir la densité de $\cD(\bR^d)$, il est nécessaire de considérer des approximations ayant un support borné. Pour cela, considérons $\chi_n \in C^\infty (\bR^d , [0,1])$, pour tout $n \in \bN^*$, une fonction radiale, telle que $\chi_n = 1$ sur $B'(0,n)$ et $\supp(\chi_n) \subset B(0,n+1)$, et dont les dérivées sont bornées indépendamment de $n$. (Par exemple $\chi_n(x) = 1-\sigma(|x|-n)$ où $\sigma \in C^\infty( \bR , [0,1])$ satisfaisant $\sigma_{|]-\infty,0]}=0$ et $\sigma_{|[1,\infty[} = 1$ est une approximation de la fonction Heaviside, voir la \cref{prop:cinfHeaviside}.) Alors $\|\partial_\alpha (f \chi_n) - \partial_\alpha f\|_p \leq C(\alpha) \sum_{|\beta|_1 \leq |\alpha|_1} \| \partial_\beta f\|_{L^p(\bR^d \sm B(0,n))}$, pour tous $\alpha \in \bN^d$ et $f \in W^{s,p}(\bR^d)$, où la constante $C$ est indépendante de $n$ (utiliser la formule \eqref{eq:deriv_binom} du binôme pour calculer $\partial_\alpha (f \chi_n)$). Ainsi $\|f \chi_n -  f\|_{W^{s,p}} \to 0$ lorsque $n \to \infty$. Par ailleurs $\|(f \chi_n)\star \rho_\ve - f \chi_n\|_{W^{s,p}} \to 0$ lorsque $\ve\to 0$ d'après le premier point établi, et comme $(f \chi_n)\star \rho_\ve \in \cD(\bR^d)$ cet ensemble est dense dans $W^{s,p}(\bR^d)$, ce qui conclut.
\end{proof}



%\todoi{$W^{1,\infty} = $ Lipschitz}

%\todoi{Ici ? Utile ?}
\notei{Supprimé le support d'une distribution, assez trivial et pas utilisé dans le cours.}
%\begin{definition}[Support d'une distribution] Pour $T \in \cD(\Omega)^*$ on définit
%\begin{align*}
%	\supp(T) := \{ x\in \Omega \mid \forall \ve>0,\, \exists \phi \in \cD(B(x,\ve)),\, \<T,\phi\> \neq 0\}.
%\end{align*}
%Cette définition prolonge le support des fonctions, et $\supp(T \star \rho) \subset \supp(T) + \supp(\rho)$.
%\end{definition}

%Un théorème de prolongement, admis et s'appliquant , permet d'étendre le résultat de densité

Un résultat de prolongement, admis, permet d'étendre à $\bR^d$ les éléments d'un espace de Sobolev $W^{s,p}(\Omega)$ défini sur un ouvert régulier, et donc de leur appliquer les techniques de convolution et les résultats de densité précédents.


\begin{definition}[Ouvert régulier]
\label{def:ouvertReg}
	Un ouvert $\Omega \subset \bR^d$ est dit régulier (de classe $C^\infty$) si pour tout $x \in \partial \Omega$ il existe des ouverts $U\ni x$ et $V \ni 0$ et un difféomorphisme $\phi \in C^\infty(U,V)$, tels que $\phi(x) = 0$ et $\phi(U \cap \Omega) = V \cap (]0,\infty[ \times \bR^{d-1})$. 
\end{definition}

\begin{theorem}[Prolongement dans les espaces de Sobolev, admis]
\label{th:prolSobolev}
Soient $\Omega \subset \bR^d$ un ouvert régulier, $s\in \bN$ et $p \in [1,\infty]$. Alors il existe un opérateur linéaire continu $P : W^{s,p}(\Omega)\to W^{s,p}(\bR^d)$ tel que $(P u)_{|\Omega} = u$ pour tout $u \in W^{s,p}(\Omega)$.  
\end{theorem}

\begin{corollary}
\label{corol:testDenseWspO}
	Soient $\Omega \subset \bR^d$ un ouvert régulier, $s \in \bN$ et $p \in [1,\infty[$. Alors $\{\phi_{|\Omega} \mid \phi \in \cD(\bR^d)\}$ est dense dans $W^{s,p}(\Omega)$. 
\end{corollary}

\begin{proof}
	Soit $u\in W^{s,p}(\Omega)$, $Pu\in W^{s,p}(\bR^d)$ son extension par le \cref{th:prolSobolev}, et $(\phi_n)_{n \in \bN}$ une suite dans $\cD(\bR^d)$ convergeant vers $Pu$ dans $W^{s,p}(\bR^d)$ par la \cref{prop:testDenseWsp}. Alors $\phi_{n|\Omega} \to u$ dans $W^{s,p}(\Omega)$ comme souhaité. 
\end{proof}

Nous terminons cette section par une présentation du problème de Poisson. La preuve choisie, particulièrement directe, est fondée sur le théorème de représentation de Riesz.
Cependant, elle se généralise difficilement, car elle requiert une forme bilinéaire $Q$ définissant un produit scalaire et une norme équivalents à ceux de $H^1(\Omega)$. 
Le théorème classique de Lax-Milgram, qui ne requiert pas la symétrie de $Q$, aurait également pu s'appliquer.
Une autre preuve alternative utilisant la compacité faible est présentée dans \cite{desvilettes2015analyse}, voir aussi l'\nref{ex:nonLinMinH1}. 


\begin{remark}[Conditions au bord de Dirichlet]
\label{rem:bcDirichlet}
L'adhérence de $\cD(\Omega)$ dans $W^{1,p}(\Omega)$ est notée $W^{1,p}_0(\Omega)$. Ce sous espace, strict en général, est donc séparable et complet. On dit qu'il implémente les conditions au bord nulles de Dirichlet.
\end{remark}


\begin{theorem}[Problème de Poisson] % avec conditions au bord de Neumann ou de Dirichlet]
\label{th:pbPoisson}
Soit $\Omega \subset \bR^d$ un ouvert borné, et soit $f \in L^2(\Omega)$. %Alors~:
\begin{itemize}
	\item (Problème de Neumann)
Il existe une unique $u \in H^1(\Omega)$ telle que 
\begin{align*}
	\forall v \in H^1(\Omega),\, Q(u,v) &= L(v), &
	\text{où }\
	Q(u,v) &:= \int_\Omega \<\nabla u,\nabla v\> + u v &
	\text{et }\ L(v) &:= \int_\Omega f v.
\end{align*}
\item (Problème de Dirichlet) Il existe une unique $u \in H^1_0(\Omega)$ t.q.\ $\forall v \in H^1_0(\Omega)$, $Q(u,v)=L(v)$. 
\end{itemize}
\end{theorem}

\begin{proof}
	La forme bilinéaire symétrique définie positive $Q$ définit une norme $\sqrt{Q(u,u)}$ équivalente à celle de $H^1(\Omega)$, et la forme linéaire $L$ est continue pour la norme $L^2$ donc aussi $H^1$. Le résultat découle donc de la complétude de $H^1$ et du \nref{th:reprRiesz}.	
%	\todoi{Mettre ailleurs la preuve alternative. Peut-être après la séparabilité de $W^{s,p}(\Omega)$, $p<\infty$.}
%	Noter qu'une preuve alternative est possible (cf polycopié de \cite{desvilettes2015analyse}), en considérant la minimisation de la fonctionnelle
%	\begin{equation*}
%		J(u) := \frac 1 2 \int_\Omega \Big(\| \nabla u\|^2 + u^2\Big) - \int_\Omega f u,
%	\end{equation*}
%	et en notant qu'une suite minimisante est nécessairement bornée. Par séparabilité de $L^2$, on a des sous-suites faiblement convergentes pour $u$ et ses dérivées $\partial_i u$, $1 \leq i \leq d$.  
\end{proof}

Le problème de Poisson admet plusieurs reformulations. En termes d'optimisation, il correspond à la minimisation de la fonctionnelle 
\begin{equation*}
	J(u) := \frac 1 2 \int_\Omega \Big(\| \nabla u\|^2 + u^2\Big) - \int_\Omega f u,
\end{equation*}
sur les espaces $H^1(\Omega)$ et $H^1_0(\Omega)$ respectivement pour les conditions au bord de Neumann et Dirichlet.

Une intégration par parties formelle (i.e.\ non justifiée a priori au regard de la régularité des termes), permet également de reformuler le problème de Neumann sous la forme de l'équation aux dérivées partielles (EDP) d'ordre deux 
%[Si $\Omega$ est un ouvert régulier, il s'agit de la solution faible du problème de Neumann
\begin{align*}
	-\Delta u + u &= f \text{ sur }\Omega, &
	\<\nabla u,n_{\partial \Omega}\> = 0 \text{ sur } \partial \Omega,
\end{align*}
où $n_{|\partial \Omega}$ désigne la normale extérieure. De même, le problème de Dirichlet est formellement équivalent à l'EDP
\begin{align*}
	-\Delta u + u &= f \text{ sur }\Omega, &
	u= 0 \text{ sur } \partial \Omega.
\end{align*}

\subsection{Injections de Sobolev - cas sous-critique}
\label{subsec:injSousCrit}

Le théorème d'injection de Sobolev sous-critique montre que les éléments de $W^{s,p}(\bR^d)$, où $s>0$ et $p\in [1,\infty[$, sont intégrables à une puissance $p^*>p$. En d'autres termes, abaisser l'exigence de régularité $s$ permet d'augmenter l'intégrabilité $p$. 

\begin{retenir}
Le théorème d'injection sous-critique entre espaces de Sobolev.

\noindent
(\Cref{th:injSousCrit}.)	
\end{retenir}

La preuve de l'injection de Sobolev sous-critique repose sur l'inégalité de convolution dite de Young faible, faisant intervenir une extension dite ``faible'' (weak) des espaces de Lebesgue  $L^p$. Nous admettons ce résultat, mais prouvons néanmoins l'inégalité de Young usuelle, plus élémentaire et faisant intervenir les espaces de Lebesgue standard, voir le \cref{th:Young}.

Soient $(X,\mu)$ un espace mesuré, $f : X \to \bR$ mesurable et $p \in [1,\infty[$. On définit 
\begin{equation*}
	\|f\|^p_{p,w} := \sup_{t > 0} \mu\{|f|\geq t\} t^p.
\end{equation*}
C'est une \emph{quasi-norme} sur l'espace $L^{p,w} := \{f \text{ mesurable}\mid \|f\|_{p,w} < \infty\}$, et on a $\|f\|_{p,w} \leq \|f\|_p$. 
Cet espace a l'avantage de contenir une fonction radiale puissance : 
\begin{equation}
\label{eq:powLpw}
	\frac 1 {|x|^\alpha} \in L^{d/\alpha,\, w} (\bR^d),
\end{equation}
pour tout $0< \alpha \leq d$. Par contraste, ces fonctions n'appartiennent à aucun espace $L^p (\bR^d)$.


\begin{theorem}[Inégalité de Young faible, admis]
\label{th:YoungFaible}
	Soient $p,q,r'\in ]1,\infty[$ tels que $1+ \frac 1 {r'} = \frac 1 p + \frac 1 q$. 
	On a 
	$\|f\star g \|_{r'} \leq C(p,q,d) \|f\|_p \|g\|_{q,w}$ pour tous $f \in L^p(\bR^d)$ et $g\in L^{q,w}(\bR^d)$.
\end{theorem}

La convolution dans l'inégalité de Young doit s'interpréter comme l'unique prolongement continu (pour les normes $\|\cdot\|_p$ et $\|\cdot\|_{r'}$) de la convolution usuelle $f\in \cD(\bR^d) \mapsto f\star g \in C^\infty(\bR^d)$, qui est bien définie car $g \in L^{q,w} \subset L^1_{\loc}$ puisque $q>1$.


Le second ingrédient de la preuve de l'injection de Sobolev sous-critique est la proposition suivante, exprimant toute fonction régulière et à support compact comme la convolution de son gradient avec un noyau vectoriel radial. Elle est liée au \nref{corol:noyauGreen}, donnant la solution fondamentale de l'opérateur Laplacien dans $\bR^d$.

Rappelons le théorème fondamental de l'analyse, qui assure que pour toute $u \in C^1(\bR^d)$ et tous $x,y \in \bR^d$ on a 
\begin{equation}
\label{eq:int_seg}
	u(y)-u(x) = \int_0^1 \<\nabla u((1-t) x+t y),(y-x)\> \diff t	.
\end{equation}
La sphère unité de $\bR^d$ est notée $\bS_{d-1} := \{x \in \bR^d \mid \|x\|=1\}$, et est munie de la mesure d'angle usuelle définie par $\mu(A) / d := \Leb\{ta \mid a\in A, t \in ]0,1]\}$ pour tout borélien $A \subset \bS_{d-1}$.

\begin{proposition}
\label{prop:convolPow} 
Pour tous $f \in C^1_c(\bR^d)$ et $x \in \bR^d$ on a 
\begin{equation*}
	f(x) = \frac 1 {|\bS_{d-1}|} \int_{\bR^d} \big\<\nabla f(x-y), \frac y {\|y\|^d}\big\> \diff y.
\end{equation*}
\end{proposition}

\begin{proof}
On peut supposer que $x=0$, sans restreindre la généralité, quitte à considérer la fonction translatée $f(x+\cdot)$. Notons $\bS := \bS_{d-1}$, et calculons
\begin{align*}
	-|\bS| f(0) 
	&= \int_\bS \big(f(r z) - f(0)\big) \diff z
	= \int_\bS \int_0^r \big\<\nabla f(s z), z\big\> \diff s \diff z 
	= \int_\bS \int_0^r \frac{\<\nabla f(s z), sz\>}{\|s z\|^d} s^{d-1}\diff s \diff z \\
	&= \int_{B(0,r)} \big\<\nabla f(y),\frac{y}{\|y\|^d}\big\> \diff y,
\end{align*}
où l'on a (i) choisi $r>0$ tel que $\supp(f) \subset B(0,r)$, (ii) utilisé \eqref{eq:int_seg}, (iii) utilisé que $\|z\|=1$, (iv) fait un changement de variables en coordonnées polaires. Le résultat s'ensuit.
\end{proof}


Avant d'attaquer la preuve du théorème d'injection sous-critique de Sobolev, nous discutons de l'homogénéité de l'estimation établie, c'est-à-dire de son invariance aux changements d'échelle. 
\begin{remark}[Homogénéité]
\label{rem:homogSobolev}
Soit $u \in \cD(\bR^d)$, $p \in [1,\infty]$,  $\alpha \in \bN^d$ et $s := |\alpha|_1$. Posons 
$u_\lambda(x) := u(\lambda x)$ pour tout $\lambda\in ]0,\infty[$. Alors
\begin{align}
\label{eq:homog_lp}
	\|u_\lambda\|_{p} &= \lambda^{-d/p} \|u\|_p, &
	\|\partial_\alpha u_\lambda\|_p &= \lambda^{s-d/p}\|\partial_\alpha u\|_p,
\end{align}
	par changement de variable dans l'intégrale et comme $\partial_\alpha u_\lambda (x) = \lambda^s \partial_\alpha u(\lambda x)$. 
	En insérant $u_\lambda$ dans l'estimation \eqref{eq:inj_ss_critique} ci-dessous, et en considérant les limites $\lambda\to 0^+$ et $\lambda\to \infty$, on voit que la relation imposée entre $s$, $d$, $p$ et $p^*$ est la seule envisageable.
\end{remark}

Dans le résultat suivant la notation ``$C(s,p,d)$'' désigne ``une constante dépendant seulement des paramètres $s,p,d$'', et de la même manière dans la suite, avec d'autres paramètres.

\begin{theorem}[Injection sous critique $W^{s,p}(\bR^d) \subset L^{p^*}(\bR^d)$, $s p < d$]
\label{th:injSousCrit}
	Soient $d\in \bN^*$, $s\in \bN$, $p \in [1,\infty[$. Supposons $s p<d$ (hypothèse de sous-criticité) et définissons $p^* \in [p,\infty[$ par 
	\begin{equation*}	
	\frac 1 {p^*} = \frac 1 p - \frac s d.
	\end{equation*}
	Alors $W^{s,p}(\bR^d) \subset L^{p^*}(\bR^d)$ et pour tout $u \in W^{s,p}(\bR^d)$ on a 
	\begin{equation}
	\label{eq:inj_ss_critique}
		\|u\|_{L^{p^*}} \leq C(p,s,d) \sum_{|\alpha| = s} \|\partial_\alpha u\|_{L^p}.
	\end{equation}
\end{theorem}

\begin{proof}
On prouve seulement le cas $p>1$. Noter que le cas $p=1$, dû à Gagliardo-Nirenberg, est plus difficile en général, sauf si $d=2$ auquel cas on dispose d'une preuve élémentaire \cite{desvilettes2015analyse}.

Par une récurrence immédiate, il suffit de considérer le cas où $s=1$. 
Définissons $h(x) := \frac 1 {|x|^{d-1}}$ sur $\bR^d$, et notons que $h \in L^{q,w}(\bR^d)$ avec $q := \frac d {d-1}$, comme observé dans \eqref{eq:powLpw}. 
Par la \cref{prop:convolPow} et le \nref{th:YoungFaible}, on a pour toute $u \in \cD(\bR^d)$
\begin{align*}
	\|u\|_{r'} &\leq C'(p,q,d) \|\nabla u\|_p \| h \|_{q,w},  & \text{avec } \frac 1 {r'} &= \frac 1 p + \frac 1 q -1 = \frac 1 p -\frac 1 d = \frac 1 {p^*}.
\end{align*}
L'estimation annoncée s'ensuit par densité de $\cD(\bR^d)$ dans $W^{1,p}(\bR^d)$, voir la \cref{prop:testDenseWsp}.
\end{proof}


\begin{corollary}
\label{corol:injSousCritOm}
	Soient $\Omega\subset \bR^d$ un ouvert régulier, $s\in \bN$ et $p \in [1,\infty[$. Supposons $s p <d$ (hypothèse de sous-criticité) et définissons $p^* \in [p,\infty[$ par $\frac 1 {p^*} = \frac 1 p - \frac s d$.
	Alors $W^{s,p}(\Omega) \subset L^{p^*}(\Omega)$ et pour tout $u \in W^{s,p}(\Omega)$ on a 
	\begin{equation*}
		\|u\|_{L^{p^*}(\Omega)} \leq C'(p,s,\Omega) \|u\|_{W^{s,p}(\Omega)} %\sum_{|\alpha| = s} \|\partial_\alpha u\|_{L^p}.
	\end{equation*}
\end{corollary}

\begin{proof}
	C'est une conséquence du \cref{th:prolSobolev} de prolongement, et du cas de $\bR^d$ établi dans le \cref{th:injSousCrit}.
\end{proof}

Comme annoncé, nous prouvons l'inégalité de Young usuelle, dont la variante faible \cref{th:YoungFaible} est au coeur de la preuve du \cref{th:injSousCrit}.

\begin{theorem}[Inégalité de Young]
\label{th:Young}
	Soient $p,q,r \in [1,\infty]$ tels que $2 = \frac 1 p + \frac 1 q +\frac 1 r$. 
	Alors 
	\begin{align*}
		I(f,g,h) &\leq \|f\|_p \|g\|_q \|h\|_r, &
		\text{où }
		I(f,g,h) &:= \iint f(x) g(y-x) h(y) \, \diff y \, \diff x,
	\end{align*}
	pour toutes $f \in L^p(\bR^d), g \in L^q(\bR^d), h \in L^r(\bR^d)$. On en déduit l'estimation de convolution
	%Ainsi $I : L^p \times L^q \times L^r$ est une forme trilinéaire continue. On en déduit que  et on a la forme bilinéaire continue
	\begin{align*}
		\|f\star g\|_{r'} &\leq \|f\|_p \|g\|_q, &
		\forall f \in L^p(\bR^d),\, \forall g \in L^q(\bR^d),
		\quad \text{avec }
		1+\frac 1 {r'} = \frac 1 p + \frac 1 q.
	\end{align*}
	Ainsi $I : L^p \times L^q \times L^r\to \bR$ et $\star : L^p\times L^q \to L^{r'}$ sont tri- et bi-linéaires continues respectivement. 
\end{theorem}

\begin{proof}
	Introduisons les exposants conjugués $p',q',r' \in [1,\infty]$, définis par $\frac 1 p + \frac 1 {p'}=1$, $\frac 1 q + \frac 1 {q'}=1$, $\frac 1 r + \frac 1 {r'}=1$.
	Pour toutes fonctions mesurables positives $f,g,h$ sur $\bR^d$ on a 
	\begin{align*}
	I(f,g,h) &= 
	\iint (f(x)^p g(y-x)^q)^{\frac 1 {r'}}  (f(x)^p h(y)^r)^{\frac 1 {q'}} (g(y-x)^q h(y)^r)^{\frac 1 {p'}} \diff x \diff y \\
	&\leq \Big(\iint f(x)^p g(y-x)^q\diff x \diff y\Big)^{\frac 1 {r'}} 
	\Big(\iint  (f(x)^p h(y)^r)^{\frac r {q'}} (g(y-x)^q h(y)^r)^{\frac r {p'}} \Big)^\frac 1 r
	\\
	&\leq \Big(\iint f(x)^p g(y-x)^q\diff x \diff y\Big)^{\frac 1 {r'}}
	\Big(\iint f(x)^p h(y)^r\diff x \diff y\Big)^{\frac 1 {q'}}
	\Big(\iint g(y-x)^q h(y)^r\diff x \diff y\Big)^{\frac 1 {p'}}
	\\
	&= \Big(\int f(x)^p \diff x\Big)^{\frac 1 {q'}+\frac 1 {r'}} 
	\Big(\int g(x)^q \diff x\Big)^{\frac 1 {r'}+\frac 1 {p'}} 
	\Big(\int h(x)^r \diff x\Big)^{\frac 1 {p'}+\frac 1 {q'}} = \|f\|_p \|g\|_q \|h\|_r.
	\end{align*}
	On a utilisé (i) les identités $1 = p(2-\frac 1 q-\frac 1 r) = p (\frac 1 {q'} + \frac 1 {r'})$ et de même par symétrie $1= q (\frac 1 {r'} + \frac 1 {p'})=r (\frac 1 {p'} + \frac 1 {q'})$,
	(ii) l'inégalité de Hölder avec les exposants conjugués $r$ et $r'$, (iii) l'inégalité de Hölder avec les exposants $\frac {q'} r$ et $\frac {p'} r$ qui sont bien conjugués car $\frac r {q'} + \frac r {p'} = r (2-\frac 1 p -\frac 1 q) = 1$, (iv) le théorème de Fubini.
	La bonne définition et la continuité de la forme trilinéaire $I : L^p \times L^q \times L^r \to \bR$ s'ensuit. 
	
	Étant données $f,g,h\in C^0_c$, on a $I(f,g,h) = \int (f\star g)(x) h(x) \diff x$ donc $\|f\star g\|_{r'} \leq \|f\|_p \|g\|_q$, en choisissant $h$ correspondant au cas d'égalité dans le
	\nref{th:HolderIneq}. Si $p,q\in [1,\infty[$, alors par densité de $C^0_c$ dans $L^p$ et $L^q$ (\cref{corol:testDenseLp}) la convolution admet une unique extension continue $L^p \times L^q \to L^{r'}$, comme annoncé. Finalement, dans le cas où $p=\infty$, on a forcément $q=1$ et $r'=\infty$, donc le résultat est trivial par intégrabilité et d'ailleurs déjà vu dans la \nref{prop:convolLp}, ce qui conclut.
\end{proof}


Le noyau de Green, solution fondamentale de l'opérateur laplacien sur $\bR^d$, permet d'exprimer une fonction régulière et à support compact comme une convolution de son Laplacien, de manière analogue à la \cref{prop:convolPow} qui l'exprimait comme convolution de son gradient. 

Dans le résultat suivant, la notation $\delta_0$ désigne la masse de Dirac à l'origine, qui est une distribution d'ordre $0$ définie par $\<\delta_0,\vp\> = \vp(0)$ pour toute $\vp \in \cD(\bR^d)$. La convolution doit être comprise au sens de la \nref{def:convolDistrib}.

\begin{corollary}[Noyau de Green]
\label{corol:noyauGreen}
La fonction $G : \bR^d \to \bR$ définie par
\begin{equation*}
	G(x) := c(d)
	\begin{cases}
	-|x| &\text{si } d=1,\\
	- \ln |x| &\text{si } d=2,\\
	|x|^{2-d} & \text{si } d\geq 3.
	\end{cases}
\end{equation*}
satisfait $-\Delta G = \delta_0$ au sens des distributions sur $\bR^d$, où $c(d) := 1/(\max \{1,d-2\} |\bS_{d-1}|)$. 
En particulier, $\vp = -\Delta (G \star \vp)$ pour toute $\vp \in \cD(\bR^d)$. 
\end{corollary}

\begin{proof}
\def\bn{\boldsymbol{n}}
Notons $V(x) := |\bS_{d-1}|^{-1} x/\|x\|^d$, le champ de vecteurs intervenant dans la \cref{prop:convolPow}.
Observons que $G \in C^\infty(\bR^d \sm \{0\})$, et que $\nabla G = - V$ sur $\bR^d \sm \{0\}$.
	%Tout d'abord notons que $\nabla g = -\alpha(d) V \in C^\infty(\bR^d\sm \{0\})$, avec $\alpha(d) = d-2$ si $d\geq 3$ et $\alpha(d)=1$ si $d\in\{1,2\}$. 
	La première étape de la preuve consiste à montrer que l'égalité $\nabla G = - V$ est également vraie en tant que distributions sur $\bR^d$. Étant donnée une fonction test (vectorielle) $\phi \in \cD(\bR^d,\bR^d)$ on calcule
	\begin{equation*}
		- \<\nabla G,\phi\>_{\cD' \times \cD} 
		= \<G,\diver \phi\>_{\cD' \times \cD} 
		= \int_{\bR^d} G \diver \phi 
		= \lim_{\delta \to 0^+} \int_{B(R) \sm B(\delta)} G \diver \phi,
	\end{equation*}
	où $R>0$ est tel que $\supp (\phi) \subset B(R)$. On a utilisé (i) la \nref{prop:deriv_distrib}, (ii) la définition de la distribution associée à une fonction $G \in L^1_\loc(\bR^d)$, voir la \cref{prop:distribL1Loc}, (iii) le fait que $\phi$ est à support compact, et l'intégrabilité de $G$ sur un voisinage de $0$. Par ailleurs
	\begin{equation*}
		\int_{B(R) \sm B(\delta)} G \diver \phi = \int_{B(R) \sm B(\delta)} \diver(G \phi) - \<\nabla G,\phi\> = \int_{\partial B(R)} G\<\phi,\bn\> - \int_{\partial B(\delta)} G\<\phi,\bn\> - \int_{B(R) \sm B(\delta)} \<\nabla G,\phi\>,
	\end{equation*}
	pour tout $\delta>0$, et avec $\bn(x):=x/|x|$ le vecteur unitaire radial, en utilisant (i) la formule de Leibniz de dérivation d'un produit, et (ii) la formule de Stokes d'intégration d'une divergence sur un domaine lisse borné.
	Dans le membre de droite, l'intégrale sur $\partial B(R)$ est nulle par choix de $R$, celle sur $\partial B(\delta)$ est de l'ordre de $\cO(\delta)$ (ou de $\cO(\delta |\ln \delta|)$ si $d=2$), et celle sur $B(R) \sm B(\delta)$ converge vers $-\int_{\bR^d} \<\nabla G,\phi\>$ lorsque $\delta\to 0$ par intégrabilité de cette fonction au voisinage de $0$.
	
	Nous avons établi que $\<\nabla G,\phi\>_{\cD' \times \cD} = \int_{\bR^d} \<\nabla G,\phi\>$, et donc $\nabla G = - V$ au sens des distributions. Par ailleurs $\diver( V ) = \delta_0$ au sens des distributions par la \cref{prop:convolPow}, donc $-\Delta G = -\diver(\nabla G) = \delta_0$ par composition, comme annoncé.
	
	Finalement, pour toute $\vp \in \cD(\bR^d)$ on a $-\Delta(G \star \vp) = (-\Delta G) \star \vp =  \delta_0 \star \vp = \vp$ par la \cref{prop:convolDistrib} comme annoncé, ce qui conclut.
%	\todoi{Pas forcément nécessaire de passer par la fonction test vectorielle. On peut travailler directement sur $\<\Delta g, \phi\>_{\cD'\times \cD}$. 	}
\end{proof}



\subsection{Injections de Sobolev - cas sur-critique}
\label{subsec:injSurCrit}

Le Théorème d'injection sur-critique de Sobolev complète le précédent, et montre que les éléments de $W^{s,p}(\bR^d)$ où $s\in \bN$ et $p \in [1,\infty[$ satisfont $s p>d$, sont des fonctions continues.
Ces deux résultats laissent ouverte la question des injections dites \emph{critiques}, pour lesquelles  $s p = d$, qui doivent être traitées au cas par cas et pour lesquelles on se réfèrera à la littérature. 

\begin{retenir}
	Le théorème d'injection sur-critique entre espaces de Sobolev.
	(\Cref{th:injSurCrit}.)
\end{retenir}

La preuve est fondée sur une variante de l'identité de convolution établie à la \cref{prop:convolPow}.
\begin{proposition}
\label{prop:convolPow2}
Pour toute $f \in C^1(\bR^d)$, et tous $x \in \bR^d$ et $r>0$, on a 
\begin{equation*}
	\int_{B(x,r)} |f(x) -  f(y)| \diff y \leq \frac{r^d}{d}\int_{B(0,r)} \|\nabla f(x-y)\| \frac{\diff y}{\|y\|^{d-1}}.
\end{equation*}
\end{proposition}

\begin{proof}
\todoi{
On peut pas mal simplifier la preuve si on laisse tomber le facteur $1/d$, qui est inutile.

On pourrait aussi faire rentrer ici l'estimation Hölder, comme conséquence naturelle.
}

	%Les preuves des deux points sont semblables. On se concentre sur le second, et sans nuire à la généralité on peut supposer $x=0$. 
	On suppose $x=0$, sans nuire à la généralité, quitte à considérer la fonction translatée $f(x+\cdot)$. On note $\bS = \bS_{d-1}$ la sphère unité de $\bR^d$. Soit $z \in \bS$, et posons $g(s) := f(s z)$ pour tout $s\in \bR_+$, de sorte que $|g'(s)| = |\<\nabla f(s z),z\>| \leq \|\nabla f(s z)\|$ par dérivation composée et par l'inégalité de Cauchy-Schwarz. 
	Étant donné $r>0$, on calcule
\begin{align*}
	&\int_0^r |f(s z)-f(0)| s^{d-1} \diff s
	= \int_0^r |g(s)-g(0)| s^{d-1} \diff s 
	\leq \int_0^r \int_0^1 |g'(st)| \diff t \, s^d \diff s \\
	&= \int_{s=0}^r \int_{\alpha=0}^r |g'(\alpha)| \bbone_{\alpha \leq s} \diff \alpha \, s^{d-1} \diff s
	=\int_{\alpha=0}^r |g'(\alpha)| \frac{r^d-\alpha^d}{d} \diff \alpha \leq \frac{r^d}{d} \int_{s=0}^r |g'(s)| \diff s.
\end{align*}
On a utilisé (i) la définition de $g$, (ii) l'inégalité des accroissements finis, (iii) le remplacement de la variable $t \in [0,1]$ par $\alpha := s t \in [0,s]$, (iv) le théorème de Fubini et l'intégration explicite de $\int_\alpha^r s^{d-1} \diff s = (r^d-\alpha^d)/d$, (v) la minoration $\alpha^d \geq 0$. 
Calculons maintenant % \belowdisplayskip=-12pt
\begin{align*}
	&\int_{B(0,r)} |f(y)-f(0)| \diff y = \int_{\bS} \int_0^r |f(s z)-f(0)| s^{d-1} \diff s \, \diff z
	\leq \frac{r^d}{d} \int_\bS \int_0^r\|\nabla f(s z)\| \diff s \, \diff z\\
	&= \frac{r^d}{d} \int_\bS \int_0^r\frac{\|\nabla f(s z)\|}{\|s z\|^{d-1}} s^{d-1}\diff s \, \diff z
	= \frac{r^d}{d} \int_{B(0,r)} \| \nabla f(y)\| \frac{\diff y} {\|y\|^{d-1}}. 
\end{align*}
On a utilisé (i) un changement de variables en coordonnées polaires, (ii) l'inégalité établie précédemment, (iii) l'égalité $\|z\|=1$ et (iv) un nouveau changement de variables en coordonnées polaires. Ceci conclut la preuve.
\end{proof}

Pour exploiter la proposition précédente, nous aurons besoin du calcul suivant de l'intégrale de la fonction radiale puissance $h : x \in \bR^d \mapsto 1/|x|^{d-1}$. Si $0 \leq q < d/(d-1)$, alors elle satisfait
%Avant de commencer, remarquons l'intégrale des fonctions radiale puissance. faisons la remarque que si $1 \leq q < d/(d-1)$ alors $h : x \in \bR^d \mapsto 1/|x|^{d-1}$ satisfait 
\begin{equation}
\label{eq:normLqPow}
	\|h\|^q_{L^q(B(0,r))} = |\bS| \int_0^r s^{-q(d-1)} s^{d-1} \diff s = |\bS| \frac{r^{(d-1)(1-q)+1}}{(d-1)(1-q)+1}
\end{equation}
pour tout $r>0$, où $\bS := \bS_{d-1}$ désigne la sphère unité de $\bR^d$.
On désigne par $C^\alpha$ l'espace des fonctions de régularité $\alpha$-Hölder, où $\alpha\in ]0,1[$, et on note $C^\alpha_0$ le sous espace de celles qui tendent vers $0$ à l'infini. 

\begin{theorem}[Injection sur-critique $W^{1,p} \subset C^{1-d/p}_0$, $p>d$, inégalité de Morrey]
\label{th:injSurCrit}
Soit $d \in \bN^*$ la dimension, et $p \in ]d,\infty[$ un exposant. Toute fonction $u\in W^{1,p}(\bR^d)$ admet un représentant continu satisfaisant
\begin{align}
\label{eq:inj_sur_critique}
	\|u\|_\infty &\leq C(p,d) \|u\|_{W^{1,p}},&
	|u(x)-u(y)| \leq C'(p,d) |x-y|^\alpha \| \nabla u\|_p, 
\end{align}
 pour tous $x,y \in \bR^d$, avec $\alpha := 1-d/p$. Ainsi, on a l'injection continue $W^{1,p}(\bR^d) \subset C_b^\alpha(\bR^d)$. De plus, pour toute $u \in W^{1,p}(\bR^d)$, on a $u(x) \to 0$ lorsque $|x| \to \infty$.
\end{theorem}

\begin{proof}
On montre ci-dessous les inégalités annoncées sous l'hypothèse supplémentaire $u \in \cD(\bR^d)$. 
Il s'ensuit que l'application linéaire $(\cD(\bR^d), \|\cdot\|_{W^{1,p}}) \to (C^\alpha_0, \|\cdot\|_{C^\alpha})$ est continue, donc admet un prolongement unique à $W^{1,p}(\bR^d)$ par densité de $\cD(\bR^d)$, voir la \cref{prop:testDenseWsp}.
On peut alors passer à la limite les inégalités (\ref{eq:inj_sur_critique}, gauche) puis (\ref{eq:inj_sur_critique}, droite) en notant qu'elles ne font intervenir que des quantités continues vis à vis des normes $\|\cdot\|_{W^{1,p}}$ et $\|\cdot\|_{C^\alpha}$, ce qui conclut.

\emph{Preuve de (\ref{eq:inj_sur_critique}, droite)}. Soit $r := |x-y|$, soit $B$ la boule de rayon $r/2$ centrée en $(x+y)/2$, et soit $u_B$ la valeur moyenne de $u$ sur $B$. Comme $B \subset B(x,r)$ on a avec $h : x \in \bR^d \mapsto |x|^{-(d-1)}$
\begin{equation*}
	|B| |u(x)-u_B| \leq \int_{B} |u(x)-u(z)| \diff z \leq  \int_{B(x,r)} |u(x)-u(z)| \diff z \leq \frac{r^d}{d} \|\nabla u\|_{L^p(B(x,r))} \| h\|_{L^q(B(0,r))},
\end{equation*} 
la \cref{prop:convolPow2} et le \nref{th:HolderIneq}, en notant $q$ l'exposant conjugué de $p$, i.e.\ $\frac 1 p + \frac 1 q =1$. 
En utilisant \eqref{eq:normLqPow} et en notant que $[(d-1)(1-q)+1]/q = 1-d/p>0$, on conclut la preuve de (\ref{eq:inj_sur_critique}, droite) 
\begin{equation*}
	|u(x)-u(y)| \leq |u(x)-u_B|+|u(y)-u_B| \leq C r^{1-d/p} \| \nabla u\|_{L^p}, 
\end{equation*}
où $C = C(d,p)$. 
Pour établir (\ref{eq:inj_sur_critique}, gauche) on note $\tilde B := B(x,1)$ la boule unité centrée en $x$, puis
\begin{equation}
\label{eq:surCritBd}
	|u(x)| 
	\leq |u(x)-u_{\tilde B}| + |u_{\tilde B}| 
	\leq C \|\nabla u\|_{L^p(\tilde B)} + \frac 1 {|\tilde B|} \|u\|_{L^1(\tilde B)} 
	\leq C \|\nabla u\|_{L^p(\tilde B)} + C'\|u\|_{L^p(\tilde B)},
\end{equation}
où l'on a utilisé (i) la moyenne $u_{\tilde B}$, (ii) l'argument précédent pour contrôler l'écart à la moyenne, (iii) \cref{corol:lpInclusLq} pour comparer les normes $L^1$ et $L^p$ sur l'espace de mesure finie $\tilde B$, où $C = C(d,p)$ et $C'=C'(d,p)$. 

Finalement, reste à établir la convergence vers $0$ à l'infini de toute $u \in W^{1,p}(\bR^d)$, $p>d$. 
Étant donné $\ve>0$, soit $R$ tel que $\|\nabla u\|_{L^p(\bR^d \sm B(0,R))} \leq \ve $ et $\| u\|_{L^p(\bR^d \sm B(0,R))} \leq \ve$. Alors $|u(x)| \leq (C+C') \ve$ dès que $|x| \geq R+1$ par \eqref{eq:surCritBd}, d'où la convergence vers $0$, ce qui conclut. 
\end{proof}

Pour compléter le résultat précédent, mentionnons que $W^{1,\infty}$ est l'espace des fonctions bornées et Lipschitz. (L'énoncé précédent ne s'applique pas au cas $p=\infty$ : en général ces fonctions ne sont pas $C^1$, et ne tendent pas vers $0$ à l'infini.)

\begin{corollary}
\label{corol:ingSurCritOm}
Soit $\Omega \subset \bR^d$ un ouvert régulier, et $p\in ]d,\infty[$. Alors $W^{1,p}(\Omega) \subset C^\alpha_b(\Omega)$ avec $\alpha := 1-d/p$.
Toute $u\in W^{1,p}(\Omega)$ admet un représentant continu satisfaisant
\begin{align*}
	\|u\|_\infty &\leq C(p,\Omega) \|u\|_{W^{1,p}},&
	|u(x)-u(y)| \leq C'(p,\Omega) |x-y|^\alpha \|u\|_{W^{1,p}}.
\end{align*}
\end{corollary}

\begin{proof}
	C'est une conséquence du \cref{th:prolSobolev} de prolongement, et du cas de $\bR^d$ entier établi dans le \cref{th:injSurCrit}.
\end{proof}


%\begin{theorem}
%\label{th:injSousCrit}
%	Soient $d\in \bN^*$, $s\in \bN$, $p \in [1,\infty[$. Supposons $p > d/s$ (hypothèse de sur-criticité).
%	Alors $W^{s,p}(\Omega) \subset C^0_b(\Omega)$ et pour tout $u \in W^{s,p}(\bR^d)$ on a 
%	\begin{equation*}
%		\|u\|_\infty \leq C(p,s,d) \sum_{|\alpha| = s} \|u\|_{W^{s,p}}.
%	\end{equation*}
%\end{theorem}


On peut composer les cas sous-critique et sur-critique des injections de Sobolev. 
La preuve ci-dessous évite le cas $p=1$, plus technique, et ne donne pas l'exposant de Hölder des fonctions continues obtenues; se référer aux inégalités de Sobolev générales pour plus de précisions.

\begin{remark}
\label{rem:SobSurCritIter}
Soit $\Omega \subset \bR^d$ un ouvert régulier, et soient $s \in \bN^*$ et $p \in ]1,\infty]$ tels que $s p  > d$. Alors $W^{s,p}(\Omega) \subset C^0(\overline \Omega)$. 
\end{remark}

\begin{proof}
	Par hypothèse, il existe $0 \leq s_0 < s$ tel que $p s_0 \leq d$ et $p(s_0+1) > d$. 
	Pour commencer, supposons que $p<\infty$ et $ps_0 < d$, et considérons $u \in W^{s,p}(\Omega) \subset W^{s_0+1,p}(\Omega)$. Alors $u \in W^{1,q}(\Omega)$ par le \cref{th:injSousCrit} (injection sous-critique) avec $\frac 1 p = \frac 1 q + \frac{s_0} d$. 
	L'hypothèse $p(s_0+1) > d$ est équivalente à $q>d$, donc $u$ est continue par le \cref{th:injSurCrit} (injection sur-critique), comme annoncé. 
	
	Reste à traiter le cas $p=\infty$, qui est trivial puisque $W^{1,p}$ est constitué de fonctions Lipschitz, ainsi que le cas critique où $p s_0 = d$. Pour ce dernier, on note que $u \vp_{|\Omega} \in W^{s,p'}(\Omega)$ pour toute $\vp \in \cD(\bR^d)$ et tout $p'\in [1,p[$ (par dérivation d'un produit, et par le \cref{corol:lpInclusLq} qui s'applique puisque $\supp(\vp)$ est de mesure finie, et en se rappelant que $p>1$). En choisissant $p'$ assez proche de $p$, on obtient $p's_0 < d < p'(s_0+1)$, donc $u \vp_{|\Omega}$ est continue par l'argument précédent, et donc $u$ est continue puisque $\vp$ est arbitraire.
\end{proof}


\subsection{Injection compacte de Sobolev - théorème de Rellich}
\label{subsec:injCompacte}

Le théorème de Rellich énonce que certaines injections de Sobolev définissent des opérateurs compacts pour les topologies associées. Cela permet, entre autres applications, de montrer l'existence de minimiseurs à des fonctionnelles très non-linéaires, voir l'\cref{ex:nonLinMinH1}.
La preuve est fondée sur le critère de compacité de Fréchet-Kolmogorov, et fait intervenir deux résultats préliminaires. %Seul le cas $p=1$ du résultat suivant est utilisé dans la preuve. 

\begin{retenir}
Le Théorème de Rellich et son corollaire.
(\Cref{th:Rellich,corol:Rellich}.)
\end{retenir}


\begin{proposition}
\label{prop:trans_grad_lp}
Soient $\Omega\subset \bR^d$ un ouvert, $p\in [1,\infty[$, $u \in W^{1,p}(\Omega)$, et $h\in \bR^d$.
Soit $\omega \csubset \Omega$ compact tel que $\omega_h := \omega+B'(0,|h|) \subset \Omega$. Alors 
%$\|\tau_h u - u\|_{L^p(\omega)} \leq |h| \|\nabla u\|_{L^p(\Omega)}$.
	\begin{equation}
	\label{eq:uhdup}
		\int_\omega |u(x+h)-u(x)|^p \diff x \leq |h|^p \int_\Omega |\nabla u(x)|^p \diff x.
	\end{equation}
\end{proposition}

\begin{proof}
	Pour toute $u \in C^1(\bR^d)$ et tout $x \in \bR^d$
	\begin{equation*}
		|u(x+h)-u(x)|^p = \Big|\int_0^1 \<\nabla u(x+t h),h\> \diff t\Big|^p \leq |h|^p \int_0^1 \|\nabla u(x+th)\|^p \diff t,
	\end{equation*}
	par les inégalités (i) des accroissements finis, et (ii) de Cauchy-Schwarz et de Jensen.
		Ainsi 
	\begin{equation*}
		\int_\omega |u(x+h)-u(x)|^p \diff x \leq |h|^p \int_0^1 \int_\omega \|\nabla u(x+t h) \|^p \diff x \diff t \leq |h|^p \int_{\omega_h} \|\nabla u(x)\|^p \diff x,
	\end{equation*}
	par (i) intégration sur $\omega$ du résultat précédent, et interversion de Fubini, puis (ii) par inclusion du domaine translaté $\omega + t h \subset \omega_h$. On conclut par densité de $C^1_c(\Omega)$ dans $W^{1,p}(\omega_h)$, puisque $\omega_h \csubset \Omega$. 
\end{proof}

\begin{proposition}[Inégalité d'interpolation $L^p$]
\label{prop:interpLp}
	Soit $(X,\mu)$ un espace mesuré, et $f: X \to \bR$ mesurable, et $p_0,p_1\in [1,\infty]$ des exposants. Alors pour tout $\theta \in ]0,1[$ on a
	\begin{align*}
		\|f\|_{p_\theta} &\leq \|f\|_{p_0}^{1-\theta} \|f\|_{p_1}^\theta, &
		\text{où } \frac 1 {p_\theta} := \frac {1-\theta} {p_0} + \frac \theta {p_1}.
	\end{align*}
\end{proposition}

\begin{proof}
	Si $p_0<\infty$ et $p_1 = \infty$, le résultat revient à la majoration triviale $\int_X f^{p_\theta} \leq \int_X f^{p_0} \|f\|_\infty^{p_\theta-p_0}$. On peut donc supposer dans la suite que $p_0,p_1 \in [1,\infty[$.
	Posons $g := |f|^{p_\theta}$, qui est positive. Alors 
	\begin{equation*}
		\|f\|_{p_\theta}^{p_\theta} = \int_X f^{p_\theta} = \int_X g^{1-\theta} g^\theta  \leq \|g^{1-\theta} \|_r \| g^\theta \|_{r'} = \|f\|_{p_0}^{(1-\theta) p_\theta} \|f\|_{p_1}^{\theta p_\theta},
	\end{equation*}
	d'après l'inégalité de Hölder, 
	où l'on a noté $r := p_0/{((1-\theta) p_\theta)}$ et $r':=p_1/(\theta p_\theta)$, qui sont par hypothèse des exposants conjugués i.e.\ $\frac 1 r + \frac 1 {r'} = 1$. Le résultat est prouvé.
\end{proof}

\begin{theorem}[Rellich]
\label{th:Rellich}
	Soit $\Omega\subset \bR^d$ un ouvert borné régulier, et $p \in [1,\infty]$ un exposant. La boule unité fermée de $W^{1,p}(\Omega)$ est : 
	\begin{itemize}
		\item (cas sous-critique, $p<d$) précompacte dans $L^q(\Omega)$ fort, pour tout $q \in [1,p^*[$ où $\frac 1 {p^*} = \frac 1 p - \frac 1 d$.
		\item (cas critique, $p=d$) précompacte dans $L^q(\Omega)$ fort pour tout $q \in [1,\infty[$. 
		\item (cas sur-critique, $p>d$) précompacte dans $C(\overline \Omega)$. 
	\end{itemize}
\end{theorem}

\begin{proof}
	Le cas sur-critique découle du théorème d'Ascoli, grâce à l'estimation de régularité Hölderienne du \cref{th:injSurCrit}. Le cas critique découle du cas sous-critique et de l'injection continue $W^{1,p}(\Omega) \subset W^{1,q}(\Omega)$ lorsque $p>q$ et comme $\Omega$ est un domaine borné.
	Reste le cas sous-critique, pour lequel nous vérifions les critères du \nref{th:FrechetKolmogorov}.
	
	(Masse évanescente au bord) Soit $K \csubset \Omega$ compact. Comme $q < p^*$ on a 
	\begin{equation*}
		|\Omega\sm K|^{-\frac 1 q} \|u\|_{L^q(\Omega\sm K)} 
		\leq |\Omega\sm K|^{-\frac 1 {p^*}} \|u\|_{L^{p^*}(\Omega\sm K)}
		\leq C |\Omega\sm K|^{-\frac 1 {p^*}} \|u\|_{W^{1,p}(\Omega)}
	\end{equation*}
	en utilisant (i) le \cref{corol:lpInclusLq}, et (ii) l'inclusion de Sobolev sous-critique \cref{corol:injSousCritOm} où $C = C(\Omega,p)$.
	Étant donné $\ve>0$, il suffit de choisir $K \csubset \Omega$ tel que $C|\Omega\sm K|^{\frac 1 q-\frac 1 {p^*}} < \ve$ pour satisfaire le critère requis.
	
	(Régularité sous translation) Soit $K \csubset \Omega$, et $h\in \bR^d$ tel que $K + B'(0,|h|) \subset \Omega$. Notons $\tau_h$ l'opérateur de translation, et soit $\alpha \in ]0,1[$ tel que $\frac 1 q = \alpha + \frac{1-\alpha}{p^*}$, puisque $q<p^*$. Alors 
	\begin{equation*}
		\|\tau_h u-u\|_{L^q(K)} 
		\leq \|\tau_h u-u\|^\alpha_{L^1(K)} \|\tau_h u-u\|^{1-\alpha}_{L^{p^*}(K)} 
		\leq C |h|^\alpha \| \nabla u\|^\alpha_{L^1(\Omega)} \|u\|^{1-\alpha}_{W^{1,p}(\Omega)}
	\end{equation*}
	en utilisant (i) l'inégalité d'interpolation $L^p$ \cref{prop:interpLp}, et (ii) la \cref{prop:trans_grad_lp} avec l'exposant $1$. 
	(Noter qu'on a possiblement $q>p$, d'où la nécessité de l'inégalité d'interpolation.)
	Le résultat s'ensuit.
\end{proof}

En combinant les résultats de compacité du \nref{th:Rellich} et du \nref{th:BanachAlaoglu}, on s'assure de la convergence forte d'une suite extraite de fonctions, et de la convergence faible de ses dérivées, dans les espaces adéquats, comme l'illustrent le corollaire et l'exemple suivants.

%, ce qui permet de montrer l'existence de solutions à des problèmes d'optimisation 

\begin{corollary}
\label{corol:Rellich}
	Soit $\Omega \subset \bR^d$ un ouvert borné régulier et $p \in ]1,\infty[$ un exposant.
	Pour toute suite $(u_n)_{n \in \bN}$ bornée dans $W^{1,p}(\Omega)$, il existe $u_* \in W^{1,p}(\Omega)$ et une extractrice $\vp$ tels que $\nabla u_{\vp(n)} \weakto \nabla u_* \in L^p(\Omega,\bR^d)$ faiblement, et de plus $u_{\vp(n)} \to u_*$ fortement dans $L^q(\Omega)$ ou $C(\overline \Omega)$ selon les conditions du \nref{th:Rellich}. 
\end{corollary}

\begin{proof} 
	L'espace $W^{1,p}(\Omega)$ est un Banach, séparable au vu de l'isométrie $u\in W^{1,p} \mapsto (u,\nabla u) \in (L^p)^{d+1}$ à valeurs dans un espace séparable, et réflexif par les inégalités de Clarkson (voir la \cref{prop:ClarksonIneq} et la note afférente). Par \cref{th:BanachAlaoglu}, on peut donc extraire une suite faiblement convergente $u_{\vp(n)} \weakto u_* \in W^{1,p}(\Omega)$. La convergence faible de $u_{\vp(n)}$ dans $W^{1,p}$ implique la convergence faible de $\nabla u_{\vp(n)}$ dans $L^p$. 
	Par le \nref{th:Rellich}, l'injection de $W^{1,p}$ dans $L^q$ ou $C(\overline \Omega)$ est un opérateur compact, et donc d'après \nref{prop:opCompactComplCont} on a la convergence forte $u_{\vp(n)} \to u_*$ dans la topologie associée.	
\end{proof}



\begin{exemple}[Minimisation d'une fonctionnelle non-linéaire sur $H^1$]
\label{ex:nonLinMinH1}
	Soit $\Omega\subset \bR^d$ un ouvert borné régulier, $F\in C^1_b(\bR)$ (i.e.\ $F$ et sa dérivée sont bornées), et $f \in L^2(\Omega)$. Alors il existe $u\in H^1(\Omega)$ minimisant 
	\begin{equation*}
		J(u) := \int_\Omega \Big(\frac 1 2 |\nabla u|^2 + \frac 1 2 u^2 - f u + F(u) \Big).
	\end{equation*}
\end{exemple}

\begin{proof}
	Posons $J_0(u) := \frac 1 2 \|u\|_{H^1(\Omega)}^2$, qui est convexe continue sur $H^1(\Omega)$, donc faiblement s.c.i.\ par le \cref{corol:cvxFermeBdComp}.
	Posons $J_1(u) := \int_\Omega f u$ et $J_2(u) := \int_\Omega F(u)$, et notons que 
	\begin{align*}
		|J_1(u)| &\leq \|f\|_2 \|u\|_2, &
		|J_2(u)-J_2(v)| &\leq \|F'\|_\infty \|u-v\|_1 \leq \|F'\|_\infty |\Omega|^\frac 1 2  \|u-v\|_2.
	\end{align*}
	Ainsi $J(u) \geq \frac 1 2 (\|\nabla u \|^2_2 + \|u\|^2_2) - C\|u\|_2 -C'$, pour des constantes $C,C'$ appropriées. Si $(u_n)_{n \in \bN}$ est une suite minimisante pour $J$, alors $\|u_n\|_2$ et $\|\nabla u_n\|_2$ sont bornées par l'estimation précédente, et donc il existe une extractrice telle que $\nabla u_{\vp(n)}$ converge faiblement dans $L^2$ et $u_{\vp(n)}$ converge fortement dans $L^2$, par le \cref{th:Rellich} de Rellich. Comme $J_0$ est faiblement s.c.i., et $J_1$ et $J_2$ sont (fortement) continues dans $L^2$, on peut passer à la limite et obtenir un minimiseur. 
	
	Alternativement, les conditions sont remplies pour appliquer le \cref{corol:minFctCvxRefl}. 
\end{proof}

Le théorème de Rellich soulève naturellement la question de la fermeture, et donc de la compacité, de la boule unité de $W^{1,p}$ dans $L^q$.  

\begin{remark}[Fermeture dans $L^q$ de la boule unité de $W^{1,p}$]
\label{rem:bouleW1pFerme}
Soit $\Omega$ un ouvert régulier, et $p \in ]1,\infty]$ et $q \in [1,\infty]$ des exposants tels qu'on ait une inclusion continue $W^{1,p}(\Omega) \subset L^q(\Omega)$. Alors la boule unité de $W^{1,p}(\Omega)$ est fermée dans $L^q(\Omega)$.  	
\end{remark}

\begin{proof}
	Pour $p \in ]1,\infty[$, l'espace $W^{1,p}(\Omega)$ est uniformément convexe par les inégalités de Clarkson -- la \cref{prop:ClarksonIneq} établit la première, s'appliquant lorsque $p\geq 2$, et mentionne en note la seconde, couvrant le cas $1<p<2$. Il est donc réflexif, et donc sa boule unité est faiblement compacte par le \cref{th:BanachAlaoglu} de Banach-Alaoglu. L'image directe de cette boule unité par l'inclusion de Sobolev dans $L^q$, qui est une application linéaire continue donc faiblement continue, est donc également faiblement compacte par la \nref{prop:compactImage}, donc faiblement fermée, donc fermée. 
	
	Dans le cas $p=\infty$, la boule unité de $W^{1,\infty}(\Omega)$ est constituée des fonctions $u : \Omega \to \bR$ telles que $|u(x)|\leq 1$ et $|u(x)-u(y)| \leq |x-y|$ pour tous $x,y\in \Omega$ tels que $[x,y] \subset \Omega$.
	Toute suite de telles fonctions admet une sous-suite uniformément convergente sur tout compact $K \subset \Omega$ par le \nref{th:AscoliLip}, ce qui permet de passer ces inégalités à la limite et donc d'établir la fermeture.
\end{proof}

En revanche, la boule unité de $W^{1,1}(]-1,1[)$ n'est pas fermée dans $L^p(]-1,1[)$, $p<\infty$. Par exemple, la suite d'approximations de la fonction de Heaviside définie par $u_n(x) := \max\{0,\min\{1,n x\}\}$ est bornée\footnote{
On vérifie que $\|u_n\|_{W^{1,1}} = \|u_n\|_{L^1} + \|u'_n\|_{L^1} = (1-1/(2n)) + 1$, donc $u_n/2 \in B'_{W^{1,1}}(0,1)$ pour tout $n \geq 1$.
} dans $W^{1,1}(]-1,1[)$, converge dans $L^p(]-1,1[)$ pour tout $p \in [1,\infty[$ (cas critique), mais sa limite a une discontinuité en $0$ donc n'appartient pas à $W^{1,1}$ (la fonction de Heaviside appartient en fait à l'espace $\mathrm{BV}(\Omega)$ des fonctions à variation bornée). 


%\section{Transformée de Fourier}
%\label{sec:Fourier}
%\subsection{Transformation de Fourier}

\todoi{Définition sur $L^1(\bR)$}

\todoi{Remarque : définition sur un groupe commutatif}

\subsection{Distributions tempérées}

\todoi{Rappel de la définition}

\todoi{Extension de Fourier par dualité}

\subsection{Théorème de Plancherel}

\todoi{Fourier définit une isométrie de $L^2(\bR)$.}

\bibliographystyle{alpha}
\bibliography{Include/Papers.bib}

\appendix

\section{Fonctions bosse $C^\infty$}
\label{sec:bosse}
La classe des fonctions $C^\infty$ est assez souple pour construire toutes sortes d'approximations, de la fonction de seuil de Heaviside, de la masse de Dirac, ou de diverses fonctions indicatrices, dans cette section. Ces outils interviennent dans la régularisation par convolution, voir la \cref{prop:testDenseWsp}, ou la preuve de certains résultats d'analyse comme la \nref{prop:distribCont}. La souplesse des fonctions $C^\infty$ est à contraster avec la rigidité des fonctions analytiques, voir la \nref{prop:StoneCtrEx}.


\begin{retenir}
L'existence et les propriétés recherchées d'un noyau de convolution. (\Cref{prop:cinfBosse}.)	
\end{retenir}

\inventaire{
\Nref{sec:bosse}\\
\Nref{lem:cinfExpInv}\\
\Nref{prop:cinfHeaviside}\\
\Nref{prop:cinfBosse}\\
\Nref{prop:cinfSupport}\\
\Nref{prop:cinfPartition}\\
\Nref{prop:cinfExhaust}
}


\begin{lemma}[Prolongement lisse de la fonction nulle]
\label{lem:cinfExpInv}
	La fonction $\vp : \bR\to \bR$ suivante est $C^\infty$: 
	\begin{equation}
		\vp(x) := \begin{cases}
			\exp(-1/x) & \text{si } x>0, \\
			0 & \text{sinon}.
		\end{cases}		
	\end{equation}
\end{lemma}

\begin{proof}
	Par construction, $\vp$ est $C^\infty$ sur $\bR^*$. Par récurrence, $\vp^{(n)}(x) = P_n(x)\exp(-1/x)/x^{2n}$ pour tout $n \in \bN$ et tout $x>0$, où $P_n$ un polynôme de degré au plus $n$. On en déduit que $\vp^{(n)}(x) \to 0$ lors que $x\to 0$, pour tout $n \in \bN$, et donc que $\vp\in C^\infty$. 
\claudereport{Typos : ``où $P_n$ un polynôme'' $\to$ ``où $P_n$ est un polynôme'' ; ``lors que'' $\to$ ``lorsque''.}
\end{proof}




\begin{proposition}[Approximation de Heaviside]
\label{prop:cinfHeaviside}
	Il existe $\vp \in C^\infty(\bR)$, croissante, positive, telle que $\vp=0$ sur $]-\infty,0]$ et $\vp=1$ sur $[1,\infty[$. 
\end{proposition}

\begin{proof}
	Soit $\vp_0$ la fonction du \cref{lem:cinfExpInv}.
	On peut choisir $\vp := \vp_1/\vp_1(1)$ où 
	\begin{equation*}
		\vp_1(t) := \int_0^t \vp_0(s) \vp_0(1-s) \diff s.
	\qedhere
	\end{equation*}
\end{proof}


\begin{proposition}[Noyau de convolution]	
\label{prop:cinfBosse}
	Il existe $\vp \in C^\infty(\bR^d)$, radiale, positive, d'intégrale unité, telle que $\vp= 1$ sur $B(0,1/3)$ et $\supp(\vp) \subset B'(0,2/3)$. En particulier $\supp(\vp) \csubset B(0,1)$.
\claudereport{Erreur : $\vp = 1$ sur $B(0,1/3)$ est incompatible avec l'intégrale unité (la preuve donne $\vp = c$ sur $B(0,1/3)$).}
	%de support contenu dans la boule unité.
\end{proposition}

\begin{proof}
	Notons $\vp_\mathrm{heav}$ la fonction de la \cref{prop:cinfHeaviside}, et observons que $\vp_1(t) := \vp_\mathrm{heav}(2-3t)$ satisfait $\vp_1=1$ sur $]-\infty,1/3]$ et $\vp_1=0$ sur $[2/3,\infty[$. 
	On pose $\vp(x) := c \vp_1(|x|)$ où la constante $c$ est choisie telle que $\int \vp = 1$, qui convient clairement. 
\end{proof}


Les résultats suivants visent à approcher des fonctions indicatrices de manière lisse. Rappelons que $C^\infty_c(\Omega)$ désigne l'ensemble des fonctions indéfiniment dérivables et à support compact sur un domaine ouvert $\Omega \subset \bR^d$.


\begin{proposition}[Approximation d'un support]	
\label{prop:cinfSupport}
	Soit $\Omega\subset \bR^d$ ouvert et $K \csubset \Omega$ compact. 
	Alors il existe $\vp\in C^\infty_c(\Omega,[0,1])$ telle que $\vp=1$ sur $K$. 
\end{proposition}

\begin{proof}
Soit $r := \min\{\|x-y\| \mid x \in K, y \in \bR^d \sm \Omega\}$. Notons que $r>0$ par compacité de $K$ et fermeture de $\bR^d \sm \Omega$. 
L'ensemble $K$ étant compact, il admet une couverture $\cup_{i=1}^I B(x_i,r/3)$ par des boules ouvertes de rayon $r/3$ centrées en des points $x_1,\cdots,x_I \in K$. 

Notons $\vp_\mathrm{heav}$ et $\vp_\mathrm{conv}$ les fonctions des \cref{prop:cinfHeaviside,prop:cinfBosse}. Posons 
\begin{equation}
	\vp_1(x) := \sum_{1 \leq i \leq I} \vp_\mathrm{conv}\Big(\frac{x-x_i} {r/3}\Big).
\end{equation}
Alors $\vp_1\in C^\infty$, $\supp(\vp_1)\csubset \Omega$, et $\vp_1 > 0$ sur $K$.
\claudereport{Problème : $\vp_\mathrm{conv}((x-x_i)/(r/3))$ n'est strictement positive que sur $B(x_i,2r/9)$, ce qui ne recouvre pas forcément $K$ ; utiliser $(x-x_i)/r$ (support dans $B'(x_i,2r/3) \subset \Omega$, positivité sur $B(x_i,2r/3) \supset B(x_i,r/3)$).}
Posons $c := \inf_K \vp_1$ et notons que $c>0$ par compacité. 
On définit finalement $\vp(x) := \vp_\mathrm{heav}(\vp_1(x)/c)$ qui convient.
\end{proof}

\begin{proposition}[Partition de l'unité]
\label{prop:cinfPartition}
	Soit $K \csubset \bR^d$ compact, dont on considère une couverture $K \subset \cup_{i=1}^I \Omega_i$ par des ouverts. Alors il existe $(\vp_i)_{i=1}^I$ telles que $\vp_i \in C^\infty(\bR^d)$, $\supp \vp_i \csubset \Omega_i$, et $\sum_{i=1}^I \vp_i = 1$ sur $K$.
\end{proposition}

\begin{proof}
	Définissons $r := \min_{x \in K} \max_{i=1}^I d(x,\bR^d \sm \Omega_i)$, appelé nombre de Lebesgue de la couverture, et notons que $r>0$ par compacité de $K$.  Posons $K_i := \{x \in K \mid d(x,\bR^d \sm \Omega_i) \geq r\}$ et notons que $K_i \csubset \Omega_i$ et que $\cup_{i=1}^I K_i = K$.
	Choisissons par la \cref{prop:cinfSupport} une fonction $\psi_i \in C^\infty_c(\Omega_i,[0,1])$ telle que $\psi_i=1$ sur $K_i$, pour tout $1 \leq i \leq I$, étendue par $0$ hors de $\Omega_i$. Les fonctions $\vp_i := \psi_i / \sum_{j=1}^I \psi_j$, où $1 \leq i \leq I$, conviennent. 
\claudereport{Problème : $\sum_j \psi_j$ s'annule loin de $K$, donc $\vp_i := \psi_i/\sum_j\psi_j$ n'est pas définie (ni $C^\infty$) sur $\bR^d$ ; par ex. multiplier par une fonction $\chi \in C^\infty_c$ égale à $1$ sur $K$ et à support dans $\{\sum_j \psi_j > 0\}$.}
\end{proof}

\begin{proposition}[Exhaustion lisse] 
\label{prop:cinfExhaust}
Soit $\Omega \subset \bR^d$ un ouvert, et et soit $(K_n)_{n \in \bN}$ une suite exhaustive de compacts, voir la \cref{prop:exhaustComp}; posons $K_{-1} := \emptyset$.
\claudereport{Typo : ``et et soit'' $\to$ ``et soit''.}
Alors il existe une suite de fonctions $\vp_n \in C^\infty_c(\bR^d,[0,1])$ telles que $\supp(\vp_n) \subset \mathring K_{n+1}\sm K_{n-1}$ et $\sum_{n \in \bN} \vp_n = 1$ sur $\Omega$. 
\end{proposition}

\begin{proof}
Pour tout $n \in \bR$, introduisons $K_{n+1/2}$ tel que $K_n \csubset \mathring K_{n+1/2}$ et $K_{n+1/2} \csubset K_{n+1}$. Par exemple $K_{n+1/2} := \{x \in \bR^d \mid d(x,K_n) \leq r/2\}$ où $r := \inf \{d(x,\bR^d \sm K_{n+1}) \mid x \in K_n\}$. 
\claudereport{Typo : ``Pour tout $n \in \bR$'' $\to$ ``Pour tout $n \in \bN$''.}

	Choisissons $\psi_n \in C^\infty_c(\mathring K_{n+1} \sm  K_{n-1}, [0,1])$ telle que $\psi_n=1$ sur $K_{n+1/2}\sm \mathring K_{n-1/2}$, par la \cref{prop:cinfSupport}, pour tout $n \in \bN$ et avec $K_{-1/2} := \emptyset$. 
	Alors $\vp_n := \psi_n / \sum_{m \in \bN} \psi_m$ convient; noter qu'en chaque point $x$, la somme contient au plus deux termes non-nuls. 
\end{proof}


\section{Espaces Hilbertiens}
\label{sec:Hilbert}

La structure Hilbertienne permet de définir la projection sur un convexe fermé, et d'identifier l'espace avec son dual via le théorème de Riesz, voir les \cref{th:projConvexe,th:reprRiesz}. Ces résultats puissants trivialisent pour l'essentiel les questions de réflexivité et de séparation de convexes étudiées dans le cas des espaces de Banach, voir les \cref{subsec:reflexivite,subsec:HahnBanachGeom}.

Dans cette section, $\bK$ désigne le corps $\bR$ ou $\bC$.

\begin{retenir}
	L'inégalité de Cauchy-Schwarz, la définition et la caractérisation de la projection sur un convexe fermé, le théorème de Représentation de Riesz, de la base Hilbertienne.
	(\Cref{prop:CauchySchwarz,th:projConvexe,th:reprRiesz,th:baseHilbert}.)
\end{retenir}

\inventaire{
\Nref{def:prehilbert} \\
\Nref{prop:CauchySchwarz} \\
\Nref{def:Hilbert} \\
\Nref{th:projConvexe} \\
\Nref{th:reprRiesz} \\
\Nref{ex:l2Hilbert} (Hilbert $L^2(X,\mu)$)\\
\Nref{prop:decompOrth} \\
\Nref{prop:decompCone} \\
\Cref{ex:coneSdp}(Cône dual de $S_d^+$) \\
\Nref{th:baseHilbert}
}

\begin{definition}[Espace préhilbertien]
\label{def:prehilbert}
	Une forme sesquilinéaire sur un $\bK$-espace vectoriel $E$ est une application $\vp : E \times E \to \bK$ qui satisfait pour tous $x,x',y,y' \in E$, $\lambda \in \bK$ : 
	\begin{itemize}
		\item (Linéarité à droite) $\vp(x,\lambda y) = \lambda \vp(x,y)$, $\vp(x,y+y')=\vp(x,y)+\vp(x,y')$. %, $\forall x,y,y' \in E$, $\lambda \in \bK$
		\item (Antilinéarité à gauche) $\vp(\lambda x,y) = \overline \lambda \vp(x,y)$, $\vp(x+x',y)=\vp(x,y)+\vp(x',y)$.  
	\end{itemize}
	Elle est dite 
	\begin{itemize}
		\item \emph{symétrique} si $\vp(x,y) = \overline {\vp(y,x)}$, $\forall x,y \in E$. 
		\item \emph{positive} si $\vp(x,x) \in \bR_+$, $\forall x \in E$.
		\item \emph{définie} si $\vp(x,x) \neq 0$, $\forall x \in E\sm \{0\}$.
	\end{itemize}
	Un \emph{produit scalaire} $\<\cdot,\cdot\>$ est une forme sesquilinéaire symétrique définie positive. Dans ce cas, l'espace est dit \emph{préhilbertien}, et on note $\|x\| := \sqrt{\<x,x\>}$ pour tout $x \in E$. 
\end{definition}

Tous les éléments $x,y$ d'un espace préhilbertien, satisfont les égalités suivantes : 
\begin{align*}	
	\|x+y\|^2 &= \|x\|^2 + \|y\|^2 + 2 \Re(\<x,y\>), &
	\|x+y\|^2 + \|x-y\|^2 &= 2 (\|x\|^2 + \|y\|^2),
\end{align*}
la seconde étant appelée identité du parallélogramme. 

\begin{proposition}[Inégalité de Cauchy-Schwarz]
\label{prop:CauchySchwarz}
Soit $(E,\<\cdot ,\cdot\>)$ un espace préhilbertien. Alors pour tous $x,y \in E$, 
\begin{equation}
\label{eq:CauchySchwarz}
	|\<x,y\>| \leq \|x\| \|y\|, 
\end{equation}
et $\|\cdot\|$ définit une norme sur $E$. L'égalité a lieu ssi $x$ et $y$ sont colinéaires. 
\end{proposition}

\begin{proof}
	L'égalité a clairement lieu dans \eqref{eq:CauchySchwarz} si $x$ et $y$ sont colinéaires. On  suppose désormais qu'ils ne le sont pas. 
	Soit $\alpha\in \bK$, $|\alpha|=1$, tel que $\Re(\alpha\<x,y\>) = |\<x,y\>|$.
	Considérons le polynôme %du second degré 
	\begin{equation*}
		P(t) := \|x+\alpha t y\|^2 
		= \<x+\alpha t y, x+\alpha t y\>
		= \<x,x\> + 2 t \Re(\alpha\<x,y\>)+ t^2\<y,y\> 
		= \|x\|^2 + 2t|\<x,y\>| + t^2\|y\|^2,
	\end{equation*}
	d'indéterminée $t \in \bR$.
	Comme $P$ est du second degré et ne s'annule pas, son discriminant est strictement négatif, d'où \eqref{eq:CauchySchwarz}. 
	
	La fonction $\|\cdot\|$ est clairement définie, positive et homogène, conformément aux axiomes de la \nref{def:norme}. Il reste donc à vérifier l'inégalité triangulaire : $\|x+y\|^2 = \|x\|^2 + 2\Re(\<x,y\>) + \|y\|^2 \leq \|x\|^2 + 2 \|x\| \|y\| + \|y\|^2 = (\|x\|+\|y\|)^2$, $\forall x,y \in E$, ce qui conclut. 
\end{proof}

\begin{definition}[Hilbert]
\label{def:Hilbert}
%Un espace vectoriel normé dont la norme découle d'un produit scalaire préhilbertien. 
Un espace préhilbertien complet est un Hilbert. 
\end{definition}

\begin{theorem}[Projection sur un convexe fermé]
\label{th:projConvexe}
Soit $H$ Hilbert et $K\subset H$ un convexe fermé non-vide. 
Pour tout $x \in H$, le problème suivant admet une unique solution
\begin{equation}
\label{eqdef:proj_cvx_ferme}
	P_K(x) := \argmin_{y \in K} \|x-y\|,
\end{equation}
définissant la projection $P_K : H \to K$. Elle est caractérisée par la condition d'angle obtus : pour tout $y \in K$
\begin{equation}
\label{eq:proj_obtus}
	\Re\<x-P_K(x), y-P_K(x)\> \leq 0.
\end{equation}
Enfin, $P_K$ est $1$-Lipschitz, et plus précisément pour tous $x,y \in H$
\begin{equation}
\label{eq:proj_moyennee}
	\|P_K(x)-P_K(y)\|^2 \leq \Re\<x-y,\ P_K(x)-P_K(y)\>.
\end{equation} 
\end{theorem}

\begin{proof}
Notons $\delta:=d_K(x)$ l'infimum de \eqref{eqdef:proj_cvx_ferme}, et 
soit $(x_n)_{n \in \bN}$ une suite minimisante, i.e.\ $x_n\in K$ et $\ve_n := \|x_n - x\|-d_K(x) \to 0$ lorsque $n \to \infty$.
Etant donnés $p,q \geq N$ on a 
\begin{equation*}
	\delta^2 \leq \big\| x - \frac{x_p+x_q} 2\big\|^2 
	= \frac 1 2 \big(\|x-x_p\|^2 + \|x-x_q\|^2\big) - \frac 1 4\|x_p-x_q\|^2
	\leq \frac 1 2 \big((\delta+\ve_p)^2 + (\delta+\ve_q)^2\big)	- \frac 1 4\|x_p-x_q\|^2,
\end{equation*}
par l'identité du parallélogramme, d'où $\|x_p-x_q\| \leq 2\sqrt{\delta(\ve_p+\ve_q)+ \frac 1 2 (\ve_p^2+\ve_q^2)}$. Ainsi $(x_n)_{n \in \bN}$ est de Cauchy, donc convergente par complétude de $H$, et par continuité de la norme et fermeture de $K$ sa limite est un minimiseur de \eqref{eqdef:proj_cvx_ferme}.
On déduit aussi de cette preuve l'unicité du minimiseur de \eqref{eqdef:proj_cvx_ferme}.

Considérons $x_* := P_K(x)$, et $y\in K$ arbitraire. On a $(1-t) x_*+ ty \in K$ pour tout $t \in [0,1]$, par convexité de $K$, puis
\begin{equation}
\label{eq:caract_proj_proof}
	\|x_*-x\|^2 \leq \|(1-t) x_* + t y - x\|^2 = \|x-x_*\|^2 - 2 t \Re(\<x-x_*,y-x_*\>) + t^2 \|y-x_*\|^2
\end{equation}
par minimalité de $x_*$ dans \eqref{eqdef:proj_cvx_ferme}. On en déduit $\Re(\<x-x_*,y-x_*\>)\leq 0$ par un développement limité en $t=0^+$, comme annoncé. 
Réciproquement, considérons $x_* \in K$ tel que $\Re(\<x-x_*,y-x_*\>)\leq 0$ pour tout $y \in K$.
Alors $\|y-x\|^2 = \|x-x_*\|^2 - 2 \Re(\<x-x_*,y-x_*\>) + \|y-x_*\|^2 \geq \|x-x_*\|^2$ pour tout $y \in K$, donc $x_*$ minimise \eqref{eqdef:proj_cvx_ferme}, comme annoncé.

Finalement \eqref{eq:proj_moyennee} s'obtient en sommant les inégalités $\Re\<x-P_K(x), P_K(y)-P_K(x)\> \leq 0$ et $\Re\<y-P_K(y), P_K(x)-P_K(y)\> \leq 0$, issues de \eqref{eq:proj_obtus}. La Lipschitzianité de $P_K$ en découle par la \nref{prop:CauchySchwarz}.
\end{proof}

La propriété \eqref{eq:proj_moyennee} est équivalente à la $1$-Lipschitzianité de $Q := 2P_K-\Id$. On dit que $P_K := (\Id+Q)/2$ est un \emph{opérateur moyenné}: c'est la moyenne de l'identité et d'un opérateur $1$-Lipschitz.


Le résultat suivant montre que l'on peut identifier un Hilbert à son dual. Les normes se correspondent isométri\-quement, ainsi que les topologies faible et *faible. 

\begin{theorem}[Représentation de Riesz]
\label{th:reprRiesz}
Soit $H$ un espace de Hilbert. Alors $H \to H^* : x \mapsto \vp_x := \<x,\cdot\>$, est une bijection isométrique anti-linéaire.  
\end{theorem}

\begin{proof}
	\emph{Preuve de l'isométrie}. On a $\|\vp_x\|_{H^*} \geq \vp_x( x/\|x\|) = \|x\|$ pour tout $x \neq 0$, et par ailleurs $|\vp_x(y)| = |\<x,y\>| \leq \|x\| \|y\|$ par Cauchy-Schwarz, d'où $\|\vp_x\|_{H^*} = \|x\|$, comme annoncé.
	
	\emph{Preuve de la surjectivité}. Soit $\vp \in H^*\sm \{0\}$, et soit $K := \vp^{-1}\{0\}$ son noyau, qui est convexe en tant qu'hyperplan, et fermé car $\vp$ est continue. 
	Soit $x_0\in H$ tel que $\vp(x_0)=1$, qui existe par homogénéité, et soit $x_* := x_0 - P_K(x_0)$, de sorte que $\vp(x_*) = \vp(x_0) = 1$ et $\<x_*,y\> = 0$ pour tout $y \in K$. Alors pour tout $x\in H$ on a $\vp(x-\vp(x) x_*) = 0$ par linéarité, donc $x-\vp(x) x_* \in K$, donc $\<x_*,x-\vp(x) x_*\> = 0$, donc $\vp(x) = \<x_*,x\>/\|x_*\|^2$ ce qui conclut. 
\end{proof}

\begin{exemple}
\label{ex:l2Hilbert}
	Soit $(X,\mu)$ un espace mesuré. Alors $L^2(X,\mu)$ est un Hilbert, muni du produit scalaire $\<f,g\> := \int_X \overline f(x) g(x) \diff\mu(x)$, voir l'\cref{ex:LpBanach}. En particulier, toute forme linéaire continue s'y écrit $g \mapsto \int_X \overline f(x) g(x) \diff\mu(x)$, pour un certain $f \in L^2(X, \mu)$.
\end{exemple}

%Les résultats suivants illustrent les décompositions d'élements d'un Hilbert obtenues grâce 
La projection orthogonale permet de décomposer les éléments d'un Hilbert.

\begin{proposition}[Décomposition orthogonale]
\label{prop:decompOrth}
	Soit $H$ Hilbert. 
	La projection sur un sous-espace vectoriel fermé $F \subset H$ est caractérisée par la condition $\<x-P_F(x),y\> = 0$, $\forall y \in F$.
	De plus $F \oplus F^\perp = H$, $P_F+P_{F^\perp} = \Id$, et $P_F$ et $P_{F^\perp}$ sont linéaires continues, où $F^\perp := \{y \in H \mid \forall x \in F,\, \<x,y\>=0\}$ désigne le sous-espace orthogonal.
\end{proposition}

\begin{proof}
	Soit $x \in H$ et $p := P_F(x) \in F$. D'après \eqref{eq:proj_obtus} on a $\Re\<x-p,y-p\> \leq 0$ pour tout $y \in F$. Donc $\Re\<x-p,y\> = 0$ pour tout $y \in F$, par stabilité additive de $F$, et donc $\<x-p,y\> = 0$, par stabilité de $F$ par multiplication par un scalaire, comme annoncé. La linéarité de cette caractérisation de $P_F(x)$ implique la linéarité du projecteur $P_F$, qui est également continu par le \cref{th:projConvexe}, comme annoncé.
	
	Notons $z := x-p$, et observons que $z \in F^\perp$ par la discussion précédente. Par ailleurs $x-z = p \in F \subset (F^\perp)^\perp$, donc par la caractérisation \eqref{eq:proj_obtus} on a $z = P_{F^\perp} x$. Ainsi $x = P_F x + P_{F^\perp} x$, comme annoncé, et ainsi $F + F^\perp = H$. Remarquons finalement que $F \cap F^\perp = \{0\}$ (si $x \in F \cap F^\perp$, alors $\|x\|^2 = \<x,x\> = 0$), donc $F \oplus F^\perp = H$ comme annoncé, ce qui conclut.
\end{proof}

Une variante du résultat classique précédent traite des cônes convexes fermés. Rappelons qu'un cône convexe est une partie $K$ d'un espace $\bR$-vectoriel telle que $\lambda x+\mu y \in K$ pour tous $x,y \in K$ et tous $\lambda,\mu \geq 0$. 

\begin{proposition}[Décomposition orthogonale associée à un cône \facultatif]
\label{prop:decompCone}
	Soit $H$ un $\bR$-Hilbert, soit $K \subset H$ un cône convexe fermé, et soit $K^\circ := \{y\in H\mid \forall x \in K,\ \<x,y\> \leq 0\}$ le cône polaire. %; noter que $K \cap K^\circ = \{0\}$.
	On a la caractérisation suivante des projections sur $K$ et $K^\circ$ : %sont caractérisées par %$P_K+P_K^\perp = \Id$ et $\<P_K,P_{K^\circ}\> = 0$, 
	\begin{equation}
	\label{eq:proj_cone}
		(y = P_K(x),\ z=P_{K^\circ} (x)) \ \Leftrightarrow \ (y\in K,\, z \in K^\circ,\, \<y,z\>=0 \text{ et } x = y + z). 
	\end{equation}
\end{proposition}

\begin{proof}
	Supposons (\ref{eq:proj_cone}, droite). Alors pour tout $y' \in K$ on a $\<x-y,y'-y\> = \<z,y'-y\> = \<z,y'\> \leq 0$, donc $y = P_K(x)$ par la caractérisation \eqref{eq:proj_obtus}. De même $z = P_{K^\circ} (x)$, comme annoncé.
	
	Réciproquement, considérons $y:=P_K(x)$ et posons $z := x-y$. On a $\<x-y,y'-y\> =\<z,y'-y\> \leq 0$, $\forall y' \in K$, par \eqref{eq:proj_obtus}. Par la propriété de cône on obtient $t \<z,y'\> \leq \<z,y\>$ pour tout $t \geq 0$, donc $\<z,y'\> \leq 0$, et donc $z \in K^\circ$. 
	A l'inverse, en choisissant $y'=0$ on obtient $\<z,0-y\>\leq 0$, donc $\<z,y\>=0$. Finalement $\<x-z,z'-z\> = \<y,z'-z\> = \<y,z'\> \leq 0$ pour tout $z' \in K^\circ$, donc $z = P_{K^\circ}(x)$, ce qui conclut.  
\end{proof}

\begin{exemple}[Cône dual de $\bR_+^d$ et $\cS^+_d$ \facultatif]
\label{ex:coneSdp}
	On note parfois $K^* := - K^\circ  = \{y \in H\mid \forall x \in K, \<x,y\> \geq 0\}$ le cône dual, d'un cône $K$ donné. 
	Un cône est dit auto-dual si $K = K^*$. Par exemple $\bR_+^d \subset \bR^d$ et $\cS_d^+ \subset \cS_d$ sont des cônes auto-duaux. La décomposition de la \cref{prop:decompCone} associée au cône $\bR_+^d$ correspond aux parties positive et négative de $x=(x_i)_{1 \leq i \leq d} \in \bR^d$~: 
	\begin{align*}
		x &= x^+ - x^-, &\text{où } x^+_i &:= \max\{0,x_i\},\ x^-_i := \max\{0,-x_i\}
	\end{align*}
	De même, la décomposition de la \cref{prop:decompCone} associée au cône $\cS_+^d$ correspond aux parties positive et négative du spectre : pour $D \in \cS_d$, diagonalisée sous la forme $U^\top \Diag(\lambda_1,\cdots,\lambda_d) U$ avec $U$ une matrice orthogonale, on a 
	\begin{equation*}
		D  
		= U^\top \Diag(\max \{\lambda_1,0\},\cdots,\max \{\lambda_d,0\}) U - U^\top \Diag(\max \{-\lambda_1,0\},\cdots,\max \{-\lambda_d,0\}) U.
	\end{equation*}
\end{exemple}

%En optimisation, le cône normal est le dual du cône intérieur, en un point $x$ d'un convexe $K$.
%\begin{align*}
%	\cC_K^\mathrm{int}(x) &= \{v \in H \mid \exists t >0, x+t v\in K\},&
%	\cC_K^\mathrm{norm}(x) &= \{v \in H \mid \forall y \in K, \<v,y-x\> \leq 0\}.
%\end{align*}
%Noter que le cône normal est fermé, mais pas forcément le cône intérieur.
%Dans un Banach, des notions analogues sont discutées après \nref{def:convexePolaire}.

%\todoi{
%Si nécessaire, évoquer le fait que les Hilbert séparables de dimension infinie sont isomorphes à $L^2(\bN,\bK)$.
%}

La structure Hilbertienne est si forte que $\ell^2 := L^2(\bN,\bK)$ est, à isométrie linéaire près, son seul représentant intéressant, comme le montre le théorème suivant qui conclut cette introduction. 

\begin{theorem}[Base Hilbertienne]
\label{th:baseHilbert}
Soit $H$ un espace de Hilbert. Sont équivalents : 
\begin{enumerate}[(i)]
	\item $H$ est séparable et de dimension infinie.
	\item $H$ admet une \emph{base hilbertienne}, c'est-à-dire une famille $(e_n)_{n\in \bN}$ de vecteurs unitaires deux à deux orthogonaux, telle que $\Vect\{e_n \mid n \in \bN\}$ est dense dans $H$.
	\item Il existe une bijection linéaire et isométrique entre $H$ et $\ell^2 := L^2(\bN,\bK)$.
\end{enumerate}
En supposant (ii), la bijection de (iii) est donnée par $x \in H \mapsto (\<e_n,x\>)_{n \in \bN} \in \ell^2$, et sa réciproque par $(x_n)_{n \in \bN} \mapsto \sum_{n \in \bN} x_n e_n$.
\end{theorem}

\begin{proof}
	\emph{Preuve de (i) $\Rightarrow$ (ii).} Soit $(x_n)_{n \in \bN}$ une suite dense dans $H$, et soit $F_N := \Vect\{x_n \mid n < N\}$ pour tout $N \geq 0$, avec la convention $F_0:=\{0\}$.
	Comme $F_n \subset F_{n+1}$ pour tout $n \geq 0$, comme l'espace vectoriel $\cup_{n \in \bN} F_n$ est dense dans $H$ par hypothèse, et comme $H$ est de dimension infinie, on a $\dim(F_n) \to \infty$ lorsque $n \to \infty$; de plus, $\dim(F_{n+1}) \leq \dim(F_n)+1$ pour tout $n \in \bN$ par construction. 
	Ainsi, il existe $\vp : \bN \to \bN$ extractrice telle que $\dim(F_{\vp(n)}) = n$ pour tout $n \in \bN$; quitte à extraire, on peut supposer $\dim(F_n) = n$. 
	On vérifie alors facilement que $e_n := (x_n - P_{F_n}(x_n)) / \|x_n - P_{F_n}(x_n)\|$, pour tout $n \in \bN$, définit une base Hilbertienne.
	
	\emph{Preuve de (ii) $\Rightarrow$ (iii).}
	La projection orthogonale de $x \in H$ sur $F_N$ est donnée par $P_{F_N}(x) = \sum_{n< N} \<e_n,x\>e_n$, pour tout $N \geq 0$, d'après le critère de la \cref{prop:decompOrth}, et donc $\|P_{F_N}(x)\|^2 = \sum_{n<N} |\<e_n,x\>|^2$. 
	Ainsi $\|x\|^2 - \sum_{n<N} |\<e_n,x\>|^2 = \|x-P_{F_N}(x)\|^2 \to 0$ lorsque $N \to \infty$, ce qui montre la bonne définition et l'isométrie.
	Réciproquement, si $(x_n)_{n\in \bN} \in \ell^2$, alors les sommes partielles $s_N := \sum_{n<N} x_n e_n$ satisfont le critère de Cauchy, puisque $\|s_n-s_m\|^2 = \sum_{n\leq k <m} |x_k|^2$, donc convergent vers un élément $x\in H$ par complétude, et celui-ci satisfait $\<e_n,x\> = \lim_{N \to \infty} \<e_n,s_N\> = x_n$ pour tout $n \in \bN$ par continuité du produit scalaire.
	
	\emph{Preuve de (iii) $\Rightarrow$ (i).} Les éléments de $\ell^2$ à support fini et coefficients dans $\bQ$ forment une famille dénombrable et dense, d'où la séparabilité. De plus $\ell^2$ contient la famille libre infinie  $(\delta^N)_{N\geq 0}$, sa base hilbertienne canonique, définie par $\delta^N_n = 1$ si $n=N$ et $\delta^N_n = 0$ sinon.
\end{proof}



\section{Dualité de Legendre-Fenchel}
\label{sec:LegendreFenchel}
Cet appendice ouvre sur la thématique de l'optimisation convexe, en tant qu'application du \nref{th:HahnBanach}, et parce que de nombreuses EDPs peuvent s'interpréter dans ce cadre, comme le \nref{th:pbPoisson} ou l'\cref{ex:nonLinMinH1}.
%En physique, la dissipation de l'énergie d'un système, ou la formulation Lagrangienne de ses lois d'évolution, amènent naturellement à des problèmes d'optimisation. 

\begin{retenir}
Cet appendice doit être considéré comme une ouverture, dont le contenu n'est pas exigible pour ce cours d'analyse fonctionnelle. La dualité de Legendre-Fenchel, la notion de sous-gradient d'une fonction convexe, et le théorème de Fenchel-Rockafellar sont toutefois fondamentaux en optimisation convexe. (\Cref{th:LegendreFenchel,th:ssGrad,th:FenchelRockafellar}.)
\end{retenir}

\inventaire{
\Nref{sec:LegendreFenchel}\\
\Nref{def:legendre_fenchel}\\
\Nref{def:domaineFct} (Domaine d'une fonction)\\
\Nref{lem:FenchelPropre} ($f^*$ propre ssi $f$ a un minorant affine)\\
\Nref{lem:minorantCvxSci} (Existence d'un minorant affine d'une fonction cvx, sci)\\
\Nref{th:LegendreFenchel} \\
\Nref{lem:cvxImplBorne} (Convexe implique localement borné)\\
\Nref{lem:cvxImplLip} (Convexe implique localement Lipschitz)\\
\Nref{th:ssGrad}\\
\Nref{lem:cvxSciAlt} (Alternative pour l'enveloppe sci d'une fonction convexe)\\
\Nref{th:FenchelRockafellar}
}

Étant donné un élément $x\in E$ d'un evn, et une forme linéaire continue $\vp \in E^*$, on note parfois $\<\vp,x\> := \vp(x)$ le \emph{crochet de dualité}. 

\begin{definition}[Conjuguée de Legendre-Fenchel]
\label{def:legendre_fenchel}
\leancheck{def:legendre_fenchel}%
Soit $E$ un evn et $f : E \to \overline \bR$ une fonction. Sa conjuguée de Legendre-Fenchel est définie pour toute $\vp \in E^*$ par 
	\begin{equation}
	\label{eqdef:lengedre_fenchel}
		f^*(\vp) := \sup_{x \in E} \<\vp,x\> - f(x).
	\end{equation}
\end{definition}

Par exemple, si $f(x) = \frac 1 2 \|x\|^2_E$, alors on vérifie que $f^*(\vp) = \frac 1 2 \|\vp\|_{E^*}^2$. 

\begin{definition}
\label{def:domaineFct}
\leancheck{def:domaineFct}%
	Étant donnée une fonction $f : X \to ]-\infty,\infty]$ définie sur un ensemble $X$, on note $\Dom(f) := \{x \in X \mid f(x) < \infty\}$ son domaine, et on dit qu'elle est propre si $\Dom(f) \neq \emptyset$. 
\end{definition}


On montre dans le lemme suivant que $f^*$ est propre ssi $f$ admet un minorant affine continu. Le lemme ultérieur est dédié à la construction de tels minorants affines, pour les fonctions convexes sci.


\begin{lemma}
\label{lem:FenchelPropre}
\leancheck{lem:FenchelPropre}%
	Soit $f : E \to ]-\infty,\infty]$ une fonction propre sur un evn. Alors  $f^* : E^* \to ]-\infty,\infty]$ est convexe et sci. De plus, $f^*$ est propre ssi $f$ admet un minorant affine continu; plus précisément $f \geq \vp + \alpha$ ssi $f^*(\vp) \leq -\alpha$, pour tous $\alpha \in \bR$, $\vp \in E^*$.
\end{lemma}

\begin{proof}
	La fonction $f^*$ est convexe et sci comme supremum de la famille non-vide de fonctions $\vp \in E^* \mapsto \<\vp,x\> - f(x) \in \bR$, associées à chaque $x \in \Dom(f)$, qui sont affines et continues, donc convexes et sci. Le second point est immédiat compte tenu de la définition \eqref{eqdef:lengedre_fenchel} de $f^*$.
\end{proof}

\begin{lemma}
\label{lem:minorantCvxSci}
\leancheck{lem:minorantCvxSci}%
	Soit $f : E \to ]-\infty,\infty]$ convexe et sci, et soient $x_0 \in E$ et $t_0 < f(x_0)$. Alors il existe $\alpha \in \bR$, $\vp \in E^*$ tels que $f > \alpha + \vp$ sur $E$, et $\alpha + \vp(x_0) > t_0$. 
\end{lemma}

\begin{proof}
	Le sur-graphe $C := \{(x,t) \mid f(x) \leq t\} \subset E \times \bR$ est convexe et fermé, voir la \cref{prop:surGrapheSCI}. 
	Par ailleurs une forme linéaire sur $E \times \bR$ est continue ssi elle prend la forme $(x,t) \mapsto \vp(x) - \lambda t$ avec $\vp \in E^*$ et $\lambda \in \bR$. 

	Par le \nref{th:separFermeComp} il existe $(\vp,\lambda) \in E^* \times \bR$ et $\delta>0$ tels que $\vp(x)-\lambda t +\delta \leq  \vp(x_0) - \lambda t_0$ pour tout $x\in E$ et $t \geq f(x)$. 
	Supposons d'abord $f(x_0) < \infty$.  
	En choisissant $x=x_0$ on obtient $\lambda >0$, et donc le minorant affine continu $x\in E \mapsto \vp(x-x_0)/\lambda +t_0$.
		
	Supposons maintenant $f(x_0) = \infty$. En considérant un point $x\in \Dom(f)$, on obtient que $\lambda \geq 0$. Si $\lambda>0$ c'est terminé, et sinon on a obtenu $\vp \in E^*$ tel que $\vp(x) +\delta \leq \vp(x_0)$ pour tout $x\in \Dom(f)$. Choisissons un minorant affine $f_1$ de $f$, qui existe par l'argument précédent puisqu'on peut supposer $f$ propre (sinon le résultat est trivial). 
	Alors $f_1(x) + \mu (\vp(x-x_0)+\delta)$ est un minorant affine de $f$ pour tout $\mu \geq 0$, et prend la valeur $f_1(x_0)+\mu \delta$ en $x_0$, qui est arbitrairement grande lorsque $\mu \to \infty$, ce qui conclut. 
\end{proof}


Le théorème de dualité de Legendre-Fenchel assure que $f = f^{**}$ (moyennant l'identification $E=E^{**}$) pour toute fonction $f : E \to ]-\infty,\infty]$ propre, convexe et sci sur un espace réflexif $E$. On peut faire le lien avec la \nref{prop:cvxFaibleFerme} affirmant similairement que $C = C^{\circ\circ}$ pour tout convexe, fermé contenant l'origine d'un espace réflexif.

\begin{theorem}[Dualité de Legendre-Fenchel]
\label{th:LegendreFenchel}
\leancheck{th:LegendreFenchel}%
	Soit $f : E \to ]-\infty,\infty]$ une fonction propre et admettant un minorant affine sur un evn (c'est toujours le cas si $f$ est convexe et sci). 
	Alors $f^{**}$ est convexe, sci, propre, et $f^{**} \leq f$ sur $E$ avec égalité ssi $f$ est convexe et sci. 
\end{theorem}

\begin{proof}
	Par le lemme précédent, $f^*$ est convexe et sci, donc admet un minorant affine par le \cref{lem:minorantCvxSci}, et propre. 
	Ainsi $f^{**}$ est convexe, sci, et propre comme annoncé. 
	Par construction, on a $f(x) \geq \<\vp,x\> - f^*(\vp)$ pour tout $\vp \in E^*$, donc $f \geq f^{**}$ en prenant le supremum. 
	
	Supposons $f$ convexe et sci. Soit $x_0 \in E$ et soit $t_0 < f(x_0)$. Par le \cref{lem:minorantCvxSci}, on a un minorant affine $f \geq \alpha + \vp$ tel que $\alpha + \vp(x_0) \geq t_0$. Alors $f^*(\vp) \leq - \alpha$, et donc $f^{**} (x_0) \geq \vp(x_0) + \alpha \geq t_0$. Il s'ensuit que $f^{**} \geq f$, et donc $f=f^{**}$ ce qui conclut.
\end{proof}

La dualité de Legendre-Fenchel permet d'exprimer une fonction convexe comme enveloppe de fonctions affines. 
On peut remarquer que si $f\leq g$ alors $f^*\geq g^*$ et donc $f^{**} \leq g^{**}$; on en déduit que  $f^{**}$ est la plus grande fonction convexe sci bornée supérieurement par $f$ (supposée admettre un minorant affine).

Dans le \cref{th:ssGrad} suivant, on montre l'existence d'une fonction affine tangente à une fonction convexe, aux points de l'intérieur de son domaine de définition. 

\begin{lemma}
\label{lem:cvxImplBorne}
\leancheck{lem:cvxImplBorne}%
\claudereport{Problème (preuve, relevé lors de la formalisation Lean) : l'union des $\{f \leq n\}$ est $\Dom(f)$ et non $E$, donc le \cref{corol:BaireFerme} ne s'applique pas tel quel. Appliquer Baire dans la boule fermée (complète) $B'(0,\ve) \subset \Dom(f)$, recouverte par les $\{f\leq n\} \cap B'(0,\ve)$ ; on obtient alors $x_0 \in B'(0,\ve)$, ce qui assure aussi $-\ve x_0 \in \Dom(f)$ (sinon $f(-\ve x_0)$ peut valoir $+\infty$).}
	Soit $f$ une fonction convexe s.c.i.\ sur un Banach. Alors $f$ est localement bornée sur $\mathring \Dom(f)$. 
\end{lemma}

\begin{proof}
	Supposons $B'(0,\ve) \subset \Dom(f)$, quitte à considérer une translation $x \mapsto f(x+x_*)$ où $x_* \in \mathring \Dom(f)$, et montrons que $f$ est bornée au voisinage de $0$.
	Comme $f$ est sci, les ensembles $(\{f \leq n\})_{n \in \bN}$ sont fermés. Comme leur union est $\Dom(f)$, l'un d'entre eux est d'intérieur non-vide. Par le \cref{corol:BaireFerme}, il existe donc $x_0 \in E$, $r>0$ et $n \in \bN$ tels que $f\leq n$ sur $B'(x_0,r)$. 
	Par convexité, $f(\ve h/(1+\ve)) \leq \ve f(x_0+h)/(1+\ve) + f(-\ve x_0)/(1+\ve)$, donc $f \leq M_0 := \ve n /(1+\ve) + f(-\ve x_0)/(1+\ve)$ sur $B'(0,r_0)$ avec $r_0 := \ve r/(1+\ve)$.
\end{proof}

\begin{lemma}
\label{lem:cvxImplLip}
\leancheck{lem:cvxImplLip}%
\claudereport{Détail (formalisation Lean) : supposer $f(x) < \infty$ (sinon l'hypothèse $f \leq f(x) + M$ est vide).}
	Soit $f : E \to ]-\infty,\infty]$ une fonction convexe sur un evn. Supposons $f \leq f(x) + M$ sur $B'(x,r)$, où $x \in E$, $r>0$, $M > 0$. 
	Alors $|f(x)-f(y)|\leq (M/r) \|x-y\|$ pour tout $y \in B'(x,r)$. On en déduit que $f$ est $2 M/r$-Lipschitz sur $B'(x,r/3)$. 
\end{lemma}

\begin{proof}
	Supposons d'abord $f(0)=0$ et $f \leq 1$ sur $B'(0,1)$. Alors pour tout $x \in B'(0,1)$ et $t \geq 1$, on a par convexité $f(t x)-f(0) \geq t(f(x)-f(0))$, donc $f(x) \leq \|x\|$ en choisissant $t = 1/\|x\|$. Également $f(x)+f(-x) \geq 0$ par convexité, donc $f(x) \geq -f(-x) \geq -\|x\|$. L'estimation annoncée s'en déduit en considérant $(f( (\cdot+x)/r) - f(x))/M$.
	
	Passons à la propriété de régularité Lipschitz. Par le résultat précédent, on a $|f(y)-f(x)|\leq M/3$ pour tout $y\in B'(x,r/3)$, et donc $f \leq f(y) + \frac 4 3 M$ sur $B_y := B(y,2 r/3) \subset B'(x,r)$. Par la première partie, $|f(y)-f(z)|/\|y-z\| \leq (\frac 4 3 M) / (\frac 2 3 r)  = (2 M/r)$ pour tout $z \in B_y\sm \{y\}$. On conclut en notant que $B'(x,r/3) \subset B_y$. 
\end{proof}

%\notei{La notion d'intérieur relatif a été supprimée du cours}


\begin{theorem}[Sous-différentiel et sous-gradient d'une fonction convexe \facultatif]
\label{th:ssGrad}
\leancheck{th:ssGrad}%
	Soient $E$ un Banach et $f : E \to ]-\infty,\infty]$ convexe et sci. Alors $f$ est localement Lipschitz sur $\mathring\Dom(f)$. De plus %Si $x \in \Relint(\Dom(f))$ 
	\begin{equation*}
		\partial f(x) := \{\vp \in E^* \mid \forall y \in E, f(y) \geq f(x) + \vp(y-x)\},
	\end{equation*}
	est un ensemble non-vide pour tout $x \in \mathring \Dom(f)$, appelé \emph{sous-différentiel} de $f$ en $x$. Ses éléments sont appelés sous-gradients de $f$ en $x$.
\end{theorem}

\begin{proof}	
	Notons $D := \mathring{\Dom(f)}$ et soit $x_0 \in D$. 
	La restriction $f_{|D}$ est localement bornée par le \cref{lem:cvxImplBorne}, donc localement Lipschitz comme annoncé par le \cref{lem:cvxImplLip}, et donc l'ensemble $G := \{(x,t)\mid x \in D,\, f(x) < t\}$ est ouvert. On peut donc séparer $(x_0,f(x_0))$ de $G$ par le \nref{th:separOuvertPt} appliqué à l'evn $\bR \times E$: il existe $(\lambda,\vp) \in \bR \times E^*$ tels que $\lambda f(x_0) - \vp(x_0) < \lambda t - \vp(x)$ pour tous $x\in D$, $t>f(x)$. En choisissant $x=x_0$ on obtient $\lambda >0$. Alors $f(x) \geq f(x_0) + \vp(x-x_0)/\lambda$ comme annoncé pour tout $x \in D$, donc pour tout $x \in E$ puisque $D$ est un voisinage de $x_0$ et $f$ est convexe, comme annoncé. Ainsi $\partial f(x_0) \neq \emptyset$. 
%	Finalement, supposons seulement $x\in \Relint(\Dom(f))$. 
%	Quitte à translater on peut supposer que l'espace affine $F := \Aff(C)$ contient l'origine, et donc c'est un Banach car il est fermé. Par le résultat précédent, il existe $\vp_F \in F^*$ telle que $f(y) \geq f(x) + \vp_F(y-x)$ pour tout $y \in F$. Par le \cref{corol:prol_forme_evn} on a une extension continue $\vp\in E^*$ de $\vp_F$, et par hypothèse $f(y) = \infty \geq f(x) + \vp_F(y-x)$ pour tout $y \in E\sm F$, ce qui conclut.  
\end{proof}

On conclut avec un résultat fondamental d'optimisation : la dualité de Fenchel-Rockafellar. 
L'hypothèse $(A \Dom f) \cap \mathring\Dom g \neq \emptyset$ admet des affaiblissements et variantes, étudiés dans la littérature.
%Noter que l'hypothèse $(A \Dom f) \cap \mathring\Dom g \neq \emptyset$ admet des affaiblissements, comme $\Relint (A\Dom f-\Dom g) \neq \emptyset$, ainsi que plusieurs variantes ($\Relint$ désignant l'intérieur relatif).
Rappelons aussi que si $A\in L(E,F)$, alors $A^\top \in L(F^*,E^*) : \vp \mapsto \vp \circ A$ désigne l'application transposée.


\begin{lemma}
\label{lem:cvxSciAlt}
\leancheck{lem:cvxSciAlt}%
	Soit $E$ evn et $f : E \to ]-\infty,\infty]$ convexe. Posons $g(x) := \liminf_{y \to x} f(y)$, pour tout $x \in E$, son enveloppe sci (on note parfois $g = f_*$). Alors on a l'alternative : 
	\begin{enumerate}[(a)]
		\item $g: E \to ]-\infty,\infty]$ est convexe, sci, et satisfait $g^* = f^*$.
		\item $g : E \to \{-\infty,\infty\}$ prend seulement deux valeurs, et $\{g = -\infty\} = \overline{\Dom f}$. %\{f < \infty\} 
	\end{enumerate}
\end{lemma}

\begin{proof}
	Supposons qu'il existe $x_* \in E$ tel que $g(x_*) = -\infty$, et donc qu'il existe une suite convergente $x_n \to x_*$ lorsque $n \to \infty$, telle que $f(x_n) \leq -n^3$ pour tout $n \in \bN$.
	Soit $y_* \in \overline{\Dom f}$, de sorte qu'il existe une suite convergente $y_n \to y_*$ lorsque $n \to \infty$, telle que $f(y_n) \leq  n$ à partir d'un certain rang. 
	Alors $z_n := y_n (1-1/ n) + x_n /n$ satisfait $z_n \to y_*$ lorsque $n \to \infty$, et $f(z_n) \to -\infty$ par convexité. 
	Ainsi, si $g$ prend la valeur $-\infty$ en un point, alors l'alternative (b) est satisfaite. 
	Sinon, $g : E \to ]-\infty,\infty]$ est par construction clairement convexe, sci, et satisfait $g^*=f^*$, ce qui conclut. 
\end{proof}

\begin{theorem}[Fenchel-Rockafellar \facultatif]
\label{th:FenchelRockafellar}
\leancheck{th:FenchelRockafellar}%
Soient $E$ et $F$ des Banach, et $A \in L(E,F)$ un opérateur linéaire continu. Soient $f: E \to ]-\infty,\infty]$ et $g : F \to ]-\infty,\infty]$, convexes, s.c.i., et telles que $(A \Dom f) \cap \mathring\Dom\, g \neq \emptyset$. 
Alors 
\begin{equation*}
	\sup_{y \in F^*} -f^*(-A^\top y) - g^*(y) = \inf_{x \in E} f(x) + g(A x).
\end{equation*}
\end{theorem}

\begin{proof}
	Les membres de gauche et droite de l'égalité à démontrer, notés $\alpha$ et $\beta$, s'écrivent %respectivement
	\begin{align*}
		\alpha &:= \sup_{y \in F^*} \inf_{x \in E} c(x,y), & 
		\beta &:= \inf_{x \in E} \sup_{y \in F^*} c(x,y), &
		\text{où } c(x,y) := f(x) + \<A x,y\> - g^*(y).
	\end{align*}
	On a donc $\alpha \leq \beta$ par construction, et il suffit d'établir que $\alpha \geq \beta$ (i.e.\ l'absence de saut de dualité, ``duality gap'' en anglais).
	En particulier, on peut supposer $\beta \neq - \infty$, sinon il n'y a rien à démontrer; noter que $\beta \neq \infty$ par hypothèse.
	Définissons $\cF : F \to ]-\infty,\infty]$ par 
	\begin{equation*}
		\cF(u) := \inf_{x \in E} f(x) + g(A x+u). 
	\end{equation*}
	Cette fonction est convexe, et satisfait $\cF(0) = \beta \in \bR$.
	
%	Montrons que $\cF$ est s.c.i.\ en $0$. Quitte à translater, on peut supposer $0 \in (A \Dom f) \cap \Relint(\Dom g)$. Notons $\tilde F := \overline {\Aff(\Dom g)}$ et $\tilde E := A^{-1} \tilde F$, qui sont des sev fermés de $E$ et $F$ donc des Banach. Quitte à remplacer $E$ et $F$ par $\tilde E$ et $\tilde F$ on peut supposer $0 \in (A \Dom f) \cap \mathring \Dom g$. 
On dispose par hypothèse de $x_* \in E$, $\ve > 0$ et $M \in \bR$ tels que $f(x_*) + g(A x_* + u)  \leq M$ pour tout $u\in B_F(0,\ve)$, voir le \cref{lem:cvxImplBorne}. 	
	Soient $u_n \to 0$ et $x_n\in E$ telles que $f(x_n) + g(A x_n +u_n) \to \liminf_{u \to 0} \cF(u)$. Alors 
	\begin{equation*}
	\hspace{-0.5cm}
		f((1-t_n) x_n + t_n x_*) + g(A [(1-t_n) x_n + t_n x_*] )
		\leq (1-t_n) [f(x_n) + g(A x_n + u_n)] + t_n [f(x_*) + g(A x_* - u_n (\tfrac 1 {t_n}-1))]
	\end{equation*}
	par convexité de $(x,u) \mapsto f(x) + g(Ax + u)$, avec $t_n \in ]0,1]$. Choisissons une suite $(t_n)_{n\geq 0}$ telle que $t_n \to 0$ et $u_n (\frac 1 {t_n}-1) \in B_F(0,\ve)$ pour tout $n \in \bN$ assez grand. Comme $f$ et $g$ sont  s.c.i., l'inégalité précédente donne $\cF(0) = \liminf_{u \to 0} \cF(u)$, c'est-à-dire que $\cF$ est s.c.i.\ au point $0$.  

	Pour tout $y \in F^*$ on a 
	\begin{align*}
		\cF^*(y) &= \sup_{u\in F} \sup_{x \in E} \<y,u\> -f(x) - g(Ax+u)\\
		&= \sup_{x \in E} \sup_{v \in F} \<y,v-Ax\> - f(x) - g(v)\\
		&=\sup_{x \in E} \<y, -Ax\> - f(x) + g^*(y)\\
		&= f^*(-A^\top y) + g^*(y). 
	\end{align*}
	On a utilisé (i) l'expression de $\cF$, (ii) l'interversion des supremum, toujours licite, et le changement de variable $v=Ax+u$, (iii - iv) la définition de la conjuguée de Legendre-Fenchel.
	Comme $\cF$ est convexe, et sci au point $0$ avec une valeur $\cF(0) = \beta$ finie, elle correspond au cas (a) du \cref{lem:cvxSciAlt}. On a donc $\cF(0) = \cF^{**}(0) = \sup_{y \in F^*} - f^*(-A^\top y) - g^*(y)$ comme annoncé.
\end{proof}


\end{document}